\documentclass[11pt]{amsart}
\usepackage{amssymb,amsmath,amsthm,mathtools}
\usepackage[margin=1.15in]{geometry}
\usepackage{booktabs,graphicx,array}
\usepackage{enumitem}
\usepackage{cite}
\usepackage[hidelinks]{hyperref}
\usepackage{longtable}
\hypersetup{pdftitle={Convex non-ball Schiffer domains in infinitely many dimensions},pdfauthor={Jizhou Guo}}
\newtheorem{theorem}{Theorem}[section]
\newtheorem{maintheorem}{Theorem}
\renewcommand{\themaintheorem}{\Alph{maintheorem}}
\newtheorem{proposition}[theorem]{Proposition}
\newtheorem{lemma}[theorem]{Lemma}
\newtheorem{corollary}[theorem]{Corollary}
\newtheorem{conjecture}{Conjecture}
\theoremstyle{definition}
\newtheorem{definition}[theorem]{Definition}
\theoremstyle{remark}
\newtheorem{remark}[theorem]{Remark}
\numberwithin{equation}{section}
\newcommand{\R}{\mathbb{R}}
\newcommand{\C}{\mathbb{C}}
\newcommand{\Z}{\mathbb{Z}}
\newcommand{\D}{\mathbb{D}}
\newcommand{\Sc}{\mathcal{S}}
\newcommand{\Gc}{\mathcal{G}}
\newcommand{\Fc}{\mathcal{F}}
\newcommand{\Nc}{\mathcal{N}}
\newcommand{\Bc}{\mathcal{B}}
\newcommand{\one}{\mathbf{1}}
\newcommand{\eps}{\varepsilon}
\newcommand{\bw}{\overline{w}}
\newcommand{\jst}{j_*}
\newcommand{\sU}{\mathsf{U}}
\DeclareMathOperator{\re}{Re}
\DeclareMathOperator{\im}{Im}
\DeclareMathOperator{\dist}{dist}
\newcommand{\nrm}[1]{\lVert #1\rVert}
\newcommand{\Ad}{\operatorname{Ad}}
\newcommand{\kf}{\mathfrak{k}}
\newcommand{\tf}{\mathfrak{t}}
\newcommand{\Hc}{\mathcal{H}}
\newcommand{\Yc}{\mathcal Y}
\newcommand{\Ecal}{\mathcal E}
\newcommand{\PP}{\mathsf P}
\newcommand{\QQ}{\mathsf Q}
\DeclareMathOperator{\tr}{tr}
\DeclareMathOperator{\diag}{diag}
\DeclareMathOperator{\acosh}{acosh}
\newcommand{\Cc}{\mathcal C}
\newcommand{\Mmix}{\mathcal M}
\newcommand{\Ec}{\mathcal E}
\newcommand{\Xc}{\mathcal X}
\newcommand{\Dom}{\operatorname{Dom}}
\newcommand{\Euler}{\mathsf E}
\newcommand{\SeedA}{\mathfrak A}
\newcommand{\Hu}{\widehat H_\Phi}
\newcommand{\HP}{\mathcal H_{\mathrm{ext}}}
\newcommand{\KD}{K_{\!D}}
\newcommand{\ab}{\alpha_b}
\newcommand{\bb}{\beta_b}
\newcommand{\wt}{\varpi_p}
\newcommand{\Fnl}{\mathcal F}
\newcommand{\Feff}{F_{\mathrm{eff}}}
\newcommand{\Qmat}{\mathcal Q_{\mathrm{mat}}}
\newcommand{\AN}{\mathsf T_{\mathrm{Newt}}}
\DeclareMathOperator{\spec}{spec}
\DeclareMathOperator{\supp}{supp}
\DeclareMathOperator{\spanop}{span}
\DeclareMathOperator{\rank}{rank}
\title{Convex non-ball Schiffer domains in infinitely many dimensions}
\author{Jizhou Guo}
\address{Dots Studio, Rednote}
\email{mitsuha2021b@gmail.com, sjtu18640985163@sjtu.edu.cn}
\date{October 8, 2026}
\begin{document}
\begin{abstract}
We construct bounded strictly convex non-ball Schiffer domains with
real-analytic boundary in two arithmetic families of dimensions $2p$,
with $p$ integral or half-integral. The bounded-detuning families use
a cofactor inverse and include infinitely many dimensions in every
residue class, with a detuning bound depending on the modulus. A second family
contains every sufficiently large half-lattice input within
$\log p/(8p)$ of a primary crossing and has unbounded detuning in both
parities. Its mixed conditioning estimate closes finite-order centres
at accuracy twenty. In each family all but a sublinear number of
interior integer block splits yield pairwise non-similar domains.
The constructions use one graded Jacobi coefficient space and a quartic
angular seed. The thresholds and splits are implicit. Each fixed-width
candidate family and the logarithmic family have density zero.
Green's identity gives failure
of the Pompeiu property for the resulting domains.
\end{abstract}
\maketitle
\tableofcontents
\section{Results and scope}\label{inf:sec-introduction}

In the smooth sphere-boundary formulation, the Schiffer conjecture
states that a bounded domain $\Omega\subset\R^n$ with boundary
homeomorphic to $S^{n-1}$ must be a ball if it supports a nonconstant
solution of $\Delta u+\lambda u=0$, for some $\lambda>0$, with
constant boundary value and zero normal derivative
\cite{inf:Wil76,inf:Zal92}. We normalize the frequency to one and
the boundary value to one. We construct domains with these
overdetermined data in unbounded families of both odd and even ambient
dimensions, while retaining strict convexity.

Write $J_\nu$ for the Bessel function of the first kind of order $\nu$.
The real primary crossing orders $\widehat p_m$, for all sufficiently
large indices $m\ge m_0$, satisfy
\[
 J_{\widehat p_m}(2\sqrt{(\widehat p_m+1)(\widehat p_m+2)})=0.
\]
They are specified uniquely by the expansion in
Proposition~\ref{inf:primary-expansion}. The arithmetic construction
uses the fixed irrational number
\[
 \alpha_{\mathrm{ar}}=\frac{\pi}{2(\sqrt3-\pi/3)}
\]
and continued-fraction approximations with denominator one modulo
four. Define $\mathcal P_{\mathrm{bd}}$ to consist of all positive
$p\in\tfrac12\Z$ with $|p-\widehat p_m|<31/(10\widehat p_m)$
for some $m\ge m_0$. The detuning is \eqref{inf:detuning-definition}
at the Neumann zero continued from the nearby crossing.
For every constructed domain the frequency-one
Schiffer data mean a nonconstant classical solution of
\begin{equation}\label{inf:main-equations}
 \Delta u+u=0\quad\hbox{in }\Omega,\qquad
 u=1,\quad\nabla u=0\quad\hbox{on }\partial\Omega.
\end{equation}
\begin{maintheorem}[Bounded detuning]\label{inf:main}\label{inf:theorem-A}
For every sufficiently large $p\in\mathcal P_{\mathrm{bd}}$, there
is an integer $a\in[p/2,3p/4]$ and a bounded strictly convex non-ball
domain $\Omega\subset\R^{2p}$ with real-analytic boundary,
invariant under $O(a)\times O(2p-a)$ and carrying
\eqref{inf:main-equations}. Its detuning satisfies $0<|\theta|\le20$.
There are infinitely many such inputs in each of $\Z$ and $\Z+1/2$.

For every fixed positive integer $L_{\mathrm{ap}}$ and integer
$b_{\mathrm{ap}}$, there is also an unbounded selection
$n\equiv b_{\mathrm{ap}}\pmod{4L_{\mathrm{ap}}}$ with
$|n/2-\widehat p_m|<32L_{\mathrm{ap}}^2/\widehat p_m$ and
$0<|\theta|<256L_{\mathrm{ap}}^2$, whose sufficiently large members
have the same domain conclusion. The threshold may depend on the
fixed modulus and residue. Hence every residue modulo
$L_{\mathrm{ap}}$ contains infinitely many such dimensions.
For even $L_{\mathrm{ap}}$ a residue modulo $L_{\mathrm{ap}}$ permits
only its compatible parity; for odd $L_{\mathrm{ap}}$ each residue
modulo $L_{\mathrm{ap}}$ contains infinitely many dimensions of both parities.

In each of these bounded families, fix
$0<\varepsilon_{\mathrm{sp}}<1/4$ and $0<\eta_{\mathrm{cnt}}<1/3$.
For every sufficiently large member $n=2p$, all but
$O(n^{2/3+\eta_{\mathrm{cnt}}})$ integer splits
$\varepsilon_{\mathrm{sp}}n\le a\le(1/2-\varepsilon_{\mathrm{sp}})n$
give such domains, pairwise non-similar for distinct retained splits.
The constant and threshold may depend on the family and these fixed
parameters. Let $\mathcal N_{\mathrm{sim}}(n)$ be the number of
integer splits $0<a<n/2$ for which the construction on some fixed
compact interior interval gives a domain. Choosing one domain at
each such split gives this many similarity classes, independently
of the choices.
Then along each family's dimensions
\[
 \liminf_{n\to\infty}\mathcal N_{\mathrm{sim}}(n)/n\ge1/2.
\]
The closure uses the cofactor estimate and a fixed centre order
$N=4M+40$, with $M$ fixed after the bounded family and compact
split interval have been chosen and before $N$.
\end{maintheorem}

\begin{maintheorem}[Logarithmic window]\label{inf:theorem-B}
Let
\[
 \mathcal P_{\log}=\{p\in\tfrac12\Z:p>1,\quad
 \exists m\ge m_0:\ |p-\widehat p_m|\le\log p/(8p)\}.
\]
For every sufficiently large $p\in\mathcal P_{\log}$, with $n=2p$, there is an
integer split $0<a<n/2$ and a bounded strictly convex non-ball domain
$\Omega\subset\R^{2p}$ with real-analytic boundary, invariant under
$O(a)\times O(2p-a)$ and carrying \eqref{inf:main-equations}.
For every fixed $0<\varepsilon_{\mathrm{sp}}<1/4$ and
$0<\eta_{\mathrm{cnt}}<1/3$, all but
$O_{\mathrm{sp},\eta_{\mathrm{cnt}}}(n^{2/3+\eta_{\mathrm{cnt}}})$
interior integer splits in the interval of Theorem~\ref{inf:theorem-A}
give pairwise non-similar domains. The same similarity-class lower
limit is at least $1/2$. This family contains sequences of both
parities with $|\theta|\to\infty$. The closure uses mixed
conditioning and the fixed centre order $N=20$.
\end{maintheorem}

\begin{proof}[Proofs of Theorems~\ref{inf:theorem-A} and \ref{inf:theorem-B}]
Corollary~\ref{inf:bounded-basic} proves the first bounded construction.
Lemmas~\ref{inf:many-splits} and \ref{inf:similarity-classes},
Proposition~\ref{inf:bounded-assembly}, and
Theorem~\ref{inf:residue-arithmetic} give all the further assertions
of Theorem~\ref{inf:theorem-A} through the cofactor route.
Proposition~\ref{inf:log-assembly} gives the logarithmic construction
and multiplicity, and Lemma~\ref{inf:log-multiplier} gives unbounded
detuning in both parities. These prove Theorem~\ref{inf:theorem-B}.
\end{proof}

Both assertions concern sufficiently large inputs, with implicit
thresholds and splits. Each fixed-width family and the logarithmic
family have density zero
by Proposition~\ref{inf:arithmetic-density}. The quartic profile and
its Hermite limit are described in Section~\ref{inf:sec-shape}.

The Pompeiu property means that a continuous function on $\R^n$
whose integral over every rigid copy of a domain is zero must
itself be zero. The Pompeiu conjecture states that, among bounded
domains with Lipschitz boundary homeomorphic to $S^{n-1}$,
failure of this property forces the domain to be a ball
\cite{inf:Wil76,inf:Zal92}. The domains in both theorems fail this property
because their indicator transforms vanish on the unit sphere.
A single nonzero plane wave, or its real part, then gives the
required nonzero function. This implication is proved directly
by Green's identity in Proposition~\ref{inf:Pompeiu}.
Consequently the Schiffer and Pompeiu conjectures fail within
the class of convex domains in infinitely many dimensions.

Brown, Schreiber and Taylor developed the spectral-synthesis
approach to the Pompeiu problem and its Fourier--Laplace
characterization \cite{inf:BST73}. Williams related failure of the
Pompeiu property on simply connected Lipschitz domains to the
overdetermined eigenvalue problem \cite{inf:Wil76};
Zalcman's survey gives the classical
bibliography \cite{inf:Zal92}. Recent planar counterexamples were
constructed by Cao-Labora and de Dios Pont \cite{inf:CLdDP26}
and by Colbrook and Stepaniants \cite{inf:CS26}.
Convex counterexamples in the plane, in every dimension from three
to eighteen, and in dimensions twenty and twenty-one were obtained
by computer-assisted proofs in \cite{inf:Guo26}.
The construction below uses a two-block split parameter at fixed
ambient dimension and closes finite-order centres in a graded
Jacobi coefficient space.

\subsection{The construction and its estimates}

The block split supplies a parameter while the ambient dimension
and the ball radial spectrum remain fixed. Write an ambient point as
$(X,Y)\in\R^a\times\R^{2p-a}$. With $n=2p$ and
$r=a/(2p)$, invariant functions depend on radius and
$x=|X|^2/(|X|^2+|Y|^2)$. Their angular basis is Jacobi orthogonal
for the beta measure of parameters $pr,p(1-r)$.
The monic quadratic $\mathsf H_2(x)$ has spherical degree four and is the
shape used to start the construction.
At a continued radial Neumann zero $J_p(R)=0$, define
\[
 \chi=\frac{J_{p+3}(R)}{J'_{p+3}(R)-(p-1)J_{p+3}(R)/R},
 \qquad \theta=4Rp\chi.
\]
For all sufficiently large members of the bounded base family,
$R=2p+3+O(p^{-1})$ and $0<|\theta|<20$.
The estimates are uniform for $0<|\theta|\le20$.

The normalized quartic radius is
$1-\theta \mathsf H_2/(2R^2pg_2)$, where
$g_2=[\mathsf H_2]\mathsf H_2^2\asymp p^{-1}$.
Vary the real split through $\varsigma=(r-1/2)^2\in[1/64,1/16]$.
The exact weighted Dirichlet eigenvalue curves have signed speed
at least $c|\theta|/p$ when their offsets from $R^2$ lie in
$(-|\theta|/(16p),|\theta|/(16p))$.
Only $O(p^{2/3})$ ordered curves can enter this band.
Avoiding their small intervals on the integer split grid gives an
actual invariant Dirichlet gap $c|\theta|p^{-5/3}$ for at least
one admissible $a$.

For each fixed accuracy order $N$, finite-order centre construction
at that split solves the interior
Helmholtz equation exactly and makes both boundary defects
$O_N(\theta^2p^{-N})$ on every prescribed finite complex disc,
with the constant and threshold allowed to depend on that disc.
The recurrence uses only fixed nonzero leading divisors, including
the seed divisor $-1/8$ after its natural scaling.
The centre differs from the quartic radius by
$O_N(|\theta|p^{-4})$ in normalized coordinates, so it inherits
half the Dirichlet gap by domain monotonicity.

The nonlinear argument uses one fixed coefficient space
$\Cc_{p,2}$ for sources and a boundary space $\Ec_1$ with one
angular degree moment. The clamped source subspace $\Gc_p$
contains exactly the sources whose Poisson fields have zero
Dirichlet and normal traces. Positive Jacobi multiplication and a
parameter-raising coupling give a high-shape-degree gain on these
clamped fields. This makes the scalar pullback analytic on the
fixed small-shape chart in $\Gc_p\times\Ec_1$, with a
second-derivative bound $C_Np^6$ on the stated centre neighbourhood.
An exact material change of variables has inverse bounded by
$C_Np^6$ in those same spaces.

The remaining inverse estimate passes through
$\Mmix$, which has radial $L^2$ norms and angular $\ell^1$
coefficients. In a ball window of width $E_{\mathrm{inv}}|\theta|$,
the exact Schur matrix has entries of size $O(|\theta|)$,
exponential angular decay in two matrix norms, and the actual
Hilbert gap. Its logarithmic angular clusters have bounded
\emph{cardinality} $M$, independent of $p$ and of the separation
constant. A determinant and cofactor estimate on each cluster
therefore gives
\[
 \nrm{F_{\mathrm{eff}}^{-1}}_{\ell^1(2^j)}
 \le C|\theta|^{-1}p^{5M/3}.
\]
Radial graph reconstruction and a finite-head/infinite-tail
argument give the full coefficient inverse for the scalar Dirichlet
graph operator $J_D=B_{\Phi_N,D}\KD$, with $B_{\Phi,D}$ defined after
\eqref{inf:scalar-pullback-definition},
\begin{equation}\label{inf:intro-inverse}
 \nrm{J_D^{-1}}_{\Cc_{p,2}\to\Cc_{p,2}}
 \le C|\theta|^{-1}p^{5/2+5M/3}.
\end{equation}

The constants determining $M$ and the inverse exponent are fixed
before the centre accuracy is chosen. With
$N=4M+40$, the residual, corrected derivative and second derivative
satisfy the three gates of Theorem~\ref{inf:newton-threshold} on a
ball of radius $|\theta|p^{-2M-62/3}$.
The exact zero remains strictly convex and keeps a nonzero
physical monic seed coefficient. Its coefficient equation is
identified with a weak equation on the completed spaces, then
with a classical solution by elliptic regularity.

The logical order of the main steps is
\[
 \begin{gathered}
 \text{uniform Bessel expansion and congruence approximation}
 \ \Longrightarrow\ \mathcal P_{\mathrm{bd}},\ R,\ \theta,\\
 \text{angular positivity and exact local spectral flow}
 \ \Longrightarrow\ \text{integer-split gap},\\
 \text{finite centres and exact clamping}
 \ \Longrightarrow\ \text{residuals and material map},\\
 \text{relative graph bounds and bounded-cardinality clusters}
 \ \Longrightarrow\ \text{coefficient inverse},\\
 \text{Newton threshold}
 \ \Longrightarrow\ \text{classical convex non-ball domain}.
 \end{gathered}
\]
For Theorem~\ref{inf:theorem-B}, one radial pole per angular degree
is retained first. A finite-shift Lommel recurrence isolates four
critical phase ratios. Their short components have subpolynomial
coefficient-to-Hilbert conditioning, while the complementary
couplings have two inverse powers of $p$, up to a chosen small
loss. The exact centre is restored in the mixed norm before the
single radial conversion factor $\sqrt p$.
Section~\ref{inf:sec-sharp} gives material and nonlinear exponents
two and three; Section~\ref{inf:sec-tame} records their dependence
on $\tau\le\log p$. The three gates then close at order twenty
and radius $\tau p^{-11}$ as in
Theorem~\ref{inf:fixed-order-newton}.

The notation index in Appendix~\ref{inf:sec-notation} distinguishes
the fixed flow window, the local-count window, and the shrinking
inverse window.

\subsection{Classical inputs}

The proof uses the following classical results; the specialized
estimates needed in the construction
are proved in the body of the paper.
\begin{enumerate}[label=\textup{(\arabic*)}]
\item The integer-order common-zero theorem of Bourget, as a
consequence of Siegel's theorem, excludes integer primary
crossings \cite[\S15.28]{inf:watson}; see also
\cite{inf:siegel}. At half-integer orders the sine--cosine
representation and Hermite--Lindemann theorem
\cite[\S1.2]{inf:BakerWustholz} give non-coincidence as proved in
Lemma~\ref{inf:half-integer-noncoincidence}.
The finite Lommel polynomials \cite{inf:watson} give the
shift resonance estimate in Lemma~\ref{inf:lommel-resonance}.
The transcendence of $\pi$
\cite{inf:lindemann} makes $\alpha_{\mathrm{ar}}$ irrational.
\item Bessel power series and recurrences, the fixed-order outgoing Hankel
normalization and Watson's integral for the derivative of a
positive zero with respect to order are used in their classical
forms \cite[\S\S10.2, 10.6, 10.17; Eq.~10.21.17]{inf:dlmf}.
The uniform large-order remainders, derivative bounds, spacing,
and radial normalization used here are then proved in
Section~\ref{inf:sec-bessel}.
\item Gasper's positive linearization theorem is used with
\eqref{inf:gasper-hypothesis} and normalization at the endpoint of
the larger Jacobi exponent \cite[Theorem~1]{inf:gasper}.
The parameter substitution is given before
Lemma~\ref{inf:radial-stencils}.
\item Closed Dirichlet forms, relatively bounded perturbations,
compact resolvents, the min--max principle and finite Riesz
reductions are used as in \cite[Chs.~II and IV--VII]{inf:kato}.
The common domains, norm differentiability, eigenvalue derivative
argument and Schur reconstructions for this family are written
out in Section~\ref{inf:sec-flow}.
\item Dirichlet $H^2$ and higher regularity and Sobolev traces are
used on smooth bounded domains \cite[Chs.~6--9]{inf:GT}; analytic
boundary regularity is used for analytic boundaries and constant
Dirichlet data \cite{inf:MorreyNirenberg,inf:MorreyBoundary}.
The convex Dirichlet Hessian inequality is proved by its boundary
identity in \eqref{inf:convex-Hessian}. Polynomial core density
and the zero-capacity cutoff statements at the axes and origin
are included in Proposition~\ref{inf:flow-form}.
\item Hadamard's convexity theorem for a compact hypersurface
with positive second fundamental form is used in
Lemma~\ref{inf:geometry-convexity}; see
\cite{inf:doCarmoLima,inf:sacksteder}.
Bernstein's inequality for trigonometric polynomials, applied on
great circles, gives the dimension-independent shape derivative
bounds \cite[Vol.~II, Ch.~X, Theorem~3.13]{inf:zygmund}. The Jacobi endpoint bound itself is
proved in Lemma~\ref{inf:angular-sup}.
The rotation argument in Proposition~\ref{inf:arithmetic-density}
uses trigonometric polynomial approximation \cite{inf:zygmund}.
\item Banach's contraction theorem \cite{inf:banach} is applied to
a closed ball in the stated complete spaces.
Taylor's formula, continued-fraction determinant identities and
elementary lattice determinant facts are used in the explicit
proofs of their specialized forms.
\item The analytic implicit function theorem is used for the
scaled radius and seed equations in Theorem~\ref{inf:centres}.
Cauchy's estimates and the Hermite--Genocchi integral formula for
divided differences give the analytic coefficient bounds in
Lemma~\ref{inf:analytic-coefficients}; Cauchy's estimates also
give the centre remainder bounds.
The Dirichlet--Poincar\'e (Wirtinger) inequality on an interval is
used in Lemma~\ref{inf:bessel-spacing} and
Theorem~\ref{inf:local-count}.
\end{enumerate}

\section{Two-block reduction and normalizations}\label{inf:sec-reduction}

Fix $p\ge64$ with $2p\in\Z$ and an integer $a$ with
$p/2\le a\le3p/4$. Write
\begin{equation}\label{inf:block-parameters}
 n=2p,\qquad b=2p-a,\qquad
 \ab=\frac a2=pr,\qquad \bb=\frac b2=p(1-r),
 \qquad \frac14\le r\le\frac38.
\end{equation}
For the spectral deformation argument the parameters $\ab,\bb$ will also
be allowed to be positive real numbers with sum $p$. All constructions
that define Euclidean domains are made at integer block sizes.
In particular $a,b\ge3$ and $\ab,\bb>1$ in this range.
We use $y=(X,Y)\in\R^a\times\R^b$ and, away from the origin, set
\begin{equation}\label{inf:coordinates}
 \varrho=|y|,\qquad t=|y|^2,\qquad
 x=\frac{|X|^2}{|X|^2+|Y|^2}.
\end{equation}
Here $t$ always denotes squared radius. A nonsquared reference radial
coordinate used later will be denoted by $\zeta$.

\begin{lemma}[Invariant differential operator]\label{inf:reduction}
For a twice differentiable invariant function $u(\varrho,x)$,
\begin{equation}\label{inf:invariant-laplacian}
 \Delta u=u_{\varrho\varrho}+\frac{2p-1}{\varrho}u_\varrho
 +\frac4{\varrho^2}\left\{x(1-x)u_{xx}+(\ab-px)u_x\right\}.
\end{equation}
Normalized spherical measure induces on $x$ the probability measure
\begin{equation}\label{inf:beta-measure}
 d\mu_r(x)=\frac{x^{\ab-1}(1-x)^{\bb-1}}{B(\ab,\bb)}\,dx.
\end{equation}
Normalized invariant volume measure on the unit ball is
$p t^{p-1}\,dt\,d\mu_r(x)$. Each angular polynomial occurs with
multiplicity one in this invariant sector.
\end{lemma}
\begin{proof}
Put $u(y)=v(q_X,q_Y)$, where $q_X=|X|^2$ and $q_Y=|Y|^2$.
Differentiating in each Euclidean coordinate and summing gives
\[
 \Delta u=4q_Xv_{q_Xq_X}+2av_{q_X}
          +4q_Yv_{q_Yq_Y}+2bv_{q_Y}.
\]
Use $q_X=\varrho^2x$ and $q_Y=\varrho^2(1-x)$ in this identity.
The mixed second derivatives in $\varrho,x$ cancel, and the remaining
coefficients are those in \eqref{inf:invariant-laplacian}.
Alternatively, the coefficients of $u_{xx}$ and $u_x$ follow directly
from $|\nabla_S x|^2=4x(1-x)$ and
$\Delta_S x=4(\ab-px)$.

Polar integration in the two blocks has volume element proportional to
$q_X^{\ab-1}q_Y^{\bb-1}\,dq_X\,dq_Y$. The determinant of the change
$(t,x)\mapsto(tx,t(1-x))$ has absolute value $t$. Its angular integral
is $B(\ab,\bb)$; its radial integral over $0<t<1$ is $1/p$.
Division by these two integrals proves the measure assertions.
An invariant function is a function of $\varrho,x$, so each angular
polynomial below occurs once.
\end{proof}

\subsection{The angular polynomials}

Let $\mathsf H_j(x)$ be the monic degree-$j$ orthogonal polynomial for $\mu_r$,
and define
\begin{equation}\label{inf:angular-normalization}
 G_j(x)=\frac{\mathsf H_j(x)}{\mathsf H_j(1)},\qquad
 \eta_j^2=\int_0^1G_j(x)^2\,d\mu_r(x),\qquad
 P_j(x)=\eta_j^{-1}G_j(x).
\end{equation}
The sign of $P_j$ is chosen to give a positive leading coefficient.
The two uses of Jacobi polynomials in this paper are distinguished by
letters: $G_j$ is angular, whereas $J_s^{\beta}$ below is radial.
The symbol $J_\nu$ without an upper index denotes the Bessel function.

\begin{lemma}[Angular formulas]\label{inf:angular-formulas}
For $j\ge0$,
\begin{align}
 G_j(x)&=\frac{P_j^{(\bb-1,\ab-1)}(2x-1)}{(\bb)_j/j!},
 &\mathsf H_j(1)&=\frac{(\bb)_j}{(p+j-1)_j},\label{inf:angular-jacobi}\\
 \eta_j^2&=\frac{(p-1)(\ab)_j j!}
 {(2j+p-1)(p-1)_j(\bb)_j},
 &\Lambda_rG_j&=\lambda_jG_j,\label{inf:angular-norm}\\
 \Lambda_r&=-4x(1-x)\partial_x^2-4(\ab-px)\partial_x,
 &\lambda_j&=4j(j+p-1).\nonumber
\end{align}
The convention for a rising factorial of length zero is $(z)_0=1$.
In the orthonormal basis multiplication by $x$ is
\begin{equation}\label{inf:jacobi-matrix}
 xP_j=\mathfrak a_{j+1}P_{j+1}+\mathfrak b_jP_j
       +\mathfrak a_jP_{j-1},
\end{equation}
where $\mathfrak a_0=0$ and
\begin{align}
 \mathfrak a_j^2&=
 \frac{j(j+\ab-1)(j+\bb-1)(j+p-2)}
 {(2j+p-2)^2(2j+p-1)(2j+p-3)},\quad j\ge1,
 \label{inf:recurrence-a}\\
 \mathfrak b_j&=\frac12+
 \frac{(r-1/2)p(p-2)}{(2j+p-2)(2j+p)}.
 \label{inf:recurrence-b}
\end{align}
All these identities hold also for real positive $\ab,\bb$ with $p\ge64$.
\end{lemma}
\begin{proof}
The shifted Rodrigues identity is
\[
 x^{\ab-1}(1-x)^{\bb-1}P_j^{(\bb-1,\ab-1)}(2x-1)
 =\frac{(-1)^j}{j!}\frac{d^j}{dx^j}
       \{x^{\ab+j-1}(1-x)^{\bb+j-1}\}.
\]
Repeated integration by parts against a polynomial of degree less than
$j$ proves orthogonality. The endpoint value is $(\bb)_j/j!$ and the
leading coefficient is $(p+j-1)_j/j!$. Their quotient proves the
second identity in \eqref{inf:angular-jacobi}. The integrations by parts
are valid: at each intermediate stage the endpoint powers in the
boundary terms have positive exponent. Differentiation of the Rodrigues
identity gives the displayed angular differential equation.

A further integration by parts against the degree-$j$ polynomial gives
\[
 \int_0^1\bigl(P_j^{(\bb-1,\ab-1)}(2x-1)\bigr)^2\,d\mu_r
 =\frac{(p-1)(\ab)_j(\bb)_j}
        {(2j+p-1)(p-1)_j j!}.
\]
Dividing by the squared endpoint value proves
\eqref{inf:angular-norm}. Orthogonality makes multiplication by $x$
tridiagonal and symmetric. The ratio of the leading coefficients of
$P_{j-1}$ and $P_j$, followed by the norm formula, gives
\eqref{inf:recurrence-a}. Comparing the coefficients of $x^j$ in the
monic recurrence gives \eqref{inf:recurrence-b}. These computations
hold for these positive real parameters.
\end{proof}

\begin{lemma}[Endpoint bound and angular derivatives]\label{inf:angular-sup}
For the parameters in \eqref{inf:block-parameters},
\[
 |G_j(x)|\le1\quad(0\le x\le1),\qquad 0<\eta_j\le1.
\]
On the unit sphere the homogeneous extension
$S_j(y)=|y|^{2j}G_j(x)$ satisfies
\begin{equation}\label{inf:angular-derivatives}
 \nrm{G_j}_{L^\infty(S^{2p-1})}\le1,\qquad
 \nrm{\nabla_S G_j}_\infty\le2j,\qquad
 \nrm{\nabla_S^2G_j}_{\mathrm{op},\infty}\le4j^2.
\end{equation}
The constants are independent of $p$ and of the admissible split.
\end{lemma}
\begin{proof}
For a Jacobi polynomial $v(z)=P_j^{(\alpha,\beta)}(z)$ with
$\alpha=\bb-1\ge\beta=\ab-1\ge0$ and $j>0$, write
$\kappa=j(j+\alpha+\beta+1)$. Its equation implies
\[
 \frac d{dz}\left(v^2+\frac{1-z^2}{\kappa}(v')^2\right)
 =\frac2\kappa\{\alpha-\beta+(\alpha+\beta+1)z\}(v')^2.
\]
The factor in braces is increasing. The expression in parentheses
therefore first decreases and then increases, with either interval
possibly empty. Its maximum is attained at an endpoint. The absolute
endpoint values are $(\alpha+1)_j/j!$ and $(\beta+1)_j/j!$; the former
is at least the latter. Division by the former proves the supremum
bound. The case $j=0$ is constant. The bound for $\eta_j$ follows by
integration against a probability measure.

The restriction to any great circle of the polynomial $S_j$ is a
trigonometric polynomial of degree at most $2j$. The first and second
Bernstein inequalities give bounds $2j$ and $4j^2$ for its derivatives
along that circle \cite[Vol.~II, Ch.~X, Theorem~3.13]{inf:zygmund}.
Great circles with prescribed unit tangent vector
compute the spherical gradient and the diagonal values of the
spherical Hessian. A real symmetric bilinear form has operator norm
equal to the supremum of its absolute diagonal values on unit vectors.
This proves \eqref{inf:angular-derivatives}.
\end{proof}

\begin{lemma}[Solid harmonics and radial profiles]\label{inf:solid-harmonics}
The function
\begin{equation}\label{inf:solid-definition}
 S_j(y)=t^jG_j(x)
\end{equation}
is an invariant homogeneous harmonic polynomial of degree $2j$.
For a twice differentiable radial profile $v$,
\begin{equation}\label{inf:profile-laplacian}
 \Delta\{S_jv(t)\}=4S_j\{tv''+(2j+p)v'\}.
\end{equation}
The regular solutions of $(\Delta+1)u=0$ in angular degree $2j$ are
multiples of
\begin{equation}\label{inf:regular-bessel-mode}
 \varrho^{1-p}J_{\nu_j}(\varrho)G_j(x),\qquad
 \nu_j=p-1+2j.
\end{equation}
Their apparent singularities at the origin are removable.
\end{lemma}
\begin{proof}
As $G_j$ is a polynomial of degree $j$, multiplication by $t^j$
expresses $S_j$ as a polynomial in $q_X,q_Y$. Applying
\eqref{inf:invariant-laplacian} and using its angular eigenvalue proves
harmonicity. For the product with $v(t)$, use
$\nabla t=2y$, $\Delta t=4p$, and $y\cdot\nabla S_j=2jS_j$.
The product rule gives \eqref{inf:profile-laplacian}.
Substituting $\varrho^{1-p}w(\varrho)$ in the radial Helmholtz equation
gives Bessel's equation of order $\nu_j$. Its regular solution is
$J_{\nu_j}$. Its convergent expansion is a constant times
$\varrho^{2j}$ multiplied by a power series in $\varrho^2$.
The factor $\varrho^{2j}G_j$ is the polynomial just constructed, which
proves regularity at the origin and across both block axes.
\end{proof}

\subsection{The ball field and the invariant Dirichlet spectrum}

Write $j_{\nu,k}$ for the $k$th positive zero of $J_\nu$ and set
$z_{j,k}=j_{\nu_j,k}$. All radial zero indices $k$ start at $1$.
Choose a positive zero $R=j_{p,k_0}$, where $k_0$ will be specified by
the arithmetic construction. The radial field with boundary value one is
\begin{equation}\label{inf:ball-field}
 u_0(\varrho)=
 \frac{\varrho^{1-p}J_{p-1}(\varrho)}{R^{1-p}J_{p-1}(R)}.
\end{equation}
Consecutive orders have no common positive zero: their recurrences
would give both the value and derivative of one solution equal to zero,
contradicting uniqueness for the Bessel equation. Thus the denominator
in \eqref{inf:ball-field} is nonzero.
The derivative identity
$(\varrho^{1-p}J_{p-1}(\varrho))'
=-\varrho^{1-p}J_p(\varrho)$ proves
\begin{equation}\label{inf:ball-traces}
 u_0(R)=1,\qquad u_0'(R)=0,\qquad u_0''(R)=-1.
\end{equation}
The last equality follows from the radial equation.

\begin{proposition}[Invariant Dirichlet realization]\label{inf:ball-spectrum}
The Friedrichs Dirichlet Laplacian $H_0=-\Delta_D$ in the invariant
subspace of the unit ball has eigenvalues $z_{j,k}^2$, one radial
sequence for each $j\ge0$, with repetitions retained. It has compact
resolvent. After the unitary radial substitution
$v(\zeta)=\sqrt{2p}\,\zeta^{p-1/2}u(\zeta)$ and the angular basis $P_j$, its
$j$th radial operator is
\begin{equation}\label{inf:radial-schrodinger}
 -\frac{d^2}{d\zeta^2}+\frac{\nu_j^2-1/4}{\zeta^2},
 \qquad 0<\zeta<1,
\end{equation}
with Friedrichs condition at zero and Dirichlet condition at one.
The spectrum does not depend on the split $r$.
\end{proposition}
\begin{proof}
The orthogonal group action commutes with the Dirichlet form. Averaging
over it is an orthogonal projection preserving $H^1_0$, so restriction
to the invariant subspace gives a closed Dirichlet form. Compactness of
$H^1_0(B)\hookrightarrow L^2(B)$ implies compact resolvent for this
restriction. The polynomials $P_j$ form a complete angular orthonormal
basis; this follows from density of polynomials in the continuous
functions on $[0,1]$ and density of continuous functions in $L^2(\mu_r)$.
Separation of variables gives \eqref{inf:radial-schrodinger}. Its
regular solution at energy $z^2$ is a multiple of
$\sqrt\zeta J_{\nu_j}(z\zeta)$. The other independent solution is not
in the Friedrichs form domain, as $\nu_j\ge p-1>1$.
The boundary condition at one is $J_{\nu_j}(z)=0$.
Neither $\nu_j$ nor the separated eigenvalue equation contains $r$.
\end{proof}

\begin{lemma}[Radial polynomial orthogonality]\label{inf:radial-polynomials}
Set $\beta_j=2j+p-1$ and
\begin{equation}\label{inf:radial-basis}
 J_s^\beta(t)=P_s^{(0,\beta)}(2t-1),\qquad
 \Psi_{j,s}(y)=S_j(y)J_s^{\beta_j}(t),\quad s\ge0.
\end{equation}
Then $J_s^\beta(1)=1$ and
\begin{align}
 J_s^\beta(t)&=\sum_{h=0}^s(-1)^{s-h}\binom sh
                         \binom{s+\beta+h}{s}t^h,
                         \label{inf:radial-expansion}\\
 \int_0^1t^\beta J_s^\beta J_v^\beta\,dt
 &=\frac{\delta_{sv}}{2s+\beta+1},\qquad
 \nrm{\Psi_{j,s}}_{L^2(B)}^2
 =\eta_j^2\frac p{p+2j+2s}.
 \label{inf:radial-norm}
\end{align}
Here volume is normalized as in Lemma~\ref{inf:reduction}. All invariant
polynomials are finite linear combinations of these basis elements.
\end{lemma}
\begin{proof}
Apply the Rodrigues computation of Lemma~\ref{inf:angular-formulas}
with Jacobi parameters $0,\beta$. It gives the coefficient formula,
endpoint value, and integral norm. Multiplication by $S_j=t^jG_j$
changes the radial weight $p t^{p-1}$ to $p t^{\beta_j}$; angular
orthogonality then proves the last identity.
A monomial in $q_X,q_Y$ is $t^d$ times a polynomial in $x$ of degree
at most $d$. Expand that polynomial in the $G_j$, $j\le d$, and expand
the remaining power $t^{d-j}$ in the $J_s^{\beta_j}$. This proves the
spanning assertion.
\end{proof}

\section{Bessel estimates with uniform remainders}\label{inf:sec-bessel}

The large-order estimates needed here are uniform as the ratio of
argument to order varies. We derive them from an outgoing integral
equation. The only asymptotic normalization used in that derivation is
the fixed-order, large-argument normalization of the Hankel function.
The recurrence identities and that normalization are as in
\cite[\S\S10.6, 10.17]{inf:dlmf}; the order derivative formula used
below is \cite[Eq.~10.21.17]{inf:dlmf}.

We write $Y_\nu$ for the Bessel function of the second kind and
$H^{(1)}_\nu=J_\nu+iY_\nu$ for the outgoing Hankel function.
The modified Bessel functions of the second kind are denoted by
$\mathrm K_\nu$.

\subsection{Zeros, radial norms, and differentiation in the order}

\begin{lemma}[Spacing and boundary normalization]\label{inf:bessel-spacing}
Let $\nu>1/2$. Every positive zero of $J_\nu$ is simple,
$j_{\nu,k}>\nu$, and
\begin{equation}\label{inf:bessel-spacing-bound}
 j_{\nu,k+1}-j_{\nu,k}>\pi.
\end{equation}
At a zero $z=j_{\nu,k}$ one has
\begin{equation}\label{inf:lommel}
 \int_0^1\zeta J_\nu(z\zeta)^2\,d\zeta
 =\frac12J_\nu'(z)^2.
\end{equation}
Consequently a radial Dirichlet eigenfunction, normalized in its radial
$L^2$ space and with its sign chosen appropriately, has boundary
normal derivative $\sqrt2\,z$.
\end{lemma}
\begin{proof}
A simultaneous zero of $J_\nu$ and $J'_\nu$ would give the zero solution
of its nonsingular ordinary differential equation at that point.
Let $v(x)=\sqrt{x}J_\nu(x)$. On an interval between two positive zeros,
\[
 \int |v'|^2=\int\left(1-\frac{\nu^2-1/4}{x^2}\right)|v|^2
 <\int|v|^2.
\]
The Dirichlet Poincar\'e inequality on an interval of length $L$ gives
$\pi^2/L^2<1$, proving \eqref{inf:bessel-spacing-bound}.
On $(0,1)$ the lowest eigenvalue of
$-\partial_\zeta^2+(\nu^2-1/4)/\zeta^2$ is at least
$\pi^2+\nu^2-1/4>\nu^2$, by the same Poincar\'e inequality and
$\zeta^{-2}\ge1$. Every eigenvalue is at least this lower bound, so
$j_{\nu,k}>\nu$.

For the integral identity put $w(\zeta,z)=J_\nu(z\zeta)$.
Subtract the equation for $w$ multiplied by $\partial_z w$ from its
$z$ derivative multiplied by $w$. Integration gives
\[
 [\zeta\{w(\partial_z w)_\zeta-w_\zeta\partial_z w\}]_0^1
 =-2z\int_0^1\zeta w^2\,d\zeta.
\]
At a zero $z$, $w(1,z)=0$, $w_\zeta(1,z)=zJ'_\nu(z)$ and
$\partial_z w(1,z)=J'_\nu(z)$. The term at zero vanishes by the
regular Bessel series. This proves \eqref{inf:lommel}. The radial
mass-unitary function is $\sqrt\zeta J_\nu(z\zeta)$, whose derivative
at one is $zJ'_\nu(z)$ and whose squared norm is the integral in
\eqref{inf:lommel}. Their ratio proves the last assertion.
\end{proof}

\begin{lemma}[Order derivative]\label{inf:bessel-order}
For $\nu\ge1/2$ and $x=j_{\nu,k}$,
\begin{equation}\label{inf:order-derivative}
 0\le \frac{\arccos(\nu/x)}{\sqrt{1-(\nu/x)^2}}
       -\frac{\partial j_{\nu,k}}{\partial\nu}
 \le\frac{\pi}{16x^2}.
\end{equation}
In particular,
\begin{equation}\label{inf:order-lipschitz}
 1-\frac\pi{16x^2}\le\frac{\partial j_{\nu,k}}{\partial\nu}
 \le\frac\pi2.
\end{equation}
The estimates are independent of $k$. If $x/\nu\to2$ and
$x\to\infty$, then
$\partial_\nu j_{\nu,k}\to2\pi/(3\sqrt3)$, also when $k$ varies.
\end{lemma}
\begin{proof}
Watson's order-derivative identity uses the modified Bessel function
$\mathrm K_0$ and reads
\[
 \frac{\partial j_{\nu,k}}{\partial\nu}
 =2x\int_0^\infty \mathrm K_0(2x\sinh v)e^{-2\nu v}\,dv.
\]
It applies to the simple positive zeros and positive orders in question.
At $\nu=1/2$ the regular Schr\"odinger solution is a multiple of
$\sin x$, so its positive zeros also satisfy $x>\nu$. For larger
orders this follows from Lemma~\ref{inf:bessel-spacing}.
Replacing $\sinh v$ by $v$ increases the integral. The replacement
integral is
\[
 2x\int_0^\infty \mathrm K_0(2xv)e^{-2\nu v}\,dv
 =\frac{\arccos(\nu/x)}{\sqrt{1-(\nu/x)^2}}.
\]
Insert the integral representation
\[
 \mathrm K_0(u)=\int_0^\infty e^{-u\cosh w}\,dw
\]
and integrate first in $v$. The result is
$\int_0^\infty(\cosh w+\nu/x)^{-1}\,dw$.
The substitution $z=\tanh(w/2)$ evaluates this integral to the
quotient above.

Since $-\mathrm K'_0=\mathrm K_1$ and $\mathrm K_1$ is positive and decreasing,
\[
 0\le \mathrm K_0(2xv)-\mathrm K_0(2x\sinh v)
 \le2x(\sinh v-v)\mathrm K_1(2xv).
\]
Taylor's integral formula gives $\sinh v-v\le v^3\cosh v/6$.
The hypothesis $\nu\ge1/2$ implies
$e^{-2\nu v}\cosh v\le1$. The error in the order derivative is
therefore at most
\[
 \frac{4x^2}{6}\int_0^\infty v^3\mathrm K_1(2xv)\,dv
 =\frac1{24x^2}\int_0^\infty u^3\mathrm K_1(u)\,du
 =\frac\pi{16x^2}.
\]
To evaluate the last integral, use
$\mathrm K_1(u)=\int_0^\infty e^{-u\cosh w}\cosh w\,dw$.
Tonelli's theorem and the elementary gamma integral give
$6\int_0^\infty\operatorname{sech}^3w\,dw=3\pi/2$.
The quotient in \eqref{inf:order-derivative}, written as
$\beta/\sin\beta$ for $0<\beta<\pi/2$, lies in $[1,\pi/2]$.
This proves \eqref{inf:order-lipschitz}. If $x/\nu\to2$ and
$x\to\infty$, the error $\pi/(16x^2)$ tends to zero, and the
quotient in \eqref{inf:order-derivative} tends to
$2\pi/(3\sqrt3)$, proving the limit.
\end{proof}

\subsection{The outgoing integral equation}

For $z>1$ define
\begin{align}
 \mathfrak s(z)&=\sqrt{z^2-1}-\arccos(1/z),
 &A_{\mathrm{Deb}}(z)&=(z^2-1)^{-1/4},\label{inf:debye-definitions}\\
 I(z)&=\frac1{4\sqrt{z^2-1}}+
                \frac5{12(z^2-1)^{3/2}},
 &V(z)&=\frac{z^2(z^2+4)}{4(z^2-1)^3}.\nonumber
\end{align}
The letters $I,V$ in this subsection denote scalar functions only.
The perturbation operator in the inverse argument will be denoted by
$\mathcal V$.

\begin{lemma}[First uniform correction]\label{inf:debye-first}
Let $\nu\ge1$, $z\ge3/2$, and set
\[
 \phi=\nu\mathfrak s(z)-\frac\pi4,\qquad
 \mathcal A_\nu(z)=\sqrt{\frac2{\pi\nu}}A_{\mathrm{Deb}}(z),
 \qquad b_\nu(z)=\frac{I(z)}{2\nu}.
\]
Then
\begin{align}
 \left|\frac{J_\nu(\nu z)}{\mathcal A_\nu(z)}
           -\cos\phi-b_\nu(z)\sin\phi\right|&\le\mathsf R_\nu^{\mathrm{Deb}}(z),
 \label{inf:debye-J}\\
 \left|\frac{Y_\nu(\nu z)}{\mathcal A_\nu(z)}
           -\sin\phi+b_\nu(z)\cos\phi\right|&\le\mathsf R_\nu^{\mathrm{Deb}}(z),
 \label{inf:debye-Y}\\
 \mathsf R_\nu^{\mathrm{Deb}}(z)&=
 \frac1{2\nu^2}\left\{V(z)+\frac{I(z)^2}{1-I(z)/\nu}\right\}
 <\frac6{5\nu^2}.\label{inf:debye-remainder}
\end{align}
Thus the same bound holds on every compact interval about $z=2$,
including a value of $z$ depending on $\nu$.
\end{lemma}
\begin{proof}
Put $\xi=\mathfrak s(z)$ and write
$H_\nu^{(1)}(\nu z)=A_{\mathrm{Deb}}(z)w(\xi)$.
Here $\mathfrak s'(z)=\sqrt{z^2-1}/z$ and
\[
 2\frac{A_{\mathrm{Deb}}'}{A_{\mathrm{Deb}}}
 +\frac{\mathfrak s''}{\mathfrak s'}+\frac1z=0,
 \qquad
 A_{\mathrm{Deb}}''+z^{-1}A_{\mathrm{Deb}}'
 =\frac{z^2+4}{4(z^2-1)^{9/4}}.
\]
Thus substitution in Bessel's equation cancels the first derivative
and gives
\[
 w_{\xi\xi}+\nu^2w=\psi(\xi)w,\qquad
 \psi(\mathfrak s(z))=-V(z).
\]
The outgoing Hankel normalization for fixed $\nu$ at infinity is
$w\sim\sqrt{2/(\pi\nu)}e^{i\nu\xi-i\pi/4}$, with the differentiated
normalization as well. This fixes the outgoing solution integrated
from infinity.
Define $g$ by
$w=\sqrt{2/(\pi\nu)}e^{i\nu\xi-i\pi/4}g$. Variation of constants yields
\begin{equation}\label{inf:volterra}
 g(\xi)=1+\frac1{2i\nu}\int_\xi^\infty
       \{e^{2i\nu(v-\xi)}-1\}\psi(v)g(v)\,dv.
\end{equation}
The potential is integrable. Indeed, with $u=\sqrt{z^2-1}$,
\begin{equation}\label{inf:potential-integral}
 \int_{\mathfrak s(z)}^\infty|\psi(v)|\,dv
 =\int_z^\infty\frac{y(y^2+4)}{4(y^2-1)^{5/2}}\,dy=I(z).
\end{equation}
The iterated integral of order $k$ in \eqref{inf:volterra} has norm at
most $(I(z)/\nu)^k/k!$. The series converges absolutely and gives the
unique bounded outgoing amplitude. Integrating first to a finite
endpoint and then using the Hankel normalization identifies that
solution with the given Hankel function. It follows that
$|g-1|\le e^{I/\nu}-1$.

The first iterate has nonoscillatory term $-iI/(2\nu)$. Its oscillatory
term has absolute value at most $V/(2\nu^2)$. To see this, integrate by
parts once and use
\[
 V'(z)=-\frac{z(z^4+10z^2+4)}{2(z^2-1)^4}<0,
 \qquad
 \int_\xi^\infty|\psi'(v)|\,dv=V(z).
\]
The contribution of the remaining iterates is bounded by
\[
 \frac1\nu\int_\xi^\infty |\psi(v)|
                       (e^{I(v)/\nu}-1)\,dv
 =e^{I/\nu}-1-I/\nu
 \le\frac{I^2}{2\nu^2(1-I/\nu)}.
\]
The final inequality follows by expanding the exponential and bounding
$1/k!\le1/2$ for $k\ge2$.
Thus $|g-(1-iI/(2\nu))|\le\mathsf R_\nu^{\mathrm{Deb}}$.
Taking real and imaginary parts after multiplication by $e^{i\phi}$
proves \eqref{inf:debye-J} and \eqref{inf:debye-Y}.

For $z\ge3/2$, $I(z)\le7/(6\sqrt5)<21/40$ and $V(z)\le9/5$.
In particular $I(z)/\nu<1$. Substitution gives
\[
 \mathsf R_\nu^{\mathrm{Deb}}(z)
 <\frac1{2\nu^2}\left(\frac95+
          \frac{(21/40)^2}{1-21/40}\right)
 =\frac{1809}{1520\nu^2}<\frac6{5\nu^2}.
\]
\end{proof}

\begin{lemma}[Any fixed order]\label{inf:debye-finite}
Let $z>1$, $\nu>0$, and $v=(z^2-1)^{-1/2}$. Define polynomials
$U_k(v)$ by $U_0=1$ and
\begin{equation}\label{inf:debye-recursion}
 U_{k+1}(v)=\frac{v^2(1+v^2)}2U_k'(v)
           +\frac18\int_0^v(1+5w^2)U_k(w)\,dw.
\end{equation}
For every fixed integer $L\ge0$, the amplitude in
\eqref{inf:volterra} satisfies
\begin{equation}\label{inf:debye-all-orders}
 \left|g-\sum_{k=0}^L\frac{U_k(v)}{(i\nu)^k}\right|
 \le\frac{2U_{L+1}(v)}{\nu^{L+1}}e^{I(z)/\nu}.
\end{equation}
If $I(z)/\nu<1$, the exponential can be replaced by
$(1-I(z)/\nu)^{-1}$. In particular,
\begin{equation}\label{inf:debye-polynomials}
 U_1(v)=\frac v8+\frac{5v^3}{24},\qquad
 U_2(v)=\frac{81v^2+462v^4+385v^6}{1152}.
\end{equation}
\end{lemma}
\begin{proof}
One has $dv/d\xi=-v^2(1+v^2)$ and
$V=v^2(1+v^2)(1+5v^2)/4$. The recursion is equivalent to
\[
 2\frac{dU_{k+1}}{d\xi}
 =\psi U_k-\frac{d^2U_k}{d\xi^2},\qquad U_{k+1}(\infty)=0.
\]
Every coefficient in \eqref{inf:debye-recursion} is nonnegative.
It follows that $U_k\ge0$, $dU_k/d\xi\le0$ and
$d^2U_k/d\xi^2\ge0$. Substitution of the truncated sum $g_L$ into
$g''+2i\nu g'-\psi g=0$ leaves only
$(U_L''-\psi U_L)/(i\nu)^L$. This numerator is nonnegative and
\[
 \int_\xi^\infty(U_L''+VU_L)\,d\xi=2U_{L+1}(v).
\]
The integral equation for $g-g_L$ has a forcing term bounded by
$2U_{L+1}/\nu^{L+1}$. In each iterated potential integral, its value
at a later point is no greater than its value at the initial point.
The same ordered-simplex estimate as for \eqref{inf:volterra}
therefore multiplies this bound by at most $e^{I/\nu}$. This proves
\eqref{inf:debye-all-orders}. The inequality
$e^w\le(1-w)^{-1}$ for $0\le w<1$ proves the alternative bound.
The two polynomials follow by integrating the recursion twice.
\end{proof}

\subsection{Fixed spectral windows}

\begin{lemma}[One radial state and no adjacent angular states]
\label{inf:window-isolation}
For each fixed $W>0$ there is $P_W$ such that for real $p\ge P_W$ and
$R\in[2p-1,2p+4]$, the set
\begin{equation}\label{inf:ball-window}
 \mathcal W_p(W)=\{(j,k):|z_{j,k}^2-R^2|\le W\}
\end{equation}
contains at most one state for each angular index $j$ and has no pair
of angular indices differing by one. The threshold is independent of
the split.
\end{lemma}
\begin{proof}
The argument interval has length $O_W(p^{-1})$, which is less than
$\pi$ for large $p$. Lemma~\ref{inf:bessel-spacing} proves the first
assertion. For the second we prove an explicit elementary separation.
Let $z\ge8$ be a zero of $J_\nu$, with $1/3\le\nu/z<1$, and set
$y(v)=J_{\nu+2}(z+v)/J_{\nu+1}(z)$. Consecutive orders have no common
positive zero, so this normalization is valid. The recurrences give
\[
 y(0)=\frac{2(\nu+1)}z\in[2/3,9/4],\qquad
 y'(0)=1-\frac{2(\nu+1)(\nu+2)}{z^2},\quad |y'(0)|\le2.
\]
For $|v|\le1/10$, both coefficients
$(z+v)^{-1}$ and $1-(\nu+2)^2/(z+v)^2$ of the differential equation
have absolute value below one. In fact
$(\nu+2)/(z-1/10)\le(z+2)/(z-1/10)\le100/79<\sqrt2$.
The first-order system for $(y,y')$ has maximum row norm at most two.
Its integral equation gives
$\max(|y(v)|,|y'(v)|)\le(9/4)e^{1/5}<3$. Therefore
$y(v)\ge2/3-3/10>0$ on this interval.

Two roots in \eqref{inf:ball-window} differ in argument by at most
$2W/R$ for large $p$, which is below $1/10$ after increasing $P_W$.
Their orders are at least $p-1$ and smaller than their roots.
Thus the smaller order and its root satisfy $1/3\le\nu/z<1$.
The separation just proved excludes adjacent angular indices.
\end{proof}

\begin{lemma}[Thin phase strip]\label{inf:phase-strip}
For each fixed $W>0$ there are $C_W,P_W<\infty$ with the following
property. Suppose $p\ge P_W$, $R\in[2p-1,2p+4]$, and
$p-1\le\nu\le R-p^{2/3}/2$. Define
\begin{align}
 \kappa(\nu,z)&=\sqrt{z^2-\nu^2},\nonumber\\
 \Xi(\nu,z)&=\kappa(\nu,z)-\nu\arccos(\nu/z)-\pi/4,\label{inf:phase-def}\\
 b(\nu,z)&=\frac1{8\kappa(\nu,z)}+
                       \frac{5\nu^2}{24\kappa(\nu,z)^3}.\nonumber
\end{align}
If $J_\nu(z)=0$ and $|z^2-R^2|\le W$, then
\begin{equation}\label{inf:phase-strip-bound}
 \dist\bigl(\Xi(\nu,R)-\arctan b(\nu,R),
                   \pi(\Z+1/2)\bigr)\le C_W/p.
\end{equation}
For $\nu=p-1+2j$, put
\begin{equation}\label{inf:phase-graph}
 \varphi_p(j)=\pi^{-1}\{\Xi(p-1+2j,R)
                    -\arctan b(p-1+2j,R)\}-\tfrac12.
\end{equation}
On a band $D_0\le R-\nu\le2D_0$ with
$D_0\ge p^{2/3}/2$ one has
\begin{equation}\label{inf:phase-curvature}
 \frac{c}{\sqrt{pD_0}}\le\varphi_p''(j)
 \le\frac{C}{\sqrt{pD_0}},
\end{equation}
where $c,C>0$ are absolute and the threshold may depend on $W$.
\end{lemma}
\begin{proof}
Write $D=R-\nu$. In the stated region and for
$|z-R|\le W/R$, one has $\nu\asymp p$, $z+\nu\asymp p$, and
$\kappa^2\asymp pD$, with absolute comparison constants. The case
$L=1$ of Lemma~\ref{inf:debye-finite} gives
\[
 \frac{J_\nu(z)}{\sqrt{2/\pi}\,\kappa^{-1/2}}
 =\cos\Xi+b\sin\Xi+e,\qquad
 |e|\le\frac{2U_2(\nu/\kappa)}{\nu^2(1-I(z/\nu)/\nu)}.
\]
Here
$I(z/\nu)/\nu=1/(4\kappa)+5\nu^2/(12\kappa^3)
=O(p^{1/2}D^{-3/2})=O(p^{-1/2})$.
The explicit polynomial $U_2$ gives
$|e|=O(pD^{-3})=O(p^{-1})$. Thus the denominator is bounded away
from zero uniformly. At a zero, divide by $\sqrt{1+b^2}$ and use
$\arcsin u\le2u$ for $0\le u\le1/2$ to obtain a distance
$O(p^{-1})$ from $\Xi-\arctan b$ to $\pi(\Z+1/2)$.
Now $\Xi_z=\kappa/z\le1$ and
\[
 b_z=-\frac z{8\kappa^3}-\frac{5\nu^2z}{8\kappa^5}=o(1).
\]
Replacing $z$ by $R$ changes the phase by $O_W(p^{-1})$, proving
\eqref{inf:phase-strip-bound}.

At fixed $R$, $\Xi_{\nu\nu}=1/\kappa$ and direct differentiation gives
\[
 b_\nu=\frac{13\nu}{24\kappa^3}+\frac{5\nu^3}{8\kappa^5},\qquad
 b_{\nu\nu}=\frac{13}{24\kappa^3}
               +\frac{7\nu^2}{2\kappa^5}
               +\frac{25\nu^4}{8\kappa^7}.
\]
It follows that
\[
 \kappa\,| (\arctan b)_{\nu\nu}|
 \le\kappa\{|b_{\nu\nu}|+2|b||b_\nu|^2\}
 =O(pD^{-3})+O(p^2D^{-6})=O(p^{-1}).
\]
The constants are absolute throughout the region. Therefore
$\varphi_p''=(4/\pi)\kappa^{-1}(1+O(p^{-1}))$.
Since $\kappa\asymp\sqrt{pD_0}$ on the band, this proves
\eqref{inf:phase-curvature}.
\end{proof}

\begin{lemma}[Convex lattice strip]\label{inf:lattice-strip}
Let $f\in C^2(I)$, where $I$ has length at most $L$, and suppose
$0<m_0\le f''\le M_0$. Let $\delta>0$ satisfy
\[
 \frac12M_0L^3+2\delta L<1.
\]
Then the integer points $(j,k)$ with $j\in I$ and
$|k-f(j)|\le\delta$ are collinear. If there is at most one such point
at each $j$, their number is at most
$2+4\sqrt{\delta/m_0}$, with an integer upper bound understood.
\end{lemma}
\begin{proof}
For three such points, the determinant with rows $(1,j,k)$ is an
integer. The determinant of the three graph values is
$(x_2-x_1)(x_3-x_1)(x_3-x_2)f[x_1,x_2,x_3]$ when the abscissae
are distinct. The integral formula for the second divided
difference bounds its absolute value by $M_0/2$, so this
contribution is at most $M_0L^3/2$. If two abscissae coincide,
the graph determinant is zero and the same bound holds. The three vertical errors contribute at most
$2\delta L$: for ordered abscissae their coefficient absolute values
sum to twice the largest separation. Thus the determinant has
absolute value less than one and is zero. If there are at least three
points, fixing the two extreme ones shows that every point lies on
their line.

Take a middle point with abscissa $x_2$ and extreme abscissae
$x_1,x_3$. The chord of $f$ satisfies
\[
 \operatorname{chord}(x_2)-f(x_2)
 \ge\frac{m_0}{2}(x_2-x_1)(x_3-x_2).
\]
One proof subtracts $m_0x^2/2$ from $f$ and uses convexity of the
result. Since the three integer points are collinear, their errors
bound the same chord deficit by $2\delta$. Among $n$ distinct integer
abscissae choose a middle one; its two separations have product at
least $(n-2)^2/4$. Therefore $(n-2)^2\le16\delta/m_0$.
The assertions for zero, one or two points need no determinant.
\end{proof}

\begin{theorem}[Sublinear window count]\label{inf:global-count}
For every fixed $W>0$ there are $C_W,P_W<\infty$ such that, for
$p\ge P_W$ and $R\in[2p-1,2p+4]$,
\begin{equation}\label{inf:global-count-bound}
 \#\mathcal W_p(W)\le C_Wp^{2/3}.
\end{equation}
The bound holds for real $p$, and is independent of the split.
\end{theorem}
\begin{proof}
There is at most one radial state per angular index by
Lemma~\ref{inf:window-isolation}. Also $z_{j,k}>\nu_j$, so the
region $\nu_j>R-p^{2/3}$ has only $O_W(p^{2/3})$ possible indices.
In the remaining region Lemma~\ref{inf:phase-strip} assigns to each
state a lattice point $(j,m)$ with
$|m-\varphi_p(j)|\le\delta=C_W'/p$. This assignment is injective:
there is at most one radial state at a given $j$, and $2\delta<1$.

Divide $D=R-\nu_j$ into dyadic bands $D_0\le D\le2D_0$, from
$p^{2/3}$ up to $O(p)$. On each band
$f''\asymp(pD_0)^{-1/2}$ for $f=\varphi_p$. Divide its horizontal
interval into intervals of length
$L=c_*(pD_0)^{1/6}$, with $c_*>0$ fixed so small that
$M_0L^3/2\le1/4$. Since $D_0=O(p)$, the other term
$2\delta L=O_W(p^{-2/3})$ is below $1/4$ for large $p$.
Lemma~\ref{inf:lattice-strip} bounds the number of points per interval
by an absolute multiple of $1+\sqrt{C'_W}$, because
$\delta/m_0=O_W(\sqrt{D_0/p})=O_W(1)$.

The band has horizontal length $O(D_0)$, so it requires
$O(p^{-1/6}D_0^{5/6}+1)$ such intervals. The sum of $D_0^{5/6}$ over
the dyadic bands is $O(p^{5/6})$; the sum of the additive ones is
$O(\log p)$. This proves \eqref{inf:global-count-bound} after adding
the turning region. Every constant and threshold in this proof is
fixed once $W$ is fixed.
\end{proof}

\begin{corollary}[Domains in a thin shell]\label{inf:shell-count}
Fix $W,C_{\mathrm{shell}}>0$ and integer blocks $a,b$ with
$a+b=2p$ in the admissible range. If a two-block invariant domain in the unit-ball
normalization satisfies
$B_{1-C_{\mathrm{shell}}/p^2}\subset\Omega\subset B_{1+C_{\mathrm{shell}}/p^2}$, then its invariant
Dirichlet spectrum has at most $C_{W,C_{\mathrm{shell}}}p^{2/3}$ eigenvalues in
$[R^2-W,R^2+W]$ for large $p$ and $R\in[2p-1,2p+4]$.
The constants are uniform in the admissible integer splits.
\end{corollary}
\begin{proof}
Zero extension preserves invariance. The min--max principle and radial
dilation give, for each ordered eigenvalue index $i$,
\[
 (1+C_{\mathrm{shell}}/p^2)^{-2}\lambda_i(B)
 \le\lambda_i(\Omega)\le
 (1-C_{\mathrm{shell}}/p^2)^{-2}\lambda_i(B).
\]
Any index counted for $\Omega$ has its ball eigenvalue in
$[R^2-W',R^2+W']$, where $W'$ depends only on $W,C_{\mathrm{shell}}$ since
$R^2=O(p^2)$. Apply Theorem~\ref{inf:global-count}.
\end{proof}

\section{Primary crossings and a continued-fraction sequence}\label{inf:sec-arithmetic}

The real crossing orders and the integer orders are different objects.
We denote a real crossing by $\widehat p_m$ and use $\mathcal P$ for
the selected set of integers.

\subsection{A congruence approximation lemma}

\begin{lemma}[Congruence approximation]\label{inf:congruence}
For every positive irrational number $\alpha$ there are infinitely many
integer pairs $(P,Q)$ with
\begin{equation}\label{inf:congruence-target}
 P\equiv1\pmod2,\qquad Q\equiv1\pmod4,\qquad Q>0,
 \qquad Q|Q\alpha-P|<\frac92.
\end{equation}
They can be selected by the finite case rule in the proof from every
sufficiently late pair of consecutive continued-fraction convergents.
\end{lemma}
\begin{proof}
Write the convergents as $P_j^{\mathrm{cf}}/Q_j^{\mathrm{cf}}$ and put $v_j=(P_j^{\mathrm{cf}},Q_j^{\mathrm{cf}})$.
Use indices for which all the three preceding denominators are positive.
The determinant of consecutive vectors has absolute value one. Set
\[
 \rho_j=Q_{j-1}^{\mathrm{cf}}/Q_j^{\mathrm{cf}}\in(0,1),\qquad
 T_j^{\mathrm{cf}}=[a_{j+1};a_{j+2},\ldots]>1.
\]
The identity
$\alpha=(T_j^{\mathrm{cf}}P_j^{\mathrm{cf}}+P_{j-1}^{\mathrm{cf}})/(T_j^{\mathrm{cf}}Q_j^{\mathrm{cf}}+Q_{j-1}^{\mathrm{cf}})$ gives, for
$(P,Q)=Av_j+Bv_{j-1}$ and nonnegative integers $A,B$,
\begin{equation}\label{inf:cf-error}
 Q|Q\alpha-P|=\frac{(A+B\rho_j)|A-BT_j^{\mathrm{cf}}|}{T_j^{\mathrm{cf}}+\rho_j}.
\end{equation}
Denote this expression by $E(A,B)$ when the index is understood.
Directly from \eqref{inf:cf-error},
\begin{equation}\label{inf:cf-basic-bounds}
 E(1,0)<1,\quad E(1,1)<2,\quad E(2,1)<3,
 \quad E(1,2)<2(1+2\rho_j).
\end{equation}
For $E(2,1)$, split at $T_j^{\mathrm{cf}}=2$. On $1<T_j^{\mathrm{cf}}<2$ the quotient has
maximum at $T_j^{\mathrm{cf}}=1$ and is at most $(2+\rho_j)/(1+\rho_j)<2$.
On $T_j^{\mathrm{cf}}\ge2$ it is less than $2+\rho_j<3$.
The other three inequalities follow by bounding respectively
$1/(T_j^{\mathrm{cf}}+\rho_j)<1$, $(T_j^{\mathrm{cf}}-1)/(T_j^{\mathrm{cf}}+\rho_j)<1$, and
$(2T_j^{\mathrm{cf}}-1)/(T_j^{\mathrm{cf}}+\rho_j)<2$.

We first give a helper rule. Suppose $v_k$ has even numerator and odd
denominator, $v_{k-1}$ has odd numerator and even denominator, and
$Q_k^{\mathrm{cf}}+Q_{k-1}^{\mathrm{cf}}\equiv3\pmod4$. The vector
$S=3v_k+v_{k-1}$ has the target residues. Put $T=T_k^{\mathrm{cf}}$ and
$\rho=\rho_k$. If $T\ge2$, its error is less than four: for
$2\le T\le3$ it is at most $(3+\rho)/(2+\rho)<3/2$, and for
$T\ge3$ it is less than $3+\rho<4$.
If $1<T<2$, then $a_{k+1}=1$ and $v_{k+1}=v_k+v_{k-1}$.
The alternative $D=3v_{k+1}$ also has the target residues. Their
error products are
\[
 E_S=\frac{(3+\rho)(3-T)}{T+\rho},\qquad
 E_D=\frac{9(1+\rho)(T-1)}{T+\rho}.
\]
The inequalities $E_S\le3$ and $E_D\le3$ hold respectively when
\[
 T\ge\frac9{6+\rho},\qquad
 T\le\frac{3+4\rho}{2+3\rho}.
\]
These two ranges cover $(1,2)$, because
$(3+4\rho)(6+\rho)-9(2+3\rho)=4\rho^2\ge0$.
Choose the candidate with smaller error. This fixes the helper rule
and always gives error less than four, with positive denominator at
least $Q_{k-1}^{\mathrm{cf}}$.

Consider now $v_j,v_{j-1}$. There are two cases.

\emph{Case I: one vector has both entries odd.}
Call that vector $v_k$, where $k=j$ or $j-1$.
If $Q_k^{\mathrm{cf}}\equiv1\pmod4$, take $v_k$ itself.
Otherwise $Q_k^{\mathrm{cf}}\equiv3\pmod4$.
If $Q_{k-1}^{\mathrm{cf}}$ is even, its numerator is odd by the determinant condition.
For even $a_{k+1}$, take $3v_k$; here $T_k^{\mathrm{cf}}>2$ and its error is
$9/(T_k^{\mathrm{cf}}+\rho_k)<9/2$.
For odd $a_{k+1}$, the vector $v_{k+1}$ has even numerator and odd
denominator. Take $2v_{k+1}+v_k$. Its denominator is $1$ modulo four,
its numerator is odd, and \eqref{inf:cf-basic-bounds} at index $k+1$
gives error less than three.

If $Q_{k-1}^{\mathrm{cf}}$ is odd, its numerator is even.
When $a_k\ge2$ one has $\rho_k<1/2$; take
$v_k+2v_{k-1}$. Its denominator is $1$ modulo four and its error is
less than $2(1+2\rho_k)<4$.
When $a_k=1$, the vector $v_{k-2}=v_k-v_{k-1}$ has odd numerator
and even denominator. The pair $v_{k-1},v_{k-2}$ satisfies the helper
hypotheses, since its denominator sum is $Q_k^{\mathrm{cf}}\equiv3\pmod4$.
Apply the helper rule at index $k-1$.

\emph{Case II: neither vector has both entries odd.}
Their types modulo two must be (even, odd) and (odd, even): neither
vector is zero modulo two, and their determinant is odd.
If $Q_j^{\mathrm{cf}}+Q_{j-1}^{\mathrm{cf}}\equiv1\pmod4$, use their sum, whose error is less
than two. If that sum is $3$ modulo four and $Q_j^{\mathrm{cf}}$ is odd, apply the
helper rule directly.

It remains that $v_j$ has odd numerator and even denominator $Q_{\mathrm e}$,
$v_{j-1}$ has even numerator and odd denominator $Q_{\mathrm o}$, and
$Q_{\mathrm e}+Q_{\mathrm o}\equiv3\pmod4$.
Thus $Q_{\mathrm e}-Q_{\mathrm o}\equiv1\pmod4$.
Use $v_{j-2}=v_j-a_jv_{j-1}$ and distinguish the residue of $a_j$.
If $a_j\equiv1\pmod4$, take $v_{j-2}$.
If $a_j\equiv2\pmod4$, take $v_{j-1}+v_{j-2}$; its error is less
than two by \eqref{inf:cf-basic-bounds} at index $j-1$.
If $a_j\equiv3\pmod4$, the denominator of $v_{j-2}$ is $3$ modulo
four; take $2v_{j-1}+v_{j-2}$, whose error is less than three.
In these three choices the numerator is odd and the denominator is
$1$ modulo four. If $a_j\equiv0\pmod4$, the pair
$v_{j-1},v_{j-2}$ has the helper types and denominator sum $3$ modulo
four, so apply that rule at index $j-1$.

Every selected denominator is positive and at least $Q_{j-3}^{\mathrm{cf}}$.
These lower bounds tend to infinity. Thus infinitely many distinct
pairs have been constructed, and all their errors are strictly below
$9/2$.
\end{proof}

\subsection{The crossing expansion}

Define the fixed constants
\begin{equation}\label{inf:arithmetic-constants}
 c_{\mathrm{ar}}=\sqrt3-\frac\pi3,
 \qquad \alpha_{\mathrm{ar}}=\frac\pi{2c_{\mathrm{ar}}},
 \qquad d_{\mathrm{ar}}=\frac5{9\sqrt3c_{\mathrm{ar}}},
 \qquad \delta_0=1-\frac\pi{3\sqrt3}.
\end{equation}
Also write
\begin{equation}\label{inf:algebraic-radius}
 R_{\mathrm{alg}}(v)=2\sqrt{(v+1)(v+2)},\qquad
 A_m=\frac{\alpha_{\mathrm{ar}}(4m+1)-3}{2}.
\end{equation}
These constants are unrelated to the block parameters $\ab,\bb$.

\begin{proposition}[Primary-crossing expansion]\label{inf:primary-expansion}
For every sufficiently large integer $m$, there is a unique primary
crossing $\widehat p_m$ in a fixed $O(A_m^{-1})$ neighbourhood of $A_m$,
meaning
\[
 J_{\widehat p_m}(R_{\mathrm{alg}}(\widehat p_m))
 =J_{\widehat p_m+3}(R_{\mathrm{alg}}(\widehat p_m))=0.
\]
It satisfies
\begin{equation}\label{inf:root-expansion}
 \widehat p_m=A_m-\frac{d_{\mathrm{ar}}}{A_m}+O(A_m^{-2}),
\end{equation}
with an absolute remainder constant. No positive integer or
half-integer is a primary crossing.
\end{proposition}
\begin{proof}
At a zero $R$ of $J_v$, the three-term recurrence gives
\begin{equation}\label{inf:crossing-recurrence}
 \frac{J_{v+3}(R)}{J_{v-1}(R)}
 =1-\frac{4(v+1)(v+2)}{R^2}.
\end{equation}
The denominator is nonzero. Thus the crossing equation is precisely
$J_v(R_{\mathrm{alg}}(v))=0$.

At $R=R_{\mathrm{alg}}(v)$, direct Taylor expansion gives
\begin{align*}
 \sqrt{R^2-v^2}&=\sqrt3v+2\sqrt3-\frac2{\sqrt3v}+O(v^{-2}),\\
 v\arccos(v/R)&=\frac{\pi v}{3}+\frac{\sqrt3}{2}
                       -\frac{11}{4\sqrt3v}+O(v^{-2}).
\end{align*}
Consequently the phase of Lemma~\ref{inf:debye-first} is
\[
 \Xi(v,R)=c_{\mathrm{ar}}v+\frac{3\sqrt3}{2}-\frac\pi4
                         +\frac{\sqrt3}{4v}+O(v^{-2}),
\]
and its first sine coefficient is
$b(v,R)=7/(36\sqrt3v)+O(v^{-2})$.
Lemma~\ref{inf:debye-first} applies uniformly since $R/v\to2$.
In particular, the remainder is $O(v^{-2})$ for these varying
arguments.

The identity
$c_{\mathrm{ar}}A_m+3\sqrt3/2-\pi/4=\pi(m+1/2)$ is exact.
At $v=A_m\pm K/A_m$, for a sufficiently large fixed $K$, the leading
cosine changes sign with magnitude exceeding the first correction
and its remainder. The intermediate value theorem gives a zero.
Every such zero obeys
\[
 \Xi(v,R)=\pi(m+1/2)+\frac7{36\sqrt3v}+O(v^{-2}).
\]
Substitution of the phase expansion proves
\eqref{inf:root-expansion}, since
$\{7/(36\sqrt3)-\sqrt3/4\}/c_{\mathrm{ar}}=-d_{\mathrm{ar}}$.

For uniqueness, the zero lies on one of the analytic curves
$R=j_{v,k}$, because the zeros are simple. On the interval under
consideration, Lemma~\ref{inf:bessel-order} gives
$\partial_vj_{v,k}\le\pi/2$, whereas
$R'_{\mathrm{alg}}(v)=(2v+3)/\sqrt{(v+1)(v+2)}>2$.
Thus each difference $R_{\mathrm{alg}}(v)-j_{v,k}$ is strictly
increasing. Two different indices cannot cross the algebraic radius
in an interval of length $O(A_m^{-1})$: their roots are separated by
more than $\pi$, while both the algebraic radius and the zero curves
move by $O(A_m^{-1})$ there. This proves uniqueness for large $m$.

Finally, for integer $v$, two distinct nonnegative integer-order
Bessel functions have no common positive zero (Bourget's theorem).
This classical consequence of Siegel's theorem is stated in
\cite[\S10.21(i)]{inf:dlmf}, with the source
\cite[\S15.28]{inf:watson}. Applying it to orders $v$ and $v+3$
excludes an integer primary crossing. For positive half-integers,
Lemma~\ref{inf:half-integer-noncoincidence} gives the same conclusion.
\end{proof}

\begin{definition}[Selected integer sequence]\label{inf:integer-sequence}
Apply the case rule in Lemma~\ref{inf:congruence} to
$\alpha=\alpha_{\mathrm{ar}}$, with convergents indexed starting at
$0$, for each index $j\ge6$. In its helper choose the candidate with
smaller error. Remove repeated pairs and list the resulting integers
$(P-3)/2$ in increasing order. Their set is $\mathcal P$; its increasing
enumeration is denoted by $(p_m)$.
The integer $m$ in $(p_m)$ is an enumeration index, whereas the primary
crossing associated to a selected pair has index $(Q-1)/4$.
\end{definition}

\begin{definition}[Near-primary half-lattice inputs]\label{inf:input-families}
Fix an integer $m_0$ beyond which Proposition~\ref{inf:primary-expansion}
holds. For $C_{\mathrm{ap}}>0$ set
\begin{gather*}
 \mathcal P_{\mathrm{fix}}(C_{\mathrm{ap}})=
 \left\{p\in\tfrac12\Z:p>0,\ 
 \exists m\ge m_0:\ |p-\widehat p_m|<C_{\mathrm{ap}}/\widehat p_m\right\},\\
 \mathcal P_{\mathrm{bd}}=\mathcal P_{\mathrm{fix}}(31/10).
\end{gather*}
The continued-fraction set $\mathcal P$ remains the selected integer
subsequence of Definition~\ref{inf:integer-sequence}. The symbol
$p_0$ denotes the crossing associated to the current input. It is
unique for sufficiently large inputs in each fixed-width family,
since successive crossings have separation $2\alpha_{\mathrm{ar}}+o(1)$.
\end{definition}

\begin{theorem}[Near-integer primary crossings]\label{inf:near-integer}
The set $\mathcal P$ is unbounded. For every sufficiently large selected
pair $(P,Q)$, the integer $p=(P-3)/2$ and the primary crossing
$p_0=\widehat p_{(Q-1)/4}$ satisfy
\begin{equation}\label{inf:near-integer-bound}
 0<|p-p_0|<\frac{31}{10p_0}.
\end{equation}
Thus the assertion holds for every sufficiently large member of
$\mathcal P$, with its associated crossing. The threshold is implicit.
\end{theorem}
\begin{proof}
The number $\alpha_{\mathrm{ar}}$ is irrational. Indeed, a rational
value would imply
$\pi(1+2\alpha_{\mathrm{ar}}/3)=2\alpha_{\mathrm{ar}}\sqrt3$,
contradicting the transcendence of $\pi$.
For $Q=4m+1$ and $p=(P-3)/2$, the expansion
\eqref{inf:root-expansion} gives
\[
 p_0|p_0-p|\le\frac{p_0}{2}|Q\alpha_{\mathrm{ar}}-P|
              +d_{\mathrm{ar}}\frac{p_0}{A_m}+O(A_m^{-1}).
\]
Here $p_0/Q\to\alpha_{\mathrm{ar}}/2$ and $p_0/A_m\to1$.
The congruence error bound therefore yields
\begin{equation}\label{inf:arithmetic-limsup}
 \limsup p_0|p_0-p|
 \le\frac98\alpha_{\mathrm{ar}}+d_{\mathrm{ar}}<\frac{31}{10}.
\end{equation}
The final comparison follows from rational bounds. From
$\pi<22/7$ and $\sqrt3>19/11$ one obtains
$c_{\mathrm{ar}}>157/231$, $\alpha_{\mathrm{ar}}<363/157$, and
\[
 \frac98\alpha_{\mathrm{ar}}+d_{\mathrm{ar}}
 <\frac{220099}{71592}<\frac{31}{10},
 \qquad 31\cdot71592-10\cdot220099=18362>0.
\]
Bourget's theorem gives $p\ne p_0$. The denominators tend to infinity
and $P/Q\to\alpha_{\mathrm{ar}}>0$, so the selected integers are
unbounded. Distinct sufficiently large $Q$ cannot give the same
integer $p$: the corresponding numerators would be equal, while
$\alpha_{\mathrm{ar}}|Q-Q'|\le9/(2Q)+9/(2Q')$ would tend to zero.
This also removes any ambiguity in the associated crossing for large
members of the sequence.
\end{proof}

\begin{corollary}[Sparsity]\label{inf:sequence-sparsity}
The selected integer set satisfies $\#(\mathcal P\cap[1,X])=O(\log X)$.
\end{corollary}
\begin{proof}
The rule selects at most one pair at each continued-fraction index $j$,
and its denominator is at least $Q_{j-3}^{\mathrm{cf}}$. These denominators grow at
least as the Fibonacci sequence. If the selected integer is at most
$X$, its error bound gives $Q=O(X)$, since
$P=2p+3=\alpha_{\mathrm{ar}}Q+O(Q^{-1})$.
Hence only $O(\log X)$ indices can contribute. Removing duplicates
cannot increase the count.
\end{proof}

\subsection{The detuning parameter}

Continue the zero at a primary crossing $p_0$ as $R(p)=j_{p,k_0}$.
In this paragraph $p$ is real on the continuation interval; its
selected endpoint belongs to $\tfrac12\Z$. Put
\begin{equation}\label{inf:detuning-definition}
 \chi=\frac{J_{p+3}(R)}{J'_{p+3}(R)-(p-1)J_{p+3}(R)/R},
 \qquad \theta=4Rp\chi,\qquad \tau=|\theta|.
\end{equation}

\begin{lemma}[Uniform bounded detuning]\label{inf:detuning}
At every sufficiently large $p\in\mathcal P_{\mathrm{bd}}$,
\begin{equation}\label{inf:detuning-range}
 R=2p+3+O(p^{-1}),\qquad 0<|\theta|<20.
\end{equation}
More precisely, uniformly for $|p-p_0|\le31/(10p_0)$,
\begin{equation}\label{inf:detuning-asymptotic}
 \theta=(16\delta_0+o(1))p_0(p-p_0).
\end{equation}
The denominator in \eqref{inf:detuning-definition} is nonzero in this
interval for large $p_0$.
\end{lemma}
\begin{proof}
Repeated use of the Bessel recurrences at $J_p(R)=0$ gives the exact
identity
\begin{equation}\label{inf:trace-ratio-exact}
 \chi=-\frac R{2(p+1)}\frac{R^2-4(p+1)(p+2)}{2R^2-4(p+1)(p+2)}.
\end{equation}
The omitted common factor in numerator and denominator is
$J_{p-1}(R)$, which is nonzero. The denominator in the rational
expression is $4p^2+O(p)$ throughout the interval, so the physical
normal-trace denominator is also nonzero.
Lemma~\ref{inf:bessel-order}, $R/p\to2$, and $R\to\infty$ give
\[
 R'(p)=\frac{2\pi}{3\sqrt3}+o(1),\qquad
 R'_{\mathrm{alg}}(p)=2+o(1),
\]
uniformly on that interval. Since the two radii agree at $p_0$,
\[
 R(p)-R_{\mathrm{alg}}(p)=-(2\delta_0+o(1))(p-p_0).
\]
Substitute this in \eqref{inf:trace-ratio-exact} and multiply by
$4Rp$ to prove \eqref{inf:detuning-asymptotic}.
The radius assertion follows from the same comparison and
$R_{\mathrm{alg}}(p)=2p+3+O(p^{-1})$.
The inequality $\delta_0<2/5$ and the defining distance bound
give eventually
$\tau<16(2/5)(31/10)+o(1)=19.84+o(1)<20$.
The difference of the two radii has a strictly negative derivative,
because $R'\le\pi/2<2<R'_{\mathrm{alg}}$. It vanishes only at
$p=p_0$. Proposition~\ref{inf:primary-expansion} excludes that equality
on $\tfrac12\Z$. Hence $\chi\ne0$ and $\theta\ne0$ at every input.
\end{proof}

\subsection{Low angular degrees at near-primary inputs}

\begin{lemma}[A radial pole estimate]\label{inf:riccati-gap}
Let $p\ge64$, $R=R_{\mathrm{alg}}(p)$, and
$\nu=p+\ell-1<2p+3$. If $J_\nu(R)\ne0$, set
$M_\ell=J'_\nu(R)/J_\nu(R)-(p-1)/R$. Then every positive zero $z$
of $J_\nu$ satisfies
\begin{equation}\label{inf:riccati-distance}
 |z^2-R^2|\ge\frac{2(R-1)}{|M_\ell|+2}.
\end{equation}
\end{lemma}
\begin{proof}
For positive $\lambda$ away from poles define
$w(\lambda)=\sqrt\lambda J'_\nu(R\sqrt\lambda)/
J_\nu(R\sqrt\lambda)$. Bessel's equation implies
\[
 w'(\lambda)=-\frac R{2\lambda}
             \left(w(\lambda)^2+\lambda-\frac{\nu^2}{R^2}\right).
\]
Put $\eta=1/R$ and $A=\sqrt{1+\eta}$. On
$[1-\eta,1+\eta]$ one has
$|\lambda-\nu^2/R^2|\le1+\eta$, because $\nu^2<R^2+1$.
Consequently
\[
 \left|\frac d{d\lambda}\arctan\frac{w(\lambda)}A\right|
 \le\frac{RA}{2(1-\eta)}.
\]
At the first pole on either side of one the arctangent tends to an
endpoint $\pm\pi/2$. Its distance to that endpoint at $\lambda=1$
is at least $\arctan(A/|w(1)|)\ge A/(|w(1)|+A)$.
(The latter inequality follows by differentiating
$\arctan u-u/(1+u)$.) A pole in this interval is therefore at
$\lambda$-distance at least
$2(1-\eta)/(R(|w(1)|+A))$ from one.
Now $|w(1)|+A<|M_\ell|+2$, since $(p-1)/R<1/2$ and
$A<3/2$. Multiplication by $R^2$ proves
\eqref{inf:riccati-distance} for these poles. A pole outside the
interval gives $|z^2-R^2|\ge R$, which is at least the claimed
right-hand side. Poles beyond the first on either side are farther
away. This proves the bound for the entire radial spectrum.
\end{proof}

\begin{lemma}[Uniform recurrence remainder]\label{inf:low-recurrence}
At a real primary crossing $p\ge10^6$, put $s_0=R_{\mathrm{alg}}(p)/2$
and define
\[
 q_{-1}=1,\quad q_0=0,\quad
 q_{k+1}=\frac{p+k}{s_0}q_k-q_{k-1}.
\]
Let $\mathsf a_{-1}=1,\mathsf a_0=0$,
$\mathsf a_{k+1}=\mathsf a_k-\mathsf a_{k-1}$, and
$\mathsf b_{-1}=\mathsf b_0=0$,
$\mathsf b_{k+1}=\mathsf b_k-\mathsf b_{k-1}+(k-3/2)\mathsf a_k$.
Set $\mathsf e_k=q_k-\mathsf a_k-\mathsf b_k/p$ and
$\nrm{(u,v)}_\triangle=\max(|u|,|v|,|u-v|)$. Then
\begin{equation}\label{inf:low-remainder}
 \nrm{(\mathsf e_k,\mathsf e_{k-1})}_\triangle
 \le e^{k(k+3)/p}\frac{(k+2)^4}{4p^2},\qquad
 \mathsf b_{3h}=(-1)^h3h(h-1).
\end{equation}
\end{lemma}
\begin{proof}
The transformation $(u,v)\mapsto(u-v,u)$ preserves the norm
$\max(|u|,|v|,|u-v|)$. The sequence $\mathsf a_k$ has period six and norm one.
Write $\delta_k=(p+k)/s_0-1$. From
$p+3/2-s_0=1/[4(p+3/2+s_0)]$ one obtains
\[
 \left|\delta_k-\frac{k-3/2}p\right|
 \le\frac{3(k+2)}{2p^2},\qquad
 |\delta_k|\le\frac{2(k+2)}p.
\]
Summing the first-order recurrence in the invariant norm bounds its
vector by $k(k+3)/2$. Subtracting the first-order approximation from
the exact recurrence then bounds the next remainder vector by
$(1+|\delta_k|)$ times the current vector plus $(k+2)^3/p^2$.
Iteration and $1+u\le e^u$ give an exponential no greater than
$e^{k(k+3)/p}$; the sum of the cubics is at most $(k+2)^4/4$.
This proves the first bound. Three successive steps of the recurrence
for $\mathsf b_k$ give
$\mathsf b_{3(h+1)}=-\mathsf b_{3h}+6h(-1)^{h+1}$. Its initial value is zero,
so induction gives the second formula.
\end{proof}

\begin{proposition}[Low-degree exclusion]\label{inf:low-exclusion}
Let $p_0\ge10^6$ be a sufficiently large primary crossing and
$R_0=R_{\mathrm{alg}}(p_0)$. Apart from the degree-four state at $R_0$,
every even degree $\ell\ge0$ with
$\ell-1\le\sqrt{p_0}/4$ has its ball Dirichlet spectrum at distance
greater than $64$ from $R_0^2$.
At any real $p$ with $|p-p_0|\le31/(10p_0)$ and continued radius
$R=j_{p,k_0}$, those nonseed degrees have their spectrum outside
$[R^2-32,R^2+32]$. The seed's other radial states are outside this
window as well.
\end{proposition}
\begin{proof}
At the crossing the recurrence is exactly
$q_k=J_{p_0+k}(R_0)/J_{p_0-1}(R_0)$.
Put $k=\ell-1$, initially $k\ge7$ with $k^2\le p_0/16$.
The first-order contribution in the invariant vector norm is at most
$5/112$ and the remainder in \eqref{inf:low-remainder} is at most
$81/8192$; their sum is less than $1/16$.
If $3\nmid k$, this gives $|q_k|>15/16$, $|q_{k-1}|<17/16$,
and hence
\[
 |M_{k+1}|=
 \left|\frac{q_{k-1}}{q_k}-\frac{2p_0+k-1}{R_0}\right|<4.
\]
If $3\mid k$ and $k\ge15$, then
$|\mathsf b_k|=k(k-3)/3\ge4k^2/15$. The remainder is at most
$648k^4/(625p_0^2)$, which is less than $|\mathsf b_k|/(4p_0)$.
It follows that
\[
 |q_k|\ge\frac{k^2}{5p_0},\qquad |q_k|<\frac1{32},
 \qquad |M_{k+1}|\le\frac{10p_0}{k^2}.
\]
For the last estimate use $|q_{k-1}|<17/16$ and
$(2p_0+k-1)/R_0<2$.
Lemma~\ref{inf:riccati-gap} now bounds the inverse distance by
\[
 \frac{|M_{k+1}|+2}{2(R_0-1)}
 \le\frac{27}{8k^2}\le\frac{27}{1800}<\frac1{64}
\]
in the resonant range. In the nonresonant range it gives a smaller
bound for $p_0\ge10^6$.

The remaining resonant index is $k=9$.
Equation~\eqref{inf:low-remainder} gives
$|q_9|\ge(18-3661/p_0)/p_0>17.99/p_0$ and
$|q_8|\le1+44/p_0+2501/p_0^2<1.001$.
Thus $|M_{10}|<p_0/17+2$, and its inverse distance is less than
$(p_0/17+4)/(4p_0-2)<1/64$.
For $\ell=2$ use $q_1=-1,q_0=0$, giving $|M_2|<1$.
For $\ell=6$ the exact values
$q_4=(p_0+1)/s_0$ and $q_5=(p_0+4)/(p_0+2)$ give $|M_6|<2$.
For $\ell=0$ one has $M_0=0$ by \eqref{inf:ball-traces}.
These exhaust the low even degrees except $\ell=4$.
Its remaining radial states are separated by Lemma~\ref{inf:bessel-spacing}.

To pass to $p$, continue each radial zero in its order.
Both relevant orders exceed eight along the interval, so
\eqref{inf:order-lipschitz} bounds the difference of any two such
order derivatives by less than $3/5$.
Consequently $j_{p+\ell-1,k}-R(p)$ changes by less than
$93/(50p_0)<2/p_0$.
If its squared-energy offset at $p$ had absolute value at
most $32$, then $R(p)>2p_0$ and
$j_{p+\ell-1,k}+R(p)\ge(39/10)p_0$ for large $p_0$.
Its radius offset would have absolute value below $9/p_0$,
and at $p_0$ below $11/p_0$. Its squared offset at the crossing
would then be less than
\[
 \frac{11}{p_0}\left(4p_0+6+\frac{11}{p_0}\right)<55,
\]
contradicting the bound $64$. The other radial seed states remain
separated by the spacing and the same Lipschitz estimate.
\end{proof}

\section{Half-integer orders and the Euclidean chain}\label{inf:sec-parity}

The order parameter belongs to $\tfrac12\Z$ at a Euclidean input.
The block dimensions $a,b$ and the angular and radial degree indices
are integers; the Jacobi parameters and the Bessel orders may be
half-integers.

\begin{lemma}[Half-integer non-coincidence]\label{inf:half-integer-noncoincidence}
For $\nu\in\{1/2,3/2,\ldots\}$, $J_\nu$ has no positive algebraic
zero. Two distinct positive half-integer orders have no common positive
zero. In particular $J_p$ and $J_{p+3}$ have no common positive zero
for $p\in\Z+1/2$, $p>0$.
\end{lemma}
\begin{proof}
Write $\nu=l+1/2$, $l\ge0$. The half-order identities and recurrence
\cite[\S\S10.6, 10.16]{inf:dlmf} give
\[
 \sqrt{\pi x/2}\,J_{l+1/2}(x)
 =\mathsf F_l(x)\sin x+\mathsf G_l(x)\cos x,
\]
where $\mathsf F_l,\mathsf G_l\in\mathbb Q[x^{-1}]$ and
\[
 (\mathsf F_{-1},\mathsf G_{-1})=(0,1),\qquad
 (\mathsf F_0,\mathsf G_0)=(1,0),\qquad
 (\mathsf F_{l+1},\mathsf G_{l+1})
 =\frac{2l+1}{x}(\mathsf F_l,\mathsf G_l)
       -(\mathsf F_{l-1},\mathsf G_{l-1}).
\]
The determinant of the successive rows $(l,l-1)$ is one: the
recurrence preserves it and the initial determinant is one.
Thus the two entries in row $l$ never vanish together at $x\ne0$.
If $x>0$ were algebraic and a zero, the real algebraic numbers
$\mathsf F_l(x),\mathsf G_l(x)$ would satisfy
\[
 e^{2ix}=-\frac{\mathsf G_l(x)+i\mathsf F_l(x)}
                  {\mathsf G_l(x)-i\mathsf F_l(x)}.
\]
The denominator is nonzero. This contradicts the Hermite--Lindemann
theorem, which states that $e^z$ is transcendental for each nonzero
algebraic complex number $z$ \cite[\S1.2]{inf:BakerWustholz}.

At a zero $x$ of $J_\nu$, $J_{\nu+1}(x)\ne0$: otherwise the derivative
recurrence and uniqueness for Bessel's equation give $J_\nu=0$
identically. Define the Lommel polynomials by
\[
 \mathsf L^{\mathrm{Lom}}_{0,\mu}(x)=1,\qquad
 \mathsf L^{\mathrm{Lom}}_{1,\mu}(x)=2\mu/x,\qquad
 \mathsf L^{\mathrm{Lom}}_{k+1,\mu}(x)
 =\frac{2(\mu+k)}x\mathsf L^{\mathrm{Lom}}_{k,\mu}(x)
     -\mathsf L^{\mathrm{Lom}}_{k-1,\mu}(x).
\]
The Bessel recurrence implies
$J_{\nu+k}(x)=\mathsf L^{\mathrm{Lom}}_{k-1,\nu+1}(x)J_{\nu+1}(x)$
for $k\ge1$; see \cite{inf:watson}.
For $k\ge2$ this is a nonzero polynomial in $x^{-1}$ with rational
coefficients and leading coefficient $\prod_{h=1}^{k-1}2(\nu+h)$.
A common zero would therefore make $x$ algebraic, contrary to the
first assertion. For $k=1$ the non-coincidence was already proved.
For the primary pair the identity specializes to
\[
 J_{p+3}(x)=\left\{\frac{4(p+1)(p+2)}{x^2}-1\right\}J_{p+1}(x)
 \quad\text{when }J_p(x)=0.
\]
\end{proof}

\begin{lemma}[Shifted congruence approximation]\label{inf:shifted-arithmetic}
The set $\mathcal P_{\mathrm{bd}}$ has infinitely many members in each
of $\Z$ and $\Z+1/2$. For the pairs produced by
Lemma~\ref{inf:congruence} applied to $\alpha_{\mathrm{ar}}+1$, set
$\widetilde P=P-Q$ and $p=(\widetilde P-3)/2$. Along this particular
half-integer selection,
\[
 \limsup\widehat p_{(Q-1)/4}|p-\widehat p_{(Q-1)/4}|
 \le\frac98\alpha_{\mathrm{ar}}+d_{\mathrm{ar}}<31/10.
\]
\end{lemma}
\begin{proof}
Both $P,Q$ are odd, so $\widetilde P$ is even and $2p$ is odd.
Moreover $Q\equiv1\pmod4$ and
\[
 Q|Q\alpha_{\mathrm{ar}}-\widetilde P|
 =Q|Q(\alpha_{\mathrm{ar}}+1)-P|<9/2.
\]
The denominators tend to infinity. Put $m=(Q-1)/4$ and use
\eqref{inf:root-expansion}; multiplication by $\widehat p_m$ gives
\[
 \widehat p_m|p-\widehat p_m|
 \le\tfrac12\widehat p_m|Q\alpha_{\mathrm{ar}}-\widetilde P|
       +d_{\mathrm{ar}}\widehat p_m/A_m+O(A_m^{-1}).
\]
Now $\widehat p_m/Q\to\alpha_{\mathrm{ar}}/2$ and
$\widehat p_m/A_m\to1$. The rational comparison in
\eqref{inf:arithmetic-limsup} proves the bound. Positivity of $p$ and
$m\ge m_0$ hold eventually. Lemma~\ref{inf:half-integer-noncoincidence}
excludes zero distance. The integer selection is
Theorem~\ref{inf:near-integer}.
\end{proof}

\begin{proposition}[The half-lattice analytic chain]\label{inf:parity-chain}
The analytic statements listed in the table below that use a
near-primary input hold for
$p\in\mathcal P_{\mathrm{bd}}$, uniformly for $0<\tau\le20$ and
$r\in[1/4,3/8]$, with integer $a,b$ for each Euclidean conclusion.
The remaining analytic statements hold under their displayed real
parameter hypotheses.
\end{proposition}
\begin{proof}
The chain is given below, with its integer objects and the identities
that determine the permissible parameters. Each row is applied under
the hypotheses of its indicated statements.
\begingroup\small\normalfont
\renewcommand{\arraystretch}{1.12}
\begin{longtable}{@{}>{\raggedright\arraybackslash}p{.32\textwidth}>{\raggedright\arraybackslash}p{.62\textwidth}@{}}
\toprule Statements & Parameter dependence\\\midrule
\endfirsthead
\toprule Statements & Parameter dependence\\\midrule
\endhead
Lemma~\ref{inf:reduction} &
The coordinate Laplacian and beta measure use integer $a,b$ and
$p=(a+b)/2$. The invariant angular multiplicity is one.\\
Lemma~\ref{inf:angular-formulas} &
Rodrigues integration and rising factorials use positive real
$\ab,\bb$; $j$ is the integer polynomial degree.\\
Lemma~\ref{inf:angular-sup} &
The endpoint comparison needs $\bb-1\ge\ab-1\ge0$.
Bernstein applies on great circles in every integer dimension.\\
Lemma~\ref{inf:solid-harmonics} &
$t^jG_j$ is a polynomial in $q_X,q_Y$.
The regular Bessel mode is this polynomial times an entire power
series in $\varrho^2$, including the origin and axes.\\
Proposition~\ref{inf:ball-spectrum} &
The radial order $p-1+2j>1$ is real; its singular solution is
outside the Friedrichs form domain. Integer blocks identify the
Euclidean invariant realization.\\
Lemma~\ref{inf:radial-polynomials} &
The radial parameter $2j+p-1$ is real. Rodrigues and generalized
binomial identities require only integer polynomial indices.\\
Lemmas~\ref{inf:bessel-spacing}, \ref{inf:bessel-order} &
The radial ODE, Poincar\'e inequality and Watson integral have
their stated real-order hypotheses.\\
Lemmas~\ref{inf:debye-first}, \ref{inf:debye-finite} &
The outgoing Volterra equation is for real $\nu>0$;
only its finite expansion order is integral.\\
Lemma~\ref{inf:window-isolation} &
The recurrence shifts the real Bessel order by two.
Its frequency spacing concerns all radial indices.\\
Lemma~\ref{inf:phase-strip}, Theorem~\ref{inf:global-count} &
The lattice coordinates are the integer angular index and
integer phase index; their offsets may depend on real $p$.\\
Lemma~\ref{inf:congruence}, Theorem~\ref{inf:near-integer},
Lemma~\ref{inf:shifted-arithmetic} &
Apply the same irrational approximation to $\alpha_{\mathrm{ar}}$
or $\alpha_{\mathrm{ar}}+1$. The latter changes $P$ to $P-Q$
and preserves the error product.\\
Proposition~\ref{inf:primary-expansion} &
The root expansion and uniqueness use real-order continuation.
Bourget and Lemma~\ref{inf:half-integer-noncoincidence}
exclude both half-lattice coincidences.\\
Lemma~\ref{inf:detuning} &
The exact rational trace identity and the real-order derivative
give the same multiplicative detuning relation.\\
Lemmas~\ref{inf:riccati-gap}, \ref{inf:low-recurrence} &
The finite integer order shifts leave the base order real.
The period-six sequence comes from the limiting recurrence.\\
Proposition~\ref{inf:low-exclusion} &
All relevant zeros are continued on a real interval in their
orders; the selected endpoint may be a half-integer.\\
Lemma~\ref{inf:angular-weighted} &
The entry and derivative formulas are rational identities for
positive real beta parameters with sum $p$.\\
Lemma~\ref{inf:angular-positive} &
The coefficient certificates hold for all real $p\ge64$.
The recurrence-ratio bounds use $\ab,\bb\ge16$.\\
Proposition~\ref{inf:flow-form} &
Axis cutoffs use $\ab,\bb>1$ and the origin cutoff $p>1$.
Integer blocks identify the weighted and Euclidean closures.\\
Lemmas~\ref{inf:flow-relative}, \ref{inf:transport} &
The band estimates and Green identity use real Jacobi parameters
and radial Bessel norms; both angular parities are retained.\\
Lemma~\ref{inf:flow-complement} &
The complementary inverse is radial spectral calculus and a
relative-form Neumann series on every complementary radial state.\\
Lemma~\ref{inf:flow-low-block} &
The finite recurrence gives seed offset $-\theta/(2p)$ and
index-five offset $-72$ in either ambient parity.\\
Lemma~\ref{inf:ordered-eigenvalues}, Theorem~\ref{inf:flow-theorem} &
Finite Riesz reduction and the weighted derivative margin apply
on the common real-parameter form domain, in the stated narrow band.\\
Theorem~\ref{inf:integer-gap} &
Integers in $[p/2,3p/4]$ number at least $p/4-1$ and have
$\varsigma$ spacing at least $1/(8p)$.
$2p\in\Z$ ensures integer $b=2p-a$.\\
Theorem~\ref{inf:centres} &
The finite recursion uses $\eps=1/p$, fixed degree indices,
the divisor $-1/8$, positive $c_j$ and powers of $r(1-r)$.\\
Lemma~\ref{inf:analytic-coefficients},
Corollary~\ref{inf:centre-coefficient-defect} &
The analytic coefficient transfer uses any positive measure and
$0<\mathsf H_j(1)=(\bb)_j/(p+j-1)_j\le1$.\\
Proposition~\ref{inf:centre-gap-transfer} &
Indexwise min--max and radial dilation hold in every integer dimension;
the normalized error is $O_N(\tau p^{-4})$.\\
Definition~\ref{inf:spaces}, Lemma~\ref{inf:radial-stencils},
Proposition~\ref{inf:module} &
The grade exponent is real. Gasper uses real Jacobi parameters;
the integer raising/lowering shifts are angular degree differences.\\
Lemma~\ref{inf:poisson-columns}, Proposition~\ref{inf:poisson-bounds} &
The radial columns are rational in the real parameter $\beta$;
the separate harmonic source column remains included.\\
Lemmas~\ref{inf:raising-coupling}, \ref{inf:clamped-smoothing} &
Positive radial stencils and their coupling use real $\beta\ge0$;
the smoothing gain acts on the clamped layers $s\ge1$.\\
Lemma~\ref{inf:gradient-profiles},
Proposition~\ref{inf:grouped-shapes} &
The solid product identity uses integer solid degrees and real $p$;
the grouped high-shape estimates have clamped source domains.\\
Lemmas~\ref{inf:double-lift}, \ref{inf:derivative-moments} &
The two lift columns and derivative columns use real radial
parameters; their reference normal factor is $R=O(p)$.\\
Lemma~\ref{inf:radial-centre-norm},
Proposition~\ref{inf:centre-source-norm} &
Rodrigues--Sonine, Lommel--Parseval and the positive Bessel ratio
hold for real orders. The head cutoff satisfies $\nu+1+2S\ge2R$;
composition uses the finite list of known centre modes.\\
Lemma~\ref{inf:pullback}, Proposition~\ref{inf:nonlinear-map} &
The Jacobian is $q^{2p-1}d=q^{n-1}d$ at integer $n$.
The nonlinear map has clamped source domain $\Gc_p\times\Ec_1$.\\
Proposition~\ref{inf:residual-estimates},
Theorem~\ref{inf:material-isomorphism} &
The double lift, normal Hessian and material identities use the
same real parameters; the $L^2$ completion is at fixed integer $n$.\\
Lemmas~\ref{inf:mixed-algebra}, \ref{inf:mixed-Poisson},
\ref{inf:mixed-domain}, Theorem~\ref{inf:relative-graph} &
The radial Hilbert blocks use real weights; density uses real
binomial powers. Hilbert and coefficient estimates retain their
respective graph domains and all radial layers.\\
Theorem~\ref{inf:local-count}, Corollary~\ref{inf:clusters} &
The integer phase lattice and real turning-point ODE give the
same local count and cluster cardinality.\\
Lemma~\ref{inf:cluster-inverse},
Proposition~\ref{inf:mixed-inverse} &
The cofactor argument is finite linear algebra, followed by the
radial graph reconstruction in both spaces.\\
Lemma~\ref{inf:head-tail}, Theorem~\ref{inf:coefficient-inverse} &
The head has at most $K_{\mathrm{head}}p+1$ radial indices;
the infinite tail uses the same all-layer Poisson columns.\\
Theorem~\ref{inf:newton-threshold}, Corollary~\ref{inf:Newton-zero} &
The contraction uses the displayed powers of $p,\tau$ in the same
real Banach spaces; its centre order is fixed first.\\
Lemmas~\ref{inf:geometry-convexity}, \ref{inf:geometry-nonball},
Propositions~\ref{inf:classical-solution}, \ref{inf:Pompeiu} &
The sphere and elliptic arguments require integer $n,a,b$.
The regular solid expansion covers the axes and origin in each parity.\\
\bottomrule
\end{longtable}
\endgroup
The only lattice changes are thus the non-coincidence theorem,
the arithmetic choice of $p$, the integer split grid, and the
integer block realization. All estimates in the table retain their
stated constants on this half-lattice.
\end{proof}

\section{The quartic domain and finite-order centres}\label{inf:sec-centres}

In this section $r$ ranges over the fixed compact interval
$[1/4,3/8]$. All finite-order constants are uniform on this interval
and for $|\theta|\le20$, including $\theta=0$ when analytic
extensions are discussed. The accuracy order is always fixed before
$p$ tends to infinity.

\subsection{The quartic profile}

The monic degree-two angular polynomial is
\begin{equation}\label{inf:H2}
 \mathsf H_2(x)=x^2-\frac{2(\ab+1)}{p+2}x
                 +\frac{\ab(\ab+1)}{(p+1)(p+2)}.
\end{equation}
Its homogeneous spherical degree is four. Define its self-coupling by
\begin{equation}\label{inf:seed-coupling}
 g_2=[\mathsf H_2]\mathsf H_2^2
 =\frac{\int \mathsf H_2^3\,d\mu_r}{\int \mathsf H_2^2\,d\mu_r}.
\end{equation}
The bracket denotes the coefficient in the monic orthogonal expansion.

\begin{lemma}[Quartic normalization]\label{inf:quartic-normalization}
Put $\varsigma=(r-1/2)^2\in[1/64,1/16]$. Then
\begin{equation}\label{inf:g2-exact}
 g_2=\frac{p-1}{(p+1)(p+5)}
 -\frac{4p^2(p^2-9p-16)\varsigma}{(p+1)(p+2)^2(p+4)(p+5)}.
\end{equation}
For $p\ge64$,
\begin{equation}\label{inf:g2-bounds}
 \frac9{16}<pg_2<1,\qquad
 \nrm{\mathsf H_2}_\infty<\frac9{16},\qquad
 \nrm{T_\varsigma}_\infty<1,\quad T_\varsigma(x)=\frac{\mathsf H_2(x)}{pg_2}.
\end{equation}
Also $g_2=4r(1-r)p^{-1}+O(p^{-2})$ uniformly in $r$.
\end{lemma}
\begin{proof}
The beta moments are $\int x^h\,d\mu_r=(\ab)_h/(p)_h$.
They give the two orthogonality conditions for \eqref{inf:H2}.
Multiplying the three-term recurrence twice gives
\[
 \int \mathsf H_2P_j^2\,d\mu_r
 =\mathfrak a_j^2+\mathfrak a_{j+1}^2+\mathfrak b_j^2
 -\frac{2(\ab+1)}{p+2}\mathfrak b_j
 +\frac{\ab(\ab+1)}{(p+1)(p+2)}.
\]
At $j=2$, substitute \eqref{inf:recurrence-a}--\eqref{inf:recurrence-b}
and $\ab=p/2+p(r-1/2)$. The terms linear in $r-1/2$ cancel;
the constant term is $(p-1)/((p+1)(p+5))$, and the coefficient of
$(r-1/2)^2$ is the second coefficient in \eqref{inf:g2-exact}.
Since $P_2$ is a multiple of $\mathsf H_2$, this gives $g_2$.

Write $g_2=c_2^{\mathrm{diag}}+d_2^{\mathrm{diag}}\varsigma$
with the coefficients in \eqref{inf:g2-exact}. For $p\ge64$,
$d_2^{\mathrm{diag}}<0$, $|pd_2^{\mathrm{diag}}|<4$, and
$pc_2^{\mathrm{diag}}\ge4032/4485>13/16$.
The last inequality persists for larger $p$ because the derivative of
$p(p-1)/((p+1)(p+5))$ has numerator $7p^2+10p-5>0$.
Thus $pg_2>13/16-4/16=9/16$, while $pg_2<pc_2^{\mathrm{diag}}<1$.
The two positive endpoint values of $\mathsf H_2$ are less than $9/16$,
using $\ab,\bb\le3p/4$. Its negative minimum has magnitude
$(\ab+1)(\bb+1)/((p+1)(p+2)^2)\le1/(4(p+1))$.
This proves \eqref{inf:g2-bounds}. Expansion of the rational formula
gives $1-4\varsigma=4r(1-r)$ as its leading coefficient.
\end{proof}

The physical quartic radius and its normalized version are
\begin{equation}\label{inf:quartic-domain}
 R+h_{\mathrm q}(x),\qquad
 h_{\mathrm q}=-\frac{2\chi}{g_2}\mathsf H_2,\qquad
 1+f_{\mathrm q}(x)=1-\frac{\theta}{2R^2}T_\varsigma(x).
\end{equation}
We denote the corresponding physical domain by $\Omega_{\mathrm q}$.
For real block parameters the same formula denotes its weighted
meridian domain. The scalar $\theta$ and the radius $R$ stay fixed
when $\varsigma$ varies.

\subsection{An analytic finite-order construction}

\begin{theorem}[Finite-order centres]\label{inf:centres}
Fix an integer $N\ge4$. For every prescribed finite radius
$R_{\mathrm{an}}>0$, put $U=\{x\in\C:|x|\le R_{\mathrm{an}}\}$.
There exist $C_{N,U},P_{N,U}<\infty$ such that the
following assertion holds. Suppose $p\in\tfrac12\Z$, $p\ge P_{N,U}$,
$J_p(R)=0$, $R\in[2p-1,2p+4]$, and
$0<|\theta|\le20$ with \eqref{inf:detuning-definition}.
For every $r\in[1/4,3/8]$ one can choose a polynomial $h_N(x)$ and
a finite sum $u_N$ of regular Helmholtz modes. For real block
parameters the equation means the weighted operator
\eqref{inf:invariant-laplacian}; at integer blocks it is the
Euclidean equation. These choices satisfy
$(\Delta+1)u_N=0$ exactly and
\begin{equation}\label{inf:centre-complex-defect}
 \sup_{x\in U}\left(|u_N(R+h_N(x),x)-1|
       +|\partial_\varrho u_N(R+h_N(x),x)|\right)
 \le C_{N,U}\theta^2p^{-N}.
\end{equation}
The number and degrees of the modes depend on $N$, not on $p,r,\theta$.
For every fixed polynomial coefficient norm and every fixed derivative
order $k$ on $U$, with constants $C_{N,U,k}$ and threshold
$P_{N,U,k}$,
\begin{align}
 h_N&=-\frac{2\chi}{g_2}\mathsf H_2+O_{N,U}(|\theta|p^{-3}),
 \label{inf:centre-profile}\\
 [\mathsf H_2]h_N&=-\frac{\theta}{16r(1-r)p}
                                  +O_{N,U}(|\theta|p^{-2}).
 \label{inf:centre-witness}
\end{align}
The radius $R_{\mathrm{an}}$ can be chosen arbitrarily large before
the constants and threshold are fixed. In particular the statement
applies as follows: fix $\rho'>2$ and $d>\rho'$, and then
choose one $R_{\mathrm{an}}>1+d$.
\end{theorem}

\begin{proof}
We construct a finite jet with more terms than the requested accuracy.
Put $\eps=1/p$, write $R=\rho/\eps$, and set
$Z=\rho^2$ and $S_0=4(1+3\eps+2\eps^2)$.
Equation~\eqref{inf:trace-ratio-exact} is equivalent to
\begin{equation}\label{inf:rho-equation}
 2Z(Z-S_0)+\theta\eps^2(1+\eps)(2Z-S_0)=0.
\end{equation}
At $(\eps,Z)=(0,4)$ its derivative in $Z$ is eight.
The analytic implicit-function theorem gives a unique root near four,
uniformly for $|\theta|\le21$. The supplied radius satisfies
$(R/p)^2\to4$ and, by \eqref{inf:trace-ratio-exact},
equation~\eqref{inf:rho-equation} with $Z=(R/p)^2$.
Thus $Z$ is this root for sufficiently large $p$. The other quadratic root is
$O(\theta\eps^2)$. Take the positive square root near two.
Then
\begin{equation}\label{inf:chi-analytic}
 \rho=2+3\eps+O(\eps^2),\qquad
 \chi=\frac{\theta\eps^2}{4\rho}
      =\frac18\theta\eps^2+O(\theta\eps^3).
\end{equation}
Here and below analyticity means on an open complex neighbourhood of
the stated compact real parameter set.

For finitely many angular indices, the boundary traces of regular
Bessel modes are computed from
\begin{equation}\label{inf:finite-q-recurrence}
 q_{-1}=1,\quad q_0=0,\quad
 q_{k+1}=\frac{2(1+k\eps)}\rho q_k-q_{k-1}.
\end{equation}
This finite recurrence is analytic in $\eps,\theta$.
Its limiting zero indices give the resonant angular set
$\mathcal I_{\mathrm{res}}=\{2,5,8,\ldots\}$.
For $j\notin\mathcal I_{\mathrm{res}}$ normalize the mode $f_j$ by $f_j(R)=1$;
for $j\in\mathcal I_{\mathrm{res}}$ normalize it by $f_j'(R)=1$.
The latter convention gives $f_2(R)=\chi$. For $j\ge5$ in
$\mathcal I_{\mathrm{res}}$,
\begin{equation}\label{inf:resonant-trace}
 f_j(R)=c_j\eps+O_j(\eps^2),\qquad
 c_j=\frac{(2j-1)(2j-4)}3>0.
\end{equation}
Indeed the first correction of \eqref{inf:finite-q-recurrence} at
$k=3h$ is $(-1)^h3h(h-1)\eps$, while the corresponding normal
trace tends to $(-1)^h$. Put $3h=2j-1$ to obtain
\eqref{inf:resonant-trace}. The change in $\rho$ due to $\theta$
starts at order $\eps^2$, so this coefficient is independent of
$\theta$. The nonresonant normalization denominators tend to $1$ or
$-1$ and are units in the analytic coefficient ring.

All collar derivatives are supplied by
\begin{equation}\label{inf:collar-ode}
 f_j''+\frac{(2-\eps)}{\rho+\eps y}f_j'
 +\left(1-\frac{2j(2j\eps^2+2\eps-2\eps^2)}
                    {(\rho+\eps y)^2}\right)f_j=0,
 \qquad y=\varrho-R.
\end{equation}
Its coefficients are analytic near $\eps=0$ on a fixed collar.
The ball mode has traces $u_0(R)=1$, $u'_0(R)=0$, $u''_0(R)=-1$.
The collar construction is a local analytic family.
At the actual inputs $J_p(R)=0$ its traces
are precisely those of the regular Bessel modes; uniqueness of the
ordinary differential equation identifies them on the collar and
then throughout their regular continuation.

We next record the angular filtration that makes the recursion finite.
For fixed $i,j$ write
\begin{equation}\label{inf:monic-product-valuation}
 \mathsf H_i\mathsf H_j=\sum_{k\le i+j}\mathfrak c_{ij}^k\mathsf H_k,\qquad
 \mathfrak c_{ij}^k=O\left(\eps^{\lceil(i+j-k)/2\rceil}\right).
\end{equation}
To prove this, use the monic recurrence
$\mathsf H_{j+1}=(x-\mathfrak b_j)\mathsf H_j-\mathfrak a_j^2\mathsf H_{j-1}$.
For fixed indices \eqref{inf:recurrence-a}--\eqref{inf:recurrence-b}
give $\mathfrak b_j-r=O(\eps)$ and
$\mathfrak a_j^2=O(\eps)$.
Induction shows that the coefficient of $(x-r)^k$ in $\mathsf H_j$ has
valuation at least $\lceil(j-k)/2\rceil$. Descending elimination
proves the same valuation for the inverse monic triangular
connection. Multiplying and composing these two connections proves
\eqref{inf:monic-product-valuation}, since a sum of ceilings is at
least the ceiling of the sum. In particular,
\begin{equation}\label{inf:seed-square}
 \mathsf H_2^2=\mathsf H_4+O(\eps)\mathsf H_3+g_2\mathsf H_2+O(\eps^2)\mathsf H_1+O(\eps^2)\mathsf H_0,
 \qquad g_2=4r(1-r)\eps+O(\eps^2).
\end{equation}

Write the formal field and boundary as
\[
 u=u_0+\sum_j a_jf_j\mathsf H_j,\qquad \varrho=R+h(x).
\]
Every coefficient of $\eps$ is a finite polynomial in $x$.
Eliminate $h$ from the radial Neumann condition first. Its derivative
in $h$ at zero is $u_0''(R)=-1$, so at each order the new coefficient
of $h$ is determined by the already known coefficients and the new
field coefficients, using the unit divisor $u_0''(R)=-1$.
In the Dirichlet equation eliminate $a_j$ for $j\notin\mathcal I_{\mathrm{res}}$.
Its diagonal linear coefficient is $f_j(R)=1$. Thus its new
coefficients are also determined by unit divisions.
This elimination includes the $j=0$ field coefficient.

Parametrize the remaining coefficients by
\begin{equation}\label{inf:resonant-scaling}
 a_2=\theta\eps \SeedA,\qquad
 a_j=\theta^2\eps^3 Z_j\quad(j\in\mathcal I_{\mathrm{res}},\ j\ge5).
\end{equation}
After the two eliminations the residual Dirichlet coefficients
$F_j$ have the form
\begin{align}
 \frac{F_2}{\theta^2\eps^3}
 &=\frac{\SeedA}{8}+2r(1-r)\SeedA^2+O(\eps),\label{inf:seed-equation}\\
 \frac{F_j}{\theta^2\eps^4}
 &=c_jZ_j+B_j(\SeedA,\theta,r)+O(\eps),
 \qquad j\in\mathcal I_{\mathrm{res}},\ j\ge5.\label{inf:other-equations}
\end{align}
Here $B_j$ is a finite polynomial and is zero except possibly at
$j=5,8$. These divisions are identities of formal series, including
at $\theta=0$.

The linear seed term is
$a_2\chi=\theta^2\eps^3 \SeedA/8+O(\theta^2\eps^4)$.
Eliminating $h$ gives quadratic Dirichlet term
$a_2^2\mathsf H_2^2/2$; its seed coefficient is the second term of
\eqref{inf:seed-equation} by \eqref{inf:seed-square}.
A pure cubic term has angular degree at most six; reaching index two
loses four degrees and contributes two further powers of $\eps$.
The leading nonresonant coefficient is $a_4=-a_2^2/2+O(\eps^3)$,
and its product with the seed has the same loss when projected to
index two. Thus neither contributes to the leading seed equation.
Quadratic terms have no nonseed resonant index. Cubic terms first
reach index five after losing one angular degree, which supplies one
extra $\eps$. At total order four the only possible new resonant
indices are five and eight. A new $Z_j$ contributes linearly through
\eqref{inf:resonant-trace}, giving $c_jZ_j$.
Its product with the seed has top index $j+2$, which is nonresonant;
reaching a resonant index at most $j$ loses at least two degrees and
therefore at least one extra $\eps$. Contributions through a newly
eliminated nonresonant coefficient have the same extra factor after
its unit division. At order zero in \eqref{inf:other-equations},
$Z_j$ enters linearly as $c_jZ_j$. Every other term has at least two
amplitudes or the factor $a_2\chi$, so it is divisible by $\theta^2$.

Choose the nonzero root
$\SeedA_0=-1/(16r(1-r))$ of \eqref{inf:seed-equation}.
The derivative of its left side with respect to $\SeedA$ at this root is
$-1/8$. After choosing $\SeedA_0$, determine
$Z_{j,0}=-B_j(\SeedA_0,\theta,r)/c_j$; these have finite support.
Work in the ring of formal
$\eps$ series whose coefficients are polynomials in $x$, with
coefficients analytic in $r,\theta$ on the indicated compact sets.
A series is required to have a finite polynomial coefficient at
every order.
For a fixed list of resonant amplitude coefficients the two unit
eliminations have linearization, in the variables $(h,a_{\rm nr})$,
\[
 (\dot h,\dot a_{\rm nr})\longmapsto
 \left(-\dot h+\sum_{j\notin\mathcal I_{\mathrm{res}}}
       f'_j(R)|_{\eps=0}\,\dot a_j\mathsf H_j|_{\eps=0},
       \dot a_{\rm nr}\right).
\]
Its inverse is triangular and sends a finite polynomial forcing
to finite polynomial coefficients. At each order, the remaining
terms involve only coefficients at lower orders. This proves
existence and uniqueness of those unit eliminations by ordinary
coefficient induction. It also proves preservation of finite
support: multiplication and the monic connection use only finite
polynomials, and the inverse displayed here adds no angular
indices not present in its forcing. Dependence of the exact
collar traces on $\eps$ contributes only previously determined
lower-order coefficients in this calculation.

After these eliminations let
$\mathcal F_2=F_2/(\theta^2\eps^3)$ and
$\mathcal F_j=F_j/(\theta^2\eps^4)$ for nonseed resonant $j$.
Their constant terms are the right sides of
\eqref{inf:seed-equation}--\eqref{inf:other-equations} without the
remainders. At $(\SeedA,Z)=(\SeedA_0,Z_0)$ their leading Jacobian is
\begin{equation}\label{inf:formal-Jacobian}
 \mathcal J_0(\dot{\SeedA},\dot Z)=
 \left(-\frac18\dot{\SeedA},
       \{\partial_{\SeedA}B_j(\SeedA_0,\theta,r)\dot{\SeedA}+c_j\dot Z_j\}_{j\ge5,
                                                    j\in\mathcal I_{\mathrm{res}}}\right).
\end{equation}
The off-diagonal column in this matrix has finite support,
since $B_j=0$ outside $j=5,8$. Each $c_j$ is positive and fixed.
Thus its inverse preserves finite support, using only the fixed
nonzero divisors $-1/8$ and $c_j$.

Suppose all coefficients of $\SeedA,Z$ below order $v\ge1$ have
already been selected and the reduced equations vanish below
that order. Substitute those truncated series with the unknown
order-$v$ coefficients set to zero, and denote the coefficients
of $\eps^v$ in $\mathcal F_2,\mathcal F_j$ by
$e_{2,v},e_{j,v}$. They are explicitly defined by finite
coefficient extraction from the collar Taylor formula, the
polynomial product connection, and the two unit eliminations.
Only finitely many $e_{j,v}$ are nonzero.
The new coefficients occur only through
\eqref{inf:formal-Jacobian}: every other occurrence is multiplied
by an additional $\eps$, or by a coefficient of positive order
of $\SeedA-\SeedA_0$ or $Z-Z_0$. Therefore the unique update is
\begin{equation}\label{inf:formal-update}
 \SeedA_v=8e_{2,v},\qquad
 Z_{j,v}=-\frac{e_{j,v}
       +\partial_{\SeedA}B_j(\SeedA_0,\theta,r)\SeedA_v}{c_j}.
\end{equation}
Set the other coefficients to zero. The update is finite, cancels
exactly the order-$v$ coefficient, and leaves every lower order
unchanged. Recompute the unit-eliminated coefficients to that
order. This completes an induction producing every finite jet.

All operations in \eqref{inf:formal-update} are additions,
multiplications, coefficient extractions and divisions by the
nonzero quantities already listed. The resulting coefficients
are analytic in a common neighbourhood of the compact real
parameter set for each finite jet. The neighbourhood may shrink
with its order. At a fixed order only finitely many collar
normalizations and monic connections occur, so a positive common
neighbourhood exists for each finite order.

Through order $\eps^4$ the seed equation is the expansion of
$\chi a_2+g_2a_2^2/2$. All cubic seed projections start at order
$\eps^5$ by the degree loss just described. Hence the chosen root
agrees with $-2\chi/g_2$ through order $\eps^2$.
The order-$\theta^2\eps^2$ shape coefficient vanishes by the
following cancellation.
At that order the index-four amplitude is $-a_2^2/2$, whose
physical normal derivative tends to $-1$.
The quadratic Neumann correction is
$-(2p-1)a_2^2\mathsf H_2^2/(2R)$ and $(2p-1)/R\to1$.
Their index-four contributions cancel.
All lower indices of $\mathsf H_2^2$ gain at least one $\eps$, and the other
resonant amplitudes start at order $\theta^2\eps^3$.
It follows that $h-a_2\mathsf H_2=O(\theta^2\eps^3)$ at these orders,
proving \eqref{inf:centre-profile} for any sufficiently long finite
jet with an $O_N(|\theta|\eps^3)$ remainder.
A finite truncation of the linear-in-$\theta$ part is included in
this latter bound. Inserting \eqref{inf:chi-analytic} and
\eqref{inf:g2-exact} gives \eqref{inf:centre-witness}.

Choose the recursion order greater than $N+4$ and retain its finite
set of modes. Evaluate the truncated coefficients in the exact
finite-$p$ polynomials $\mathsf H_j$ and exact normalized radial functions.
Restoring the factors $\theta^2\eps^3$ and $\theta^2\eps^4$
in \eqref{inf:seed-equation} and \eqref{inf:other-equations}
recovers the original resonant Dirichlet coefficients. The two
unit eliminations preserve the chosen Taylor jet of the Neumann
and nonresonant Dirichlet equations. Thus a jet of order greater
than $N+4$ makes every original defect coefficient through order
$N$ vanish.
Choose the truncations so that $h-a_2\mathsf H_2$ has a factor
$\theta^2$, using the identical truncated $a_2$ in the field and
the boundary. The higher powers of $\theta$ already present in
$a_2$ are retained in this convention.
For $x$ in any fixed bounded complex neighbourhood $U'\supset U$,
this polynomial has size $O_{N,U'}(|\theta|\eps)$.
For sufficiently small $|\eps|$ its values lie in a fixed collar,
where the integral equation for \eqref{inf:collar-ode} gives jointly
analytic solutions with uniform bounds. The denominators
$\rho+\eps y$ remain away from zero, and only finitely many mode
normalizations occur. Thus both boundary defects are analytic on a
common neighbourhood of the compact parameter set and $U'$.
Their Taylor coefficients in $\eps$ vanish through order $N$ by the
finite induction, as polynomial identities in $x$.

Both defects have an analytic factor $\theta^2$.
At $\theta=0$ the field and domain are the ball. The
linear Neumann defect is
$-\partial_\theta h+(\partial_\theta a_2)\mathsf H_2=0$ by the identical
seed choice; the linear Dirichlet defect is zero since
$u'_0(R)=0$ and $f_2(R)=\chi$ itself has a factor $\theta$.
All other terms have at least two such factors.
Divide the analytic defects by $\theta^2$ and apply Taylor's formula
with remainder in $\eps$ on a smaller fixed complex disk. Compactness
in $r,\theta,x$ bounds the remainder by $C_{N,U}|\eps|^N$.
Cauchy's formula on $U'$ gives every stipulated fixed angular
derivative bound. This proves \eqref{inf:centre-complex-defect} and
the remainder statements. At actual integer blocks the finite regular field is
analytic across axes and origin by Lemma~\ref{inf:solid-harmonics}.
\end{proof}

\begin{lemma}[Analytic defects give weighted coefficients]
\label{inf:analytic-coefficients}
Let $\mu$ be any positive measure on $[0,1]$ with infinite support,
$\mathsf H_j$ its monic orthogonal polynomials, and suppose $0<\mathsf H_j(1)\le1$.
For a complex-valued function $F$ analytic on a neighbourhood of
$U_d=\{z:\dist(z,[0,1])\le d\}$, its coefficients in
$G_j=\mathsf H_j/\mathsf H_j(1)$ satisfy
\begin{equation}\label{inf:cauchy-coefficients}
 |F_j|\le d^{-j}\sup_{U_d}|F|.
\end{equation}
Consequently, for each fixed $k\ge0$ and $1<\rho'<d$,
\begin{equation}\label{inf:weighted-analytic-coefficients}
 \sum_{j\ge0}(1+j)^k(\rho')^j|F_j|
 \le\left(\sum_{j\ge0}(1+j)^k(\rho'/d)^j\right)\sup_{U_d}|F|.
\end{equation}
The constants do not depend on the measure. These assertions apply
to all the beta measures used here.
\end{lemma}
\begin{proof}
For $j=0$ the coefficient is the normalized integral of $F$, so
\eqref{inf:cauchy-coefficients} follows directly. Let $j\ge1$.
The zeros $x_1,\ldots,x_j$ of a monic orthogonal polynomial are simple
and lie in $(0,1)$. Interpolate $F$ at those nodes by a polynomial
$L$ of degree at most $j-1$. The divided-difference identity gives
\[
 F(x)-L(x)=F[x_1,\ldots,x_j,x]\mathsf H_j(x).
\]
The integral formula for divided differences, valid also for
complex-valued functions, gives
\[
 |F[x_1,\ldots,x_j,x]|\le\frac1{j!}
                                  \sup_{[0,1]}|F^{(j)}|.
\]
Orthogonality to $L$ therefore gives
\[
 \frac{|\langle F,\mathsf H_j\rangle|}{\nrm{\mathsf H_j}_{L^2(\mu)}^2}
 \le\frac{\sup_{[0,1]}|F^{(j)}|}{j!}
 \le d^{-j}\sup_{U_d}|F|,
\]
where the last inequality is Cauchy's estimate on each radius-$d$
disk centred on the interval. Multiplication by $\mathsf H_j(1)\le1$
proves \eqref{inf:cauchy-coefficients}; summation gives
\eqref{inf:weighted-analytic-coefficients}.
For our measure \eqref{inf:angular-jacobi} shows
$\mathsf H_j(1)=(\bb)_j/(p+j-1)_j\le1$.
\end{proof}

\begin{corollary}[The coefficient form of the centre defects]
\label{inf:centre-coefficient-defect}
Fix $N,k$ and $\rho'>1$. The centre in Theorem~\ref{inf:centres}
can be chosen so that its physical Dirichlet and radial Neumann
boundary defects $D_N,N_N^{\mathrm{phys}}$, expanded in $G_j$, obey
\begin{equation}\label{inf:centre-coeff-defect-bound}
 \sum_j(1+j)^k(\rho')^j
       (|D_{N,j}|+|N_{N,j}^{\mathrm{phys}}|)
 \le C_{N,k,\rho'}\theta^2p^{-N}.
\end{equation}
The constants and threshold are uniform in $r$ and $0<|\theta|\le20$.
\end{corollary}
\begin{proof}
Choose $d>\rho'$ and then one fixed $R_{\mathrm{an}}>1+d$.
The set $U_d$ lies in the disc of radius $R_{\mathrm{an}}$.
Use the analytic remainder established by the finite-jet construction
in Theorem~\ref{inf:centres} on that disc, and apply
Lemma~\ref{inf:analytic-coefficients} separately to the two defects.
The radius choice changes only the constants and threshold for
this fixed jet. For the Newton application choose $\rho'>2$ once.
\end{proof}

\begin{proposition}[Transfer of an invariant spectral gap]
\label{inf:centre-gap-transfer}
Fix the centre order $N$ and suppose the actual quartic domain has
invariant Dirichlet gap at least $c_{\mathrm q}\tau p^{-5/3}$ about
$R^2$ in unit-ball normalization, with fixed $c_{\mathrm q}>0$.
Then, for all sufficiently large $p$ depending on $N,c_{\mathrm q}$,
the centre domain $\Omega_N/R$ has gap at least
$(c_{\mathrm q}/2)\tau p^{-5/3}$. The threshold is independent of
the admissible split and of $0<\tau\le20$.
\end{proposition}
\begin{proof}
By \eqref{inf:centre-profile}, the two normalized radius functions
differ by $O_N(\tau p^{-4})$ and are both $1+O_N(\tau p^{-2})$.
Thus, for $\delta\le C_N\tau p^{-4}$,
\[
 (1-\delta)\Omega_{\mathrm q}/R\subset\Omega_N/R
 \subset(1+\delta)\Omega_{\mathrm q}/R.
\]
Zero extension preserves the invariant sector, so the min--max
principle gives, for every ordered index $i$,
\[
 (1+\delta)^{-2}\lambda_i(\Omega_{\mathrm q}/R)
 \le\lambda_i(\Omega_N/R)
 \le(1-\delta)^{-2}\lambda_i(\Omega_{\mathrm q}/R).
\]
For any eigenvalue within a fixed window of $R^2$, its counterpart
is $O(p^2)$, and the difference is bounded by $C_N\tau p^{-2}$.
This is smaller than $(c_{\mathrm q}/2)\tau p^{-5/3}$ once
$p^{1/3}>2C_N/c_{\mathrm q}$. An eigenvalue inside the proposed
centre gap would then put its ordered quartic counterpart strictly
inside the original gap, a contradiction. The factor $\tau$ cancels
from the required threshold.
\end{proof}

\section{The quartic spectral flow}\label{inf:sec-flow}

Throughout this section $p\ge64$, $\varsigma\in I=[1/64,1/16]$,
$r=1/2-\sqrt \varsigma$, $R\in[2p-1,2p+4]$, and $0<\tau=|\theta|\le20$.
Unless indicated otherwise, constants are absolute and independent of
$p,\varsigma,R,\theta$. The fixed window used for the flow argument is
$W_{\mathrm{flow}}$. The coefficient inverse uses the separately
defined windows in Appendix~\ref{inf:sec-notation}.

\subsection{Angular matrices and their weighted variation}

Identify the angular Hilbert spaces for different $\varsigma$ by their
orthonormal bases $P_j(\varsigma)$. In these coordinates the angular operator
is the fixed diagonal matrix $\Lambda=\diag(\lambda_j)$, and
multiplication by $x$ is the matrix $X_\varsigma$ in
\eqref{inf:jacobi-matrix}. Define
\[
 T_\varsigma=\frac{\mathsf H_2(X_\varsigma)}{pg_2},\qquad
 w_j=\frac{j+1}{p+j+1},\qquad \mathsf W=\diag(w_j).
\]
The same notation $T_\varsigma(x)$ denotes its scalar polynomial before
identification. All estimates involving $\mathsf W^{-1}$ initially
mean estimates on finite sequences, extended by boundedness.

\begin{lemma}[Weighted angular bounds]\label{inf:angular-weighted}
For $k=0,1,2$, the matrix $\partial_\varsigma^kT_\varsigma$ has bandwidth two and
\begin{equation}\label{inf:T-entry-bound}
 |(\partial_\varsigma^kT_\varsigma)_{ij}|\le C\min(w_i,w_j),\qquad
 \nrm{(\partial_\varsigma^kT_\varsigma)\mathsf W^{-1}}+
 \nrm{\mathsf W^{-1}(\partial_\varsigma^kT_\varsigma)}\le C.
\end{equation}
The matrices $X_\varsigma$ and their first two derivatives are bounded uniformly.
If
\begin{equation}\label{inf:G-angular}
 G_\varsigma=4x(1-x)(\partial_xT_\varsigma)^2
     =\tfrac12[T_\varsigma,[\Lambda,T_\varsigma]],
\end{equation}
then, for $k=0,1,2$,
\begin{equation}\label{inf:G-weighted}
 \nrm{\mathsf W^{-1}(\partial_\varsigma^kG_\varsigma)\mathsf W^{-1}}\le Cp.
\end{equation}
\end{lemma}
\begin{proof}
From \eqref{inf:recurrence-a}, uniformly for $j\ge1$ in the stated range,
\[
 \mathfrak a_j^2\le C\frac j{p+j},\qquad
 \frac{\partial_\varsigma\mathfrak a_j^2}{\mathfrak a_j^2}
 =-\frac{p^2}{(j+p/2-1)^2-p^2\varsigma}.
\]
For $j\ge1$, write
$\mathfrak D_j=(j+p/2-1)^2-p^2\varsigma$. The exact derivatives are
\begin{equation}\label{inf:angular-second-derivative}
 \frac{\partial_\varsigma\mathfrak a_j}{\mathfrak a_j}
 =-\frac{p^2}{2\mathfrak D_j},\qquad
 \frac{\partial_\varsigma^2\mathfrak a_j}{\mathfrak a_j}
 =-\frac{p^4}{4\mathfrak D_j^2}.
\end{equation}
The factor $\mathfrak D_j$ is bounded below by $c(p+j)^2$,
so $|\partial_\varsigma^k\mathfrak a_j|\le C\mathfrak a_j$
for $k=1,2$. The coefficient $\mathfrak a_0$ and both its
derivatives are identically zero. The diagonal entries $\mathfrak b_j$ have uniformly
bounded first two derivatives: $\varsigma$ stays away from zero, so
$\partial_\varsigma^k\sqrt \varsigma$ is bounded for $k\le2$.
The two-step entry of $\mathsf H_2(X_\varsigma)$ is
$\mathfrak a_{j+1}\mathfrak a_{j+2}$ and has size, together with its
two derivatives, at most $Cw_j$.

The one-step entry is
\[
 \mathfrak a_{j+1}\left(\mathfrak b_j+\mathfrak b_{j+1}
                               -\frac{2(\ab+1)}{p+2}\right).
\]
Its parenthesis equals $(r-1/2)$ times
\[
 p(p-2)\left\{
 \frac1{(p+2j-2)(p+2j)}+
 \frac1{(p+2j)(p+2j+2)}\right\}-\frac{2p}{p+2}.
\]
The braces with their prefactor have exactly the compensating value
at $j=0$. For $j\le p$, subtract this value and differentiate the
reciprocal products in $j$ to bound the difference by $Cj/p$.
For $j\ge p$ each term is bounded by a constant. Thus the
parenthesis and both its $\varsigma$ derivatives are $O(j/(p+j))$.
Multiplication by $\mathfrak a_{j+1}$ preserves the bound $Cw_j$.

The diagonal entries of $\mathsf H_2(X_\varsigma)$ are
$M_j=c_j^{\mathrm{diag}}+d_j^{\mathrm{diag}}\varsigma$, where
\begin{align}
 c_j^{\mathrm{diag}}&=\frac{j(p-1)(j+p-1)}{2(p+1)(2j+p-3)(2j+p+1)},
 \label{inf:T-diagonal}\\
 d_j^{\mathrm{diag}}&=-\frac{2jp^2(j+p-1)(p^2+7p-8j^2-8jp+8j)}
 {(2j+p)(p+1)(p+2)(2j+p-3)(2j+p-2)(2j+p+1)}.\nonumber
\end{align}
These formulas follow by the multiplication calculation in the proof
of Lemma~\ref{inf:quartic-normalization}, now at general $j$.
The denominator factors are comparable to $p$ or $p+j$;
hence $|c_j^{\mathrm{diag}}|+|d_j^{\mathrm{diag}}|\le Cj/(p+j)$.
The scalar $pg_2$ is bounded away from zero and has bounded derivatives
of every fixed order. Division by it proves the entry bounds.
For $|i-j|\le2$ the weights have ratios between $1/3$ and $3$.
The row and column sums over the five diagonals prove
\eqref{inf:T-entry-bound}. The same row-sum argument for $X_\varsigma$
and its derivatives proves their asserted bounds.

The identity in \eqref{inf:G-angular} follows by expanding
$[\Lambda,T]= (\Lambda T)_{\mathrm{mult}}
-8x(1-x)T_x\partial_x$.
The function $G_\varsigma$ is a degree-four polynomial in $X_\varsigma$, with
coefficients and their first two derivatives uniformly bounded.
Its bandwidth is four and its unweighted norm is $O(1)$.
When both indices exceed $p-4$, the weights are bounded below and
this already proves the required bound.
For the other nonzero entries all indices are at most $p+4$, and
\[
 (G_\varsigma)_{ij}=\frac12\sum_l(2\lambda_l-\lambda_i-\lambda_j)
                                 (T_\varsigma)_{il}(T_\varsigma)_{lj}.
\]
Only finitely many neighbours occur; their eigenvalue differences
are $O(p)$ and each product is at most $Cw_iw_j$ by
\eqref{inf:T-entry-bound}. The product rule gives the same estimate
for two derivatives. Summing over the nine diagonals proves
\eqref{inf:G-weighted}.
\end{proof}

\begin{lemma}[Positive angular variation]\label{inf:angular-positive}
For every subset $\mathcal I\subset\{8,9,\ldots\}$ containing no
adjacent indices, let $T_{\mathcal I}$ denote the principal angular
submatrix with those indices. Then
\begin{equation}\label{inf:angular-positive-bound}
 T'_{\mathcal I}(\varsigma)\ge c_0\mathsf W_{\mathcal I}^2,
 \qquad c_0=10^{-4},
\end{equation}
for $p\ge64$. The inequality holds for
finite or infinite $\mathcal I$ as a bounded-operator inequality.
\end{lemma}
\begin{proof}
Differentiate
$M_j/g_2=(c_j^{\mathrm{diag}}+d_j^{\mathrm{diag}}\varsigma)/
(c_2^{\mathrm{diag}}+d_2^{\mathrm{diag}}\varsigma)$.
The numerator is
\begin{align}
 \mathcal N_j&=d_j^{\mathrm{diag}}c_2^{\mathrm{diag}}
                    -c_j^{\mathrm{diag}}d_2^{\mathrm{diag}}\nonumber\\
 &=\frac{24jp^3(j-2)(p-1)(j+p-1)(j+p+1)}
 {(2j+p)(p+1)(p+2)^2(p+4)(p+5)
  (2j+p-3)(2j+p-2)(2j+p+1)}.
 \label{inf:N-angular}
\end{align}
Define also $\partial_\varsigma(\mathfrak a_j^2/g_2)
=\mathcal E_j/g_2^2$, where
\begin{align}
 \mathcal E_j&=\frac{4jp^2(j+p-2)F(p,j)}
 {(p+1)(p+2)^2(p+4)(p+5)(2j+p-3)(2j+p-2)^2(2j+p-1)},
 \label{inf:E-angular}\\
 F(p,j)&=(p^2-9p-16)j^2+(p^3-11p^2+2p+32)j
                         -5p^3+3p^2+8p-12.\nonumber
\end{align}
These identities are obtained by substituting
\eqref{inf:T-diagonal}, \eqref{inf:recurrence-a}, and
\eqref{inf:g2-exact} and clearing their positive denominators.
For $u=p-64$, $j_{\mathrm{pos}}=j-7$ one has the explicit identity
\begin{align*}
 F(p,j)={}&(u^2+119u+3504)j_{\mathrm{pos}}^2\\
 &+(u^3+195u^2+12548u+266304)j_{\mathrm{pos}}\\
 &+2u^3+359u^2+20957u+394500.
\end{align*}
Thus $\mathcal E_j>0$ for $p\ge64,j\ge7$.
The polynomial identity in Appendix~\ref{inf:sec-positivity} gives
\begin{equation}\label{inf:half-diagonal}
 \frac{\mathcal N_j}{2}
 >\frac23(\mathcal E_{j-1}+\mathcal E_j+
                 \mathcal E_{j+1}+\mathcal E_{j+2}),\qquad j\ge8.
\end{equation}
All denominators there are positive on the stated range.

The adjacent recurrence coefficients satisfy
$3/4<\mathfrak a_{j+1}/\mathfrak a_j<4/3$ for $j\ge7$.
For the upper bound discard the denominator ratio in
\eqref{inf:recurrence-a} to obtain
\[
 \frac{\mathfrak a_{j+1}^2}{\mathfrak a_j^2}
 \le\frac87\left(\frac{23}{22}\right)^2\frac{70}{69}<\frac{16}{9}.
\]
For the inverse ratio discard the increasing numerator factors and
use $2j+p\ge78$ to get
\[
 \frac{\mathfrak a_j^2}{\mathfrak a_{j+1}^2}
 \le\frac{78^2\cdot79}{76^2\cdot75}<\frac{16}{9}.
\]
The corresponding rational bound decreases with $2j+p$, factor by
factor, so these estimates cover every larger $j,p$.
On one parity the only off-diagonal entries of $T_\varsigma$ are
$\mathfrak a_{j+1}\mathfrak a_{j+2}/(pg_2)$.
Their derivatives are positive and, for the upper neighbour, at most
\[
 \frac{2}{3pg_2^2}(\mathcal E_{j+1}+\mathcal E_{j+2}).
\]
Indeed differentiate the product as one half the sum of the two
square derivatives times the adjacent coefficient ratios.
Equation~\eqref{inf:half-diagonal} and
$2|v_jv_k|\le|v_j|^2+|v_k|^2$ now give
\[
 \langle T'_\varsigma v,v\rangle\ge
 \sum_j\frac{\mathcal N_j}{2pg_2^2}|v_j|^2
\]
on either parity and every subset thereof.
For $8\le j\le p$, the numerator and denominator in
\eqref{inf:N-angular} are respectively at least $9j^2p^6$ and at
most $3456p^9$. Since $g_2<1/p$, the margin is at least
$j^2/(768p^2)$. For $j\ge p$, the corresponding estimates
$6p^4j^4$ and $3456p^5j^4$ give a margin at least $1/1152$.
Together with $w_j\le(9/8)j/p$ for $j\ge8$ and $w_j\le1$, these
bounds imply \eqref{inf:angular-positive-bound} with $c_0=10^{-4}$.

A set with no adjacent indices has no opposite-parity off-diagonal
entry: the matrix has bandwidth two. Its principal submatrix is
therefore the direct sum of the two parity submatrices.
The bounds extend from finite vectors by density and
Lemma~\ref{inf:angular-weighted}.
\end{proof}

\subsection{The common Dirichlet form}

Let $\zeta\in(0,1)$ be nonsquared reference radius and put
\[
 \epsilon_{\mathrm q}=-\frac{\theta}{2R^2},\qquad
 \rho_\varsigma(x)=1+\epsilon_{\mathrm q}T_\varsigma(x),\qquad c_p=p-\tfrac12.
\]
Use the Hilbert space
$\Hc_{\mathrm{flow}}=L^2((0,1),d\zeta)\otimes\ell^2(\mathbb N_0)$,
with the Jacobi identification already made. The ball operator and
its positive shift are
\begin{equation}\label{inf:flow-ball-operators}
 H_0=-\partial_\zeta^2+
          \frac{c_p(c_p-1)+\Lambda}{\zeta^2},\qquad
 K_0=H_0+R^2,\qquad
 \mathsf D=\zeta\partial_\zeta+\tfrac12.
\end{equation}
This is the mass-unitary representation of the operator already
denoted by $H_0$. Its spectrum and all its spectral projections
are unchanged by the unitary map.

\begin{proposition}[Exact pullback form]\label{inf:flow-form}
For the weighted quartic domain $0<\varrho<\rho_\varsigma(x)$, let
$U(\zeta,x)=u(\zeta\rho_\varsigma(x),x)$ and
$v=\sqrt{2p}\,\zeta^{p-1/2}\rho_\varsigma^pU$. Put
$L_\varsigma=\partial_x\log\rho_\varsigma$, $\alpha_\varsigma=\rho_\varsigma^{-2}$, and
$q_x=4x(1-x)$. Its Dirichlet form is
\begin{equation}\label{inf:flow-formula}
 \mathfrak q_\varsigma[v,w]=\int\alpha_\varsigma\left\{
 (v_\zeta-c_pv/\zeta)\overline{(w_\zeta-c_pw/\zeta)}
 +\frac{q_x}{\zeta^2}(v_x-L_\varsigma\mathsf Dv)
                       \overline{(w_x-L_\varsigma\mathsf Dw)}\right\}
 \,d\zeta\,d\mu_r.
\end{equation}
Define, initially as weak expressions on the common core,
\begin{align}
 \mathcal L(M)&=MH_0+\frac1{2\zeta^2}[\Lambda,M]\mathsf D,
 \label{inf:linear-form-operator}\\
 \mathfrak g_\gamma[v,w]&=\int\zeta^{-2}\gamma
                                  (\mathsf Dv)\overline{\mathsf Dw}.
 \nonumber
\end{align}
Write $\mathcal G(\gamma)$ for the weak form operator characterized
on this core by
$\langle\mathcal G(\gamma)v,w\rangle=\mathfrak g_\gamma[v,w]$.
The operator associated with $\mathfrak q_\varsigma$ is, in the sense of forms,
\begin{equation}\label{inf:flow-exact-operator}
 H_\varsigma=H_0+\mathcal L(\alpha_\varsigma-I)+\mathcal G(\gamma_\varsigma),
 \qquad
 \gamma_\varsigma=\epsilon_{\mathrm q}^2
                   (I+\epsilon_{\mathrm q}T_\varsigma)^{-4}G_\varsigma.
\end{equation}
The exact form has the same closed domain as $H_0$ for large $p$.
The family is $C^1$ in the relative form norm and has compact
resolvent. At integer block sizes it is unitarily equivalent to the
invariant Euclidean Dirichlet Laplacian on $\Omega_{\mathrm q}/R$.
\end{proposition}
\begin{proof}
With the fixed normalization by the unit-ball volume, the radial
change of variables has density $2p\zeta^{2p-1}\rho_\varsigma^{2p}$,
proving the mass normalization.
Its angular derivative is $U_x-\zeta L_{\varsigma}U_\zeta$.
Substitution gives
\[
 U_x-\zeta L_{\varsigma}U_\zeta
 = (2p)^{-1/2}\zeta^{-c_p}\rho_\varsigma^{-p}
       \{v_x-L_\varsigma(\zeta v_\zeta+v/2)\}.
\]
The radial derivative gives
$(2p)^{-1/2}\zeta^{-c_p}\rho_\varsigma^{-p}(v_\zeta-c_pv/\zeta)$.
These identities prove \eqref{inf:flow-formula}.

To integrate the form, radial integration of its first square gives
$\alpha_\varsigma(-\partial_\zeta^2+c_p(c_p-1)/\zeta^2)$.
The purely angular square gives
\[
 \zeta^{-2}\{\alpha_\varsigma\Lambda+\tfrac12[\Lambda,\alpha_\varsigma]
                         -\tfrac12(\Lambda\alpha_\varsigma)_{\mathrm{mult}}\}.
\]
For the cross terms use $\alpha_{\varsigma}L_\varsigma=-(\alpha_\varsigma)_x/2$ and
$\mathsf D^*=-\mathsf D+2$ in radial measure $\zeta^{-2}d\zeta$.
Their sum is
\[
 \zeta^{-2}\{\tfrac12[\Lambda,\alpha_\varsigma](\mathsf D-1)
                         +\tfrac12(\Lambda\alpha_\varsigma)_{\mathrm{mult}}\}.
\]
The multiplication terms cancel. The remaining square has
coefficient $\alpha_{\varsigma}q_xL_\varsigma^2$, giving
\eqref{inf:flow-exact-operator} and \eqref{inf:G-angular}.
These are identities on finite angular sums with radial functions
compactly supported in $(0,1)$.

We specify the closure. Let $\mathcal V_0=\Dom(K_0^{1/2})$, defined
as the completion of that core in its positive form norm.
Angular spectral truncations commute with $K_0$ and converge in
this norm. In each remaining radial block the Friedrichs definition
gives approximation by compactly supported radial functions.
Thus the stated core is form dense. Its angular integration by
parts has no endpoint term: the boundary factor is
$x^{\ab}(1-x)^{\bb}$ and the test functions are polynomial there.
Cutting off a smooth angular function near $x=0$ costs
$O(\delta^{\ab-1})$ in energy, and near $x=1$ costs
$O(\delta^{\bb-1})$. Both tend to zero since $\ab,\bb>1$.
At the physical origin a radial cutoff costs $O(\delta^{2p-2})$.
Consequently the block axes and origin have zero capacity for
this weighted energy.
For integer blocks this is the invariant Euclidean $H^1_0$
realization, by averaging smooth Dirichlet approximations over the
compact symmetry group. For real parameters define the original weighted Dirichlet
realization as the closure of smooth functions compactly supported
in the open meridian rectangle, using the norm from mass plus
\[
 2p\int_0^1\!\int_0^1
 \left(|U_\zeta|^2+\zeta^{-2}q_x|U_x|^2\right)
                  \zeta^{2p-1}\,d\zeta\,d\mu_r.
\]
The polynomial angular core has the same closure. The cutoffs
just estimated put every finite angular
polynomial times a compactly supported smooth radial function
in this closure. Conversely, for a smooth function compactly
supported in the meridian rectangle, let $U_J$ be its angular
projection on $P_0,\ldots,P_J$.
Integration by parts in $x$ gives
$\langle\Lambda_r U,P_j\rangle=\lambda_j\langle U,P_j\rangle$.
The smooth compact support makes $\Lambda_rU$ square integrable,
so Parseval gives square summability of these coefficients with
weight $\lambda_j^2$. It follows that $U_J\to U$ in angular
energy as well as in angular $L^2$.
The same projection commutes with the radial derivative;
Parseval applied to $U_\zeta$ gives convergence of its radial
energy. The radial support stays in a compact subinterval of
$(0,1)$, so the factor $\zeta^{-2}$ is bounded there.
Dominated convergence in the radial integral proves convergence
in the full weighted form norm. Each projected radial coefficient
is smooth with compact support. Thus both cores have the same
closure, establishing the real-parameter domain.

For integer blocks, the identification with the Euclidean
invariant form domain can also be read in both directions.
Average a compactly supported smooth Euclidean approximation
over $O(a)\times O(b)$; averaging is a contraction for mass and
Dirichlet energy. The result is smooth and invariant.
Cut it off at the origin and at each block axis. In squared
angular distance the extra energy is bounded by
$C\delta^{\ab-1}+C\delta^{\bb-1}$; at the origin it is bounded
by $C\delta^{2p-2}$. All tend to zero. Away from those sets it
is a smooth function of the meridian coordinates with compact
support, so it belongs to the weighted closure just described.
In the reverse direction a meridian test function with compact
support away from the axes and origin lifts to a smooth
invariant Euclidean test function, extended by zero across
those sets. This proves equality of the closed domains.

The estimates proved immediately below show that the difference of
\eqref{inf:flow-formula} and the ball form has relative norm below
$1/2$ on $\mathcal V_0$. Hence its closed form norm is equivalent to
that of $K_0$, and its closure has exactly domain $\mathcal V_0$.
The angular mass factor $\rho_\varsigma^p$ also preserves the
weighted first-order domain. The required control at the origin
is Hardy's inequality in the radial weight:
\begin{equation}\label{inf:flow-Hardy}
 \int_0^1|U|^2\zeta^{2p-3}\,d\zeta
 \le\frac1{(p-1)^2}
      \int_0^1|U_\zeta|^2\zeta^{2p-1}\,d\zeta.
\end{equation}
For compactly supported $U$ this follows by integrating
$(\zeta^{2p-2})'|U|^2$ and applying Cauchy--Schwarz; completion
extends it. Integrate also against $\mu_r$.
Multiplication by $\rho_\varsigma^{\pm p}$ differentiates in $x$
by the extra term $\pm p L_\varsigma U$. Its angular energy has
the factor $p^2|U|^2/\zeta^2$, which is controlled by
\eqref{inf:flow-Hardy}. Both multipliers and their inverses are
bounded for fixed parameters. Thus the unitary mass transformation
preserves the closure, including its origin behaviour.
The radial graph map is bi-Lipschitz even at the origin: on the
sphere its radial factor and first derivatives are bounded and it
is bounded away from zero. Such a map transports Dirichlet energy
and its closure. Therefore the just described closure is the
physical one, including at the origin. The relative derivative
bounds below prove norm $C^1$ dependence. Compactness of the
embedding $\mathcal V_0\hookrightarrow\Hc_{\mathrm{flow}}$ persists
under equivalent form norms and proves compact resolvent.
\end{proof}

\begin{lemma}[Relative form estimates]\label{inf:flow-relative}
Set $B_{1,\varsigma}=-2\mathcal L(\epsilon_{\mathrm q}T_\varsigma)$ and
$\mathcal R_{2,\varsigma}=H_\varsigma-H_0-B_{1,\varsigma}$. For $k=0,1,2$ and large $p$,
\begin{align}
 \nrm{K_0^{-1/2}(\partial_\varsigma^kB_{1,\varsigma})K_0^{-1/2}\mathsf W^{-1}}
 &\le C\tau p^{-2},\label{inf:flow-linear-relative}\\
 \nrm{\mathsf W^{-1}K_0^{-1/2}(\partial_\varsigma^k\mathcal R_{2,\varsigma})
               K_0^{-1/2}\mathsf W^{-1}}
 &\le C\theta^2p^{-3}.\label{inf:flow-quadratic-relative}
\end{align}
These are estimates on the full radial spaces. They also establish
the closure and differentiability assertions in
Proposition~\ref{inf:flow-form}.
\end{lemma}
\begin{proof}
We first prove the band estimate used to sum the exact form.
If an angular matrix $M$ has bandwidth $m$ and
$|M_{ij}|\le M_*w_i^aw_j^b$, $a,b\in\{0,1\}$, then
\begin{equation}\label{inf:band-form-bound}
 \nrm{\mathsf W^{-a}K_0^{-1/2}\mathcal L(M)
                K_0^{-1/2}\mathsf W^{-b}}
 \le C(m+1)^3M_*.
\end{equation}
Indeed the radial forms
$K_{0,j}=-\partial_\zeta^2+(\nu_j^2-1/4)/\zeta^2+R^2$
are comparable with ratio at most $C(m+1)^2$ for $|i-j|\le m$.
This follows from $\nu_i/\nu_j$ and its inverse being at most
$C(m+1)$. Thus
$\nrm{K_{0,i}^{-1/2}K_{0,j}^{1/2}}\le C(m+1)$.
The block $M_{ij}H_{0,j}$ in \eqref{inf:linear-form-operator}
is bounded after normalization since
$\nrm{H_{0,j}K_{0,j}^{-1}}\le1$.
For the commutator block use the two form bounds
\[
 \nrm{\zeta^{-1}K_{0,i}^{-1/2}}\le C/(p+i),\qquad
 \nrm{\zeta^{-1}\mathsf DK_{0,j}^{-1/2}}\le C.
\]
The first follows from the centrifugal potential; the second from
$\zeta^{-1}\mathsf D=\partial_\zeta+1/(2\zeta)$ and the same
positive form. Also
$|\lambda_i-\lambda_j|\le Cm(p+i+j)\le Cm(m+1)(p+i)$.
Each normalized block is bounded by $C(m+1)^2M_*$ after its
weights are removed. The row and column sums involve at most
$2m+1$ blocks. The scalar Schur inequality proves
\eqref{inf:band-form-bound} on the form core, then by closure.

For $n\ge2$, Lemma~\ref{inf:angular-weighted} gives
\[
 |(\partial_\varsigma^kT_\varsigma^n)_{ij}|\le n^2C^nw_iw_j,\qquad k\le2,
\]
with the factor $n^2$ replaced by one when $k=0$ if desired.
Put one inverse weight on the first factor and one on the last;
each differentiated endpoint has the same one-sided bound, and
interior factors are uniformly bounded. The product rule has at
most a fixed multiple of $n^2$ terms. The bandwidth is $2n$.
The series
\[
 \alpha_\varsigma=I-2\epsilon_{\mathrm q}T_\varsigma
       +\sum_{n\ge2}(-1)^n(n+1)\epsilon_{\mathrm q}^nT_\varsigma^n
\]
therefore has a two-sided weighted remainder under
\eqref{inf:band-form-bound} of size $C\epsilon_{\mathrm q}^2$,
including two derivatives. Its dominating series is
$\sum n^6(C|\epsilon_{\mathrm q}|)^n$, which converges uniformly.

Equation~\eqref{inf:G-weighted} and the convergent series for
$(I+\epsilon_{\mathrm q}T_\varsigma)^{-4}$ give
\[
 \nrm{\mathsf W^{-1}(\partial_\varsigma^k\gamma_\varsigma)\mathsf W^{-1}}
 \le Cp\epsilon_{\mathrm q}^2,\qquad k\le2.
\]
To move the left weight through the inverse series, use boundedness
of $\mathsf W^{-1}T_\varsigma\mathsf W$ and its derivatives, which follows
from \eqref{inf:T-entry-bound}. Since
$\zeta^{-1}\mathsf D$ is controlled by $K_0^{1/2}$ and commutes
with the angular weights, the same bound holds for the normalized
form $\mathcal G(\gamma_\varsigma)$.
Now $|\epsilon_{\mathrm q}|\le C\tau/p^2$.
The first-order band estimate proves
\eqref{inf:flow-linear-relative}; the two remainder estimates prove
\eqref{inf:flow-quadratic-relative}.
Their unweighted sum is below $1/2$ for large $p$.

The $k=2$ bounds provide a uniform modulus of continuity for the
first derivative: Taylor's integral formula on the finite bands,
followed by the uniformly convergent series, bounds the difference
quotient remainder in relative form norm. Thus the family is norm
$C^1$ on the common form space. These bounds also prove the
closure and regularity assertions of Proposition~\ref{inf:flow-form}.
\end{proof}

\subsection{The boundary matrix and the transport term}

Fix a finite $W_{\mathrm{flow}}$ and set
$\PP=\one_{[-W_{\mathrm{flow}},W_{\mathrm{flow}}]}(H_0-R^2)$.
For large $p$, Lemma~\ref{inf:window-isolation} says its states have
distinct, nonadjacent angular indices. Index its orthonormal states by
$i$, let $j(i)$ be their angular indices and $z_i$ their positive
zeros, and write
\[
 D_\PP=\diag(z_i^2-R^2),\quad
 Z_\PP=\diag(z_i/R),\quad
 T_\PP^{\mathrm{ang}}=[(T_\varsigma)_{j(i),j(i')}],\quad
 \mathcal S_\varsigma=\PP(\epsilon_{\mathrm q}T_\varsigma\mathsf D)\PP.
\]
The matrix $T_\PP^{\mathrm{ang}}$ is the \emph{angular principal
submatrix}, while $\mathcal S_\varsigma$ is the actual Hilbert-space
compression with radial overlaps.

\begin{lemma}[Transport identity]\label{inf:transport}
One has $\mathcal S_\varsigma^*=-\mathcal S_\varsigma$ and the exact identity
\begin{equation}\label{inf:transport-identity}
 \PP B_{1,\varsigma}\PP=\theta Z_\PP T_\PP^{\mathrm{ang}} Z_\PP
                             -[D_\PP,\mathcal S_\varsigma].
\end{equation}
For $k=0,1$,
\begin{equation}\label{inf:transport-bounds}
 \nrm{\PP(\partial_\varsigma^kB_{1,\varsigma})\PP\mathsf W^{-1}}\le C\tau,
 \qquad
 \nrm{(\partial_\varsigma^k\mathcal S_\varsigma)\mathsf W^{-1}}\le C\tau/p,
\end{equation}
with constants independent of $W_{\mathrm{flow}}$ once
$p^2\ge W_{\mathrm{flow}}$.
\end{lemma}
\begin{proof}
As interior differential expressions,
$[H_0,\mathsf D]=2H_0$ and
$[H_0,h]=\zeta^{-2}[\Lambda,h]$ for an angular multiplier $h$.
Thus $B_{1,\varsigma}=-[H_0,\epsilon_{\mathrm q}T_\varsigma\mathsf D]$.
Green's identity applied to two ball Dirichlet eigenfunctions gives
\[
 \langle\phi_i,H_0(h\mathsf D\phi_{i'})\rangle
 =z_i^2\langle\phi_i,h\mathsf D\phi_{i'}\rangle
 +\int h\,\partial_\zeta\phi_i(1)
                 \partial_\zeta\phi_{i'}(1)\,d\mu_r.
\]
The origin term vanishes by the regular radial behaviour. The
boundary derivatives are $\sqrt2z_iP_{j(i)}$ by
Lemma~\ref{inf:bessel-spacing}. Substitution of
$h=-\theta T_\varsigma/(2R^2)$ gives
\eqref{inf:transport-identity}. Radial integration on pairs of
Dirichlet eigenfunctions gives $\mathsf D^*=-\mathsf D$; the
angular multiplier commutes with $\mathsf D$. Hence $\mathcal S_\varsigma$ is
skew-adjoint.
On the fixed window $\nrm{K_0^{1/2}\PP}\le Cp$.
Equation~\eqref{inf:flow-linear-relative} proves the first bound
in \eqref{inf:transport-bounds}. The form estimate also gives
$\nrm{\mathsf D\PP}\le Cp$; use
\eqref{inf:T-entry-bound} and
$|\epsilon_{\mathrm q}|\le C\tau/p^2$ for the second bound.
\end{proof}

\subsection{Eliminating the full complementary space}

\begin{lemma}[Complementary inverse and remainder]
\label{inf:flow-complement}
There is an absolute lower bound on $W_{\mathrm{flow}}\ge100$ such
that the following holds for $p$ sufficiently large depending on
that fixed window. For every $\varsigma\in I$ and real $|\lambda|\le1$,
the form operator
$\PP^\perp(H_\varsigma-R^2-\lambda)\PP^\perp$ is invertible. Its exact
Schur matrix is
\begin{equation}\label{inf:flow-schur}
 F_\PP(\varsigma,\lambda)=D_\PP-\lambda+
                        \PP B_{1,\varsigma}\PP+\mathcal R_\varsigma(\lambda),
\end{equation}
where, for $k=0,1$,
\begin{equation}\label{inf:flow-schur-remainder}
 \nrm{\mathsf W^{-1}(\partial_\varsigma^k\mathcal R_\varsigma)\mathsf W^{-1}}
 \le C\theta^2(W_{\mathrm{flow}}^{-1}+p^{-1}).
\end{equation}
The constant $C$ is independent of the chosen window provided
$p^2\ge W_{\mathrm{flow}}$ and the window exceeds its fixed lower
bound. A kernel vector is reconstructed uniquely from $v=\PP u$,
and its complementary part satisfies
\begin{equation}\label{inf:flow-Q-reconstruction}
 \nrm{\PP^\perp u}\le C\tau W_{\mathrm{flow}}^{-1}
                                      \nrm{\mathsf W v}.
\end{equation}
\end{lemma}
\begin{proof}
Use the positive form normalization $K_0$ and abbreviate
$Q=\PP^\perp$, $\mathcal B_\varsigma=H_\varsigma-H_0$, and
$\mathfrak C_\varsigma=K_0^{-1/2}\mathcal B_{\varsigma}K_0^{-1/2}$.
Write $K_{0,\PP}$ and $K_{0,Q}$ for the restrictions of $K_0$
to $\PP\Hc_{\mathrm{flow}}$ and $Q\Hc_{\mathrm{flow}}$.
The inverse $\mathfrak U_0$ of the normalized unperturbed complement is
radially diagonal with entries
\[
 \frac{z_{j,k}^2+R^2}{z_{j,k}^2-R^2-\lambda},
 \qquad |z_{j,k}^2-R^2|>W_{\mathrm{flow}}.
\]
Their absolute supremum is at most $Cp^2/W_{\mathrm{flow}}$.
This includes the whole infinite tail, where the ratio tends to
one; $p^2\ge W_{\mathrm{flow}}$ covers that limit.
Lemma~\ref{inf:flow-relative} gives
$\nrm{\mathfrak C_\varsigma\mathsf W^{-1}}+
\nrm{\mathfrak C'_\varsigma\mathsf W^{-1}}\le C\tau/p^2$.
Choose the fixed window large enough that
$\nrm{\mathfrak U_0Q\mathfrak C_{\varsigma}Q}\le1/2$ uniformly for $\tau\le20$.
The normalized inverse
\[
 \mathfrak U=(I+\mathfrak U_0Q\mathfrak C_{\varsigma}Q)^{-1}\mathfrak U_0
\]
has norm at most $Cp^2/W_{\mathrm{flow}}$, and
$\mathfrak U'_\varsigma=-\mathfrak UQ\mathfrak C'_{\varsigma}Q\mathfrak U$ has norm at most
$Cp^2\tau/W_{\mathrm{flow}}^2$.
The actual form inverse is $K_{0,Q}^{-1/2}\mathfrak UK_{0,Q}^{-1/2}$.

Eliminate $Q$ using this inverse. The remainder in
\eqref{inf:flow-schur} is
\[
 \mathcal R_\varsigma=\PP\mathcal R_{2,\varsigma}\PP
 -\PP\mathcal B_{\varsigma}QK_{0,Q}^{-1/2}\mathfrak UK_{0,Q}^{-1/2}
                                      Q\mathcal B_\varsigma\PP.
\]
In the second term the two exterior $K_{0,\PP}^{1/2}$ factors cost
$Cp$ each; each normalized perturbation with its exterior inverse
weight costs $C\tau/p^2$. The weighted norm is therefore
$C\theta^2/W_{\mathrm{flow}}$.
Its derivative either replaces one perturbation by $\mathfrak C'_\varsigma$ or replaces
$\mathfrak U$ by $\mathfrak U'_\varsigma$. The added bound in the latter case is
$C\tau^3/W_{\mathrm{flow}}^2$, absorbed in the same bound.
The first remainder term costs $C\theta^2/p$ by
\eqref{inf:flow-quadratic-relative}. This proves
\eqref{inf:flow-schur-remainder} with no dependence of $C$ on the
window size.

The eliminated component is
$u_Q=-K_{0,Q}^{-1/2}\mathfrak UQ\mathfrak C_\varsigma\PP K_{0,\PP}^{1/2}v$.
Since $K_0\ge cp^2I$, the factors cost respectively
$C/p$, $Cp^2/W_{\mathrm{flow}}$, $C\tau/p^2$ after extracting
$\mathsf Wv$, and $Cp$. This proves
\eqref{inf:flow-Q-reconstruction}. The reconstruction belongs to
the common form domain. The eliminated equation is exactly zero
there, so the reduced kernel equation is equivalent to the full
weak eigenvalue equation. A weak solution of that equation is in
the operator domain by the definition of the associated closed-form
operator, completing the equivalence of the reduced and full
eigenvalue equations.
\end{proof}

\subsection{The seed and the second low state}

Assume now that $R$ and $\theta$ satisfy
\eqref{inf:detuning-definition} with $J_p(R)=0$.

\begin{lemma}[Low-block elimination]\label{inf:flow-low-block}
For every fixed $W_{\mathrm{flow}}$ above the absolute lower bound in
Lemma~\ref{inf:flow-complement} and all sufficiently large
$p$, the only angular indices below eight in the flow window are
$j=2,5$, each with one radial state. Their ball offsets satisfy
\begin{equation}\label{inf:low-offsets}
 d_2=-\frac{\theta}{2p}+O(\theta^2p^{-3}),\qquad
 d_5=-72+O(p^{-1}).
\end{equation}
Let $L$ be their two-dimensional span and $\PP_H=\PP-L$.
Write $F_L,\mathsf W_L,u_L$ for the $L$ blocks of $F_\PP,\mathsf W,u$,
and $\mathcal R_{HL},\mathcal R_{LH}$ for the corresponding off-diagonal
blocks of $\mathcal R_\varsigma$ between $L$ and $\PP_H$.
Write $D_H,Z_H,\mathsf W_H$ for the restrictions of
$D_\PP,Z_\PP,\mathsf W$ to this high block, and
$T_H^{\mathrm{ang}}$ for its angular principal submatrix.
For $|\lambda|\le\omega_p=\tau/(16p)$, the $L$ block of
$F_\PP(\varsigma,\lambda)$ is invertible. Eliminating it gives
\begin{equation}\label{inf:high-schur}
 F_H(\varsigma,\lambda)=D_H-\lambda+B_{1,HH}
                                      +\widehat{\mathcal R}_\varsigma(\lambda),
\end{equation}
with
\begin{equation}\label{inf:high-schur-remainder}
 \nrm{\mathsf W_H^{-1}(\partial_\varsigma^k\widehat{\mathcal R}_\varsigma)
                         \mathsf W_H^{-1}}
 \le C\theta^2(W_{\mathrm{flow}}^{-1}+p^{-1}),\quad k=0,1.
\end{equation}
A nonzero full eigenvector in this band has $v=\PP_Hu\ne0$ and
$\nrm{u}\le C\nrm v$, uniformly also as $\tau\downarrow0$.
\end{lemma}
\begin{proof}
The trace ratio and \eqref{inf:rho-equation} give
$R=2p+3+O(p^{-1})$. Normalize the seed radial function by
$f'(R)=1$, $f(R)=\chi=\theta/(4Rp)$.
The collar equation has bounded coefficients and one bounded
derivative for fixed angular index. On $|y|\le2|\chi|$ its integral
equation yields $f'(R+y)=1+O(|\chi|)$. The endpoint values have
opposite signs, and the derivative has positive sign.
Its unique zero in the collar is
$R-\chi+O(\chi^2)$. Squaring gives
$d_2=-2R\chi+O(R\chi^2)=-\theta/(2p)+O(\theta^2p^{-3})$,
uniformly for arbitrarily small nonzero $\theta$.

For $j=5$ use \eqref{inf:resonant-trace} with $c_5=18$.
The same collar argument places its zero at
$R-18/p+O(p^{-2})$, giving the second offset in
\eqref{inf:low-offsets}. For $0\le j<8$ other than $2,5$, the
finite recurrence \eqref{inf:finite-q-recurrence} has nonzero
limiting value at $k=2j-1$. Its derivative trace stays bounded.
The collar equation excludes a zero in a fixed positive-width
collar of $R$. Such an index is therefore outside every fixed
energy window for large $p$. The radial spacing excludes a second
state in each remaining index.

The first-order operator has no $L$--$\PP_H$ coupling or coupling
between the two $L$ states. Their angular distances are at least
three, whereas both matrices in
\eqref{inf:transport-identity} have angular bandwidth two.
The seed diagonal is
\[
 \sigma_2=d_2+\frac{\theta z_2^2}{pR^2}
          =\frac{\theta}{2p}+O(\theta^2p^{-3}),
\]
because $(T_\varsigma)_{22}=1/p$ exactly. The other diagonal tends to $-72$
and its perturbation is $O(\tau/p)$.
For $|\lambda|\le\omega_p$, these two uncorrected diagonal entries
have modulus at least $\tau/(4p)$ and $36$, respectively.
The remainder block has norm at most
$C\theta^2(W_{\mathrm{flow}}^{-1}+p^{-1})p^{-2}$ because
$w_2,w_5\le6/p$. A two-dimensional Neumann series gives
\[
 \nrm{F_L^{-1}}+\nrm{\partial_{\varsigma}F_L^{-1}}\le Cp/\tau.
\]
For the derivative use
$(F_L^{-1})'=-F_L^{-1}F'_LF_L^{-1}$ and
$\nrm{F'_L}\le C\tau/p$, which follows from the diagonal structure
and the angular derivative bound. In particular,
\begin{equation}\label{inf:weighted-low-inverse}
 \nrm{\mathsf W_LF_L^{-1}\mathsf W_L}+
 \nrm{\mathsf W_L(\partial_{\varsigma}F_L^{-1})\mathsf W_L}
 \le C/(p\tau).
\end{equation}
The second elimination has remainder
\[
 \widehat{\mathcal R}
 =\mathcal R_{HH}-\mathcal R_{HL}F_L^{-1}\mathcal R_{LH}.
\]
Put $a_*=C\theta^2(W_{\mathrm{flow}}^{-1}+p^{-1})$.
The weighted new term and its derivative are at most
$Ca_*^2/(p\tau)$ by \eqref{inf:weighted-low-inverse}.
For $\tau\le20$ this is absorbed in $a_*$ after increasing the
$p$ threshold, proving \eqref{inf:high-schur-remainder}.
Reconstruction gives
\[
 \nrm{u_L}\le C\tau(W_{\mathrm{flow}}^{-1}+p^{-1})
                           \nrm{\mathsf W_Hv}.
\]
Combining this with \eqref{inf:flow-Q-reconstruction} proves
$\nrm u\le C\nrm v$ and the nonvanishing assertion.
The inverse factors $1/\tau$ are paired with the quadratic
remainders, giving uniform reconstruction for $0<\tau\le20$.
\end{proof}

\begin{lemma}[Ordered eigenvalue derivatives]\label{inf:ordered-eigenvalues}
For the common $C^1$ closed-form family above, the ordered
eigenvalues are locally Lipschitz. At an interior parameter
$\varsigma_0\in(1/64,1/16)$, suppose an eigenvalue $\lambda$ has multiplicity
$m$ and occupies ordered positions $i,\ldots,i+m-1$.
Let $\mu_1\le\cdots\le\mu_m$ be the eigenvalues of the derivative
form restricted to its Hilbert eigenspace. The one-sided
derivatives exist and satisfy
\begin{equation}\label{inf:one-sided-eigenvalues}
 \lambda'_{i+r-1,+}(\varsigma_0)=\mu_r,\qquad
 \lambda'_{i+r-1,-}(\varsigma_0)=\mu_{m-r+1},\qquad 1\le r\le m.
\end{equation}
Thus every one-sided derivative is an eigenvalue of the
compressed derivative form. An ordered eigenvalue need not
have a two-sided derivative at a multiple crossing.
\end{lemma}
\begin{proof}
A fixed positive shift has inverse that is norm $C^1$, by the
inverse identity on the common form space. Its compact resolvent
admits a contour surrounding the isolated finite cluster in
question. The contour integral gives a $C^1$ finite-rank
projection, also with values in the form domain. Local unitary
transport of this projection identifies its range with a fixed
finite-dimensional space. Restricting the resolvent there and
inverting that finite matrix gives a $C^1$ self-adjoint matrix
for the original spectral cluster; see also
\cite[Chapter~II, Theorem~5.4 and Theorem~6.8]{inf:kato}.
At $\varsigma_0$ it is $\lambda I$. Its derivative is the
restricted derivative form: differentiation of its unitary
transport contributes a commutator with $\lambda I$, which is
zero. The matrix expansion is therefore $\lambda I+hB+o(|h|)$,
with $B$ having eigenvalues $\mu_r$.
The finite-dimensional min--max principle bounds changes of
ordered eigenvalues by the operator norm of the remainder.
For $h>0$ the order of $h\mu_r$ is unchanged, and for $h<0$ it
is reversed. Division by $h$ gives
\eqref{inf:one-sided-eigenvalues}.
The same estimate for two nearby parameters gives local
Lipschitz continuity. Compactness of the parameter interval
gives absolute continuity for each fixed ordered curve, hence
a two-sided derivative almost everywhere.
\end{proof}

\begin{theorem}[Local spectral flow]\label{inf:flow-theorem}
For all sufficiently large $p\in\mathcal P_{\mathrm{bd}}$, let $R,\theta$
be as in Lemma~\ref{inf:detuning}. There are absolute constants
$c_{\mathrm{flow}}>0$ and $P_{\mathrm{flow}}<\infty$ such that
every ordered eigenvalue offset $\lambda_i(\varsigma)$ of $H_\varsigma-R^2$
satisfies
\begin{equation}\label{inf:flow-conclusion}
 \operatorname{sign}(\theta)\lambda_i'(\varsigma)
 \ge c_{\mathrm{flow}}\tau/p
 \quad\hbox{for almost every }\varsigma\hbox{ with }
 |\lambda_i(\varsigma)|<\tau/(16p),
\end{equation}
when $p\ge P_{\mathrm{flow}}$. The constants are uniform in
$0<\tau\le20$.
\end{theorem}
\begin{proof}
Fix an eigenvalue $\lambda$ in the band and a vector $u$ in its
entire eigenspace. Put $v=\PP_Hu$. The two exact eliminations show
that $v\ne0$. With $v,\lambda$ held fixed, differentiation of the
reconstructed quadratic form gives
\begin{equation}\label{inf:feshbach-derivative}
 \langle u,H'_{\varsigma}u\rangle
 =\langle v,\partial_{\varsigma}F_H(\varsigma,\lambda)v\rangle.
\end{equation}
Indeed, terms differentiating the reconstruction lie in
$\PP^\perp+L$, and their pairing with the operator vanishes by the
eliminated equations. The bounded normalized inverses above
make the reconstruction differentiable in the common form norm.

The transport identity gives
\[
 \partial_{\varsigma}F_H
 =\theta Z_H(T_H^{\mathrm{ang}})'Z_H
          -[D_H,\mathcal S'_{HH}]+\partial_\varsigma\widehat{\mathcal R}.
\]
The first term, after multiplication by $\operatorname{sign}\theta$,
is at least $c_1\tau\nrm{\mathsf W_Hv}^2$ in quadratic form.
Here Lemma~\ref{inf:angular-positive} applies to all the retained
indices, and $z_i/R\to1$ uniformly in the fixed window; $c_1>0$
is absolute.
Estimate the transport term on the actual kernel vector, using
\eqref{inf:high-schur}:
$D_Hv=\lambda v-(B_{1,HH}+\widehat{\mathcal R})v$.
Since $\mathcal S'_{HH}$ is skew-adjoint, its real pairing with $\lambda v$
vanishes. Thus
\begin{align*}
 |\langle v,[D_H,\mathcal S'_{HH}]v\rangle|
 &\le2\nrm{(B_{1,HH}+\widehat{\mathcal R})v}
                                  \nrm{\mathcal S'_{HH}v}\\
 &\le C\theta^2p^{-1}\nrm{\mathsf W_Hv}^2.
\end{align*}
The final remainder is at most
$C\theta^2(W_{\mathrm{flow}}^{-1}+p^{-1})
\nrm{\mathsf W_Hv}^2$.
Choose one constant $C_{\mathrm{flow,err}}$ dominating the
transport-error constant and the constants in
\eqref{inf:flow-schur-remainder} and
\eqref{inf:high-schur-remainder}. It is independent of
$W_{\mathrm{flow}}$, under their stated lower bound on that
window. The combined error is at most
$C_{\mathrm{flow,err}}\theta^2(W_{\mathrm{flow}}^{-1}+2p^{-1})
\nrm{\mathsf W_Hv}^2$.
Choose $W_{\mathrm{flow}}$ so large that
$20C_{\mathrm{flow,err}}/W_{\mathrm{flow}}<c_1/4$ and the
complementary-inverse conditions hold. Next choose $p$ large
enough that $40C_{\mathrm{flow,err}}/p<c_1/4$.
The other window-dependent requirements are
$p^2\ge W_{\mathrm{flow}}$, $2W_{\mathrm{flow}}/R<1/10$,
and the fixed-index asymptotic threshold in
Lemma~\ref{inf:flow-low-block}. Also impose the near-crossing
low-degree threshold in Proposition~\ref{inf:low-exclusion}. Equations
\eqref{inf:feshbach-derivative} and these bounds give
\begin{equation}\label{inf:weighted-flow-margin}
 \operatorname{sign}(\theta)\langle u,H'_{\varsigma}u\rangle
 \ge(c_1/2)\tau\nrm{\mathsf W_Hv}^2.
\end{equation}

To turn the weighted bound into a $1/p$ margin, put
$J_p=\lfloor\sqrt p/16\rfloor$.
For large $p$ every index $j<J_p$ other than the removed seed is in
the exclusion range of Proposition~\ref{inf:low-exclusion}.
Its ball offset has modulus greater than $32$.
On those rows of \eqref{inf:high-schur}, invert $D_H-\lambda$,
whose inverse norm is at most $1/31$, to get
\[
 \nrm{v_{j<J_p}}\le C\tau\nrm{\mathsf W_Hv}.
\]
On the other indices $w_j\ge c/\sqrt p$, giving
$\nrm{v_{j\ge J_p}}\le C\sqrt p\nrm{\mathsf W_Hv}$.
Lemma~\ref{inf:flow-low-block} therefore yields
$\nrm{\mathsf W_Hv}^2\ge cp^{-1}\nrm u^2$.
The further resonant low states $j=8,11,\ldots$ in a large fixed
window are included in this argument: they remain in $\PP_H$,
are covered by its positive angular variation, and are excluded
from the near-zero mass by the last row estimate.
Substitution into \eqref{inf:weighted-flow-margin} proves the
claimed derivative lower bound on the entire eigenspace.

By Lemma~\ref{inf:ordered-eigenvalues}, every one-sided derivative
of an ordered eigenvalue in the band is an eigenvalue of this
compressed derivative form. Its signed lower bound therefore
applies at every differentiability point, and almost everywhere
on each absolutely continuous curve. This proves
\eqref{inf:flow-conclusion}, including multiple crossings.
\end{proof}

\section{Choosing an integer split}\label{inf:sec-gap}

The counting window in this section is denoted by $W_{\mathrm{count}}$.
It comes from a radial shell comparison and is independent of the
wide window $W_{\mathrm{flow}}$ used in the spectral-flow proof.

\begin{lemma}[Avoiding narrow bands on a discrete grid]
\label{inf:grid-avoidance}
Let a compact interval contain at least $p/8$ grid points separated
by at least $1/(9p)$. Fix $C_*>0$ and suppose at most
$C_*p^{2/3}$ absolutely
continuous real functions can enter $[-1,1]$, and for each such
function $\lambda_i$ one has
$\epsilon_{\mathrm{sgn}}\lambda_i'\ge\mu>0$ almost everywhere on
$|\lambda_i|<\omega$, where $\epsilon_{\mathrm{sgn}}\in\{-1,1\}$ is fixed
and $0<\omega<1$. If $p$ is sufficiently large depending on $C_*$
and $\gamma=\mu/(576C_*p^{2/3})<\omega$, then some grid point
satisfies $|\lambda_i|\ge\gamma$ for every $i$.
\end{lemma}
\begin{proof}
Clip $\epsilon_{\mathrm{sgn}}\lambda_i$ to $[-\omega,\omega]$.
The chain rule for absolutely continuous functions says that its
derivative is zero almost everywhere outside the open band and is
at least $\mu$ inside. On a level set an absolutely continuous
function has derivative zero almost everywhere, so the two band
endpoints introduce no exceptional contribution. The clipped
function is therefore nondecreasing. The preimage of
$(-\gamma,\gamma)$ is an interval, and its length is at most
$2\gamma/\mu$: integrate the derivative between any two points
of that interval. It contains at most $1+18p\gamma/\mu$ grid points.
The union of all excluded grid points has cardinality at most
\[
 C_*p^{2/3}(1+18p\gamma/\mu)
 =C_*p^{2/3}+p/32\le p/16<p/8
\]
for large $p$. A remaining point has the asserted property.
Functions that never enter $[-1,1]$ exclude no point.
\end{proof}

\begin{lemma}[A fixed counting window]\label{inf:quartic-count-window}
For $p\ge64$, $R\in[2p-1,2p+4]$, $|\theta|\le20$, and every
$\varsigma\in I$, the quartic domain lies in the radial shell
$B_{1-3/p^2}\subset\Omega_{\mathrm q}/R\subset B_{1+3/p^2}$.
Every ordered eigenvalue curve that enters the unit energy
window about $R^2$ has its ball counterpart in the fixed window
$W_{\mathrm{count}}=40$. This statement holds also for real block
parameters in their weighted Dirichlet realization.
\end{lemma}
\begin{proof}
The radius differs from one by at most $10/R^2$, using
$\nrm{T_\varsigma}_\infty<1$. For $p\ge64$,
$10/(2p-1)^2\le3/p^2$, because
$2p^2-12p+3>0$. Set $\delta=3/p^2$.
The ordered min--max shell comparison implies that
$\lambda_i(\Omega_{\mathrm q}/R)\in[R^2-1,R^2+1]$ forces
\[
 (1-\delta)^2(R^2-1)\le\lambda_i(B)
              \le(1+\delta)^2(R^2+1).
\]
Since $(2p+4)^2+1\le5p^2$ for $p\ge64$, the upper displacement
from $R^2$ is at most
$1+5(6+9/p^2)<32$; the lower displacement is at most $31$.
Both are below $40$. These are indexwise inequalities for every
split, and the ball eigenvalue is split independent.
The weighted min--max comparison has the same dilation factors,
so the proof includes real block parameters.
\end{proof}

\begin{theorem}[Integer-split gap]\label{inf:integer-gap}
There are absolute constants $c_{\mathrm q}>0$ and $P_{\mathrm q}$
such that for every $p\in\mathcal P_{\mathrm{bd}}$ with $p\ge P_{\mathrm q}$,
there is an integer $a\in[p/2,3p/4]$ for which the actual quartic
domain satisfies
\begin{equation}\label{inf:quartic-gap}
 \dist\bigl(R^2,\spec(-\Delta_D|_{O(a)\times O(2p-a)};
                               \Omega_{\mathrm q}/R)\bigr)
 \ge c_{\mathrm q}\tau p^{-5/3}.
\end{equation}
For each fixed centre accuracy $N$, the same split gives centre gap
\begin{equation}\label{inf:centre-gap}
 \gamma_p=c_{\mathrm{gap}}\tau p^{-5/3},\qquad
 c_{\mathrm{gap}}=c_{\mathrm q}/2,
\end{equation}
for all $p$ above a further threshold depending on $N$.
\end{theorem}
\begin{proof}
Lemma~\ref{inf:quartic-count-window} places every ball counterpart
of a curve entering $[-1,1]$ about $R^2$ in the fixed window
$W_{\mathrm{count}}=40$, independently of $\varsigma,\theta$.
Theorem~\ref{inf:global-count} therefore bounds the number of
possible ordered indices by $C_*p^{2/3}$, where one may fix
$C_*=\max(1,C_{40})$.
This count concerns the ball spectrum, which is independent of
$\varsigma$; the single ball count bounds all curves that can enter
the band.

The integers $a\in[p/2,3p/4]$ provide at least $p/4-1$ points
$\varsigma_a=(a/(2p)-1/2)^2$ in $I$. For consecutive admissible integers, put
$d_a=1/2-a/(2p)$. Both $d_a,d_{a+1}$ are at least $1/8$, so
\[
 \varsigma_a-\varsigma_{a+1}
 =(d_a-d_{a+1})(d_a+d_{a+1})
 =\frac{d_a+d_{a+1}}{2p}\ge\frac1{8p}.
\] In particular the weaker constants
$p/8$ and $1/(9p)$ in Lemma~\ref{inf:grid-avoidance} hold for
large $p$. Apply that lemma with
\[
 \epsilon_{\mathrm{sgn}}=\operatorname{sign}\theta,\qquad
 \mu=c_{\mathrm{flow}}\tau/p,\qquad
 \omega=\tau/(16p).
\]
The flow hypothesis is Theorem~\ref{inf:flow-theorem}.
For large $p$ the resulting
$\gamma=(c_{\mathrm{flow}}/(576C_*))\tau p^{-5/3}$ is smaller
than $\omega$. Thus a grid point has the full invariant gap
\eqref{inf:quartic-gap}, with
$c_{\mathrm q}=c_{\mathrm{flow}}/(576C_*)$.
At that grid point the weighted realization is the Euclidean
invariant realization by Proposition~\ref{inf:flow-form}.
Finally Proposition~\ref{inf:centre-gap-transfer} transfers half of
the gap to every fixed-order centre after increasing its threshold.
\end{proof}

\section{Coefficient spaces and clamped Poisson estimates}\label{inf:sec-spaces}

From now on $p\in\tfrac12\Z$ and the integer split are fixed. Constants in operator
estimates are uniform over $p\ge64$ and admissible splits unless a
subscript explicitly permits dependence on the centre order.
The angular coefficient radius is fixed at $\rho_w=2$ throughout
the Newton argument. Larger radii will be used only to estimate the
known centre.

\subsection{The source, shape, and clamped subspaces}

Define the degree weight
\begin{equation}\label{inf:degree-weight}
 \wt(\ell)=(1+\ell/p^2)^p,\qquad \ell\ge0.
\end{equation}
It is increasing and obeys
\begin{equation}\label{inf:weight-properties}
 \wt(\ell_1+\ell_2)\le\wt(\ell_1)\wt(\ell_2),\qquad
 \frac{\wt(\ell+2)}{\wt(\ell)}\le e^{2/p},\qquad
 \wt(\ell)\le e^{\ell/p}.
\end{equation}
These follow from $1+u+v\le(1+u)(1+v)$ and $1+u\le e^u$.

\begin{definition}[The fixed Banach spaces]\label{inf:spaces}
For the polynomial basis \eqref{inf:radial-basis}, let $\Cc=\Cc_{p,2}$
be the completion with norm
\begin{equation}\label{inf:C-norm}
 \nrm f_\Cc=\sum_{j,s\ge0}2^j\wt(2j+2s)|f_{j,s}|.
\end{equation}
Its moment spaces have norms
\begin{equation}\label{inf:C-moments}
 \nrm f_{\Cc^{[k]}}=
 \sum_{j,s\ge0}2^j\wt(2j+2s)(1+2j+2s)^k|f_{j,s}|.
\end{equation}
Let $\Gc_p=\{g\in\Cc:g_{j,0}=0\text{ for every }j\}$ and let
$\Ec_1$ be the space of all boundary coefficient sequences with
finite norm
\begin{equation}\label{inf:E-norm}
 h=\sum_{j\ge0}h_jG_j,\qquad
 \nrm h_{\Ec_1}=\sum_{j\ge0}(1+j)2^j\wt(2j)|h_j|.
\end{equation}
Write $\HP  h=\sum_jh_jS_j$ for its harmonic extension, and
$\Xc_p=\Gc_p\times\Ec_1$ with the sum norm.
The Newton spaces are real; their complexifications are used for
analytic maps and estimates. Dot products in algebraic differential
identities are extended complex-bilinearly; Hilbert inner products
and norms retain their Hermitian meaning.
\end{definition}

\begin{lemma}[Completions and shape regularity]\label{inf:space-embedding}
The space $\Cc$ embeds injectively into invariant normalized
$L^2(B)$ with norm at most one. Polynomial partial sums are dense.
For each fixed $k$, harmonic extension sends $\Ec_1$ continuously
into $C^k(\overline B)$, with constants depending on $k$ and the
dimension. The first two boundary derivative bounds are independent
of $p$.
In particular
\begin{equation}\label{inf:shape-C2}
 \nrm h_{C^2(S^{2p-1})}\le C\nrm h_{\Ec_1},\qquad
 \nrm{D((\HP h)y)}_{\mathrm{op},\infty}\le C\nrm h_{\Ec_1}.
\end{equation}
The boundary function $h$ is real analytic.
\end{lemma}
\begin{proof}
By \eqref{inf:radial-norm}, every basis function has $L^2$ norm at
most one. The triangle inequality gives the continuous embedding
on finite sums and hence on their completion. Its injectivity
follows by taking the inner product with each orthogonal basis
function: all coefficients of a zero $L^2$ limit vanish.
Weighted $\ell^1$ completeness and truncation give the remaining
Banach-space assertions.

Lemma~\ref{inf:angular-sup} bounds the first two boundary
derivatives by $2j$ and $4j^2$. The ratios of these quantities to
$(1+j)2^j\wt(2j)$ are uniformly bounded. Summing proves the first
estimate in \eqref{inf:shape-C2}. Radial homogeneity gives
$\partial_\varrho S_j=2j\varrho^{2j-1}G_j$;
the tangential bound gives $|\nabla S_j|\le2\sqrt2j$ in the ball.
The product rule proves the second estimate.
For the full extension, choose any fixed $1<R_*<\sqrt2$.
Homogeneity and the endpoint bound give
$|S_j(y)|\le R_*^{2j}$ for real $|y|\le R_*$.
Thus the series for $\HP h$ converges uniformly on each compact
subset of $B_{\sqrt2}$, with supremum there controlled by the
shape norm. Its finite sums are harmonic polynomials. Passing
to the limit in the mean-value identity shows that the limit is
harmonic in that larger ball. The Poisson formula on a ball of
fixed radius about each point of $\overline B$ can therefore be
differentiated any fixed number of times. It gives uniform
convergence of all those derivatives on $\overline B$, with
constants depending on the derivative order and fixed dimension.
It also proves real analyticity of the extension on a
neighbourhood of $\overline B$. The dimension-independent first
two boundary bounds are the separate estimates above.

To see analyticity of the boundary directly, expand $G_j$ in
Chebyshev polynomials on $[0,1]$. Its real supremum is at most one,
so each coefficient has absolute value at most two. On any
Bernstein ellipse of parameter $1<\rho_{\mathrm{ell}}<2$ this gives
$|G_j(z)|\le2(j+1)\rho_{\mathrm{ell}}^j$.
The series for $h$ converges uniformly on that ellipse, since
$\rho_{\mathrm{ell}}<2$, and defines a holomorphic function there. Composing
with the polynomial $x=|X|^2$ on the real sphere proves real
analyticity of the boundary function.
\end{proof}

\subsection{The positive multiplier module}

For endpoint-normalized Jacobi polynomials with parameters
$\alpha,\beta>-1$, put $u=\alpha+\beta+1$ and $v=\alpha-\beta$.
Gasper's theorem \cite[Theorem~1]{inf:gasper} gives nonnegative
linearization coefficients provided
\begin{equation}\label{inf:gasper-hypothesis}
 \alpha\ge\beta,\qquad
 u(u+5)(u+3)^2\ge(u^2-7u-24)v^2.
\end{equation}
The normalization is at the endpoint corresponding to the larger
Jacobi exponent $\alpha$. In the two-block parameters
$\ab=a/2\le\bb=p-a/2$, this is $x=1$, the endpoint of the
factor $(1-x)^{\bb-1}$.
Here $\alpha=\bb-1\ge\beta=\ab-1$, $u=p-1$, and
$0\le v=\bb-\ab\le p/2$. For $p\ge64$ the margin in
\eqref{inf:gasper-hypothesis} is at least
\[
 \frac{3p^4+37p^3+64p^2-16p-64}{4}>0.
\]
Thus in our normalization
\begin{equation}\label{inf:gasper-product}
 G_iG_j=\sum_{|i-j|\le k\le i+j}c_{ij}^kG_k,
 \qquad c_{ij}^k\ge0,\qquad\sum_kc_{ij}^k=1.
\end{equation}
The upper support follows from polynomial degree; the lower support
follows from orthogonality, since a polynomial of degree below $i$
has zero inner product with $G_i$. The sum equals one by evaluating
at $x=1$.

\begin{lemma}[Positive radial stencils]\label{inf:radial-stencils}
For $\beta\ge0$ and $s\ge0$,
\begin{align}
 J_s^\beta&=\frac{s+\beta+1}{2s+\beta+1}J_s^{\beta+1}
              +\frac{s}{2s+\beta+1}J_{s-1}^{\beta+1},
 \label{inf:raising-stencil}\\
 tJ_s^{\beta+1}&=\frac{s+\beta+1}{2s+\beta+2}J_s^\beta
              +\frac{s+1}{2s+\beta+2}J_{s+1}^\beta.
 \label{inf:lowering-stencil}
\end{align}
The term with index $-1$ in the first identity is zero.
All coefficients are nonnegative and sum to one.
Multiplication by $t$ in a fixed radial basis also has three
nonnegative coefficients of sum one.
\end{lemma}
\begin{proof}
Insert the finite binomial formula
\eqref{inf:radial-expansion}. Comparing each power of $t$ and using
$\binom{z+1}{s}=\binom z s+\binom z{s-1}$ gives
\eqref{inf:raising-stencil} and \eqref{inf:lowering-stencil}.
Their coefficients have the displayed signs and sums.
For multiplication by $t$ in a fixed basis, orthogonality makes
its matrix tridiagonal. The upper coefficient is the positive
ratio of leading coefficients; symmetry and the positive squared
norms give a positive lower coefficient. The diagonal coefficient
is $\int t(J_s^\beta)^2t^\beta dt/\int(J_s^\beta)^2t^\beta dt$,
which is nonnegative. Evaluation at $t=1$ gives their sum one.
\end{proof}

\begin{proposition}[Solid multipliers and the shape algebra]
\label{inf:module}
For integers $i,h\ge0$, multiplication by $S_it^h$ is bounded on $\Cc$ and
\begin{equation}\label{inf:module-bound}
 \nrm{M_{S_it^h}}_{\Cc\to\Cc}\le2^i\wt(2i+2h).
\end{equation}
The boundary space $\Ec_1$ is a Banach algebra, with
$\nrm{h_1h_2}_{\Ec_1}\le\nrm{h_1}_{\Ec_1}\nrm{h_2}_{\Ec_1}$.
In particular
\begin{equation}\label{inf:harmonic-module}
 \nrm{M_{\HP h}}\le\nrm h_{\Ec_1},\qquad
 \nrm{M_{\Euler\HP h}}\le2\nrm h_{\Ec_1}.
\end{equation}
Here $\Euler=y\cdot\nabla$.
\end{proposition}
\begin{proof}
A product $S_it^h\Psi_{j,s}$ has, by
\eqref{inf:gasper-product}, output profiles
$t^{i+j-k+h}J_s^{\beta_j}$ with positive angular coefficients of
sum one. If $k\ge j$, raise the radial parameter by $2(k-j)$
using \eqref{inf:raising-stencil}, then multiply by the remaining
powers of $t$. If $k<j$, use $2(j-k)$ powers for lowering by
\eqref{inf:lowering-stencil}. There are enough powers because
$i+j-k\ge2(j-k)$ is equivalent to $k\ge j-i$, part of the
triangle support. The remaining powers use the fixed-parameter
positive multiplication stencil.
Every conversion has coefficient sum one. The output total degree
is at most the input degree plus $2i+2h$, and $2^k\le2^{i+j}$.
Equation~\eqref{inf:weight-properties} proves
\eqref{inf:module-bound}, first on basis columns and then by
absolute summation and completion.

For boundary products use the same angular support,
$1+k\le(1+i)(1+j)$, and
$\wt(2k)\le\wt(2i)\wt(2j)$.
This proves the algebra inequality. The harmonic multiplier
bounds follow by summing \eqref{inf:module-bound} and using
$\Euler S_i=2iS_i$.
These coefficient multipliers equal pointwise multiplication
on $\Cc$: equality holds on finite source sums, which converge
in $\Cc$ and hence in $L^2$, while multiplication by a fixed
solid polynomial is continuous in $L^2$.
For a shape extension its finite sums converge uniformly on
the ball and in multiplier norm, so both limits give the same
pointwise multiplier.
\end{proof}

\subsection{Exact Poisson columns and traces}

Let $\KD=\Delta_D^{-1}$ on the unit ball, so $\KD=-H_0^{-1}$.
Its sign is part of this convention. The same operator restricted
to $\Gc_p$ will occasionally be denoted by $K$.

\begin{lemma}[Poisson columns]\label{inf:poisson-columns}
Fix an angular input $j$, and put $B=2j$, $\beta=B+p-1$,
$b_0=\beta+1$, and $d_s=\beta+2s$.
For $s\ge1$,
\begin{align}
 \KD\Psi_{j,s}&=S_jk_s(t),\nonumber\\
 k_s&=\frac{J_{s-1}^\beta}{4d_s(d_s+1)}
 -\frac{J_s^\beta}{2d_s(d_s+2)}
 +\frac{J_{s+1}^\beta}{4(d_s+1)(d_s+2)}.\label{inf:poisson-profile}
\end{align}
For $s=0$,
\begin{equation}\label{inf:poisson-harmonic}
 k_0=\frac{t-1}{4b_0}
     =\frac{J_1^\beta-J_0^\beta}{4b_0(b_0+1)}.
\end{equation}
For $s\ge1$ the derivative identities are
\begin{align}
 k_s'&=\frac{J_s^{\beta+1}-J_{s-1}^{\beta+1}}{4(d_s+1)},
 \label{inf:poisson-derivative}\\
 \beta k_s+tk_s'&=
 \frac{J_{s+1}^{\beta-1}-J_s^{\beta-1}}{4(d_s+1)}.
 \label{inf:poisson-lowered}
\end{align}
For $s=0$ their replacements are
$k_0'=1/(4b_0)$ and
$\beta k_0+tk_0'=J_1^{\beta-1}/(4b_0)$.
\end{lemma}
\begin{proof}
Write $J_s^\beta(t)=\sum_{m=0}^s a_{s,m}t^m$, where
\eqref{inf:radial-expansion} gives
\[
 a_{s,m}=(-1)^{s-m}\binom sm\binom{s+\beta+m}s.
\]
The coefficient equation for
$4\{tk_s''+(\beta+1)k_s'\}=J_s^\beta$ is
\[
 [t^{m+1}]k_s=\frac{a_{s,m}}{4(m+1)(\beta+m+1)}.
\]
For $1\le m\le s$, the adjacent-layer ratios are
\[
 \frac{a_{s-1,m}}{a_{s,m}}=-\frac{s-m}{s+\beta+m},\qquad
 \frac{a_{s+1,m}}{a_{s,m}}=-\frac{s+\beta+m+1}{s+1-m},
\]
whereas
\[
 \frac{a_{s,m-1}}{a_{s,m}}
 =-\frac{m(\beta+m)}{(s+1-m)(s+\beta+m)}.
\]
With $d=d_s=\beta+2s$, the coefficient in the proposed formula,
divided by $a_{s,m}$, is therefore
\begin{align*}
 &-\frac{s-m}{4d(d+1)(s+\beta+m)}-\frac1{2d(d+2)}
 -\frac{s+\beta+m+1}{4(d+1)(d+2)(s+1-m)}\\
 &\hspace{35mm}=-\frac1{4(s+1-m)(s+\beta+m)}.
\end{align*}
The equality follows by multiplication by the common denominator
and substitution of $d=\beta+2s$; its right side is exactly
$a_{s,m-1}/\{4m(\beta+m)a_{s,m}\}$.
For the top coefficient use
\[
 \frac{a_{s+1,s+1}}{a_{s,s}}
 =\frac{(d+2)(d+1)}{(s+1)(s+\beta+1)}.
\]
Finally the three column coefficients sum to zero, fixing the
constant coefficient by $k_s(1)=0$. This proves
\eqref{inf:poisson-profile} for every $s\ge1$.

For \eqref{inf:poisson-derivative}, the coefficient of $t^m$ in
$J_s^{\beta+1}-J_{s-1}^{\beta+1}$, divided by $a_{s,m}$, equals
\[
 \frac{s+\beta+m+1}{\beta+m+1}
 +\frac{s-m}{\beta+m+1}
 =\frac{d+1}{\beta+m+1}.
\]
This is exactly $4(d+1)[t^m]k_s'/a_{s,m}$.
For \eqref{inf:poisson-lowered}, at $m\ge1$ the left coefficient is
$(\beta+m)[t^m]k_s=a_{s,m-1}/(4m)$.
The two binomial ratios in the right numerator give
\[
 \frac{[t^m](J_{s+1}^{\beta-1}-J_s^{\beta-1})}{a_{s,m-1}}
 =\frac{s+\beta+m}{m}+\frac{s+1-m}{m}
 =\frac{d+1}{m}.
\]
Their nonconstant coefficients agree. The previous derivative
identity gives $k_s'(1)=0$; both sides of
\eqref{inf:poisson-lowered} vanish at one, so their constants
agree as well.
For $s=0$, the elementary expression
\eqref{inf:poisson-harmonic} solves the equation and gives both
replacement derivative identities directly. In particular this
layer has the denominator $4b_0$ in \eqref{inf:poisson-harmonic}.
By \eqref{inf:profile-laplacian}, these polynomial functions solve
the Dirichlet Poisson problem. Uniqueness identifies them with
$\KD\Psi_{j,s}$.
\end{proof}

\begin{proposition}[Poisson bounds and exact clamping]
\label{inf:poisson-bounds}
On all of $\Cc$, including the layer $s=0$,
\begin{equation}\label{inf:poisson-operator-bounds}
 \nrm{\KD}\le\frac3{p(p+1)},\qquad
 \nrm{\Euler\KD}\le C/p,\qquad
 \nrm{\Euler^2\KD}\le C.
\end{equation}
More precisely, its weighted basis columns satisfy, for all $j,s\ge0$,
\begin{align}
 \nrm{\KD\Psi_{j,s}}_\Cc
 &\le\frac{C}{(p+2j+2s)^2}\nrm{\Psi_{j,s}}_\Cc,\nonumber\\
 \nrm{\Euler\KD\Psi_{j,s}}_\Cc
 &\le\frac{C}{p-1+2j+2s}\nrm{\Psi_{j,s}}_\Cc.
 \label{inf:poisson-column-bounds}
\end{align}
The normal trace is the bounded map
\begin{equation}\label{inf:poisson-trace}
 \partial_\nu\KD f=
 \sum_{j\ge0}\frac{f_{j,0}}{2(2j+p)}G_j,
 \qquad
 \nrm{\partial_\nu\KD}_{\Cc\to\Ec_1}\le\tfrac12.
\end{equation}
Consequently $U=1+\KD g$ has Dirichlet trace one and normal trace
zero exactly when $g\in\Gc_p$.
\end{proposition}
\begin{proof}
The absolute sums of the columns in
\eqref{inf:poisson-profile} and \eqref{inf:poisson-harmonic} are
$1/(d_s(d_s+2))$ and $1/(2b_0(b_0+1))$.
The only upward radial shift increases degree by two, whose
weight ratio is at most $e^{2/p}<2$. This proves the first bound.
On one profile,
\[
 \Euler\KD=Bk_s+2tk'_s,\qquad
 \Euler^2\KD=B^2k_s+4(1-p)tk'_s+tJ_s^\beta.
\]
The second equality uses the Poisson equation to eliminate $k_s''$.
The derivative column in parameter $\beta+1$ has sum at most
$1/(2(d_s+1))$; multiplication by $t$ lowers it positively to
parameter $\beta$. The separate $s=0$ column gives the same bound
up to an absolute constant. Since $B\le d_s$ and $d_s\ge p-1$,
these formulas prove the other two operator bounds and
\eqref{inf:poisson-column-bounds}.

At the boundary $S_j=G_j$ and
$\partial_\nu(S_jk_s)=2jk_s(1)G_j+2k'_s(1)G_j$.
Only $s=0$ contributes, giving \eqref{inf:poisson-trace}.
Its weighted trace bound follows from
$(1+j)/(2(2j+p))\le1/2$ and
$\wt(2j)\le\wt(2j+2s)$.
For a general coefficient source, approximate by finite sums.
The embedding in $L^2$ and Dirichlet $H^2$ regularity give
$\KD f_n\to\KD f$ in $H^2(B)$.
The displayed trace maps converge in $\Ec_1$ and thus in every
fixed boundary Sobolev norm, by exponential angular decay.
The Sobolev trace theorem identifies this limit with the actual
normal trace of the $H^2$ solution. Dirichlet traces are zero by
Poisson construction. Finally orthogonality of the $G_j$ shows
that the trace vanishes if and only if each $f_{j,0}$ vanishes.
This proves the asserted characterization on the completed spaces.
\end{proof}

\subsection{Raising coupling and the clamped gain}

\begin{lemma}[Contraction under parameter raising]\label{inf:raising-coupling}
For $\beta\ge0$ and integers $s\ge1$, $m\ge0$, the coefficient
vectors of $J_s^\beta$ and $J_{s-1}^\beta$ in parameter
$\beta+m$ are probability vectors: their entries are nonnegative
and sum to one by \eqref{inf:raising-stencil}. Their $\ell^1$
distance is at most
\begin{equation}\label{inf:coupling-bound}
 2e^2\frac z{z+m},\qquad z=2s+\beta+1.
\end{equation}
The bound is uniform in all its parameters.
\end{lemma}
\begin{proof}
One raising step in \eqref{inf:raising-stencil} keeps index
$\sigma$ or lowers it by one, with down probability
$q_\sigma(b)=\sigma/(2\sigma+b+1)$.
Couple the chains started at $s,s-1$ with one uniform random
number at each step. Since $q_\sigma$ increases in $\sigma$,
the indices remain adjacent or merge, and after merging remain
equal. Before merging, at step $h$ their upper index is at most
$s$ and their parameter is $b=\beta+h$. Their merging probability
is at least
\[
 q_\sigma(b)-q_{\sigma-1}(b)
 =\frac{b+1}{(2\sigma+b+1)(2\sigma+b-1)}
 \ge\frac1{z+h}-\frac{2s}{(z+h)^2}.
\]
This also holds when the lower index is zero.
The probability of no merger by step $m$ is at most
\[
 \exp\left(-\sum_{h<m}\frac1{z+h}
                     +2s\sum_{h<m}\frac1{(z+h)^2}\right)
 \le e^2\frac z{z+m}.
\]
Here the first sum is at least $\log((z+m)/z)$, and the second
multiplied by $2s$ is at most
$2s(z^{-1}+z^{-2})\le2$.
The total variation distance of the two distributions is no greater
than the nonmerger probability. Their $\ell^1$ distance is twice
total variation, giving \eqref{inf:coupling-bound}.
Every subsequent positive stencil of sum one is an $\ell^1$
contraction, so it preserves this estimate.
\end{proof}

\begin{lemma}[Clamped multiplication gain]\label{inf:clamped-smoothing}
For integers $c,h\ge0$ put $T=S_ct^h$ and $D=2c+2h$.
On clamped sources, that is on $\Gc_p$,
\begin{equation}\label{inf:clamped-gain}
 \nrm{M_T\Euler K}_{\Gc_p\to\Cc}
 \le C\frac{2^c\wt(D)}{p+D}.
\end{equation}
The constant is independent of $c,h$, the source indices, $p$ and
the admissible split.
\end{lemma}
\begin{proof}
Use a source $\Psi_{j,s}$ with $s\ge1$, and put $B=2j$,
$\beta=B+p-1$, $d_s=\beta+2s$.
For an angular output of degree $L=2k$ set
\[
 n_{\downarrow}=h+(2c+B-L)/2,\qquad
 n_{\uparrow}=h+(2c-B+L)/2.
\]
There are $n_{\downarrow}$ powers of $t$. First raise the source
parameter by $n_{\uparrow}$, then lower it with those
$n_{\downarrow}$ powers; the net change is
$n_{\uparrow}-n_{\downarrow}=L-B$.
The triangle support guarantees $n_{\downarrow},n_{\uparrow}\ge0$.
If $D\le8(B+p)$, the ordinary multiplier bound and the column
bound $C/d_s$ for $\Euler K$ prove the claim, since
$p+D\le9(B+p)\le9(d_s+1)$.

Suppose $D>8(B+p)$. Then $n_{\uparrow}\ge D/4$.
If $2c<D/2$, one has $h>D/4$ and $n_{\uparrow}\ge h$.
If $2c\ge D/2$, the triangle inequality gives
$n_{\uparrow}\ge h+2c-B>D/4$.
Since $s\ge1$, the Poisson column has the zero-sum form
\[
 k_s=\frac{J_{s-1}^\beta-J_s^\beta}{4d_s(d_s+1)}
       +\frac{J_{s+1}^\beta-J_s^\beta}{4(d_s+1)(d_s+2)}.
\]
After $n_{\uparrow}$ raises, Lemma~\ref{inf:raising-coupling} bounds its
coefficient sum by $C/(d_s(d_s+n_{\uparrow}))$.
The term $2tk_s'$ begins with the adjacent difference in
\eqref{inf:poisson-derivative}, in parameter $\beta+1$.
Raise it by $n_{\uparrow}$ and then lower with $n_{\downarrow}+1$
powers, including its extra $t$. Its coefficient sum is at most
$C/(d_s+n_{\uparrow})$.
Since $B\le d_s$, the combination $Bk_s+2tk_s'$ is bounded by
$C/(d_s+n_{\uparrow})\le C/(p+D)$.

All output degrees are at most the input degree plus $D+2$.
The extra factor $\wt(2)$ is bounded by two, and angular weights
are at most $2^{c+j}$. Summing the positive angular coefficients
in \eqref{inf:gasper-product}, and using
\eqref{inf:weight-properties}, proves \eqref{inf:clamped-gain}.
\end{proof}

\subsection{Gradient estimates with one shape moment}

\begin{lemma}[Gradient profiles on every source layer]
\label{inf:gradient-profiles}
Let $A=2i>0$, $B=2j$, $\ell_{\mathrm{out}}=2k$, and $h=(A+B-\ell_{\mathrm{out}})/2$ for an
angular product output in \eqref{inf:gasper-product}. Put
$Z_{\mathrm{grad}}=h+\ell_{\mathrm{out}}+p-1$.
For $\KD\Psi_{j,s}=S_jk_s$ the two output profiles
of $\nabla S_i\cdot\nabla\KD$ and
$\nabla S_i\cdot\nabla\Euler\KD$, with the coefficient
$c_{ij}^k$ suppressed, are respectively
\begin{align}
 &2hZ_{\mathrm{grad}}t^{h-1}k_s+2At^hk'_s,\label{inf:gradient-profile-one}\\
 &At^hJ_s^\beta+
 (4hZ_{\mathrm{grad}}-2AB-4A(p-1))t^hk'_s
                  +2BhZ_{\mathrm{grad}}t^{h-1}k_s.\label{inf:gradient-profile-two}
\end{align}
The term with $h=0$ and a negative power of $t$ is absent.
Consequently, including $s=0$,
\begin{align}
 \nrm{\nabla S_i\cdot\nabla\KD}_{\Cc\to\Cc}
 &\le CA2^i\wt(A)/p,\label{inf:gradient-C-one}\\
 \nrm{\nabla S_i\cdot\nabla\Euler\KD}_{\Cc\to\Cc}
 &\le CA2^i\wt(A).\label{inf:gradient-C-two}
\end{align}
\end{lemma}
\begin{proof}
Since both solid factors are harmonic,
$2\nabla S_i\cdot\nabla S_j=\Delta(S_iS_j)$.
Use $S_iS_j=\sum c_{ij}^kS_kt^h$ and
\eqref{inf:profile-laplacian}; its coefficient is $4hZ_{\mathrm{grad}}$.
The derivative of the radial profile contributes $2At^hk_s'$,
proving \eqref{inf:gradient-profile-one}.
Apply that formula to $Bk_s+2tk_s'$.
The Poisson equation gives
$(Bk_s+2tk_s')'=J_s^\beta/2-(B+2p-2)k_s'$;
substitution proves \eqref{inf:gradient-profile-two}.
The triangle support gives $0\le h\le\min(A,B)$ and
\begin{equation}\label{inf:hZ-bound}
 hZ_{\mathrm{grad}}=h(A+B-h+p-1)\le A(B+p-1).
\end{equation}
For $A\le B$ the left side is increasing up to $h=A$.
For $A\ge B$ use its value at $h=B$ and
$B(A+p-1)\le A(B+p-1)$.

In every nonsaturated output the positive parameter conversions of
Lemma~\ref{inf:radial-stencils} apply separately to the two terms
in \eqref{inf:gradient-profile-one}.
When $\ell_{\mathrm{out}}<B$, the required lowerings are $B-\ell_{\mathrm{out}}$ for the $k_s$ term
and $B-\ell_{\mathrm{out}}+1$ for the derivative term, whose parameter is $\beta+1$.
The available powers are $h-1$ and $h$, respectively.
They suffice whenever $\ell_{\mathrm{out}}\ge B-A+2$; the degrees are even.
When $\ell_{\mathrm{out}}=B$, the $k_s$ term keeps its parameter.
The derivative term starts in parameter $\beta+1$ and needs one
lowering; one of its $h=i\ge1$ powers of $t$ supplies it.
When $\ell_{\mathrm{out}}>B$, raise the $k_s$ term by
$\ell_{\mathrm{out}}-B$ and the derivative term by
$\ell_{\mathrm{out}}-B-1\ge1$, followed by their nonnegative
remaining powers of $t$. The case $h=0$ omits the first term.
Thus the only remaining case is $\ell_{\mathrm{out}}=B-A$, $h=A$, where the two
terms combine exactly to
\[
 2At^{A-1}(\beta k_s+tk_s').
\]
Use \eqref{inf:poisson-lowered} before the remaining $A-1$
lowerings. This supplies the one missing lowering step.
The separate harmonic-column identity in Lemma~\ref{inf:poisson-columns}
proves the same statement at $s=0$.

The absolute Poisson column sums are $O((B+p+2s)^{-2})$, and
the two first-derivative columns have sums $O((B+p+2s)^{-1})$.
Together with \eqref{inf:hZ-bound}, these give a per-output bound
$CA/(B+p)$ for \eqref{inf:gradient-profile-one}.
In \eqref{inf:gradient-profile-two} the three terms have bounds
$CA$, $4CA$, and $2CA$ after positive conversion.
In the saturated case use their combined expression
\[
 At^AJ_s^\beta+2ABt^{A-1}(\beta k_s+tk_s'),
\]
which gives the same bound, also at $s=0$.
In that layer the coefficient norm of $k_0$ is
$1/(2b_0(b_0+1))$, which supplies the factor $1/(B+p)$.
All final degrees are at most input degree plus $A$.
The positive angular coefficients sum to one and
$2^k\le2^{i+j}$. Applying \eqref{inf:weight-properties} proves
\eqref{inf:gradient-C-one}--\eqref{inf:gradient-C-two}.
\end{proof}

\begin{proposition}[Grouped shape derivatives]\label{inf:grouped-shapes}
For $H_i=\HP h_i$, the following bounds hold, with absolute
constants:
\begin{align}
 \nrm{M_{\nabla H_1\cdot\nabla H_2}}_{\Cc\to\Cc}
 &\le Cp\nrm{h_1}_{\Ec_1}\nrm{h_2}_{\Ec_1},
 \label{inf:gradient-multiplier}\\
 \nrm{M_{\Euler(\Euler-1)H_1}\Euler K}_{\Gc_p\to\Cc}
 &\le C\nrm{h_1}_{\Ec_1},\label{inf:second-shape-clamped}\\
 \nrm{M_{\nabla H_1\cdot\nabla\Euler H_2}\Euler K}_{\Gc_p\to\Cc}
 &\le C\nrm{h_1}_{\Ec_1}\nrm{h_2}_{\Ec_1}.
 \label{inf:mixed-shape-clamped}
\end{align}
The last two bounds are on clamped sources in $\Gc_p$.
\end{proposition}
\begin{proof}
For a pair of harmonics of degrees $A=2i,B=2j$, the product of
gradients has output coefficient $2hZ_{\mathrm{grad}}$ and solid profile
$S_kt^{h-1}$. It is zero if $h=0$; otherwise its total degree is
$A+B-2$. From \eqref{inf:hZ-bound},
$2hZ_{\mathrm{grad}}\le Cp(1+i)(1+j)$.
The module bound and positive angular summation prove
\eqref{inf:gradient-multiplier}.
For the second inequality apply Lemma~\ref{inf:clamped-smoothing}
to the harmonic degree $A$, with coefficient $A(A-1)$.
The quotient $A(A-1)/(p+A)$ is at most $A$, leaving one shape
moment.
For the last inequality its pairwise coefficient is $2BhZ_{\mathrm{grad}}$.
Apply the same clamped gain at degree $D=A+B-2$.
For $h\ge1$, $Z_{\mathrm{grad}}=A+B-h+p-1\le p+A+B-2$, so
\[
 \frac{2BhZ_{\mathrm{grad}}}{p+A+B-2}\le2B\min(A,B)\le2AB.
\]
The grading and angular weights are bounded by the product of the
two shape weights. Summation gives \eqref{inf:mixed-shape-clamped}.
\end{proof}

\subsection{Double-trace lifts and source moments}

\begin{lemma}[Exact double-trace lift]\label{inf:double-lift}
For each angular index define
\begin{equation}\label{inf:lift-columns}
 L_j^D=S_j[1-j(t-1)],\qquad L_j^N=\tfrac12S_j(t-1).
\end{equation}
Their Dirichlet and radial normal traces are $(G_j,0)$ and
$(0,G_j)$, respectively. Their coefficient bounds and Laplacians
are
\begin{align}
 \nrm{L_j^D}_\Cc&\le3\,2^j\wt(2j+2),&
 \nrm{L_j^N}_\Cc&\le\frac{2^j\wt(2j+2)}{2j+p+1},
 \label{inf:lift-column-bounds}\\
 \Delta L_j^D&=-2j(4j+2p)S_j,&
 \Delta L_j^N&=(4j+2p)S_j.\label{inf:lift-laplacians}
\end{align}
For the defects of Corollary~\ref{inf:centre-coefficient-defect},
let $f_N=h_N/R$, $q_N=1+\HP f_N$ and
$d_N=q_N+\Euler q_N$. Define
\begin{equation}\label{inf:reference-defect}
 N_N^{\mathrm{ref}}=R(d_N|_{\partial B})N_N^{\mathrm{phys}},
 \qquad
 \ell_N=\sum_j(D_{N,j}L_j^D+N_{N,j}^{\mathrm{ref}}L_j^N).
\end{equation}
For every fixed moment $k$ and $p\ge P_{N,k}$,
\begin{equation}\label{inf:lift-estimates}
 \nrm{\ell_N}_{\Cc^{[k]}}\le C_{N,k}\theta^2p^{-N},\qquad
 \nrm{\Delta\ell_N}_{\Cc^{[k]}}\le C_{N,k}\theta^2p^{2-N}.
\end{equation}
\end{lemma}
\begin{proof}
Differentiating \eqref{inf:lift-columns} at $t=1$ gives their
traces. The identity
$t-1=(J_1^{\beta_j}-J_0^{\beta_j})/(\beta_j+2)$
places both lifts in radial layers zero and one and gives
\eqref{inf:lift-column-bounds}. Formula
\eqref{inf:profile-laplacian} proves \eqref{inf:lift-laplacians}.
The reference normal trace in \eqref{inf:reference-defect} follows
by differentiating the physical composition $u_N(Rq_N(y)y)$ in
the radial direction on the sphere: $D(q_Ny)\nu=d_N\nu$.

The multiplier $d_N|_{\partial B}$ is a fixed-degree polynomial
with uniformly bounded coefficients at any fixed angular radius.
Thus \eqref{inf:centre-coeff-defect-bound} bounds its product with
the physical defect, losing only the factor $R=O(p)$.
The factor $(2j+p+1)^{-1}$ in the normal lift cancels that factor.
Choose a slightly larger fixed angular radius in
\eqref{inf:centre-coeff-defect-bound} to absorb
$\wt(2j+2)\le e^{(2j+2)/p}$ and the prescribed angular moments.
This proves the first estimate in \eqref{inf:lift-estimates}.
The Laplacians in \eqref{inf:lift-laplacians}, including the extra
$p$ in the reference normal trace, cost at most $C p^2$ after
absorbing fixed powers of $j$ by the same radius choice. This proves
the second estimate. The sums converge in all stipulated moment
norms and their traces agree with the classical lifts by the same
coefficient identities and smooth convergence.
\end{proof}

\begin{lemma}[Derivative columns and moments]\label{inf:derivative-moments}
For $\beta\ge0$, $s\ge1$, and $0\le k<s$,
\begin{equation}\label{inf:derivative-column}
 \partial_tJ_s^\beta=\sum_{k<s}d_{s,k}J_k^\beta,
 \quad d_{s,k}=(\beta+2k+1)
 \left(1-(-1)^{s+k}\frac{\binom{k+\beta}{k}}{\binom{s+\beta}{s}}\right)
 \ge0.
\end{equation}
For every $0\le m\le s$ the endpoint derivative is
\begin{equation}\label{inf:Jacobi-endpoint-derivatives}
 (J_s^\beta)^{(m)}(1)
 =\frac{(s)_{\mathrm{fall},m}(s+\beta+1)_m}{m!}.
\end{equation}
In particular, the first derivative column has sum
$s(s+\beta+1)$ and the second has sum
$s(s-1)(s+\beta+1)(s+\beta+2)/2$.
Consequently
\begin{equation}\label{inf:Euler-moment}
 \nrm{\Euler u}_{\Cc^{[k]}}\le Cp\nrm u_{\Cc^{[k+2]}}.
\end{equation}
The Laplacian has positive radial columns of sum at most
$C(p+2j+2s)^4$ and does not raise total degree. For each fixed $k\ge0$,
\begin{equation}\label{inf:Delta-moment}
 \nrm{\Delta u}_{\Cc^{[k]}}
 \le C_k\{p^4\nrm u_{\Cc^{[k]}}+\nrm u_{\Cc^{[k+4]}}\}.
\end{equation}
\end{lemma}
\begin{proof}
For $\beta>0$, integration by parts gives, since $k<s$,
\[
 \int_0^1t^\beta(J_s^\beta)'J_k^\beta\,dt
 =1-\beta J_k^\beta(0)\int_0^1t^{\beta-1}J_s^\beta\,dt.
\]
To obtain this expression, subtract $J_k^\beta(0)$ in the
integrated derivative term; its remaining factor is $t$ times a
polynomial of degree below $s$, hence orthogonal to $J_s^\beta$.
Rodrigues' formula and $s$ integrations by parts give the remaining
integral as $(-1)^sB(\beta,s+1)$.
Divide by the squared norm $(\beta+2k+1)^{-1}$ and use the endpoint
value at zero to obtain \eqref{inf:derivative-column}.
The formula at $\beta=0$ is its finite continuous limit.
Monotonicity in $k$ of $\binom{k+\beta}{k}$ proves nonnegativity.
Every further differentiation preserves coefficient positivity.
For the endpoint values, write $J_s^\beta(1+z)=\sum_m\mathfrak d_mz^m$
and use its differential equation
\[
 t(1-t)J''+\{\beta+1-(\beta+2)t\}J'
                      +s(s+\beta+1)J=0.
\]
Its coefficient of $z^m$ gives
$\mathfrak d_{m+1}=(s-m)(s+\beta+1+m)\mathfrak d_m/(m+1)^2$,
with $\mathfrak d_0=1$.
Iteration and multiplication by $m!$ prove
\eqref{inf:Jacobi-endpoint-derivatives}. Evaluating the positive
columns at one now gives their sums.

On one solid profile $\Euler=B+2t\partial_t$.
The positive $t$ stencil gives column sum
$B+2s(s+\beta+1)\le Cp(1+2j+2s)^2$ and no increase in total degree.
This proves \eqref{inf:Euler-moment}.
The profile Laplacian is $4(t\partial_t^2+(\beta+1)\partial_t)$,
so its column sum is
$4((J_s^\beta)''(1)+(\beta+1)(J_s^\beta)'(1))$,
bounded by $C(p+2j+2s)^4$.
Each output degree is at most the input degree $D=2j+2s$, so
its moment weight $(1+D_{\mathrm{out}})^k$ is at most $(1+D)^k$.
Multiplying the column bound by this weight and using
$(p+D)^4\le8(p^4+D^4)$ gives \eqref{inf:Delta-moment} after
absolute summation.
\end{proof}

\section{The nonlinear map and the material isomorphism}\label{inf:sec-material}

The chart in this section uses a harmonic extension on the whole ball.
It is different from the homogeneous radial stretching used to derive
the spectral flow. Both charts describe the same boundary domain at a
fixed centre; only its spectrum is passed from one description to the
other.

\subsection{A smooth scalar chart on the ball}

For a normalized boundary shape $f\in\Ec_1$ put
\begin{equation}\label{inf:scalar-chart}
 H_f=\HP f,\qquad q=1+H_f,\qquad d=q+\Euler q,
 \qquad \Phi_f(y)=q(y)y.
\end{equation}
For a field $U$ on the reference ball define
\begin{equation}\label{inf:scalar-pullback-definition}
 P_{\Phi_f}\Delta U=
       [\Delta(U\circ\Phi_f^{-1})]\circ\Phi_f,
 \qquad B_{\Phi_f}=P_{\Phi_f}\Delta+R^2.
\end{equation}
For a smooth real chart, $B_{\Phi,D}$ denotes the restriction of
$B_\Phi$ to the invariant Dirichlet graph domain
$H^2(B)\cap H^1_0(B)$.
Fix the analytic chart radius $\varepsilon_0=1/2$ and the
working radius $\varepsilon_{\mathrm{work}}=1/4$.
For every real $f\in\Ec_1$ with $\nrm f_{\Ec_1}<1/2$, the map
$\Phi_f$ is a diffeomorphism of the closed ball onto the region
with radial boundary $1+f$.
Indeed $|q-1|\le\nrm f_{\Ec_1}$ and
$|d-1|\le2\nrm f_{\Ec_1}<1$ on the ball, so $q,d>0$.
On each ray, $\Phi_f(\zeta\omega)=\zeta q(\zeta\omega)\omega$
keeps the direction fixed and has strictly positive radial
derivative $d$. Distinct rays remain distinct and the radial
map is one-to-one on each ray. Its Jacobian determinant is
$q^{2p-1}d>0$, including at the origin. The inverse function
theorem and the radial endpoint give the global
identification. By Lemma~\ref{inf:space-embedding}, the harmonic
extension and chart are real analytic on a neighbourhood of
the closed ball. Both axes and the origin are ordinary points
of this chart.

\begin{lemma}[Exact scalar pullback]\label{inf:pullback}
Let $w=\nabla q$, $M_q=w\cdot w$ and
\begin{equation}\label{inf:pullback-metric}
 A_q=q^{-2}\left\{I-\frac{y\otimes w+w\otimes y}{d}
                         +\frac{M_q}{d^2}y\otimes y\right\},
 \qquad C_q=A_q:D^2q.
\end{equation}
Then
\begin{align}
 P_{\Phi_f}\Delta U={}&q^{-2}\Delta U
 -\frac2{q^2d}\nabla q\cdot\nabla\Euler U
 +\frac{M_q}{q^2d^2}\Euler(\Euler+1)U
 \label{inf:scalar-pullback}\\
 &+\left(\frac{2M_q}{qd^3}-\frac{C_q}{d}\right)\Euler U,\nonumber\\
 C_q={}&q^{-2}\left\{-\frac{2(\nabla q\cdot\nabla\Euler q-M_q)}d
                +\frac{M_q\Euler(\Euler-1)q}{d^2}\right\}.
 \label{inf:Cq-identity}
\end{align}
The second identity uses harmonicity of $q$.
\end{lemma}
\begin{proof}
Since $D\Phi_f=qI+y\otimes w$, its inverse is
$q^{-1}(I-y\otimes w/d)$ and
$D\Phi_f^{-1}D\Phi_f^{-T}=A_q$.
The coordinate change formula for a scalar Laplacian is
\[
 P_{\Phi_f}\Delta U=A_q:D^2U
   -D\Phi_f^{-1}(A_q:D^2\Phi_f)\cdot\nabla U.
\]
Now $A_q:D^2\Phi_f=yC_q+2A_qw$, and multiplication by the
inverse matrix gives
\[
 D\Phi_f^{-1}(A_q:D^2\Phi_f)
 =\frac{2w}{q^2d}
   +y\left\{\frac{C_q}d-\frac{2M_q}{q^2d^2}
                             -\frac{2M_q}{qd^3}\right\}.
\]
On the other hand,
\[
 A_q:D^2U=q^{-2}\left\{\Delta U-\frac2d
 (\nabla q\cdot\nabla\Euler U-\nabla q\cdot\nabla U)
             +\frac{M_q}{d^2}\Euler(\Euler-1)U\right\}.
\]
The terms with $\nabla q\cdot\nabla U$ cancel, giving
\eqref{inf:scalar-pullback}.
Apply the last formula for $A_q:D^2U$ to $U=q$ and use
$\Delta q=0$ to get \eqref{inf:Cq-identity}.
All computations are identities of smooth functions; extension to
the coefficient spaces is established next.
\end{proof}

\begin{proposition}[Analytic map with a polynomial second derivative]
\label{inf:nonlinear-map}
For $g\in\Gc_p$ and $\nrm f_{\Ec_1}<\varepsilon_0$, define
\begin{equation}\label{inf:nonlinear-definition}
 \Fnl(g,f)=B_{\Phi_f}(1+Kg).
\end{equation}
This is an analytic map
$\Gc_p\times\{f\in\Ec_1:\nrm f_{\Ec_1}<\varepsilon_0\}\to\Cc$,
affine in $g$,
and on $\nrm f_{\Ec_1}\le\varepsilon_{\mathrm{work}}=1/4$ it satisfies
\begin{equation}\label{inf:nonlinear-second}
 \nrm{D^2\Fnl(g,f)}\le Cp(1+\nrm g_\Cc).
\end{equation}
For real $f$, its coefficient extension is exactly the weak scalar
pullback of $1+Kg\in H^2(B)$. Its complexification is the analytic
extension of the same algebraic operator formulas.
\end{proposition}
\begin{proof}
The harmonic module estimate gives
$\nrm{M_{q-1}}\le\nrm f_{\Ec_1}$. For $d-1=H_f+\Euler H_f$,
the same estimate applied to each solid coefficient gives
\[
 \nrm{M_{d-1}}
 \le\sum_{j\ge0}(1+2j)2^j\wt(2j)|f_j|
 \le2\sum_{j\ge0}(1+j)2^j\wt(2j)|f_j|
 =2\nrm f_{\Ec_1}.
\]
Their inverse powers are norm-convergent Neumann series on
$\nrm f_{\Ec_1}<1/2$. On the closed working ball of radius $1/4$,
differentiating these series any fixed number of times gives
uniformly bounded multilinear multiplier maps.
The same series converge uniformly on the closed ball, by
\eqref{inf:shape-C2}, to the indicated pointwise inverse powers.

Apply \eqref{inf:scalar-pullback} to $U=1+Kg$.
All derivatives of the constant vanish. The term $q^{-2}\Delta Kg$
is $q^{-2}g$.
Lemma~\ref{inf:gradient-profiles} controls
$\nabla q\cdot\nabla\Euler K$ with one shape moment.
The operator $\Euler(\Euler+1)K$ is bounded by
\eqref{inf:poisson-operator-bounds}, and its multiplier $M_q$
is bounded by \eqref{inf:gradient-multiplier} with norm at most
$Cp\nrm f_{\Ec_1}^2$.
For the last term group
\[
 C_q\Euler K=q^{-2}\left\{
 -2d^{-1}(\nabla q\cdot\nabla\Euler q-M_q)\Euler K
 +M_qd^{-2}[\Euler(\Euler-1)q]\Euler K\right\}.
\]
The first grouped product is bounded by
\eqref{inf:mixed-shape-clamped} and
\eqref{inf:gradient-multiplier}; the last one by
\eqref{inf:second-shape-clamped} followed by the bounded multiplier
$M_q$. The scalar multipliers commute, so the high shape derivative
is always placed immediately next to $\Euler K$ on its clamped
source, as in Proposition~\ref{inf:grouped-shapes}.
Every numerator is thus a bounded linear, bilinear, or trilinear
map from shapes to operators $\Gc_p\to\Cc$, of norm at most $Cp$.
The inverse multiplier series prove analyticity of the operator
coefficient $L_f\in\mathcal B(\Gc_p,\Cc)$ in
$\Fnl(g,f)=R^2+L_fg$.
Its first two shape derivatives have norm at most $Cp$ on the
closed working ball of radius $1/4$. The two mixed derivatives and the pure shape second
derivative give \eqref{inf:nonlinear-second}; the pure source
second derivative is zero. The term $R^2(1+Kg)$ is shape independent
and contributes no second derivative.

To identify the extension, first use polynomial sources and finite
shape sums, where the identity is classical. Approximate any
$f\in\Ec_1$ by finite sums in its norm; the charts converge in
$C^2$ and have uniformly positive $q,d$ on any fixed smaller shape
ball. The explicit inverse matrix is therefore uniformly
bounded there. For each fixed $p$, their scalar
pullbacks converge as maps $H^2(B)\to L^2(B)$ by the coordinate
formula. Next approximate $g$ in $\Cc$, hence in $L^2$.
Dirichlet regularity and uniqueness give convergence of $Kg$ in
$H^2(B)$. The bounded coefficient maps converge in $\Cc$, hence
in $L^2$, to the same pullback. This proves the assertion on the
completed spaces.
\end{proof}

\subsection{Coefficient size of the known centre}

\begin{lemma}[Fixed radial modes in coefficient norm]
\label{inf:radial-centre-norm}
Fix the finite mode set of the order-$N$ centre in
Theorem~\ref{inf:centres}. For each fixed moment $k$ and
$p\ge P_{N,k}$, its regular radial modes,
rescaled to the unit ball, have $\Cc^{[k]}$ norms at most
$C_{N,k}p^{k+1}$. The bound is uniform in $r$ and
$0<|\theta|\le20$.
\end{lemma}
\begin{proof}
For one fixed $j$ put $\nu=p+2j-1$ and write its normalized mode as
$\zeta^{1-p}J_\nu(R\zeta)G_j/d_j^{\mathrm{mode}}$.
Here $d_j^{\mathrm{mode}}$ is either $J_\nu(R)$ for value normalization or
$J'_\nu(R)-(p-1)J_\nu(R)/R$ for physical normal normalization.
The finite recurrence \eqref{inf:finite-q-recurrence} shows that
$d_j^{\mathrm{mode}}/J_{p-1}(R)$ has a nonzero analytic limit in every case.
In particular the resonant seed is normalized by its normal trace.

In the solid basis the coefficients of this mode are
\begin{equation}\label{inf:Bessel-Jacobi-coefficients}
 c_s=\frac{2(\nu+1+2s)(-1)^sJ_{\nu+1+2s}(R)}{Rd_j^{\mathrm{mode}}}.
\end{equation}
To derive them, Rodrigues' formula and $s$ integrations by parts
in $t=\zeta^2$ give the integral identity
\[
 \int_0^1\zeta^{\nu+1}J_\nu(R\zeta)J_s^\nu(\zeta^2)\,d\zeta
 =\frac{(-1)^s}{R}J_{\nu+2s+1}(R).
\]
One may evaluate the final integral by the absolutely convergent
Bessel power series and the beta integral. The derivative identity
used in the integrations is
$\partial_t^s(t^{-\nu/2}J_\nu(R\sqrt t))
=(-R/2)^st^{-(\nu+s)/2}J_{\nu+s}(R\sqrt t)$.
Dividing by the radial norm in \eqref{inf:radial-norm} proves
\eqref{inf:Bessel-Jacobi-coefficients}.
Parseval and the Lommel integral at a general argument give
\begin{equation}\label{inf:centre-parseval}
 \sum_{s\ge0}\frac{|c_s|^2}{\nu+1+2s}
 =\frac{J_\nu(R)^2-J_{\nu-1}(R)J_{\nu+1}(R)}{(d_j^{\mathrm{mode}})^2}
 \le C_{j,\mathrm{rad}}.
\end{equation}
The general Lommel identity follows by differentiating
$x^2(J_\nu(x)^2-J_{\nu-1}(x)J_{\nu+1}(x))/2$ and using the
recurrences. The bound follows since every numerator value is a
fixed recurrence shift of $J_{p-1}(R)$ and the normalization
quotient has a nonzero limit.

For $s\le S=C_{j,\mathrm{rad}}'p+O_j(1)$, with a fixed $C_{j,\mathrm{rad}}'$ large enough that
$\nu+1+2S\ge2R$, Cauchy--Schwarz in
\eqref{inf:centre-parseval} gives
\[
 \sum_{s\le S}\wt(2j+2s)|c_s|\le C_{j,\mathrm{rad}}p.
\]
There are $O_j(p)$ terms; their denominator reciprocals in the
Cauchy--Schwarz factor are inverted to quantities $O_j(p)$, and
the weights are bounded by a constant by \eqref{inf:weight-properties}.
A fixed degree moment adds at most $C_{j,k}p^k$.

For the remaining tail use the positive ratio bound
$J_{m+1}(R)/J_m(R)<R/(m+1)$ when $m\ge2R$.
Here positivity follows from $j_{m,1}>m$.
For a proof of the ratio bound, its Riccati equation is
$r'=1-(2m+1)r/x+r^2$, with
$r(x)\sim x/(2m+2)$ at zero. At a first contact with
$x/(m+1)$ for $x\le m$, the difference of derivatives would be
$-1+x^2/(m+1)^2<0$, excluding upward contact.
Applying the bound twice to \eqref{inf:Bessel-Jacobi-coefficients}
gives $|c_{s+1}/c_s|<R^2/(m(m+1))\le1/4$ for
$m=\nu+1+2s\ge2R$.
The neighbouring grading ratio is at most $e^{2/p}$, and the
moment ratio tends uniformly to one as $p$ grows for fixed $k$.
Thus the weighted tail is geometric with ratio below $1/2$ for
large $p$, and is bounded by the head estimate. The angular factor
$2^j$ is bounded for the fixed list of modes.
\end{proof}

\begin{proposition}[The deformed centre and its source]
\label{inf:centre-source-norm}
Let $f_N=h_N/R$, $\Phi_N=\Phi_{f_N}$, and
\begin{equation}\label{inf:clamped-centre}
 U_*=u_N\circ(R\Phi_N),\qquad U_0=U_*-\ell_N,
 \qquad g_0=\Delta U_0,\qquad z_0=(g_0,f_N).
\end{equation}
For every fixed moment $k$ and $p\ge P_{N,k}$,
$g_0\in\Gc_p$, $U_0=1+Kg_0$, and
\begin{equation}\label{inf:centre-norm-bounds}
 \nrm{U_*}_{\Cc^{[k]}}+\nrm{U_0}_{\Cc^{[k]}}
 \le C_{N,k}p^{k+1},\qquad
 \nrm{g_0}_{\Cc^{[k]}}\le C_{N,k}p^{k+5}.
\end{equation}
For $p\ge P_{N,0}$, on the fixed product region
\begin{equation}\label{inf:nonlinear-working-region}
 \{g\in\Gc_p:\nrm{g-g_0}_{\Cc}\le1\}
 \ \times\ \{f\in\Ec_1:\nrm f_{\Ec_1}\le1/4\},
\end{equation}
\begin{equation}\label{inf:p6-second}
 \sup\nrm{D^2\Fnl}\le C_Np^6.
\end{equation}
\end{proposition}
\begin{proof}
For each known centre mode, fix $p$ and put $\nu=p+2j-1$.
Its radial factor in the solid basis is the entire function of
the complex squared-radius variable $t$,
\[
 F_{\nu,j}(t)=\frac1{d_j^{\mathrm{mode}}}
 \sum_{h\ge0}\frac{(-1)^h(R/2)^{\nu+2h}}{h!\,\Gamma(\nu+h+1)}t^h,
 \qquad t\in\C,
\]
where $\Gamma$ is Euler's gamma function
\cite[Eq.~10.2.2]{inf:dlmf}. For real $t>0$, this is
$t^{-\nu/2}J_\nu(R\sqrt t)/d_j^{\mathrm{mode}}$.
The radial expansion in Lemma~\ref{inf:radial-centre-norm} is
$\sum_{s\ge0}c_sJ_s^\nu(t)$ on $0\le t\le1$ in weighted $L^2$.
For $m=\nu+1+2s\ge2R$, the ratio bound proved in that lemma gives
$|c_{s+1}|\le R^2|c_s|/[m(m+1)]$. Since $m\ge2s$, iteration
from a fixed $s_0\ge1$ with $\nu+1+2s_0\ge2R$ gives
\[
 |c_s|\le |c_{s_0}|(R^2/4)^{s-s_0}
                 \left(\frac{(s_0-1)!}{(s-1)!}\right)^2,
 \qquad s\ge s_0.
\]
For every fixed $A>0$, Taylor's formula at one and
\eqref{inf:Jacobi-endpoint-derivatives} give
\[
 \sup_{|t-1|\le A}|J_s^\nu(t)|
 \le\sum_{h=0}^s\frac{[A s(2s+\nu+1)]^h}{(h!)^2}
 \le\exp\{2\sqrt{A s(2s+\nu+1)}\}.
\]
For fixed $p,j,A$, this envelope is at most exponential in $s$,
so the factorial bound makes the radial series locally uniformly
convergent on $\C$. Its sum is entire. On $[0,1]$ uniform
convergence and the original $L^2$ expansion identify that sum
with $F_{\nu,j}$; continuity turns the almost-everywhere equality
into equality on the interval. The identity theorem gives
equality on $\C$, including at the deformed radial arguments.

Write $q=1+\HP f_N$ and $e=q^2-1$.
The finite centre shape and the module estimate give
$\nrm{M_e}\le C_N\tau p^{-2}=:\epsilon_{\mathrm{cmp},p}$.
For one radial basis polynomial the exact Taylor identity is
\[
 J_s^\beta(q^2t)=\sum_{m=0}^s\frac{e^m}{m!}
                                      t^m\partial_t^mJ_s^\beta(t).
\]
The derivative columns and the positive $t$ stencils show that
$t^m\partial_t^mJ_s^\beta$, whose degree is no greater than $s$,
has coefficient sum
$(s)_{\mathrm{fall},m}(s+\beta+1)_m/m!$.
It follows that its relative composition norm is at most
\begin{equation}\label{inf:composition-bound}
 \sum_{m=0}^s\epsilon_{\mathrm{cmp},p}^m
 \frac{(s)_{\mathrm{fall},m}(s+\beta+1)_m}{(m!)^2}
 \le\exp\{2\sqrt{\epsilon_{\mathrm{cmp},p}s(2s+\beta+1)}\}.
\end{equation}
Indeed bound the product by $[s(2s+\beta+1)]^m$ and use
$\sum z^m/(m!)^2\le e^{2\sqrt z}$, which follows by comparing
with the even terms of the exponential using
$\binom{2m}m\le4^m$.
The extra solid factor is $q^{2j}$, a bounded multiplier because
$j$ belongs to the fixed finite set of centre modes.

On the $s=O_N(p)$ head the right side of
\eqref{inf:composition-bound} is bounded by $C_N$.
On the tail it is at most $\exp(C_N(s+p)/p)$.
Multiply the geometric Bessel tail in
Lemma~\ref{inf:radial-centre-norm} by this envelope and by the
weight ratio $e^{2/p}$. For large $p$ the resulting geometric
ratio is still below $1/2$. Thus each composed mode series
converges in $\Cc$ with norm at most $C_Np$.
The finite polynomial chart increases each total degree by at
most a fixed multiple depending on $N$. Therefore each output
moment is bounded by a constant times the corresponding input
moment; the same argument gives convergence and the bounds in
all fixed moments $k$.
For this fixed $p$, the radial entire-function identity makes
the composed partial sums converge uniformly on the closed ball
to $S_j(qy)F_{\nu,j}(q^2t)$: the arguments $q^2t$ range over a
compact set, and $S_j(qy)$ is a fixed polynomial. The polynomial
partial sums also converge in $\Cc$ by the preceding coefficient
estimate, hence in $L^2(B)$ by Lemma~\ref{inf:space-embedding}.
Uniform convergence on the closed ball gives the same $L^2$ limit.
Uniqueness of that limit and injectivity of $\Cc\hookrightarrow L^2(B)$
identify the coefficient sum with the composed mode.
Adding the finite list of centre modes proves the asserted
identification and bounds for $U_*$.
The finite field amplitudes are uniformly bounded by the centre
construction. Subtracting the lift uses \eqref{inf:lift-estimates}
and proves the assertion for $U_0$.
The Laplacian bound \eqref{inf:Delta-moment}, with moments included,
then gives $\nrm{g_0}_{\Cc^{[k]}}\le C_{N,k}p^{k+5}$.

By construction the two traces of $U_0$ are exactly $1,0$.
Dirichlet Poisson uniqueness yields $U_0=1+\KD g_0$, and
\eqref{inf:poisson-trace} gives $g_0\in\Gc_p$.
Finally \eqref{inf:nonlinear-second} and the $p^5$ source bound
give \eqref{inf:p6-second} on the product region
\eqref{inf:nonlinear-working-region}. Its exponent does not depend
on $N$.
\end{proof}

\subsection{Residuals in the same coefficient space}

\begin{proposition}[The residual and its material derivative]
\label{inf:residual-estimates}
At the centre $z_0$ of \eqref{inf:clamped-centre}, put
$E_0=\Fnl(z_0)$. For every fixed moment $k$ and $p\ge P_{N,k}$,
\begin{equation}\label{inf:residual-moments}
 \nrm{E_0}_{\Cc^{[k]}}\le C_{N,k}\theta^2p^{2-N},\qquad
 \nrm{\Euler E_0}_{\Cc^{[k]}}\le C_{N,k}\theta^2p^{3-N}.
\end{equation}
For a shape variation $\eta\in\Ec_1$, define
\begin{equation}\label{inf:residual-material}
 \mathcal V_{\mathrm{res}}(\delta g,\eta)
       =(\HP\eta)d_0^{-1}\Euler E_0,
 \qquad d_0=1+\HP f_N+\Euler\HP f_N.
\end{equation}
Then, for $p\ge P_N$,
\begin{equation}\label{inf:residual-material-bound}
 \nrm{\mathcal V_{\mathrm{res}}}_{\Xc_p\to\Cc}
 \le C_N\theta^2p^{3-N}.
\end{equation}
\end{proposition}
\begin{proof}
The interior equation of the finite centre is exact, so
$B_{\Phi_N}U_*=0$ and $E_0=-B_{\Phi_N}\ell_N$.
The lift occupies only radial layers zero and one.
On these layers the operators $\Delta$, $\Euler^2$ and
$\nabla a\cdot\nabla\Euler$, for any fixed solid polynomial $a$,
have coefficient bounds $C_a p^2$ times a fixed number of angular
moments.
On layer one,
\[
 \Euler\Psi_{j,1}=(2j+2)\Psi_{j,1}
                            +2(\beta_j+1)\Psi_{j,0}.
\]
On a solid monomial $S_l t^h$ of bounded radial power,
$\Delta(S_lt^h)=4h(h+2l+p-1)S_lt^{h-1}$ costs one factor $p+l$.
Writing $J_1^{\beta_j}=(\beta_j+2)t-(\beta_j+1)$ adds at most one
such factor. Multiplication by the fixed polynomial $a$ has bounded
angular shifts and bounded radial powers.
Finally
$2\nabla a\cdot\nabla U=\Delta(aU)-a\Delta U-(\Delta a)U$.
In the displayed Euler formula the large coefficient $p+j$ occurs
on the harmonic layer, whose Laplacian product costs only one
additional factor $p+j$; the layer-one coefficient is only $2j+2$.
The product therefore costs at most $C_a(p+j)^2$.
The direct Euler bound also gives $Cp^2$ after two fixed moments.

At the fixed finite chart, $q_0-1,d_0-1$ have coefficient multiplier
norm $O_N(\tau p^{-2})$. Put
\[
 \epsilon_{\mathrm{rec}}
 =\max\{\nrm{M_{q_0-1}}_{\Cc\to\Cc},
                 \nrm{M_{d_0-1}}_{\Cc\to\Cc}\},
\]
and increase the threshold until $\epsilon_{\mathrm{rec}}<1/2$.
Their inverse powers preserve all fixed moments: multiplication
by a polynomial of degree $D_N$ increases
degree by at most $D_N$, and
$\sum_{v\ge0}(1+vD_N)^k\epsilon_{\mathrm{rec}}^v<\infty$ for
$\epsilon_{\mathrm{rec}}<1/2$.
The constants may depend on $N,k$ but not on $p,r,\theta$.
Apply the two-layer estimates to every term of
\eqref{inf:scalar-pullback} and then these inverse multipliers.
Equation~\eqref{inf:lift-estimates} and $R^2=O(p^2)$ give the
first estimate in \eqref{inf:residual-moments}.
Apply \eqref{inf:Euler-moment}, using that first estimate at moment
$k+2$, to obtain the second.
The module bound for $\HP\eta$ and the fixed inverse multiplier
$d_0^{-1}$ then prove \eqref{inf:residual-material-bound}.
\end{proof}

\subsection{The exact material map and its inverse}

The radial harmonic chart specifies the interior extension as well
as the boundary displacement.

\begin{theorem}[Material isomorphism]\label{inf:material-isomorphism}
Fix $N\ge4$ and take $p$ sufficiently large for the centre above.
Put
\begin{equation}\label{inf:material-definitions}
 D_p=-\partial_{\nu\nu}U_0|_{\partial B},\qquad
 \mathfrak f_0=\frac{\Euler U_0}{d_0},\qquad
 T_\eta=(\HP\eta)\mathfrak f_0.
\end{equation}
The map
\begin{equation}\label{inf:Qmat-definition}
 \Qmat(\delta g,\eta)=\Delta(K\delta g-T_\eta)
                 :\Xc_p\longrightarrow\Cc
\end{equation}
is a bounded isomorphism, with
\begin{equation}\label{inf:Qmat-bounds}
 \nrm{\Qmat}\le C_Np^8,\qquad
 \nrm{\Qmat^{-1}}\le C_Np^6.
\end{equation}
If $J_D=B_{\Phi_N,D}\KD$, then the exact factorization is
\begin{equation}\label{inf:material-factorization}
 D\Fnl(z_0)=J_D\Qmat+\mathcal V_{\mathrm{res}}.
\end{equation}
All constants are uniform in the admissible split and
$0<|\theta|\le20$; the exponents are independent of $N$.
\end{theorem}
\begin{proof}
The centre has constant Dirichlet trace and zero gradient, so its
boundary Hessian is purely normal-normal. Evaluating its equation
at the boundary gives
\begin{equation}\label{inf:normal-hessian}
 D_p=\frac{R^2-E_0|_{\partial B}}{\nu^TA_{q_0}\nu}.
\end{equation}
Indeed all drift terms multiply the zero gradient, and
$D^2U_0=(\partial_{\nu\nu}U_0)\nu\otimes\nu$ there.
The finite centre polynomial and the boundary algebra show
$d_0|_{\partial B}=1+O_N(\tau p^{-2})$ and
$\nu^TA_{q_0}\nu=1+O_N(\tau p^{-2})$ in $\Ec_1$.
For the latter assertion insert \eqref{inf:pullback-metric};
the product formula in the proof of
\eqref{inf:gradient-multiplier}, with the extra boundary weight
$1+k\le(1+i)(1+j)$ on each angular product, gives
\[
 \nrm{M_{q_0}|_{\partial B}}_{\Ec_1}
 \le Cp\left(\sum_j(1+j)^2 2^j\wt(2j)|f_{N,j}|\right)^2
 \le C_N\tau^2p^{-3}.
\]
Here all fixed moments of the finite centre shape are
$O_N(\tau p^{-2})$ by \eqref{inf:centre-profile}.
The trace of $E_0$ in $\Ec_1$ is bounded by its first source
moment, since $J_s^{\beta_j}(1)=1$.
To identify this coefficient trace, take polynomial sums $F_m$ of
$E_0$ and put $h_m=F_m|_{\partial B}$. The first source moment
gives $h_m\to h$ in $\Ec_1$, while the fourth source moment in
Proposition~\ref{inf:residual-estimates} and
\eqref{inf:Delta-moment} gives $\Delta F_m\to\Delta E_0$ in $L^2$.
Since $F_m-\HP h_m$ has zero Dirichlet trace,
\[
 F_m-\HP h_m=\KD\Delta F_m.
\]
At fixed $p$, harmonic extension and Dirichlet Poisson regularity
make both terms on the right of
$F_m=\HP h_m+\KD\Delta F_m$ converge in $H^2$.
Their $L^2$ limit is $E_0$, hence
$E_0=\HP h+\KD\Delta E_0\in H^2$ and its Sobolev trace is $h$.
Thus \eqref{inf:residual-moments} and the boundary algebra give
\begin{equation}\label{inf:normal-hessian-inverse}
 \nrm{D_p^{-1}}_{\Ec_1}\le C_Np^{-2}
\end{equation}
by the convergent inverse series in \eqref{inf:normal-hessian}.
In particular $D_p$ has no zero.

The moment bounds for the fixed centre and its reciprocal $d_0$
give
$\nrm{\mathfrak f_0}_{\Cc^{[k]}}\le C_{N,k}p^{k+4}$ by
\eqref{inf:Euler-moment}.
Using \eqref{inf:Delta-moment} with its sum of two powers yields
\begin{equation}\label{inf:fixed-material-field}
 \nrm{\Delta \mathfrak f_0}_\Cc\le C_Np^8.
\end{equation}
Since $\mathfrak f_0$ has zero Dirichlet trace, Dirichlet uniqueness implies
$\mathfrak f_0=\KD\Delta \mathfrak f_0$.
Harmonicity of $\HP\eta$ and
\eqref{inf:gradient-C-one} therefore give
\begin{equation}\label{inf:material-field-bound}
 \nrm{\Delta T_\eta}_\Cc
 \le\nrm{M_{\HP\eta}\Delta \mathfrak f_0}_\Cc+
       2\nrm{\nabla(\HP\eta)\cdot\nabla\KD\Delta \mathfrak f_0}_\Cc
 \le C_Np^8\nrm\eta_{\Ec_1}.
\end{equation}
One moment of the unknown shape enters. This proves
boundedness of \eqref{inf:Qmat-definition}.

For an arbitrary source $\psi\in\Cc$ set
\begin{equation}\label{inf:material-reconstruction}
 W=\KD\psi,\quad N_W=\partial_\nu W,\quad
 \eta=(d_0|_{\partial B})N_W/D_p,\quad
 \delta g=\psi+\Delta T_\eta.
\end{equation}
Equations \eqref{inf:poisson-trace} and
\eqref{inf:normal-hessian-inverse} imply
$\nrm\eta_{\Ec_1}\le C_Np^{-2}\nrm\psi_\Cc$;
then \eqref{inf:material-field-bound} gives
$\nrm{\delta g}_\Cc\le(1+C_Np^6)\nrm\psi_\Cc$.
At the boundary,
\[
 T_\eta=0,\qquad \partial_\nu T_\eta=-D_p\eta/d_0.
\]
This follows by differentiating $\mathfrak f_0=(\Euler U_0)/d_0$:
$\Euler U_0=0$ there and
$\partial_\nu\Euler U_0=\partial_{\nu\nu}U_0=-D_p$.
Hence $W+T_\eta$ has both traces zero.
It is the Dirichlet Poisson solution with source $\delta g$;
\eqref{inf:poisson-trace} implies $\delta g\in\Gc_p$ and
$K\delta g=W+T_\eta$.
This proves that \eqref{inf:material-reconstruction} is a right
inverse to $\Qmat$ and gives its $p^6$ bound.
Conversely, for any input to $\Qmat$, the field
$W=K\delta g-T_\eta$ has zero Dirichlet trace and is recovered
uniquely from its Laplacian. Its normal trace is $D_p\eta/d_0$,
so \eqref{inf:material-reconstruction} recovers $\eta$ and then
$\delta g$. This proves the left inverse as well.

It remains to prove \eqref{inf:material-factorization}.
For the chart variation $\Phi_N+t(\HP\eta)y$, its reference material
vector is
$v_\eta=D\Phi_N^{-1}((\HP\eta)y)=y(\HP\eta)/d_0$.
Differentiation of the two compositions in
\eqref{inf:scalar-pullback-definition} gives the exact scalar
commutator identity
\[
 \delta(P_\Phi\Delta)U_0
 =-P_{\Phi_N}\Delta(v_\eta\cdot\nabla U_0)
       +v_\eta\cdot\nabla(P_{\Phi_N}\Delta U_0).
\]
It follows either by differentiating the inverse chart or by
applying the chain rule first to a smooth physical test function.
Since $v_\eta\cdot\nabla U_0=T_\eta$, adding the frequency term
and the source variation yields
\[
 D\Fnl(z_0)(\delta g,\eta)
 =B_{\Phi_N}(K\delta g-T_\eta)
       +(\HP\eta)d_0^{-1}\Euler E_0.
\]
The field in parentheses has zero Dirichlet trace, so it equals
$\KD\Qmat(\delta g,\eta)$. This proves
\eqref{inf:material-factorization} on finite variations.
To pass to the completion, fix $p$ and the smooth real centre chart.
Dirichlet regularity on the ball gives a bounded map
$\KD:L^2(B)\to H^2(B)\cap H^1_0(B)$; the smooth coefficients in
\eqref{inf:scalar-pullback} give a bounded map
$B_{\Phi_N,D}:H^2(B)\cap H^1_0(B)\to L^2(B)$.
These continuity bounds use only ball regularity and the fixed
chart, independently of the Dirichlet inverse estimates below.
For $\dot z\in\Xc_p$, take finite variations
$\dot z_m\to\dot z$ in $\Xc_p$.
The already proved boundedness of $\Qmat$ gives
$\Qmat\dot z_m\to\Qmat\dot z$ in $\Cc$, hence in $L^2(B)$ by
Lemma~\ref{inf:space-embedding}. Consequently
$J_D\Qmat\dot z_m\to B_{\Phi_N,D}\KD\Qmat\dot z$ in $L^2(B)$.
The finite-variation identity also says
$J_D\Qmat\dot z_m=(D\Fnl(z_0)-\mathcal V_{\mathrm{res}})\dot z_m$.
Its right side converges in $\Cc$ to
$(D\Fnl(z_0)-\mathcal V_{\mathrm{res}})\dot z$, by
Proposition~\ref{inf:nonlinear-map} and
\eqref{inf:residual-material-bound}. The continuous injective
embedding $\Cc\hookrightarrow L^2(B)$ identifies the two limits.
Thus $J_D\Qmat\dot z$ belongs to $\Cc$ and satisfies
\eqref{inf:material-factorization} for every $\dot z\in\Xc_p$.
Surjectivity of $\Qmat$ also gives boundedness of $J_D$ on all
of $\Cc$, with
$J_D=(D\Fnl(z_0)-\mathcal V_{\mathrm{res}})\Qmat^{-1}$.
\end{proof}

\section{Mixed norms and the exact relative graph bound}\label{inf:sec-mixed}

The intermediate space keeps the radial variable Hilbertian and the
angular coefficients summable and preserves radial spectral
projections in endpoint-normalized angular coordinates.

\begin{definition}[Mixed space and weighted matrix norms]
\label{inf:mixed-spaces}
Let $\Mmix$ be the space of all families $(u_j)_{j\ge0}$ with
$u_j\in L^2((0,1),pt^{\beta_j}dt)$ and finite norm below,
represented as $u=\sum_jS_ju_j(t)$. Put
\begin{equation}\label{inf:mixed-norm}
 \nrm u_{\Mmix}=\sum_j2^j\nrm{u_j}_{\beta_j},\qquad
 \nrm v_\beta^2=p\int_0^1t^\beta|v(t)|^2\,dt.
\end{equation}
The actual invariant Hilbert space $\Hc$ has norm
\begin{equation}\label{inf:true-H-norm}
 \nrm u_\Hc^2=\sum_j\eta_j^2\nrm{u_j}_{\beta_j}^2.
\end{equation}
For an angular block matrix $T$ define, at fixed $\sigma>0$,
\begin{align}
 \nrm T_{\Mmix,\sigma}
 &=\sup_j\sum_i2^{i-j}e^{\sigma|i-j|}
                            \nrm{T_{ij}}_{\beta_j\to\beta_i},
 \label{inf:M-matrix-norm}\\
 \nrm T_{\Hc,\sigma}
 &=\max\left\{
 \sup_i\sum_je^{\sigma|i-j|}\nrm{T^H_{ij}},
 \sup_j\sum_ie^{\sigma|i-j|}\nrm{T^H_{ij}}\right\},
 \quad T^H_{ij}=\frac{\eta_i}{\eta_j}T_{ij}.
 \label{inf:H-matrix-norm}
\end{align}
The block norms in the second line use the same radial Hilbert spaces;
only the angular basis has been normalized.
\end{definition}

\begin{lemma}[Embedding and matrix algebra]\label{inf:mixed-algebra}
There are norm-at-most-one embeddings
$\Cc\hookrightarrow\Mmix\hookrightarrow\Hc$.
The space $\Mmix$ is complete, and finite polynomial sums are
dense in all three spaces.
The two matrix norms in
\eqref{inf:M-matrix-norm}--\eqref{inf:H-matrix-norm} are
submultiplicative, have identity norm one, and dominate their
respective operator norms. The spaces of matrices with finite such
norms are complete. For a block-diagonal operator both reduce
to the supremum of its radial block norms.
\end{lemma}
\begin{proof}
The radial squared basis norm is
$p/(p+2j+2s)\le1$. The triangle inequality gives the first
embedding from \eqref{inf:C-norm}; the inequality $\eta_j\le1$
and $\ell^2\le\ell^1$ give the second.
Finite angular truncations and polynomial density for each positive
radial weight give density in $\Mmix$ and $\Hc$; density in
$\Cc$ is part of its definition. Completeness of $\Mmix$ follows
from its description as a weighted $\ell^1$ sum of the complete
radial Hilbert spaces: a Cauchy family has a limit in every
radial block, and Fatou's inequality bounds the sum of the
blockwise differences by its Cauchy bound.
For products of blocks use
$e^{\sigma|i-j|}\le e^{\sigma|i-k|}e^{\sigma|k-j|}$ and
$2^{i-j}=2^{i-k}2^{k-j}$, then sum the positive block norm bounds.
The ratios $\eta_i/\eta_j$ telescope through the intermediate
index. This proves submultiplicativity for both row and column
sums. The mixed operator bound follows from summing columns;
the Hilbert bound is the block Schur inequality, obtained by
Cauchy--Schwarz with these positive row and column sums.
For completeness, a Cauchy sequence in either matrix norm is
Cauchy in every block operator norm. Take its blockwise limits.
Fatou's inequality for the positive row and column sums bounds
the norm of each difference from that limit by the corresponding
Cauchy bound, uniformly over the exterior index. The sequence
therefore converges in the matrix norm, and the preceding operator
bounds identify the limit with a bounded operator on the indicated
space. For diagonal blocks all off-diagonal weights are one and
only one term remains in each sum.
\end{proof}

\begin{lemma}[Solid multipliers in the two matrix norms]
\label{inf:mixed-multipliers}
For integers $i,h\ge0$,
\begin{equation}\label{inf:mixed-multiplier-bounds}
 \nrm{M_{S_it^h}}_{\Mmix,\sigma}\le2^ie^{\sigma i},\qquad
 \nrm{M_{S_it^h}}_{\Hc,\sigma}\le(2i+1)e^{\sigma i}.
\end{equation}
\end{lemma}
\begin{proof}
An angular product output has factor
$t^{i+j-k+h}u_j$. Its squared radial weight has exponent
\[
 \beta_k+2(i+j-k+h)=\beta_j+2i+2h\ge\beta_j.
\]
Multiplication by that power of $t$ is a contraction between the
indicated radial spaces. The positive angular coefficients sum to
one, and $|k-j|\le i$; this proves the mixed bound.
In the true Hilbert space multiplication has norm at most one
because $|S_it^h|\le1$ on the ball. Every angular block is an
orthogonal compression and has norm at most one. Its bandwidth
is $i$, so each weighted row and column contains at most $2i+1$
terms, proving the Hilbert bound.
\end{proof}

\begin{lemma}[Poisson and gradient estimates in mixed and Hilbert norms]
\label{inf:mixed-Poisson}
On all source layers one has
\begin{equation}\label{inf:mixed-Poisson-bounds}
 \nrm{\KD}\le Cp^{-2},\qquad
 \nrm{\Euler\KD}\le Cp^{-1},\qquad
 \nrm{\Euler^2\KD}\le C
\end{equation}
in both matrix norms. For $A=2i>0$,
\begin{align}
 \nrm{\nabla S_i\cdot\nabla\KD}_{\Mmix,\sigma}
 &\le CA2^ie^{\sigma i}/p,&
 \nrm{\nabla S_i\cdot\nabla\Euler\KD}_{\Mmix,\sigma}
 &\le CA2^ie^{\sigma i},\label{inf:mixed-gradients}\\
 \nrm{\nabla S_i\cdot\nabla\KD}_{\Hc,\sigma}
 &\le CA(2i+1)e^{\sigma i}/p,&
 \nrm{\nabla S_i\cdot\nabla\Euler\KD}_{\Hc,\sigma}
 &\le CA(2i+1)e^{\sigma i}.
 \label{inf:Hilbert-gradients}
\end{align}
The constants are absolute for fixed $\sigma$.
\end{lemma}
\begin{proof}
In angular degree $B=2j$, let $\kf$ solve
$4(t\kf''+(\beta+1)\kf')=f$, $\kf(1)=0$, $\beta=B+p-1$.
The radial spectral lower bound gives
$\nrm\kf_\beta\le C(B+p)^{-2}\nrm f_\beta$.
Multiplication of the equation by $\overline\kf t^\beta$ and integration
by parts give
\[
 4\nrm{\kf'}_{\beta+1}^2
 =-\operatorname{Re}\langle f,\kf\rangle_\beta,
 \qquad
 \nrm{\kf'}_{\beta+1}\le C(B+p)^{-1}\nrm f_\beta.
\]
These equalities first hold for polynomial sources; their
extensions follow from the Dirichlet form and Poisson uniqueness.
The profile identities
$\Euler\KD=B\kf+2t\kf'$ and
$\Euler^2\KD=B^2\kf+4(1-p)t\kf'+tf$
then give \eqref{inf:mixed-Poisson-bounds} on each block, hence in
both matrix norms.

For the gradient operators use the two profiles
\eqref{inf:gradient-profile-one}--\eqref{inf:gradient-profile-two}
with the general source $f$ in place of $J_s^\beta$ and
$\kf$ in place of $k_s$.
When $h\ge1$, the output squared weight for $t^{h-1}\kf$ is
$t^{\beta+A-2}$, which is bounded above by $t^\beta$ since
$A\ge2$. The corresponding weight for $t^h\kf'$ is
$t^{\beta+A}\le t^{\beta+1}$.
The source term $t^hf$ is also a contraction between its two
radial spaces. Equation~\eqref{inf:hZ-bound} and the preceding
energy bounds give $CA/(B+p)$ for the first profile and $CA$ for
the second. There is no negative-power term when $h=0$.
Positive angular summation and the bandwidth weight prove
\eqref{inf:mixed-gradients}.

In Hilbert coordinates, homogeneity and
Lemma~\ref{inf:angular-sup} give
$\nrm{\nabla S_i}_\infty\le\sqrt2A$.
For $u=\KD f$, the Dirichlet spectral lower bound gives
$\nrm{\nabla u}_2\le(p-1)^{-1}\nrm f_2$.
On the ball one also has
\begin{equation}\label{inf:convex-Hessian}
 \nrm{D^2u}_2\le\nrm{\Delta u}_2=\nrm f_2.
\end{equation}
For a smooth real Dirichlet function, integrate the divergence of
$(\Delta u)\nabla u-D^2u\nabla u$.
Its interior divergence is $(\Delta u)^2-|D^2u|^2$.
At the boundary $\nabla u=(\partial_\nu u)\nu$ and
$\Delta u=\partial_{\nu\nu}u+H_{\partial B}\partial_\nu u$,
so the boundary integral is
$\int_{\partial B}H_{\partial B}|\partial_\nu u|^2\ge0$.
This proves \eqref{inf:convex-Hessian}; use real and imaginary
parts and $H^2$ approximation for the general case.
Since $\nabla\Euler u=\nabla u+(D^2u)y$, the unweighted Hilbert
bounds for the two gradient operators are $CA/p$ and $CA$.
Their angular bandwidth is $i$, as is also visible in the exact
profiles. Compressing to blocks and summing at most $2i+1$
weighted entries per row and column proves
\eqref{inf:Hilbert-gradients}. This argument covers the harmonic
source layer as well as all the others.
\end{proof}

\begin{lemma}[Closed mixed graph realization]\label{inf:mixed-domain}
Let $H_0$ have its Hilbert Dirichlet realization. Its realization
on $\Mmix$ has domain
\[
 \Dom_{\Mmix}(H_0)=\KD\Mmix
 =\{u\in\Dom_{\Hc}(H_0):H_0u\in\Mmix\}.
\]
It is closed in $\Mmix$ and its graph norm is equivalent, for
fixed $p$, to the norm of its source. Its realization on $\Cc$
has the identical description with $\Cc$ in place of $\Mmix$.
Let $V$ be defined on $\Dom_{\Hc}(H_0)$, and suppose
$\nrm{VH_0^{-1}}<1$ both on $\Hc$ and on $\Mmix$.
Then $H_0+V$ restricts to a closed operator on
$\Dom_{\Mmix}(H_0)$, and
\[
 \{u\in\Dom_{\Hc}(H_0):(H_0+V)u\in\Mmix\}
 =\Dom_{\Mmix}(H_0).
\]
The same conclusions hold with $\Cc$ in place of $\Mmix$ when
the relative norm is below one on both $\Hc$ and $\Cc$.
\end{lemma}
\begin{proof}
The Poisson estimates give a bounded map $\KD:\Mmix\to\Mmix$,
and the Hilbert identification gives $H_0\KD f=-f$.
Thus every element in its range is in the stated Hilbert domain
and has source in $\Mmix$.
Conversely, if $u$ lies in that Hilbert domain and $H_0u=f\in\Mmix$,
Dirichlet uniqueness gives $u=-\KD f$, proving equality of the
descriptions. The norm $\nrm{H_0u}_{\Mmix}$ is a complete norm
on the domain, since $\KD$ is a bijection from the complete
source space onto it. The Poisson bound makes it equivalent to
$\nrm u_{\Mmix}+\nrm{H_0u}_{\Mmix}$.
If $u_n\to u$ and $H_0u_n\to f$ in $\Mmix$, the embeddings
imply the same convergence in $\Hc$. Closedness of the Hilbert
operator gives $u\in\Dom_\Hc(H_0)$ and $H_0u=f$; hence the
description just proved puts $u$ in the mixed domain. This
proves closedness directly.
The coefficient-space proof uses the coefficient Poisson bound
and the same Hilbert embedding, so every step remains valid.

In each indicated source space, the factor
$I+VH_0^{-1}$ is bounded and boundedly invertible in the source
space. The product
$(H_0+V)=(I+VH_0^{-1})H_0$ therefore has an equivalent graph norm
on the same complete domain and is closed. Adding a bounded
scalar shift preserves closedness there.
For a source in $\Mmix$, the partial sums of the Neumann series
for $(I+VH_0^{-1})^{-1}$ agree with the Hilbert partial sums.
Convergence in $\Mmix$ implies convergence in $\Hc$, so the
two inverses coincide on these sources. If
$u\in\Dom_{\Hc}(H_0)$ and $(H_0+V)u=f\in\Mmix$, Hilbert
uniqueness gives
$u=H_0^{-1}(I+VH_0^{-1})^{-1}f\in\Dom_{\Mmix}(H_0)$.
The reverse inclusion follows by applying the mixed-space
operator. This proves the displayed equality and the restriction
property. Replacing $\Mmix$ by $\Cc$ gives the coefficient
version. The same argument applies to
a bounded angular-diagonal radial spectral projection, which
commutes with $H_0$ and preserves the graph domain.
\end{proof}

\subsection{A centre-size bound independent of the accuracy order}

\begin{lemma}[Order-independent leading size]\label{inf:uniform-centre-size}
For a fixed centre order $N$, expand $f_N=\sum_j f_{N,j}G_j$ and put
\begin{equation}\label{inf:strong-centre-norm}
 \delta=\sum_j(1+j)^{10}(2e^2)^j\wt(2j)|f_{N,j}|.
\end{equation}
There is an absolute $C_0$ such that for every fixed $N$,
\begin{equation}\label{inf:uniform-delta}
 \delta\le C_0\tau p^{-2}\qquad(p\ge P_N).
\end{equation}
Only the threshold depends on $N$. This stronger norm is used for
the known finite centre; the unknown Newton shape retains the
space $\Ec_1$ with one moment and angular radius two.
\end{lemma}
\begin{proof}
Equation~\eqref{inf:centre-profile} gives
\[
 f_N=-\frac{\theta}{2R^2pg_2}\mathsf H_2+O_N(\tau p^{-4})
\]
in every fixed polynomial coefficient norm.
The leading term has only index two. Its coefficient is bounded
by $C\tau p^{-2}$ independently of $N$, since $pg_2>9/16$ and
$\mathsf H_2(1)\le1$. Its weight in \eqref{inf:strong-centre-norm} is an
absolute constant for $p\ge64$.
For the fixed jet the remainder has bounded degree depending on
$N$, and the finite monic-to-$G_j$ connection is uniformly
bounded for large $p$, because
$\mathsf H_j(1)\to(1-r)^j>0$ for each fixed $j$.
Thus its strong norm is at most $C_N\tau p^{-4}$.
Choose $C_0$ to be the leading bound plus one, and increase $P_N$
until $C_Np^{-2}\le1$. This proves \eqref{inf:uniform-delta}.
\end{proof}

\subsection{Scalar perturbation and density conjugation}

Use the fixed centre chart $\Phi_N=qy$, $d=q+\Euler q$, and define
\begin{equation}\label{inf:density-operators}
 H_\Phi=-P_{\Phi_N}\Delta,\qquad
 m=(\det D\Phi_N)^{1/2}=q^{p-1/2}d^{1/2},\qquad
 \Hu=mH_\Phi m^{-1},\qquad \mathcal V=\Hu-H_0.
\end{equation}
Here $\Hu$ is the unitary Dirichlet realization on $\Hc$.
The target volume uses the same fixed normalization as unit-ball
volume, so multiplication by $m$ is exactly the unitary density
factor. The scalar graph operator is
$J_D=(P_{\Phi_N}\Delta+R^2)\KD$.
All full-source graph estimates in this section are at this
fixed finite centre with the strong norm
\eqref{inf:strong-centre-norm}. The Newton iteration uses its
frozen inverse, while the nonlinear map is estimated on
$\Gc_p\times\Ec_1$.

\begin{proposition}[The Hilbert relative estimate]
\label{inf:Hilbert-relative}
For the real finite-centre chart in \eqref{inf:density-operators},
let $\delta$ be \eqref{inf:strong-centre-norm}. There are absolute
constants $C,\varepsilon_H>0$ such that, if $p\ge64$ and
$p\delta\le\varepsilon_H$, then
\begin{equation}\label{inf:Hilbert-relative-bound}
 \nrm{\mathcal VK_0^{-1}}_{\Hc,1}
 \le C(\delta+p\delta^2).
\end{equation}
The scalar graph difference also satisfies
\begin{equation}\label{inf:Hilbert-scalar-bound}
 \nrm{(P_{\Phi_N}\Delta-\Delta)\KD}_{\Hc,1}
 \le C(\delta+\delta^2).
\end{equation}
These constants do not depend on $N$, the split, or the number
of nonzero centre modes. All invariant radial source layers
are included.
\end{proposition}
\begin{proof}
Use the orthonormal angular coordinates of \eqref{inf:H-matrix-norm}.
Its ingredients are the endpoint supremum bound, finite angular
bandwidth, and ball Hilbert estimates.

\emph{1. Banded operators.}
Let $P_j^H$ be the orthogonal projection onto angular degree
$2j$. If an operator $T$ has angular bandwidth $D$, each block
$P_i^HTP_j^H$ has norm at most $\nrm T_{\Hc\to\Hc}$, and each
row and column has at most $2D+1$ nonzero blocks. Therefore
\begin{equation}\label{inf:Hilbert-band-bound}
 \nrm T_{\Hc,1}\le(2D+1)e^D\nrm T_{\Hc\to\Hc}.
\end{equation}
This bounds both weighted row and column sums.

\emph{2. Multipliers.}
For an invariant polynomial multiplier $a$ of Jacobi degree
(maximum angular index) at most $D$, orthogonality gives bandwidth $D$ and multiplication
on $L^2$ has norm at most $\nrm a_\infty$. Hence
\[
 \nrm{M_a}_{\Hc,1}\le(2D+1)e^D\nrm a_\infty.
\]
In particular $|S_it^h|\le1$ gives
$(2i+1)e^i$ for that multiplier. Summing the centre coefficients
shows that multiplication by $\mathfrak h=q-1$, $\Euler\mathfrak h$,
and $\Euler(\Euler-1)\mathfrak h$ has norm at most $C\delta$.
The fixed moments in \eqref{inf:strong-centre-norm} absorb all
the displayed degree factors. The endpoint bound used here is
for $\ab\le\bb$ and $G_i(1)=1$, as fixed in
\eqref{inf:angular-jacobi}.

\emph{3. Products of gradients.}
For $A=2i$ the solid harmonic has
$|\nabla S_i|\le\sqrt2A$, by radial homogeneity and the
Bernstein bound on great circles. Thus
$|\nabla S_i\cdot\nabla S_j|\le2(2i)(2j)$.
Its angular bandwidth as a multiplier is at most $i+j$.
Equation~\eqref{inf:Hilbert-band-bound} gives
\begin{equation}\label{inf:Hilbert-gradient-product}
 \nrm{M_{\nabla S_i\cdot\nabla S_j}}_{\Hc,1}
 \le2(2i)(2j)(2i+2j+1)e^{i+j}.
\end{equation}
After summing the shape coefficients, this bounds the multipliers
$M_q=\nabla q\cdot\nabla q$,
$\nabla q\cdot\nabla\Euler q$,
$\nabla d\cdot\nabla d$, and
$\nabla q\cdot\nabla d$ by $C\delta^2$.
An Euler derivative on a shape adds only its angular degree,
which is included in the ten moments. This Hilbert multiplier
bound is $C\delta^2$.

\emph{4. Diagonal Poisson operators.}
For $u=\KD f$, the ball spectral estimate and its energy identity
give
\[
 \nrm u_2\le(p-1)^{-2}\nrm f_2,\qquad
 \nrm{\nabla u}_2\le(p-1)^{-1}\nrm f_2.
\]
The boundary identity \eqref{inf:convex-Hessian} gives
$\nrm{D^2u}_2\le\nrm f_2$ on the convex ball.
Since $\Euler u=y\cdot\nabla u$ and
$\Euler^2u=y^TD^2uy+y\cdot\nabla u$, the three Poisson operators
have norms at most $(p-1)^{-2}$, $2/(p-1)$ and $4$.
They are angular diagonal, so the same bounds hold in
$\nrm\cdot_{\Hc,1}$ on all source layers.

\emph{5. Gradient--Poisson operators.}
The preceding energy estimate gives
$\nrm{\nabla S_i\cdot\nabla\KD}_{\Hc\to\Hc}
\le\sqrt2(2i)/(p-1)$.
Using $\nabla\Euler u=\nabla u+(D^2u)y$ gives
$\nrm{\nabla S_i\cdot\nabla\Euler\KD}_{\Hc\to\Hc}
\le\sqrt2(2i)(1+(p-1)^{-1})$.
Their exact profiles \eqref{inf:gradient-profile-one}--
\eqref{inf:gradient-profile-two} show angular bandwidth $i$.
Applying \eqref{inf:Hilbert-band-bound} and summing the centre
coefficients proves
\[
 \nrm{\nabla\mathfrak h\cdot\nabla\KD}_{\Hc,1}\le C\delta/p,
 \qquad
 \nrm{\nabla\mathfrak h\cdot\nabla\Euler\KD}_{\Hc,1}\le C\delta.
\]
The same estimates hold with $\mathfrak h$ replaced by
$\Euler\mathfrak h$, and hence for
$d-1=\mathfrak h+\Euler\mathfrak h$.

\emph{6. Inverse multipliers and density.}
The Hilbert matrix norm is a Banach algebra with unit norm one
by Lemma~\ref{inf:mixed-algebra}.
For $C\delta<1$, the Neumann series for $q^{-1},d^{-1}$
converge in it; they also converge uniformly on the ball to the
pointwise multipliers. Their fixed inverse powers are bounded.
The absolute binomial majorant for $m^{\pm1}$ gives
\begin{equation}\label{inf:Hilbert-density-size}
 \nrm{M_{m^{\pm1}}}_{\Hc,1}
 \le(1-C\delta)^{-p-1}\le2
\end{equation}
when the absolute constant $\varepsilon_H$ is sufficiently small.
This choice is independent of the finite centre order.

\emph{7. The scalar graph difference.}
Compose \eqref{inf:scalar-pullback} with $\KD$ and subtract
the identity. The first term costs $C\delta$ by the inverse
multiplier bound, and the second costs $C\delta$ by step 5.
The third costs $C\delta^2$ by steps 3 and 4.
Formula~\eqref{inf:Cq-identity} bounds its contracted Hessian
multiplier by $C\delta^2(1+\delta)$; it is followed by
$\Euler\KD$, of norm $C/p$. The remaining gradient-square
term is bounded in the same way. This proves
\eqref{inf:Hilbert-scalar-bound}.

\emph{8. The density commutator.}
Both $q$ and $d$ are harmonic. With $\alpha=p-1/2$,
\[
 \frac{\nabla m^{-1}}{m^{-1}}
 =-\alpha q^{-1}\nabla q-\tfrac12d^{-1}\nabla d,
\]
\[
 \frac{\Delta m^{-1}}{m^{-1}}
 =\alpha(\alpha+1)q^{-2}\nabla q\cdot\nabla q
 +\tfrac34d^{-2}\nabla d\cdot\nabla d
 +\alpha q^{-1}d^{-1}\nabla q\cdot\nabla d.
\]
The gradient part of $[\Delta,m^{-1}]\KD$ is bounded by
$Cp(\delta/p)=C\delta$, using step 5 and the inverse multipliers.
Its potential part is bounded by
$Cp^2\delta^2(p-1)^{-2}\le C\delta^2$, using steps 3 and 4.
Thus
\begin{equation}\label{inf:Hilbert-commutator}
 \nrm{[\Delta,m^{-1}]\KD}_{\Hc,1}
 \le C(\delta+\delta^2).
\end{equation}
The potential starts at quadratic order because $q,d$ are harmonic.

\emph{9. Assembly on the graph domain.}
Put $X=(H_\Phi-H_0)H_0^{-1}
=(P_{\Phi_N}\Delta-\Delta)\KD$.
Multiplication by $m^{-1}$ preserves the Dirichlet boundary
condition, and the preceding commutator estimate bounds
$H_0m^{-1}H_0^{-1}$ by an absolute constant. Hence
\[
 \mathcal VH_0^{-1}
 =mX(H_0m^{-1}H_0^{-1})+m[H_0,m^{-1}]H_0^{-1}
\]
has Hilbert weighted norm at most $C(\delta+\delta^2)$,
which is at most $C(\delta+p\delta^2)$.
Finally $H_0K_0^{-1}$ is an angular-diagonal radial spectral
contraction, so multiplying by it proves
\eqref{inf:Hilbert-relative-bound}.
The identities hold first on the smooth Dirichlet graph core;
the bounds extend them to the Hilbert graph realization.
\end{proof}

\begin{theorem}[Exact relative graph bounds]\label{inf:relative-graph}
Fix $\sigma=1$. There is an absolute $C$ independent of $N$ such
that, for every fixed centre order $N$ and all $p\ge P_N$,
\begin{align}
 \nrm{\mathcal VK_0^{-1}}_{\Mmix,1}
       +\nrm{\mathcal VK_0^{-1}}_{\Hc,1}
 &\le C\tau p^{-2},\qquad K_0=H_0+R^2,
 \label{inf:relative-two-norms}\\
 \nrm{(P_{\Phi_N}\Delta-\Delta)\KD}_{\Cc\to\Cc}
       +\nrm{(P_{\Phi_N}\Delta-\Delta)\KD}_{\Mmix\to\Mmix}
 &\le C\tau p^{-2}.\label{inf:relative-scalar}
\end{align}
The density multipliers and their graph conjugates also satisfy,
for the same $p\ge P_N$,
\begin{align}
 \max\{\nrm{M_{m^{\pm1}}}_{\Mmix,1},
                  \nrm{M_{m^{\pm1}}}_{\Hc,1}\}&\le2,\nonumber\\
 \max\{\nrm{H_0M_{m^{-1}}H_0^{-1}}_{\Mmix,1},
                  \nrm{H_0M_{m^{-1}}H_0^{-1}}_{\Hc,1}\}&\le C.
 \label{inf:density-graph-bounds}
\end{align}
All these bounds include every source layer. The operator $\Hu$ has domain
$H^2(B)\cap H^1_0(B)$ in its invariant sector and has the actual
centre-domain spectrum.
\end{theorem}
\begin{proof}
Proposition~\ref{inf:Hilbert-relative} supplies the Hilbert half.
For the mixed and coefficient bounds, use endpoint-normalized
angular coordinates.
We keep the size $\delta$ in \eqref{inf:strong-centre-norm}
throughout the calculation. In $\Cc$ and both weighted matrix
algebras, solid multipliers and their first two Euler derivatives
are controlled by their coefficient sums with the corresponding
degree factors. The ten moments and the factor $(2e^2)^j$
dominate the angular weights, the bandwidth factor $2j+1$, and
these fixed derivatives. Proposition~\ref{inf:module} and
Lemmas~\ref{inf:mixed-multipliers} and \ref{inf:mixed-Poisson} thus give
\[
 \nrm{M_{q-1}}+\nrm{M_{d-1}}\le C\delta,
 \quad
 \nrm{\nabla(q-1)\cdot\nabla\KD}\le C\delta/p,
 \quad
 \nrm{\nabla(q-1)\cdot\nabla\Euler\KD}\le C\delta.
\]
For the coefficient-space version use
Lemma~\ref{inf:gradient-profiles}, which explicitly includes $s=0$.
The same statements hold with a fixed Euler derivative on the
harmonic shape.

Products of two harmonic gradients have multiplier norm
$Cp\delta^2$. In the mixed and coefficient spaces this follows
from the exact product coefficient $2h(h+2k+p-1)c_{ij}^k$
and its bound by $Cp(1+i)(1+j)$.
In the Hilbert matrix algebra use the pointwise product bound
$Cij$, bandwidth at most $i+j$, and the row/column estimate in
Lemma~\ref{inf:mixed-multipliers}; the available moments absorb
the resulting polynomial factors. The coarser bound $Cp\delta^2$
is therefore valid there as well. Products with one Euler
derivative on either factor obey the same bound in terms of
$\delta$.
All inverse powers of $q,d$ are convergent series in these
multiplier algebras when $C\delta<1/2$.

Subtract $\Delta$ in \eqref{inf:scalar-pullback} and compose with
$\KD$. The four terms have bounds, respectively,
\begin{equation}\label{inf:scalar-four-bounds}
 C\delta,\qquad C\delta,\qquad Cp\delta^2,
                 \qquad C\delta^2(1+\delta).
\end{equation}
For the last bound, \eqref{inf:Cq-identity} gives multipliers
$Cp\delta^2$ and $Cp\delta^3$, and
$\Euler\KD$ costs $C/p$.
The third term uses the bounded operator
$\Euler(\Euler+1)\KD$.
This proves a scalar graph bound $C(\delta+p\delta^2)$ in each
of the three realizations, including the two weighted matrix
norms.

The density requires its own estimate. Both $q$ and $d$ are
harmonic, since $\Delta\Euler q=(\Euler+2)\Delta q=0$.
Write $\alpha=p-1/2$. Direct differentiation gives
\begin{align}
 \frac{\nabla m^{-1}}{m^{-1}}
 &=-\alpha q^{-1}\nabla q-\tfrac12d^{-1}\nabla d,
 \label{inf:density-gradient}\\
 \frac{\Delta m^{-1}}{m^{-1}}
 &=\alpha(\alpha+1)q^{-2}|\nabla q|^2
       +\tfrac34d^{-2}|\nabla d|^2
       +\alpha q^{-1}d^{-1}\nabla q\cdot\nabla d.
 \label{inf:density-Laplacian}
\end{align}
The second formula is quadratic in the shape gradients.
The binomial series for $m^{\pm1}$ has norms at most
$\exp(Cp\delta)$, hence at most two after increasing the
threshold using \eqref{inf:uniform-delta}.
The commutator
\[
 [\Delta,m^{-1}]\KD
 =2\nabla m^{-1}\cdot\nabla\KD+(\Delta m^{-1})\KD
\]
has norm at most
$Cp(\delta/p)+Cp^2(p\delta^2)p^{-2}
=C(\delta+p\delta^2)$ in both weighted matrix algebras.
In particular $H_0m^{-1}H_0^{-1}$ is bounded by an absolute
constant there.

Put $X=(H_\Phi-H_0)H_0^{-1}
=(P_{\Phi_N}\Delta-\Delta)\KD$.
The exact factorization
\[
 \mathcal VH_0^{-1}
 =mX(H_0m^{-1}H_0^{-1})+m[H_0,m^{-1}]H_0^{-1}
\]
therefore has norm at most $C(\delta+p\delta^2)$.
The remaining factor $H_0K_0^{-1}$ is a block-diagonal radial
Hilbert contraction in both matrix norms. This proves the
corresponding estimate for $\mathcal VK_0^{-1}$.
Finally insert \eqref{inf:uniform-delta}, with $\tau\le20$.
The operator constants are absolute; the threshold $P_N$ absorbs
the finite-jet remainder. This proves
\eqref{inf:relative-two-norms} and \eqref{inf:relative-scalar}
with $N$-independent constants, as well as
\eqref{inf:density-graph-bounds}.

We also identify their domains. For this finite smooth chart,
composition and multiplication by the positive smooth density
map the Dirichlet $H^2$ domain to itself. The physical smooth
Dirichlet realization therefore gives the asserted domain and
unitary spectrum. Equivalently the small relative graph bound
in $\Hc$ preserves the ball domain by the closed-operator
perturbation theorem. All equalities above first hold for Poisson
images of polynomial sources. Those images are graph dense in
the Dirichlet domain, because polynomial sources are dense in
$L^2$ and $\KD:L^2\to H^2\cap H^1_0$ is continuous.
Thus the bounded graph extensions are the same geometric
operators. By Lemma~\ref{inf:mixed-domain}, on $\Mmix$ their graph domain is
$\{\KD f:f\in\Mmix\}$ with its source graph norm; the embedding
into $\Hc$ and Poisson uniqueness identify it with the restriction
of that realization. The identical argument applies to $\Cc$.
\end{proof}

\section{Local counts and cluster inversion}\label{inf:sec-localization}

The window in this section is an arbitrary fixed number
$W_{\mathrm{loc}}>0$.
The local count depends on this number, but neither on the centre
order nor on the detuning or split.

\begin{theorem}[Local logarithmic count]\label{inf:local-count}
There is an integer $M_{W_{\mathrm{loc}}}\ge1$ with the following
quantifier order: for every fixed $D>0$ there is
$P_{W_{\mathrm{loc}},D}<\infty$ such that every angular interval
of length $D\log p$ contains at most $M_{W_{\mathrm{loc}}}$
states of $\mathcal W_p(W_{\mathrm{loc}})$ whenever
$p\ge P_{W_{\mathrm{loc}},D}$ and $R\in[2p-1,2p+4]$.
In particular the count bound is independent of $D$.
\end{theorem}
\begin{proof}
In the oscillatory region $R-\nu\ge p^{2/3}/2$,
Lemma~\ref{inf:phase-strip} assigns to each state a lattice point
with $|m-\varphi_p(j)|\le C_{W_{\mathrm{loc}}}/p$.
Dividing \eqref{inf:phase-strip-bound} by $\pi$ gives the
stronger bound $C_{W_{\mathrm{loc}}}/(\pi p)$; we use the same
constant in the weaker bound above. The threshold includes those
of Lemmas~\ref{inf:window-isolation} and \ref{inf:phase-strip}
for this window and the condition $2C_{W_{\mathrm{loc}}}/p<1$.
Its curvature satisfies
\[
 c/p\le\varphi_p''(j)\le Cp^{-5/6}.
\]
The lower bound uses $\kappa\le Cp$; the upper uses
$\kappa\ge cp^{5/6}$. On an interval of length $D\log p$,
\[
 \tfrac12Cp^{-5/6}(D\log p)^3
       +2C_{W_{\mathrm{loc}}}D\log p/p<1
\]
for large $p$. Lemma~\ref{inf:lattice-strip} therefore bounds the
number of states by
$2+4\sqrt{C_{W_{\mathrm{loc}}}/c}$.
This number is independent of $D$; only the threshold of the
last inequality depends on $D$.

In the turning region $\nu\ge R-2p^{2/3}$, suppose two states
have angular distance at most $D\log p$.
Their original roots are within $O_{W_{\mathrm{loc}}}(p^{-1})$
of $R$. Transport the root of the larger order to the smaller
order $\nu_0$, keeping its radial index fixed.
Lemma~\ref{inf:bessel-order} moves it by at most
$\pi D\log p$. Every intervening root stays of order $p$.
If the two radial indices were equal, the lower derivative bound
in that lemma would make the original roots differ by at least
$2(1-o(1))$, because the angular indices differ by at least one.
That contradicts their original window width $O(p^{-1})$.
Thus there are two distinct zeros of $J_{\nu_0}$ in an interval
of length $O_{W_{\mathrm{loc}},D}(\log p)$ about $R$.

Throughout that interval the coefficient in
\[
 v''+\left(1-\frac{\nu_0^2-1/4}{x^2}\right)v=0,
 \qquad v(x)=\sqrt xJ_{\nu_0}(x),
\]
is at most $(2+o(1))p^{-1/3}$, uniformly for the fixed window
and interval constant, because
$\nu_0\ge R-2p^{2/3}$ and $x=R+O(\log p)$.
Between the two zeros, integration and the Dirichlet Poincar\'e
inequality give
$\pi^2/\ell^2\le Cp^{-1/3}$ for their separation $\ell$.
Hence $\ell\ge(\pi/\sqrt2+o(1))p^{1/6}$, in particular
$\ell\ge2.2p^{1/6}$ for large $p$, contradicting its logarithmic upper
bound for large $p$. The turning interval therefore contains
at most one state.

The two regions overlap by an order interval of length
$3p^{2/3}/2$. An angular interval of length $D\log p$ has order
length $2D\log p$, which is smaller than this overlap for large
$p$. If it is not entirely in the oscillatory region, all its
possible window points are in the turning region.
There is at most one radial state per angular index by
Lemma~\ref{inf:window-isolation}. Rounding up the oscillatory
count and increasing it to at least one proves the theorem.
\end{proof}

\begin{corollary}[Bounded cluster cardinality]\label{inf:clusters}
Fix $D>0$ and connect consecutive angular indices of
$\mathcal W_p(W_{\mathrm{loc}})$ when their separation is at most
$D\log p$. For all sufficiently large $p$, every resulting
cluster has at most $M_{W_{\mathrm{loc}}}$ states, and distinct
clusters are separated by more than $D\log p$.
A cluster may have diameter as large as
$(M_{W_{\mathrm{loc}}}-1)D\log p$.
The same cardinality bound holds for every subset of the window.
\end{corollary}
\begin{proof}
If a cluster had $M_{W_{\mathrm{loc}}}+1$ states, its first that
many states would span at most
$M_{W_{\mathrm{loc}}}D\log p$.
Apply Theorem~\ref{inf:local-count} with the fixed constant
$M_{W_{\mathrm{loc}}}D$ in place of $D$.
Its bound is still $M_{W_{\mathrm{loc}}}$, a contradiction.
The separation and diameter statements follow from the definition.
Removing states can split clusters but cannot join two clusters
of the larger set, proving the assertion for subsets.
\end{proof}

\begin{lemma}[Inversion with arbitrary diagonal Hilbert weights]
\label{inf:cluster-inverse}
Fix positive constants $c,C,q_{\mathrm{gap}},\sigma$ and an integer $M\ge1$.
Let $D$ satisfy $\sigma D>q_{\mathrm{gap}}M+q_{\mathrm{gap}}+2$.
Let $\tau>0$ and let $F$ be a finite matrix whose indices
carry angular degrees, partitioned into clusters of at most $M$
indices separated by more than $D\log p$.
Assume that:
\begin{enumerate}[label=\textup{(\roman*)}]
\item $F$ is self-adjoint for a positive diagonal Hilbert inner
product and $\nrm{Fv}_H\ge\gamma_F\nrm v_H$, where
$\gamma_F=c\tau p^{-q_{\mathrm{gap}}}$;
\item in a weighted coefficient $\ell^1$ norm,
$\nrm F_1\le C\tau$;
\item deleting all inter-cluster entries has error at most
$C\tau p^{-\sigma D}$ in both that coefficient norm and the
specified Hilbert operator norm.
\end{enumerate}
Then, for sufficiently large $p$ depending only on
$c,C,q_{\mathrm{gap}},\sigma,M,D$,
\begin{equation}\label{inf:cofactor-inverse-bound}
 \nrm{F^{-1}}_1\le C'\tau^{-1}p^{q_{\mathrm{gap}}M}.
\end{equation}
The assertion is uniform as $\tau\downarrow0$.
\end{lemma}
\begin{proof}
Let $F_{\mathrm b}$ keep the diagonal cluster blocks.
The symmetric deletion of paired entries preserves
self-adjointness in the diagonal Hilbert structure.
The Hilbert error is at most $\gamma_F/2$ for large $p$, since
$\sigma D>q_{\mathrm{gap}}$. Thus
\[
 \nrm{F_{\mathrm b}v}_H\ge
 (\gamma_F-C\tau p^{-\sigma D})\nrm v_H\ge(\gamma_F/2)\nrm v_H.
\]
Every eigenvalue of every cluster block is real and has modulus
at least $\gamma_F/2$.

Conjugate a cluster block of size $m\le M$ by its coefficient
weights, so its $\ell^1$ norm is the unweighted column norm.
All its entries have modulus at most $C\tau$.
The determinant is invariant under this conjugation and has
modulus at least $(\gamma_F/2)^m$.
Expansion of each cofactor as a sum of $(m-1)!$ products gives
its absolute value at most $(m-1)!(C\tau)^{m-1}$.
For $m=1$ the empty product is one.
Cramer's formula and summing each inverse column give
\[
 \nrm{F_{\mathrm b,m}^{-1}}_1
 \le m!(C\tau)^{m-1}(2/\gamma_F)^m
 \le C_{M,c,C}\tau^{-1}p^{q_{\mathrm{gap}}M}.
\]
The block diagonal inverse has the maximum of these norms.
Its product with the discarded coefficient error is at most
$C'p^{q_{\mathrm{gap}}M-\sigma D}$, below $1/2$ for large $p$.
A Neumann series restores that error and proves
\eqref{inf:cofactor-inverse-bound}.
In both absorptions the factor $\tau$ cancels. The inverse has
exactly one remaining factor $\tau^{-1}$.
For an empty index set the bound is vacuous.
\end{proof}

The next section uses a window scaled by $\tau$,
and applies the local count to a fixed larger window containing it.

\section{The full inverse in coefficient norm}\label{inf:sec-inverse}

Fix any centre order $N$. All constants and exponents in this section are fixed
independently of that order. Only the large-$p$ threshold may
depend on it. We assume the actual centre gap
\eqref{inf:centre-gap}, supplied by Theorem~\ref{inf:integer-gap}.
Write
\[
 \mathcal A=\Hu-R^2,\qquad \mathcal D_0=H_0-R^2,
 \qquad K_0=H_0+R^2.
\]
The actual scalar Dirichlet graph operator is
$J_D=B_{\Phi_N,D}\KD$ and includes the harmonic source layer and
the seed.

\subsection{A detuning-scaled window}

Let $C_{\mathrm{rel}}\ge1$ be a constant in
\eqref{inf:relative-two-norms}. Fix
$E_{\mathrm{inv}}\ge64C_{\mathrm{rel}}$, and define
\begin{equation}\label{inf:inverse-window}
 \Pi=\one_{[-E_{\mathrm{inv}}\tau,E_{\mathrm{inv}}\tau]}
                                    (H_0-R^2),\qquad Q=I-\Pi.
\end{equation}
This is the full ball spectral projection, including every radial
state in the interval. It is distinct from the fixed flow-window
projection. The local count will use
$W_{\mathrm{loc}}=20E_{\mathrm{inv}}$.

\begin{lemma}[Inverse on the full complement]\label{inf:inverse-complement}
For $E_{\mathrm{inv}}\ge64C_{\mathrm{rel}}$ and all $p$ above a
threshold depending on $C_{\mathrm{rel}},E_{\mathrm{inv}},N$,
the operator $\mathcal A_Q=Q\mathcal AQ$ is
invertible on its graph domain, and
\begin{equation}\label{inf:inverse-complement-bound}
 \nrm{K_0\mathcal A_Q^{-1}Q}_{\Mmix,1}
 +\nrm{K_0\mathcal A_Q^{-1}Q}_{\Hc,1}
 \le Cp^2/(E_{\mathrm{inv}}\tau).
\end{equation}
The two inverses coincide on their common sources.
\end{lemma}
\begin{proof}
The projections in \eqref{inf:inverse-window} are diagonal in
angular degree and orthogonal in each radial Hilbert space.
Their norms in both matrix algebras and both underlying spaces
are at most one.
Put $D_Q=Q\mathcal D_0Q$. Radial spectral calculus gives
\begin{equation}\label{inf:ball-complement-bound}
 \nrm{K_0D_Q^{-1}Q}\le1+\frac{2R^2}{E_{\mathrm{inv}}\tau}
\end{equation}
in all four norms.
Since $R/p\le33/16$ for $p\ge64$, its product with
\eqref{inf:relative-two-norms} is at most
\[
 \frac{20C_{\mathrm{rel}}}{p^2}
 +\frac{1089C_{\mathrm{rel}}}{128E_{\mathrm{inv}}}.
\]
The chosen window and $p^2\ge80C_{\mathrm{rel}}$ make this at
most $1/4+1089/8192<1/2$. The window is fixed before this
threshold is imposed.
Factor
\[
 \mathcal A_Q=(I+Q\mathcal VK_0^{-1}K_0D_Q^{-1}Q)D_Q.
\]
The norm-convergent Neumann series in each matrix algebra gives
\[
 \nrm{K_0\mathcal A_Q^{-1}Q}
 \le2\left(1+\frac{2R^2}{E_{\mathrm{inv}}\tau}\right).
\]
Increasing the threshold to include $p^2\ge40E_{\mathrm{inv}}$
absorbs the constant term into $Cp^2/(E_{\mathrm{inv}}\tau)$,
proving \eqref{inf:inverse-complement-bound}.
It maps sources into the graph domain because of its leading
$D_Q^{-1}$ factor and the bound after multiplication by $K_0$.
By the graph realizations of Lemma~\ref{inf:mixed-domain}, the
embedding $\Mmix\subset\Hc$, and uniqueness of the Hilbert
inverse, the two constructions agree on common sources.
\end{proof}

\begin{proposition}[The Schur matrix and its actual gap]
\label{inf:inverse-Schur}
Define the exact Schur matrix at spectral parameter zero by
\begin{equation}\label{inf:effective-matrix}
 \Feff=\Pi\mathcal A\Pi-
               \Pi\mathcal AQ\mathcal A_Q^{-1}Q\mathcal A\Pi.
\end{equation}
For large $p$, each angular index occurs at most once in its index
set. In radial unit vectors with the original angular $G_j$,
its Hilbert inner product has diagonal weights $\eta_j^2$.
It is self-adjoint for that inner product, satisfies
\begin{equation}\label{inf:effective-norm}
 \nrm{\Feff}_{\Mmix,1}+\nrm{\Feff}_{\Hc,1}
                                  \le C_{E_{\mathrm{inv}}}\tau,
\end{equation}
and has the Hilbert gap
\begin{equation}\label{inf:effective-gap}
 \nrm{\Feff v}_\Hc\ge\gamma_p\nrm v_\Hc.
\end{equation}
\end{proposition}
\begin{proof}
The containing fixed window is $20E_{\mathrm{inv}}$.
Lemma~\ref{inf:window-isolation} gives one radial state per angular
index. The normalized radial basis therefore identifies the
mixed norm on this finite space with $\ell^1(2^j)$ exactly.
Self-adjointness of \eqref{inf:effective-matrix} follows from
that of $\mathcal A$ and the complementary Hilbert inverse.
The stated diagonal Hilbert weights follow from
\eqref{inf:true-H-norm}.

On $\Pi$, $K_0$ has norm at most $Cp^2$.
Equation~\eqref{inf:relative-two-norms} gives
$\nrm{\Pi\mathcal V\Pi}\le C\tau$ in both weighted norms.
The feedback factors as
\[
 (\Pi\mathcal VK_0^{-1})Q(K_0\mathcal A_Q^{-1}Q)
                     (\mathcal VK_0^{-1})(K_0\Pi).
\]
Its four factors have sizes $C\tau/p^2$,
$Cp^2/(E_{\mathrm{inv}}\tau)$, $C\tau/p^2$, and $Cp^2$,
respectively. Its norm is at most $C\tau/E_{\mathrm{inv}}$.
The unperturbed diagonal has norm at most $E_{\mathrm{inv}}\tau$.
Submultiplicativity from Lemma~\ref{inf:mixed-algebra} proves
\eqref{inf:effective-norm}, including
exponential angular decay of the feedback.

For $v\in\Pi\Hc$ set
$u=v-\mathcal A_Q^{-1}Q\mathcal A\Pi v$.
It belongs to the graph domain and obeys
$Q\mathcal Au=0$, $\Pi\mathcal Au=\Feff v$.
Orthogonality gives $\nrm u_\Hc\ge\nrm v_\Hc$.
The actual centre-domain gap implies
$\nrm{\mathcal Au}_\Hc\ge\gamma_p\nrm u_\Hc$.
These equalities prove \eqref{inf:effective-gap} for the exact
Schur matrix.
\end{proof}

\begin{corollary}[Finite-matrix coefficient inverse]
\label{inf:effective-inverse}
Let $M=M_{20E_{\mathrm{inv}}}$ be the local-count bound from
Theorem~\ref{inf:local-count}. Then
\begin{equation}\label{inf:effective-inverse-bound}
 \nrm{\Feff^{-1}}_{\ell^1(2^j)}
                         \le C\tau^{-1}p^{5M/3}
\end{equation}
for all sufficiently large $p$.
Its constant and exponent are independent of the centre order.
\end{corollary}
\begin{proof}
Take clusters at separation $D_{\mathrm{sep}}\log p$ in the
actual window. Corollary~\ref{inf:clusters}, applied to its fixed
superset of width $20E_{\mathrm{inv}}$, bounds their cardinalities
by $M$, independently of $D_{\mathrm{sep}}$.
Deleting inter-cluster entries in \eqref{inf:effective-norm}
costs at most $C\tau p^{-D_{\mathrm{sep}}}$ in the coefficient
norm by weighted column sums and in the Hilbert norm by its
weighted row and column sums. Apply
Lemma~\ref{inf:cluster-inverse} with $q_{\mathrm{gap}}=5/3$, $\sigma=1$, and
$c=c_{\mathrm{gap}}$.
Choose once and for all
$D_{\mathrm{sep}}>5M/3+5$.
This proves \eqref{inf:effective-inverse-bound} and its stated
independence.
\end{proof}

\subsection{Recovering the full radial graph}

\begin{proposition}[Mixed-space inverse]\label{inf:mixed-inverse}
The full operator $\mathcal A$ and scalar graph $J_D$ have
inverses on the mixed space, with
\begin{equation}\label{inf:mixed-inverse-bound}
 \nrm{K_0\mathcal A^{-1}}_{\Mmix\to\Mmix}
       +\nrm{J_D^{-1}}_{\Mmix\to\Mmix}
 \le C\tau^{-1}p^{2+5M/3}.
\end{equation}
These inverses agree with their Hilbert realizations.
\end{proposition}
\begin{proof}
For a source $f\in\Mmix$, the exact block solution is
\[
 v=\Feff^{-1}(\Pi f-\Pi\mathcal VQ\mathcal A_Q^{-1}Qf),\qquad
 Qu=\mathcal A_Q^{-1}(Qf-Q\mathcal V\Pi v),\qquad u=v+Qu.
\]
The first off-diagonal factor has norm at most $C/E_{\mathrm{inv}}$,
by \eqref{inf:relative-two-norms} and
\eqref{inf:inverse-complement-bound}. The input $Q\mathcal V\Pi$
has norm at most $C\tau$, and $K_0\Pi$ costs $Cp^2$.
Combining these bounds with \eqref{inf:effective-inverse-bound}
gives $\nrm{K_0u}\le C\tau^{-1}p^{2+5M/3}\nrm f$.
By Lemma~\ref{inf:mixed-domain}, the reconstruction lies in the
mixed graph domain and solves the
full equation, whose Hilbert solution is unique by the centre gap.
This proves the first inverse and identifies it with the Hilbert one.

For the scalar equation $B_{\Phi_N}W=f$, the unitary equation is
$\mathcal A(mW)=-mf$. Equation~\eqref{inf:density-graph-bounds}
bounds multiplication by $m$ and $H_0m^{-1}H_0^{-1}$ on $\Mmix$.
Since $H_0K_0^{-1}$ is an angular-diagonal radial contraction,
it also bounds $H_0m^{-1}K_0^{-1}$ by an absolute constant.
Thus
\[
 J_D^{-1}f=\Delta W
 =(H_0m^{-1}K_0^{-1})(K_0\mathcal A^{-1})(mf),
\]
which proves the second bound in \eqref{inf:mixed-inverse-bound}.
Its signs use $\Delta=-H_0$ and $mW=-\mathcal A^{-1}mf$.
Both inverse identities follow from the graph-domain equation and
Hilbert uniqueness. Every radial and angular source, including the
harmonic layer and the seed, is part of this reconstruction.
\end{proof}

\subsection{Finite head, infinite tail, and membership}

\begin{lemma}[Radial tail and finite-head conversion]
\label{inf:head-tail}
For this lemma the scalar hypotheses on the fixed real finite-centre
chart are $R\in[2p-1,2p+4]$ and
\begin{equation}\label{inf:head-tail-hypotheses}
 \nrm{(P_{\Phi_N}\Delta-\Delta)\KD}_{\Cc\to\Cc}
 +\nrm{(P_{\Phi_N}\Delta-\Delta)\KD}_{\Mmix\to\Mmix}
 \le\frac14.
\end{equation}
There is a fixed $K_{\mathrm{head}}$ independent of the centre order
such that, with coefficient projections
\[
 \Pi_T=\one_{\{j+s\le K_{\mathrm{head}}p\}},\qquad Q_T=I-\Pi_T,
\]
the compressed operator $Q_TJ_DQ_T$ has inverse of norm at most
two in both $Q_T\Cc$ and $Q_T\Mmix$ for large $p$.
Moreover
\begin{equation}\label{inf:head-conversion}
 \nrm{\Pi_Tg}_\Cc\le C_{K_{\mathrm{head}}}\sqrt p\,
                                                   \nrm g_{\Mmix}.
\end{equation}
\end{lemma}
\begin{proof}
The Poisson columns in Lemma~\ref{inf:poisson-columns} have radial
bandwidth one and satisfy the column bound
\eqref{inf:poisson-column-bounds}, including the harmonic column.
On the indicated tail this is at most
$C/((1+2K_{\mathrm{head}})p-1)^2$.
The ratio of neighbouring coefficient weights is at most $e^{2/p}$.
Therefore
\[
 \nrm{R^2Q_T\KD Q_T}_{\Cc\to\Cc}
 \le C/(1+2K_{\mathrm{head}})^2.
\]
For the mixed bound, normalize the radial polynomials by their
squared norms $n_{j,s}=p/(p+2j+2s)$.
The neighbouring ratios of $n_{j,s}$ and their inverses are
bounded by two. After this normalization the tridiagonal
Poisson matrix has row and column sums bounded by the same
quantity up to an absolute factor: adjacent denominators differ
by two and are comparable. Removing entries by $Q_T$ does not
increase those absolute sums. The radial Schur inequality in
each angular block gives
$\nrm{R^2Q_T\KD Q_T}_{\Mmix\to\Mmix}
\le C/(1+2K_{\mathrm{head}})^2$ as well.

Choose the fixed cutoff so that both bounds are below $1/4$.
The hypothesis \eqref{inf:head-tail-hypotheses} bounds the norm of
$Q_T(P_{\Phi_N}\Delta-\Delta)\KD Q_T$ by $1/4$ in both spaces.
Since
\[
 Q_TJ_DQ_T=I_{Q_T}+R^2Q_T\KD Q_T
                   +Q_T(P_{\Phi_N}\Delta-\Delta)\KD Q_T,
\]
a Neumann series proves both tail inverses and their norm bound.
The two series agree on common sources.

On the head, $\wt(2j+2s)\le e^{2K_{\mathrm{head}}}$ and
$n_{j,s}\ge(1+2K_{\mathrm{head}})^{-1}$.
Each angular block contains at most $K_{\mathrm{head}}p+1$ radial
coefficients. Cauchy--Schwarz gives
\[
 \sum_{s\le K_{\mathrm{head}}p-j}\wt(2j+2s)|g_{j,s}|
 \le e^{2K_{\mathrm{head}}}
 \sqrt{(K_{\mathrm{head}}p+1)(1+2K_{\mathrm{head}})}
 \left(\sum_s n_{j,s}|g_{j,s}|^2\right)^{1/2}.
\]
Multiply by $2^j$ and sum in $j$ to prove
\eqref{inf:head-conversion}. The common angular coefficient
structure gives the loss $\sqrt p$.
\end{proof}

\begin{theorem}[Two-sided coefficient inverse]\label{inf:coefficient-inverse}
For every fixed centre order $N$ and all sufficiently large $p$,
$J_D$ is a bounded isomorphism of $\Cc$, and
\begin{equation}\label{inf:coefficient-inverse-bound}
 \nrm{J_D^{-1}}_{\Cc\to\Cc}
 \le C\tau^{-1}p^{5/2+5M/3}
 \le C\frac{p^{b_{\mathrm{inv}}}}{\gamma_p},
 \qquad b_{\mathrm{inv}}=2M+3.
\end{equation}
The constants and exponent are independent of $N$; its threshold
may depend on $N$. The estimates are uniform in
$0<\tau\le20$ and all admissible splits with the stated gap.
\end{theorem}
\begin{proof}
Equation~\eqref{inf:relative-scalar} of
Theorem~\ref{inf:relative-graph} bounds the sum in
\eqref{inf:head-tail-hypotheses} by $C\tau p^{-2}\le1/4$
for $0<\tau\le T$ and $p$ above the fixed-order threshold,
after increasing it depending on the fixed $T$.
Together with $R\in[2p-1,2p+4]$, this supplies the hypotheses
of Lemma~\ref{inf:head-tail}; here $T=20$, and arbitrary fixed
$T$ is covered by Lemma~\ref{inf:fixed-detuning-extension}.
For $f\in\Cc$, first solve $J_Dg=f$ in $\Mmix$ by
Proposition~\ref{inf:mixed-inverse}.
The head $\Pi_Tg$ is finite dimensional and belongs to $\Cc$,
with bound \eqref{inf:head-conversion}.
The full $J_D$ is bounded on $\Cc$ by an absolute constant:
use $J_D=I+R^2\KD+(P_{\Phi_N}\Delta-\Delta)\KD$,
\eqref{inf:poisson-operator-bounds}, $R^2=O(p^2)$, and
\eqref{inf:relative-scalar}.
Thus the tail equation with source
$Q_Tf-Q_TJ_D\Pi_Tg$ has a solution $z\in Q_T\Cc$ and
\[
 \nrm z_\Cc\le C(\nrm f_\Cc+\nrm{\Pi_Tg}_\Cc).
\]
By the embedding into $\Mmix$ and uniqueness of its tail equation,
this $z$ is exactly $Q_Tg$.
This proves membership in $\Cc$ before applying a norm estimate
to the full solution. Hence
\[
 \nrm g_\Cc\le C\nrm f_\Cc+
                      C\sqrt p\nrm g_{\Mmix}
 \le C\tau^{-1}p^{5/2+5M/3}\nrm f_\Cc.
\]
The last bound absorbs the first term since $M\ge1$ and
$\tau\le20$.
The constructed right inverse is injective on $\Cc$ by the
Hilbert realization and its gap, so both inverse identities hold.
Finally
$5/2+5M/3-5/3=5M/3+5/6\le2M+3$.
Using \eqref{inf:centre-gap} gives the second inequality in
\eqref{inf:coefficient-inverse-bound}.
Every constant used in the matrix, reconstruction and tail steps
was fixed before the centre order.
\end{proof}

\section{The Newton threshold and the exact zero}\label{inf:sec-newton}

The closing argument uses a threshold for the residual and two
polynomial operator bounds in the same spaces. It requires the
actual integer-split gap already obtained.

\begin{theorem}[Polynomial Newton threshold]\label{inf:newton-threshold}
Suppose fixed exponents $b_{\mathrm{inv}},m_{\mathrm{mat}},d_{\mathrm{nl}}\ge0$ have been chosen independently
of the centre accuracy $N$. Assume, uniformly in $0<\tau=|\theta|\le20$
and the admissible splits, for all sufficiently large $p$, that
in the real Banach spaces $\Xc_p,\Cc$ of
Definition~\ref{inf:spaces} a centre $z_0$ satisfies
\begin{align*}
 D\Fnl(z_0)&=J_D\Qmat+\mathcal V_{\mathrm{res}},&
 \nrm{\Fnl(z_0)}&\le C_N\theta^2p^{2-N},\\
 \nrm{\mathcal V_{\mathrm{res}}}&\le C_N\theta^2p^{3-N},&
 \nrm{J_D^{-1}}&\le C_Np^{b_{\mathrm{inv}}}/(c_{\mathrm{gap}}\tau p^{-5/3}),\\
 \nrm{\Qmat^{-1}}&\le C_Np^{m_{\mathrm{mat}}},&
 \sup\nrm{D^2\Fnl}&\le C_Np^{d_{\mathrm{nl}}}.
\end{align*}
Here $J_D$ and $\Qmat$ are bounded isomorphisms, and the last bound
holds on the valid closed centre ball considered below, on which
$\Fnl$ is $C^2$.
Put
\begin{equation}\label{inf:threshold-exponents}
 B_{\mathrm{Newt}}=b_{\mathrm{inv}}+m_{\mathrm{mat}}+5/3,\qquad
 L_{\mathrm{Newt}}=B_{\mathrm{Newt}}+d_{\mathrm{nl}}+4.
\end{equation}
Assume that the accuracy order satisfies
\begin{equation}\label{inf:threshold-accuracy}
 N\ge\lceil2B_{\mathrm{Newt}}+d_{\mathrm{nl}}+8\rceil.
\end{equation}
For $p$ above a threshold depending on the fixed constants and
$N$, the closed ball of radius
$\mathfrak r_p=\tau p^{-L_{\mathrm{Newt}}}$ about $z_0$ contains a
real zero of $\Fnl$, provided that ball is inside the domain of the
map. The threshold is uniform for $0<\tau\le20$.
\end{theorem}
\begin{proof}
Let $\AN=\Qmat^{-1}J_D^{-1}$. This isomorphism
$\Cc\to\Xc_p$ satisfies
$\nrm{\AN}\le C_Np^{B_{\mathrm{Newt}}}/\tau$.
The three quantities controlling the contraction are
\begin{align}
 \frac{\nrm{\AN\Fnl(z_0)}}{\mathfrak r_p}
 &\le C_Np^{B_{\mathrm{Newt}}+L_{\mathrm{Newt}}+2-N}
                                      \le C_Np^{-2},
 \label{inf:newton-gate-one}\\
 \nrm{I-\AN D\Fnl(z_0)}
 &\le C_N\tau p^{B_{\mathrm{Newt}}+3-N}\le20C_Np^{-5},
 \label{inf:newton-gate-two}\\
 \sup\nrm{\AN D^2\Fnl}\,\mathfrak r_p
 &\le C_Np^{B_{\mathrm{Newt}}+d_{\mathrm{nl}}-L_{\mathrm{Newt}}}
                                      =C_Np^{-4}.
 \label{inf:newton-gate-three}
\end{align}
The exact factorization gives
$I-\AN D\Fnl(z_0)=-\AN\mathcal V_{\mathrm{res}}$, proving
the second estimate. The exponents follow from
\eqref{inf:threshold-exponents} and \eqref{inf:threshold-accuracy}.
Increase $p$ until each quantity is at most $1/4$.
For $\mathcal T(z)=z-\AN\Fnl(z)$, Taylor's integral formula
bounds its displacement from the centre on the closed ball by
$\mathfrak r_p/4+\mathfrak r_p/4+\mathfrak r_p/8<\mathfrak r_p$.
The same formula for its derivative gives Lipschitz constant at
most $1/4+1/4=1/2$.
The ball is complete, so Banach's contraction theorem supplies
a fixed point. Since $\AN$ is injective, its fixed-point equation
implies $\Fnl(z)=0$. All the smallness conditions are uniform for
$0<\tau\le20$, as the displayed powers show.
\end{proof}

\begin{corollary}[Instantiation for the centre family]
\label{inf:Newton-zero}
Fix the constants in the following order:
\begin{equation}\label{inf:constant-order}
 \begin{aligned}
 C_0&\longrightarrow C_{\mathrm{rel}}\longrightarrow
 E_{\mathrm{inv}}\longrightarrow M=M_{20E_{\mathrm{inv}}}\\
 &\longrightarrow (D_{\mathrm{sep}},K_{\mathrm{head}})
 \longrightarrow b_{\mathrm{inv}}=2M+3
 \longrightarrow N=4M+40.
 \end{aligned}
\end{equation}
Here $C_0$ and $C_{\mathrm{rel}}$ are those of
\eqref{inf:uniform-delta} and \eqref{inf:relative-two-norms}.
For every sufficiently large $p\in\mathcal P_{\mathrm{bd}}$, choose a split
from Theorem~\ref{inf:integer-gap} and the order-$N$ centre at
that split. There is an exact real zero $z_*=(g_*,f_*)$ with
\begin{equation}\label{inf:Newton-correction}
 \nrm{z_*-z_0}_{\Xc_p}\le\tau p^{-L_{\mathrm{Newt}}},\qquad
 B_{\mathrm{Newt}}=2M+32/3,\quad L_{\mathrm{Newt}}=2M+62/3.
\end{equation}
\end{corollary}
\begin{proof}
The centre gap follows from Theorem~\ref{inf:integer-gap} and
Proposition~\ref{inf:centre-gap-transfer} at the now fixed order.
The coefficient inverse has exponent $b_{\mathrm{inv}}$ by
Theorem~\ref{inf:coefficient-inverse}.
The material inverse and second derivative have exponents
$m_{\mathrm{mat}}=d_{\mathrm{nl}}=6$ by \eqref{inf:Qmat-bounds} and \eqref{inf:p6-second}.
The residuals and factorization are
\eqref{inf:residual-moments}, \eqref{inf:residual-material-bound},
and \eqref{inf:material-factorization}.
Thus all hypotheses of Theorem~\ref{inf:newton-threshold} hold
in exactly $\Xc_p=\Gc_p\times\Ec_1$ and $\Cc_{p,2}$.
The choice $N=4M+40$ exceeds
$\lceil2B_{\mathrm{Newt}}+14\rceil$.
The three gate exponents in this instance are
\[
 B_{\mathrm{Newt}}+L_{\mathrm{Newt}}+2-N=-20/3,\qquad
 B_{\mathrm{Newt}}+3-N=-2M-79/3,\qquad
 B_{\mathrm{Newt}}+6-L_{\mathrm{Newt}}=-4.
\]
The second gate retains its factor $\tau$; the first and third
have no residual negative power of it.
Since $\nrm{f_N}_{\Ec_1}=O_N(\tau p^{-2})$ and
$\mathfrak r_p\to0$ uniformly for $\tau\le20$, the closed ball
is eventually in the product region
\eqref{inf:nonlinear-working-region}: its source displacement is
at most one and its total shape norm is at most $1/4$. Increase the threshold to meet all preceding estimates
for this one fixed $N$ and apply the threshold theorem.
The spaces, centre and operators are real by
Definition~\ref{inf:spaces} and their defining formulas, so the
contraction acts on the real closed ball.
\end{proof}

\section{Geometry, classical solutions, and the Pompeiu conclusion}
\label{inf:sec-geometry}

\begin{lemma}[Analyticity and full-dimensional strict convexity]
\label{inf:geometry-convexity}
Let $f\in\Ec_1$ be real with
$\nrm f_\infty+\nrm{\nabla_S^2f}_{\mathrm{op},\infty}<1/4$.
The radial graph $\{(1+f(\omega))\omega:\omega\in S^{2p-1}\}$
is a closed analytic hypersurface bounding a bounded strictly
convex domain. It is invariant under $O(a)\times O(b)$.
The conclusion concerns all ambient directions, including the
block-orbit directions and the axes.
\end{lemma}
\begin{proof}
By Lemma~\ref{inf:space-embedding}, $f$ is holomorphic in a
complex neighbourhood of $x\in[0,1]$. On the real sphere
$x=|X|^2$ is a polynomial. The radial parametrization is therefore
analytic, including at $x=0,1$. Its radius is positive and its
radial derivative is nonsingular, so it is an embedded closed
hypersurface. Its symmetry follows from dependence on $x$ alone.

Put $r_f=1+f$ and let $g_S$ be the sphere metric.
Differentiating the radial parametrization and its outward normal
gives the second fundamental form with positive denominator and
numerator
\begin{equation}\label{inf:curvature-numerator}
 r_f^2g_S+2\,df\otimes df-r_f\nabla_S^2f.
\end{equation}
Since $3/4<r_f<5/4$, its quadratic form on every unit tangent
vector is at least $9/16-5/16=1/4$.
This proves positivity in the full tangent space.
To make the block directions explicit, at $0<x<1$ a unit
$X$-orbit tangent vector has
$\nabla_S^2x(v,v)=2(1-x)$ and $dx(v)=0$; a unit $Y$-orbit
tangent vector has $\nabla_S^2x(v,v)=-2x$.
Their numerators in \eqref{inf:curvature-numerator} are
$r_f^2-2r_f(1-x)f'$ and $r_f^2+2r_fx f'$.
For the unit meridian direction the numerator is
\[
 r_f^2+8x(1-x)(f')^2
   -r_f\{4x(1-x)f''+2(1-2x)f'\}.
\]
The coordinate-free positive bound already controls all three
expressions and their limits at the axes.
The global convexity theorem for compact hypersurfaces with
positive second fundamental form (Hadamard's theorem, in the form
given by \cite{inf:doCarmoLima}; see also \cite{inf:sacksteder})
now makes the enclosed radial domain
strictly convex. Its hypotheses hold: the hypersurface is smooth,
compact, connected and embedded, and every principal curvature
is positive. The convex boundary contains no line segment, since
an interior point of such a segment would have zero normal
curvature in its direction. Thus the convexity is strict.
Boundedness follows from its bounded radius.
\end{proof}

\begin{lemma}[The non-ball witness survives the correction]
\label{inf:geometry-nonball}
At the zero supplied by Corollary~\ref{inf:Newton-zero}, the
normalized shape $f_*$ satisfies the hypotheses of
Lemma~\ref{inf:geometry-convexity} for large $p$, and its domain
is not a ball.
\end{lemma}
\begin{proof}
The fixed centre has normalized $C^2$ size at most
$CC_0\tau p^{-2}$ by \eqref{inf:shape-C2},
$\nrm{f_N}_{\Ec_1}\le\delta$, and \eqref{inf:uniform-delta}.
The correction has normalized $C^2$ size at most
$C\tau p^{-L_{\mathrm{Newt}}}$ by \eqref{inf:shape-C2} and
\eqref{inf:Newton-correction}. Their sum is below $1/4$ for
large $p$, uniformly for $\tau\le20$.

Consider the physical monic seed functional
$\mathfrak W_p(f)=R[\mathsf H_2]f=Rf_2/\mathsf H_2(1)$.
Its dual norm on $\Ec_1$ is at most $Cp$, since $R=O(p)$ and
$\mathsf H_2(1)=(\bb)_2/(p+1)_2$ is bounded away from zero uniformly
in the split. At the centre,
\[
 \mathfrak W_p(f_N)
 =-\frac{\theta}{16r(1-r)p}+O_N(\tau p^{-2}),
\]
so its absolute value is at least $c\tau/p$ for large $p$.
The correction changes it by at most
$Cp\tau p^{-L_{\mathrm{Newt}}}=o(\tau/p)$ since $L_{\mathrm{Newt}}>2$.
It remains nonzero.
Any ball invariant under both orthogonal block groups has centre
fixed by both groups, hence centre zero. Its radial function
would be constant and would have zero monic degree-two angular
coefficient. The nonzero functional excludes such a ball.
\end{proof}

\begin{proposition}[Identification of the zero with a classical solution]
\label{inf:classical-solution}
For a real zero $(g,f)$ of \eqref{inf:nonlinear-definition} in the
small chart, set $U=1+Kg$, $\widehat\Omega=\Phi_f(B)$, and
$w=U\circ\Phi_f^{-1}$. Then $w$ is a classical solution of
\begin{equation}\label{inf:normalized-PDE}
 \Delta w+R^2w=0\quad\hbox{in }\widehat\Omega,
 \qquad w=1,\quad\nabla w=0\quad\hbox{on }\partial\widehat\Omega.
\end{equation}
It is smooth up to the boundary and real analytic there and in
the interior. For $\Omega=R\widehat\Omega$ and $u(x)=w(x/R)$,
one has $\Delta u+u=0$, boundary value one and zero boundary
gradient. The field is nonconstant.
\end{proposition}
\begin{proof}
The source $g$ is an actual $L^2$ function by
Lemma~\ref{inf:space-embedding}. Its Poisson field belongs to
$H^2(B)\cap H^1_0(B)$, so $U\in H^2(B)$ and its Dirichlet
trace is one. The completed-space trace identity
\eqref{inf:poisson-trace} makes its normal trace zero because
$g\in\Gc_p$.
The tangential trace of its gradient is the tangential derivative
of its constant Dirichlet trace, hence zero in the Sobolev trace
sense. Equivalently, finite clamped Poisson sums have zero full-gradient
trace and converge to $U$ in $H^2$; continuity of the gradient
trace gives the same conclusion on the completion. Thus the
whole gradient trace is zero.
The map $\Phi_f$ is a smooth diffeomorphism of closed domains;
composition transports $U$ to an $H^2$ field on
$\widehat\Omega$. Proposition~\ref{inf:nonlinear-map} identifies
its coefficient equation with the weak physical equation in
\eqref{inf:normalized-PDE}.
The trace chain rule
$\nabla w\circ\Phi_f=D\Phi_f^{-T}\nabla U$ transports the zero
full-gradient trace as well.

Dirichlet elliptic regularity on a smooth bounded domain applies
to $\Delta w=-R^2w$ and the constant boundary data. Starting from
$H^2$, each application improves regularity by two derivatives;
iteration gives all Sobolev orders, hence $C^\infty$ up to the
boundary in the fixed dimension; see \cite[Chs.~6, 8, 9]{inf:GT}.
The boundary is analytic by Lemma~\ref{inf:space-embedding}, the
operator has analytic constant coefficients, and its Dirichlet
data are analytic. Analytic interior regularity \cite{inf:MorreyNirenberg} and
analytic Dirichlet boundary regularity \cite{inf:MorreyBoundary}
give real analyticity up to that boundary. The scalar Laplacian
is strongly elliptic, and the Dirichlet condition is its
complementing boundary condition; the lower-order term $R^2$
has analytic coefficients. Its now classical gradient trace agrees with the zero
Sobolev trace. This proves \eqref{inf:normalized-PDE} pointwise.

The dilation gives
$\Delta_xu=R^{-2}(\Delta w)(x/R)=-u$ and preserves both boundary
conditions. The only constant solution of $\Delta u+u=0$ is zero,
which is incompatible with boundary value one. Hence $u$ is
nonconstant.
\end{proof}

\begin{remark}[The symmetry sector]\label{inf:sector-distinction}
The Dirichlet inverse throughout the construction is only in the
$O(a)\times O(b)$ invariant sector. At the exact solution, Cartesian
derivatives of $u$ satisfy the frequency-one equation and have zero
Dirichlet trace. At least one is nonzero, so the full ambient
Dirichlet operator has eigenvalue one. Since $u$ is invariant,
its $X$-derivatives
transform as the standard $O(a)$ representation and its
$Y$-derivatives as the standard $O(b)$ representation. Averaging
any of these derivatives over the corresponding reflection group
gives zero. They therefore have zero invariant projection.
Within the invariant sector the graph inverse used above includes
every mode, in particular the degree-four seed and the harmonic
source layer.
\end{remark}

\begin{proposition}[Failure of the Pompeiu property]\label{inf:Pompeiu}
Every bounded domain carrying a classical field constructed above fails
the Pompeiu property.
\end{proposition}
\begin{proof}
For any $\xi\in\R^{2p}$ with $|\xi|=1$, put
$v(x)=e^{i\xi\cdot x}$. Both $u$ and $v$ solve the frequency-one
Helmholtz equation. Green's second identity gives
\[
 0=\int_{\partial\Omega}(u\partial_\nu v-v\partial_\nu u)
   =\int_{\partial\Omega}\partial_\nu v
   =\int_\Omega\Delta v=-\int_\Omega e^{i\xi\cdot x}\,dx.
\]
Thus the Fourier transform of the indicator vanishes on the unit
sphere. Fix a unit vector $\xi_0$.
For every orthogonal map $\mathsf O$ and translation vector $x_{\mathrm{tr}}$,
\[
 \int_{x_{\mathrm{tr}}+\mathsf O\Omega}e^{i\xi_0\cdot x}\,dx
 =e^{i\xi_0\cdot x_{\mathrm{tr}}}\int_\Omega e^{i(\mathsf O^T\xi_0)\cdot y}\,dy=0.
\]
Taking real parts provides the nonzero continuous real function
$\cos(\xi_0\cdot x)$ whose integrals over all rigid copies of
$\Omega$ vanish. This is the failure of the Pompeiu property.
\end{proof}

\begin{corollary}[The bounded base construction]\label{inf:bounded-basic}
For every sufficiently large $p\in\mathcal P_{\mathrm{bd}}$, some
integer $a\in[p/2,3p/4]$ gives a bounded strictly convex analytic
non-ball domain with $O(a)\times O(2p-a)$ symmetry and the data
\eqref{inf:main-equations}. The construction uses
$0<\tau\le20$ and the cofactor order $N=4M+40$.
\end{corollary}
\begin{proof}
Lemma~\ref{inf:shifted-arithmetic} supplies unbounded inputs of both
parities; Lemma~\ref{inf:detuning} supplies their radii and
$0<\tau<20$. Theorem~\ref{inf:integer-gap} selects an integer
split for every sufficiently large member. Fix the order of
constants \eqref{inf:constant-order} and the resulting one finite
accuracy order $N$.
The centre exists at every selected split by
Theorem~\ref{inf:centres}; its gap transfers by
Proposition~\ref{inf:centre-gap-transfer}.
Theorem~\ref{inf:coefficient-inverse},
Theorem~\ref{inf:material-isomorphism}, and the second-derivative
and residual bounds give the exact zero in
Corollary~\ref{inf:Newton-zero}.
Lemmas~\ref{inf:geometry-convexity} and
\ref{inf:geometry-nonball} give its analytic strictly convex
non-ball domain with the required block symmetry.
Proposition~\ref{inf:classical-solution} gives the classical
frequency-one field, and Proposition~\ref{inf:Pompeiu} gives the
Pompeiu assertion.
There are only finitely many thresholds after the constants and
$N$ have been fixed; all are uniform in the split and in the
nonzero detuning range. Their maximum discards only finitely
many members of $\mathcal P_{\mathrm{bd}}$. Its remaining members have
unbounded distinct dimensions $2p$, as required.
\end{proof}

\section{Constants}\label{inf:sec-constants}

After the constants and centre order have been fixed, the proof
uses the maximum of finitely many thresholds. The exponents are
\[
 b_{\mathrm{inv}}=2M+3,\qquad
 B_{\mathrm{Newt}}=2M+32/3,\qquad
 L_{\mathrm{Newt}}=2M+62/3,\qquad N=4M+40.
\]
The constants in the finite-centre remainder, the material inverse,
and the second derivative depend on this fixed $N$.
All estimates use this one fixed order.

The main cost of the coefficient inverse is the cofactor exponent
$5M/3$. Its cluster count is for the fixed containing window
$W_{\mathrm{loc}}=20E_{\mathrm{inv}}$.
The local-count proof is used at angular length
$M D_{\mathrm{sep}}\log p$.
It requires, among other conditions,
\begin{align}
 \tfrac12C p^{-5/6}(M D_{\mathrm{sep}}\log p)^3
 +\frac{2C_{20E_{\mathrm{inv}}}}p
                  M D_{\mathrm{sep}}\log p&<1,
 \label{inf:constant-lattice}\\
 c_t p^{1/6}&>\pi M D_{\mathrm{sep}}\log p
                       +O(E_{\mathrm{inv}}/p).
 \label{inf:constant-turning}
\end{align}
In the turning comparison any fixed $0<c_t<\pi/\sqrt2$ is available
after increasing its threshold, by the proof of Theorem~\ref{inf:local-count}.
The determinant step also requires
$D_{\mathrm{sep}}>5M/3+5$.
The inverse window is fixed by $E_{\mathrm{inv}}\ge64C_{\mathrm{rel}}$.
Its local-count bound is
\[
 M=\left\lceil 2+4\sqrt{C_{20E_{\mathrm{inv}}}/c}\right\rceil,
\]
with the phase-strip constant and curvature lower constant from
Theorem~\ref{inf:local-count}. Its integer $M$ is fixed before
$D_{\mathrm{sep}}$ and $N$, in the order
\eqref{inf:constant-order}.

The final threshold includes the Bessel and arithmetic estimates,
the flow and integer-grid gap, the centre construction and gap
transfer, the relative graph bounds, the local count at angular
length $MD_{\mathrm{sep}}\log p$, the cofactor absorption, the
head--tail bounds, the three Newton gates, and the geometric
smallness conditions. Each threshold is finite for the fixed
constants and order; their maximum is uniform in the admissible
split and in $0<\tau\le20$.

\section{Bounded detuning, interior splits, and residue classes}
\label{inf:sec-widened}

Fix $0<\varepsilon_{\mathrm{sp}}<1/4$ and write
\[
 I_{\mathrm{sp}}=[\varepsilon_{\mathrm{sp}}^2,
                         (1/2-\varepsilon_{\mathrm{sp}})^2],\qquad
 \varepsilon_{\mathrm{sp}}\le r\le1/2-\varepsilon_{\mathrm{sp}}.
\]
Constants in this section may depend on this interval and on a fixed
upper bound $T$ for $\tau$. The coefficient spaces and their angular
radius two remain those of Definition~\ref{inf:spaces}.

\begin{lemma}[Low-state transfer on a fixed wider window]
\label{inf:wide-low-transfer}
Fix $C_{\mathrm{ap}}>0$. The second conclusion of
Proposition~\ref{inf:low-exclusion} holds for
$|p-p_0|\le C_{\mathrm{ap}}/p_0$, after increasing the threshold
depending on $C_{\mathrm{ap}}$. More generally, its proof is uniform
for $C_{\mathrm{ap}}=O(\log p_0)$.
\end{lemma}
\begin{proof}
Suppose a nonseed low degree $\ell$ has a zero $z(p)$ with
$|z(p)^2-R(p)^2|\le32$. Continue it and the Neumann zero along the
interval to $p_0$. Equation~\eqref{inf:order-lipschitz} first gives
$|z(t)-R(t)|\le C(1+C_{\mathrm{ap}})/p_0$ there, so both radii are
comparable to $2p_0$ and both order-to-radius ratios belong to
$[1/3,2/3]$. Put
\[
 f_{\mathrm{ord}}(x)=\arccos x/\sqrt{1-x^2}.
\]
Its derivative is bounded on that compact interval.
Equation~\eqref{inf:order-derivative} therefore gives
\[
 |z'(t)-R'(t)|\le C(\ell+1)/p_0
                         +C(1+C_{\mathrm{ap}})/p_0^2.
\]
Integration, followed by multiplication by $z+R$, bounds the change
in squared-energy offset by
\begin{equation}\label{inf:wide-low-error}
 C C_{\mathrm{ap}}(\ell+1)/p_0
       +C C_{\mathrm{ap}}(1+C_{\mathrm{ap}})/p_0^2.
\end{equation}
For $\ell=O(\sqrt{p_0})$ this tends to zero in both stated ranges.
The offset at $p_0$ is then less than $64$, contradicting the first
conclusion of Proposition~\ref{inf:low-exclusion}. The other radial
seed states are excluded by their frequency spacing and the same
continuation estimate.
\end{proof}

\begin{lemma}[Extension of the split interval]\label{inf:split-extension}
The reduction, centre, flow, coefficient-space, material, mixed,
localization and cofactor estimates above hold on $I_{\mathrm{sp}}$
with the replacements in the following table. Every constant may
depend on $\varepsilon_{\mathrm{sp}}$; their powers of $p$ and their
order of choice are unchanged. The weighted reduction, centre and
flow statements allow real $r$ on this interval. The Euclidean
coefficient, material, mixed and closing estimates are applied at
integer block sizes $a,b$ with $a+b=2p$.
For the $\tau\le20$ counting estimate
one may use
\[
 K_{\mathrm{sp}}=\frac4{\varepsilon_{\mathrm{sp}}
                            (1-\varepsilon_{\mathrm{sp}})},\qquad
 W_{\mathrm{sp}}=40K_{\mathrm{sp}}.
\]
\begingroup\small\normalfont
\renewcommand{\arraystretch}{1.13}
\begin{longtable}{@{}>{\raggedright\arraybackslash}p{.29\textwidth}>{\raggedright\arraybackslash}p{.29\textwidth}>{\raggedright\arraybackslash}p{.35\textwidth}@{}}
\toprule Statement & Where the interval enters & Replacement\\\midrule
\endfirsthead
\toprule Statement & Where the interval enters & Replacement\\\midrule
\endhead
\ref{inf:reduction}, \ref{inf:angular-sup}, \ref{inf:flow-form} &
Positive beta parameters, endpoint orientation, axis cutoffs &
$\ab,\bb\ge\varepsilon_{\mathrm{sp}}p>1$ and $\bb\ge\ab$;
integer blocks for the Euclidean identification.\\
\ref{inf:quartic-normalization} & $pg_2>9/16$,
$\nrm{\mathsf H_2}_\infty<9/16$, $\nrm{T_\varsigma}_\infty<1$ &
$pg_2\ge2\varepsilon_{\mathrm{sp}}(1-\varepsilon_{\mathrm{sp}})$,
$pg_2<1$, $\nrm{\mathsf H_2}_\infty\le1$,
$\nrm{T_\varsigma}_\infty\le K_{\mathrm{sp}}$.\\
\ref{inf:angular-weighted} & Derivatives of $\sqrt\varsigma$,
lower bound on $\mathfrak D_j$ &
$\varsigma\ge\varepsilon_{\mathrm{sp}}^2$,
$\mathfrak D_j\ge c_{\mathrm{sp}}(p+j)^2$;
constants $C_{\mathrm{sp}}$ in the same weighted estimates.\\
\ref{inf:angular-positive} & Adjacent coefficient ratio;
the upper bound $g_2<1/p$ &
$\ab,\bb\ge16$ suffices for the same ratio bounds;
$\mathcal N_j,\mathcal E_j$ and their certificates are
independent of $\varsigma$.\\
\ref{inf:centres}, \ref{inf:centre-coefficient-defect} &
Compact parameter set and divisors $r(1-r)$ &
Use $r\in[\varepsilon_{\mathrm{sp}},1/2-\varepsilon_{\mathrm{sp}}]$;
the same finite recursion and each prescribed complex disc.\\
\ref{inf:flow-relative}--\ref{inf:flow-theorem} &
Quartic multiplier bounds and positive angular margin &
$C_{\mathrm{sp}}$ replaces numerical interval constants;
then choose $W_{\mathrm{flow}}$ and the threshold.\\
\ref{inf:quartic-count-window} & Supremum of $T_\varsigma$ &
Shell width $3K_{\mathrm{sp}}/p^2$ and window
$W_{\mathrm{sp}}$ for sufficiently large $p$.\\
\ref{inf:grid-avoidance}, \ref{inf:integer-gap} &
Grid size and spacing &
$\#\{a:2\varepsilon_{\mathrm{sp}}p\le a\le
(1-2\varepsilon_{\mathrm{sp}})p\}$ and spacing
at least $\varepsilon_{\mathrm{sp}}/p$.\\
\eqref{inf:gasper-hypothesis}, Section~\ref{inf:sec-spaces} &
Gasper margin for $\bb-\ab\le p/2$ &
Set $u=p-1$, $v=\bb-\ab\le u$; the lower margin is
$18u^3+63u^2+45u>0$.\\
\ref{inf:uniform-centre-size}, \ref{inf:relative-graph} &
The bounds for $pg_2$ and $\mathsf H_2$ &
$C_0,C_{\mathrm{rel}}$ depend on the split interval;
choose $E_{\mathrm{inv}}$, then $M$, thereafter.\\
\ref{inf:material-isomorphism}, \ref{inf:geometry-nonball} &
Normal Hessian and the physical seed coefficient &
$\nrm{D_p^{-1}}_{\Ec_1}\le C_{N,\mathrm{sp}}p^{-2}$ and
$\mathsf H_2(1)\ge c_{\mathrm{sp}}>0$;
the seed is $-\theta/(16r(1-r)p)+O_{N,\mathrm{sp}}(\tau p^{-2})$.\\
\bottomrule
\end{longtable}
\endgroup
\end{lemma}
\begin{proof}
The rational identity \eqref{inf:g2-exact} is valid for all these real
parameters. Its leading term is $pg_2\to4r(1-r)$ uniformly on the
compact interval, and its negative $\varsigma$ coefficient gives
$pg_2<1$ as before. The two endpoint values of $\mathsf H_2$ tend
uniformly to $r^2$ and $(1-r)^2$; its negative minimum is
$O(p^{-1})$. This proves the displayed replacement bounds and
$\mathsf H_2(1)\ge c_{\mathrm{sp}}$. The formulas for
$\mathfrak D_j$, $\mathfrak a_j$, and $\mathfrak b_j$ in
Lemma~\ref{inf:angular-weighted} then give every entry and derivative
estimate with $C_{\mathrm{sp}}$.

For positivity, impose $p\varepsilon_{\mathrm{sp}}\ge16$. The two
adjacent-ratio computations in Lemma~\ref{inf:angular-positive} use
only $\ab,\bb\ge16$, $j\ge7$, and $2j+p\ge78$.
Its numerators \eqref{inf:N-angular}, \eqref{inf:E-angular} and
\eqref{inf:half-diagonal} have no $\varsigma$. The last lower bound
uses $g_2<1/p$, which was retained. The same positive variation
therefore holds. In the form proof all series converge when
$C_{\mathrm{sp}}\tau/p^2$ is small; its two Schur reductions then
give the stated flow estimates with the changed constants.

For the counting row, the normalized radius displacement is at most
$10K_{\mathrm{sp}}/R^2\le3K_{\mathrm{sp}}/p^2$ for $\tau\le20$.
The min--max calculation of Lemma~\ref{inf:quartic-count-window}
then gives a containing window $40K_{\mathrm{sp}}$ for all sufficiently
large $p$ depending on $K_{\mathrm{sp}}$. For consecutive integers $a$,
integrating $|d\varsigma/da|=|r-1/2|/p$ gives spacing at least
$\varepsilon_{\mathrm{sp}}/p$. Use this spacing in the clipped-curve count.
For the one-split gap, take $\gamma=c_{\mathrm{sp}}'\tau p^{-5/3}$.
The clipped-curve count is at most
$C_{\mathrm{sp}}p^{2/3}+C_{\mathrm{sp}}c_{\mathrm{sp}}'p$.
Choose $c_{\mathrm{sp}}'>0$ so the second term is less than half
the leading grid cardinality, then increase $p$ to absorb the first.
This gives the $p^{-5/3}$ gap on the extended interval as well.

For the algebra put $u=\ab+\bb-1=p-1$ and $v=\bb-\ab$.
If $u^2-7u-24\le0$, \eqref{inf:gasper-hypothesis} holds directly.
Otherwise its left-minus-right side is minimized on $0\le v\le u$
at $v=u$, where it equals $18u^3+63u^2+45u$.
Positive products, endpoint bounds and real radial stencils follow.
The material and mixed proofs subsequently use these products and
the strong norm of the finite centre; their numerical use of $9/16$
is replaced by the first two rows. The recurrence for that centre
divides only by nonzero units, $-1/8$, $c_j$ and powers of $r(1-r)$.
These divisors stay uniformly away from zero on the compact interval.
Thus all listed estimates follow in their stated order.
\end{proof}

\begin{lemma}[Arbitrary fixed detuning bound]\label{inf:fixed-detuning-extension}
Fix $0<C_{\mathrm{ap}},T<\infty$ and the split interval above.
For half-lattice inputs with $|p-p_0|\le C_{\mathrm{ap}}/p_0$ and
$0<\tau\le T$, the centre, flow, cofactor inverse, and Newton
statements have the parameterized forms below. Fix
$5/3<q_{\mathrm{gap}}<2$; the base case $q_{\mathrm{gap}}=5/3$
is included with its original bounds. In the inverse and closing
rows assume a selected quartic gap $\kappa\tau p^{-q_{\mathrm{gap}}}$,
with fixed $\kappa>0$.
Use $\varsigma\in I_{\mathrm{sp}}$, equivalently
$r\in[\varepsilon_{\mathrm{sp}},1/2-\varepsilon_{\mathrm{sp}}]$;
material and closing statements use integer blocks.
Constants and thresholds may depend on these fixed data.
Choose $W_{\mathrm{flow}}(T)$ above the base lower bound and the
required multiples of $T$ and of the inverse angular margin.
After $C_0,C_{\mathrm{rel}}$ choose $E_{\mathrm{inv}}$, then
$W_{\mathrm{loc}}=TE_{\mathrm{inv}}$ and $M=M_{TE_{\mathrm{inv}}}$,
then separation and the head cutoff, and finally the centre order.
The last column identifies the proof line permitting each substitution.

\begingroup\small\normalfont
\def\UrlBreaks{\do\:\do\-}
\setlength{\tabcolsep}{4pt}
\renewcommand{\arraystretch}{1.16}
\begin{longtable}{@{}>{\raggedright\arraybackslash}p{.23\textwidth}>{\raggedright\arraybackslash}p{.16\textwidth}>{\raggedright\arraybackslash}p{.29\textwidth}>{\raggedright\arraybackslash}p{.25\textwidth}@{}}
\toprule Base label & Parameter replaced & Parameterized bound & Proof line permitting replacement\\\midrule
\endfirsthead
\toprule Base label & Parameter replaced & Parameterized bound & Proof line permitting replacement\\\midrule
\endhead
\ref{inf:detuning}, \nolinkurl{inf:detuning};
\ref{inf:low-exclusion}, \nolinkurl{inf:low-exclusion}
& Scaled distance $31/10$, base detuning bound
& Use $|p-p_0|\le C_{\mathrm{ap}}/p_0$ and
$0<\tau\le T$; the radius and detuning comparison are
\eqref{inf:log-detuning-relation}. Low exclusion retains its
original energy bounds.
& The exact quotient \eqref{inf:trace-ratio-exact} and the real
zero derivative comparison hold on this fixed window;
\ref{inf:wide-low-transfer} gives the continuation error
\eqref{inf:wide-low-error}.\\
\ref{inf:centres}, \nolinkurl{inf:centres};
\ref{inf:centre-coefficient-defect}, \nolinkurl{inf:centre-coefficient-defect}
& $|\theta|\le20$, base split interval
& \eqref{inf:centre-complex-defect} and
\eqref{inf:centre-coeff-defect-bound} are
$C_{N,T,\mathrm{sp}}\theta^2p^{-N}$;
\eqref{inf:centre-profile} has physical error
$C_{N,T,\mathrm{sp}}\tau p^{-3}$.
& In \eqref{inf:rho-equation} the $Z$ derivative at $(\eps,Z)=(0,4)$
is $8$ for every $\theta$; the recursion divides by
units, $-1/8$, $c_j$ and $r(1-r)$.\\
\ref{inf:flow-relative}, \nolinkurl{inf:flow-relative};
\ref{inf:flow-complement}, \nolinkurl{inf:flow-complement}
& $\tau\le20$, absolute flow window
& Relative bounds $C_{\mathrm{sp}}\tau p^{-2}$ and
$C_{\mathrm{sp}}\tau^2p^{-3}$;
\eqref{inf:flow-schur-remainder} is
$C_{\mathrm{sp}}\tau^2(W_{\mathrm{flow}}(T)^{-1}+p^{-1})$.
& The series is $\sum h^6(C_{\mathrm{sp}}\tau/p^2)^h$;
the complement ratio is $C_{\mathrm{sp}}\tau/W_{\mathrm{flow}}(T)$.
The derivative adds $C_{\mathrm{sp}}\tau^3/W_{\mathrm{flow}}(T)^2$.\\
\ref{inf:flow-low-block}, \nolinkurl{inf:flow-low-block};
\ref{inf:ordered-eigenvalues}, \nolinkurl{inf:ordered-eigenvalues};
\ref{inf:flow-theorem}, \nolinkurl{inf:flow-theorem}
& $\tau\le20$, $\varsigma_0\in(1/64,1/16)$
& Use $\varsigma_0\in\operatorname{int}I_{\mathrm{sp}}$;
the low inverse is $C_{T,\mathrm{sp}}/(p\tau)$.
For $|\lambda|<\tau/(16p)$ the signed slope is
at least $c_{T,\mathrm{sp}}\tau/p$.
& The proof of \ref{inf:flow-low-block} adds the first-order term
to give seed diagonal $\theta/(2p)+O_T(\theta^2p^{-3})$;
\eqref{inf:low-offsets} keeps the index-five value
$-72+O_T(p^{-1})$. Final norm recovery in
\ref{inf:flow-theorem} costs $C_{T,\mathrm{sp}}\sqrt p$.\\
\ref{inf:quartic-count-window}, \nolinkurl{inf:quartic-count-window}
& $|\theta|\le20$, $W_{\mathrm{count}}=40$
& Shell width $C_{\mathrm{sp}}\tau/p^2$ and containing
count window $C_{\mathrm{sp}}(1+T)$.
& \eqref{inf:quartic-domain} gives displacement
$\tau\nrm{T_\varsigma}_\infty/(2R^2)$;
indexwise shell bounds multiply $O(p^2)$ eigenvalues by
$1+O_{\mathrm{sp}}(\tau/p^2)$.\\
\ref{inf:grid-avoidance}, \nolinkurl{inf:grid-avoidance};
\ref{inf:integer-gap}, \nolinkurl{inf:integer-gap}
& Base split grid and $\tau\le20$
& The interior grid has $\asymp_{\mathrm{sp}}p$ points and spacing
at least $\varepsilon_{\mathrm{sp}}/p$; one split has gap
$c_{T,\mathrm{sp}}\tau p^{-5/3}$.
& The clipped-curve count is
$C_{T,\mathrm{sp}}p^{2/3}(1+C_{T,\mathrm{sp}}p\gamma/\mu)$,
with $\mu=c_{T,\mathrm{sp}}\tau/p$;
choose the gap constant before the threshold.\\
\ref{inf:uniform-centre-size}, \nolinkurl{inf:uniform-centre-size};
\ref{inf:relative-graph}, \nolinkurl{inf:relative-graph}
& Centre-size and relative constants on the base compact set
& $\delta\le C_{0,\mathrm{sp}}\tau p^{-2}$;
\eqref{inf:relative-scalar} and \eqref{inf:relative-two-norms}
are $C_{\mathrm{rel},\mathrm{sp}}\tau p^{-2}$.
These constants precede $N$.
& The leading quartic size is $C_{\mathrm{sp}}\tau p^{-2}$;
the finite-order remainder $C_{N,T,\mathrm{sp}}\tau p^{-4}$
is absorbed by the threshold. The scalar graph line is
$C(\delta+p\delta^2)$.\\
\ref{inf:centre-source-norm}, \nolinkurl{inf:centre-source-norm};
\ref{inf:residual-estimates}, \nolinkurl{inf:residual-estimates}
& Fixed-centre constants for $\tau\le20$
& \eqref{inf:centre-norm-bounds} retains powers $p^{k+1}$ and
$p^{k+5}$ with $C_{N,k,T,\mathrm{sp}}$.
Residual and material defect bounds are
$C_{N,T,\mathrm{sp}}\theta^2p^{2-N}$ and
$C_{N,T,\mathrm{sp}}\theta^2p^{3-N}$.
& The known-mode head uses the finite factor $e^{C_N\sqrt T}$;
the factorial radial tail stays geometric. The two-layer lift
and \eqref{inf:Delta-moment} give the residual powers.\\
\ref{inf:material-isomorphism}, \nolinkurl{inf:material-isomorphism};
\eqref{inf:p6-second}, \nolinkurl{inf:p6-second}
& Fixed-centre constants for $\tau\le20$
& $\nrm{\Qmat^{-1}}\le C_{N,T,\mathrm{sp}}p^6$ and
$\sup\nrm{D^2\Fnl}\le C_{N,T,\mathrm{sp}}p^6$
on \eqref{inf:nonlinear-working-region}.
& \eqref{inf:material-reconstruction} uses
$\nrm{D_p^{-1}}\le C_{N,T,\mathrm{sp}}p^{-2}$ and
$\nrm{\Delta T_\eta}\le C_{N,T,\mathrm{sp}}p^8\nrm\eta$;
\eqref{inf:nonlinear-second} costs $Cp(1+\nrm g)$.\\
\eqref{inf:inverse-window}, \nolinkurl{inf:inverse-window};
\ref{inf:inverse-complement}, \nolinkurl{inf:inverse-complement};
\ref{inf:inverse-Schur}, \nolinkurl{inf:inverse-Schur}
& Containing window $20E_{\mathrm{inv}}$
& The projection stays
$\one_{[-E_{\mathrm{inv}}\tau,E_{\mathrm{inv}}\tau]}(H_0-R^2)$;
its fixed superset has width $TE_{\mathrm{inv}}$.
The Schur norm is $C_{E_{\mathrm{inv}},T,\mathrm{sp}}\tau$.
& \eqref{inf:ball-complement-bound} and
\eqref{inf:relative-two-norms} give complementary absorption;
the four feedback factors give $C\tau/E_{\mathrm{inv}}$.\\
\ref{inf:cluster-inverse}, \nolinkurl{inf:cluster-inverse}
& $\tau\le20$, gap exponent
& For every $\tau>0$, gap $c\tau p^{-q_{\mathrm{gap}}}$ and
$\sigma D>q_{\mathrm{gap}}M+q_{\mathrm{gap}}+2$ give
$\nrm{F^{-1}}_1\le C\tau^{-1}p^{q_{\mathrm{gap}}M}$.
& The cofactor line is $m!(C\tau)^{m-1}(2/\gamma_F)^m$;
$\tau$ cancels in both error absorptions.\\
\ref{inf:effective-inverse}, \nolinkurl{inf:effective-inverse}
& $M_{20E_{\mathrm{inv}}}$, $5M/3$, $D_{\mathrm{sep}}>5M/3+5$
& $M=M_{TE_{\mathrm{inv}}}$,
$\nrm{\Feff^{-1}}_1\le C\tau^{-1}p^{q_{\mathrm{gap}}M}$;
choose $\sigma D_{\mathrm{sep}}>q_{\mathrm{gap}}M+q_{\mathrm{gap}}+2$.
& \ref{inf:clusters} applies to the fixed superset;
deleting inter-cluster entries costs
$C\tau p^{-\sigma D_{\mathrm{sep}}}$, followed by the cofactor row.\\
\ref{inf:mixed-inverse}, \nolinkurl{inf:mixed-inverse}
& Exponent $2+5M/3$
& \eqref{inf:mixed-inverse-bound} has right side
$C\tau^{-1}p^{2+q_{\mathrm{gap}}M}$.
& The block solution costs $C/E_{\mathrm{inv}}$ off diagonal,
$C\tau$ into the window and $Cp^2$ for $K_0\Pi$;
density graph multipliers are bounded.\\
\ref{inf:coefficient-inverse}, \nolinkurl{inf:coefficient-inverse}
& $\tau\le20$, exponent $5/2+5M/3$
& $\nrm{J_D^{-1}}_\Cc\le C\tau^{-1}p^{5/2+q_{\mathrm{gap}}M}
\le Cp^{2M+3}/\gamma_p$, with
$\gamma_p=c_{\mathrm{gap}}\tau p^{-q_{\mathrm{gap}}}$.
& \eqref{inf:head-conversion} costs $\sqrt p$;
\eqref{inf:relative-scalar} gives \eqref{inf:head-tail-hypotheses}.
Compare $5/2+q_{\mathrm{gap}}(M-1)\le2M+3$.\\
\ref{inf:centre-gap-transfer}, \nolinkurl{inf:centre-gap-transfer}
& Gap exponent $5/3$, $\tau\le20$
& Quartic gap $c_{\mathrm q}\tau p^{-q_{\mathrm{gap}}}$
gives centre gap $(c_{\mathrm q}/2)\tau p^{-q_{\mathrm{gap}}}$.
& The min--max displacement is $C_{N,T,\mathrm{sp}}\tau p^{-2}$;
require $p^{2-q_{\mathrm{gap}}}>2C_{N,T,\mathrm{sp}}/c_{\mathrm q}$.\\
\ref{inf:newton-threshold}, \nolinkurl{inf:newton-threshold};
\ref{inf:Newton-zero}, \nolinkurl{inf:Newton-zero}
& $\tau\le20$, gap exponent $5/3$
& Inverse hypothesis
$C_Np^{b_{\mathrm{inv}}}/(c_{\mathrm{gap}}\tau p^{-q_{\mathrm{gap}}})$;
$B=b_{\mathrm{inv}}+m_{\mathrm{mat}}+q_{\mathrm{gap}}$,
$L=B+d_{\mathrm{nl}}+4$, $N\ge\lceil2B+d_{\mathrm{nl}}+8\rceil$.
& \eqref{inf:newton-gate-one}--\eqref{inf:newton-gate-three} give
$C_Np^{B+L+2-N}$, $C_N\tau p^{B+3-N}$ and
$C_Np^{B+d_{\mathrm{nl}}-L}$; only the second uses $\tau\le T$.\\
\ref{inf:geometry-convexity}, \nolinkurl{inf:geometry-convexity};
\ref{inf:geometry-nonball}, \nolinkurl{inf:geometry-nonball}
& Fixed-centre comparison constants
& Normalized $C^2$ size
$C_{N,T,\mathrm{sp}}\tau p^{-2}+O(\tau p^{-L})$;
physical seed $-\theta/(16r(1-r)p)$ with error
$C_{N,T,\mathrm{sp}}\tau p^{-2}+O(\tau p^{1-L})$.
& \eqref{inf:curvature-numerator} applies when the shape tends
to zero in $C^2$; dividing the witness error by $\tau/p$
makes it tend to zero.\\
\bottomrule
\end{longtable}
\endgroup
In the Newton row take $b_{\mathrm{inv}}=2M+3$ and
$m_{\mathrm{mat}}=d_{\mathrm{nl}}=6$. Thus
$B=B_{\mathrm{bd}}=2M+9+q_{\mathrm{gap}}$,
$L=L_{\mathrm{bd}}=B_{\mathrm{bd}}+10$ and $N=4M+40$;
the three powers are those in \eqref{inf:bounded-gates}.
Section~\ref{inf:sec-constants} records the special case
$T=20$, $q_{\mathrm{gap}}=5/3$.
These fixed-detuning families use the cofactor estimate
with fixed local cluster cardinality.
\end{lemma}
\begin{proof}
Lemma~\ref{inf:wide-low-transfer} supplies the low-state exclusion.
The real root comparison and the exact ratio
\eqref{inf:trace-ratio-exact} give $R=2p+3+O_{C_{\mathrm{ap}}}(p^{-1})$.
For the centre choose a fixed complex $\theta$ disc containing
$[-T,T]$ before its $\eps$ disc. Its radius equation has derivative
eight at $(\eps,Z)=(0,4)$ for every $\theta$ on this fixed disc.
The finite recursion has the same nonzero leading divisors, and its
two defects have the same analytic factor $\theta^2$. Compactness
therefore proves all the stated fixed-order centre bounds with
$C_{N,T,\mathrm{sp}}$.

In Lemma~\ref{inf:flow-relative} the dominating series is
$\sum h^6(C_{\mathrm{sp}}\tau/p^2)^h$. Choose a fixed
$W_{\mathrm{flow}}$ larger than the required multiples of $T$ and
the inverse angular margin. The complementary series ratio is
$C_{\mathrm{sp}}\tau/W_{\mathrm{flow}}$, and its derivative remainder
is bounded by $C_{\mathrm{sp}}\tau^2(W_{\mathrm{flow}}^{-1}+p^{-1})$.
The extra term $C\tau^3/W_{\mathrm{flow}}^2$ is absorbed by this
bound. The seed diagonal is still
$\theta/(2p)+O_T(\theta^2p^{-3})$, the index-five diagonal is
$-72+O_T(p^{-1})$, and the weighted low inverse is $C_T/(p\tau)$.
Its two exterior remainders retain their factors $\tau^2$.
Choose the window and then $p$ so the errors are below the positive
angular margin. The last step of Theorem~\ref{inf:flow-theorem}
therefore gives $c_{T,\mathrm{sp}}\tau/p$ in its same narrow band.
The shell count uses a fixed window $C_{\mathrm{sp}}(1+T)$.

The leading quartic strong norm is bounded by
$C_{\mathrm{sp}}\tau p^{-2}$ independently of the centre order.
For each fixed $N$ its remainder is
$C_{N,T,\mathrm{sp}}\tau p^{-4}$. Increasing only the threshold
absorbs this remainder into the leading bound, as in
Lemma~\ref{inf:uniform-centre-size}. Thus $C_0,C_{\mathrm{rel}}$
are fixed before $N$. For the coefficient inverse choose
$E_{\mathrm{inv}}$ from this relative graph constant. The
containing window is now $TE_{\mathrm{inv}}$, so
$M=M_{TE_{\mathrm{inv}}}$ is fixed. The determinant proof of
Lemma~\ref{inf:cluster-inverse} retains exactly one power $\tau^{-1}$.
Choose separation and head cutoff after $M$, and the finite centre
order after those choices. The material and residual proofs use
only the fixed-centre norms and the unchanged clamped estimates.
Their constants may depend on $T$, while their exponents remain six.
The three Newton gate estimates and the geometric comparisons
retain their powers of $\tau$; replacing its upper bound twenty by
$T$ changes only the threshold.
\end{proof}

\begin{lemma}[All but a sublinear number of interior splits]
\label{inf:many-splits}
Fix $C_{\mathrm{ap}},T$, $0<\eta_{\mathrm{cnt}}<1/3$, and
$q_{\mathrm{gap}}=2-\eta_{\mathrm{cnt}}$. At every sufficiently large input of
Lemma~\ref{inf:fixed-detuning-extension}, all but
$O_{C_{\mathrm{ap}},T,\mathrm{sp},\eta_{\mathrm{cnt}}}
 (p^{2/3+\eta_{\mathrm{cnt}}})$ integers
\[
 2\varepsilon_{\mathrm{sp}}p\le a\le(1-2\varepsilon_{\mathrm{sp}})p
\]
give an actual quartic invariant gap at least
$\kappa\tau p^{-q_{\mathrm{gap}}}$, where $\kappa>0$ depends only
on the fixed data.
\end{lemma}
\begin{proof}
The enlarged shell window gives at most $C_{T,\mathrm{sp}}p^{2/3}$
possible ordered curves. Their clipped signed functions are
nondecreasing by the absolute-continuity argument of
Lemma~\ref{inf:grid-avoidance}, with slope at least
$\mu=c_{T,\mathrm{sp}}\tau/p$ inside the narrow flow band.
Since the grid spacing is at least $\varepsilon_{\mathrm{sp}}/p$,
one curve excludes at most
\[
 1+\frac{2p\kappa\tau p^{-q_{\mathrm{gap}}}}
          {\varepsilon_{\mathrm{sp}}\mu}
 \le 1+C p^{2-q_{\mathrm{gap}}}
\]
points. The gap lies inside the flow band eventually because
$q_{\mathrm{gap}}>1$. Summing the exclusions gives
$C(p^{2/3}+p^{8/3-q_{\mathrm{gap}}})\le C'p^{2/3+\eta_{\mathrm{cnt}}}$.
One can exclude twice the stated radius first, which also covers
equality at an endpoint. This proof uses $2p\in\Z$ only to obtain
integer second block sizes.
\end{proof}

\begin{proposition}[Cofactor closure]
\label{inf:bounded-assembly}
Under the hypotheses of Lemma~\ref{inf:many-splits}, every retained
split carries an analytic strictly convex non-ball domain and a
solution of \eqref{inf:main-equations} with its prescribed block symmetry.
The threshold in $p$ follows these choices.
First fix $C_{\mathrm{ap}},T,\varepsilon_{\mathrm{sp}},q_{\mathrm{gap}}$.
Next fix $C_0,C_{\mathrm{rel}},E_{\mathrm{inv}}$ and
$M=M_{TE_{\mathrm{inv}}}$, then $D_{\mathrm{sep}}$ and
$K_{\mathrm{head}}$. Finally set
\[
 b_{\mathrm{inv}}=2M+3,\qquad N=4M+40,\qquad
 B_{\mathrm{bd}}=2M+9+q_{\mathrm{gap}},\qquad
 L_{\mathrm{bd}}=B_{\mathrm{bd}}+10.
\]
The closed Newton ball has radius $\tau p^{-L_{\mathrm{bd}}}$.
\end{proposition}
\begin{proof}
Fix the data in the order stated, choosing separation with
$\sigma D_{\mathrm{sep}}>q_{\mathrm{gap}}M+q_{\mathrm{gap}}+2$
and the head cutoff required by Lemma~\ref{inf:head-tail}.
Use the parameterized statements of
Lemma~\ref{inf:fixed-detuning-extension} throughout this chain.
The centre-to-quartic normalized displacement is
$C_{N,T,\mathrm{sp}}\tau p^{-4}$. The indexwise min--max proof of
Proposition~\ref{inf:centre-gap-transfer} changes energies by
$C_{N,T,\mathrm{sp}}\tau p^{-2}=o(\tau p^{-q_{\mathrm{gap}}})$.
Thus the centre has half the selected gap.

Use Lemma~\ref{inf:cluster-inverse} with this gap exponent.
The Schur reconstruction and radial conversion in
Section~\ref{inf:sec-inverse} give
\[
 \nrm{J_D^{-1}}\le C\tau^{-1}p^{5/2+q_{\mathrm{gap}}M}
 \le C p^{2M+3}/(\tau p^{-q_{\mathrm{gap}}}),
\]
since $5/2+q_{\mathrm{gap}}(M-1)\le2M+3$.
The separation inequality makes the off-cluster restoration small.
The scalar estimate \eqref{inf:relative-scalar} supplies
\eqref{inf:head-tail-hypotheses} for fixed $T$ and sufficiently
large $p$, as in the proof of Theorem~\ref{inf:coefficient-inverse}.
The material inverse and nonlinear bound still have exponent six,
so $\AN=\Qmat^{-1}J_D^{-1}$ has norm at most
$C_N\tau^{-1}p^{B_{\mathrm{bd}}}$.

Apply the contraction proof of Theorem~\ref{inf:newton-threshold}
with these explicit exponents. The three gate powers are
\begin{equation}\label{inf:bounded-gates}
 \begin{split}
 B_{\mathrm{bd}}+L_{\mathrm{bd}}+2-N&=2q_{\mathrm{gap}}-10,\\
 B_{\mathrm{bd}}+3-N&=-2M-28+q_{\mathrm{gap}},\\
 B_{\mathrm{bd}}+6-L_{\mathrm{bd}}&=-4.
 \end{split}
\end{equation}
The second carries the additional factor $\tau\le T$; all three
tend to zero uniformly in $0<\tau\le T$. The working ball lies in
\eqref{inf:nonlinear-working-region} for large $p$. The actual inverse
identities in the same spaces give the exact real zero. The geometric
and classical-solution arguments of Section~\ref{inf:sec-geometry}
use the replacement constants in Lemma~\ref{inf:split-extension}.
Its normalized $C^2$ size is
$O_{N,T,\mathrm{sp}}(\tau p^{-2})+O(\tau p^{-L_{\mathrm{bd}}})\to0$,
so Lemma~\ref{inf:geometry-convexity} gives strict convexity.
The leading physical quartic coefficient has size comparable to
$\tau/p$; both the centre remainder and the physical Newton
correction are smaller after division by that quantity.
This proves non-sphericity as well as strict convexity, and completes
the cofactor chain uniformly for $0<\tau\le T$.
\end{proof}

\begin{lemma}[Orthogonal symmetry and similarity classes]
\label{inf:similarity-classes}
For a differentiable nonconstant radial graph about zero invariant
under $O(a)\times O(n-a)$, its orthogonal symmetry Lie algebra is
$\mathfrak{so}(a)\oplus\mathfrak{so}(n-a)$.
Consequently constructed domains at distinct splits below $n/2$
are pairwise non-similar. For a dimension family, define
$\mathcal N_{\mathrm{sim}}(n)$ as the number of integers $0<a<n/2$
for which the construction gives a domain on some fixed compact
interior split interval, with its associated fixed centre order.
Choosing any one such domain at each split gives this many
similarity classes, independently of the choices.
Suppose that, for every fixed $0<\varepsilon_{\mathrm{sp}}<1/4$
and $0<\eta_{\mathrm{cnt}}<1/3$, all but
$O(n^{2/3+\eta_{\mathrm{cnt}}})$ integer splits in
$[\varepsilon_{\mathrm{sp}}n,(1/2-\varepsilon_{\mathrm{sp}})n]$
give domains. Then along that family's dimensions,
\[
 \mathcal N_{\mathrm{sim}}(n)\ge
 (1/2-2\varepsilon_{\mathrm{sp}})n-Cn^{2/3+\eta_{\mathrm{cnt}}},\qquad
 \liminf_{n\to\infty}\frac{\mathcal N_{\mathrm{sim}}(n)}n\ge\frac12.
\]
The limit inferior is taken through that family's dimensions.
A countable exhaustion by fixed compact split intervals suffices
in the definition of $\mathcal N_{\mathrm{sim}}$.
\end{lemma}
\begin{proof}
Write the radius on the sphere as $F(x)$, $x=|X|^2$.
There is $x_0\in(0,1)$ with $F'(x_0)\ne0$.
A skew matrix has diagonal blocks in the two block Lie algebras and
off-diagonal blocks $M,-M^t$. Its infinitesimal action on $F$ is
$2F'(x)X^tMY$. Vanishing for all $X,Y$ with squared norms
$x_0,1-x_0$ forces $M=0$, by varying their directions. The block
diagonal matrices are symmetries, proving the equality.

All constructed domains have barycentre zero. A similarity between
two has zero translation and conjugates their orthogonal symmetry
Lie algebras. The dimension of the latter is
$\{a(a-1)+(n-a)(n-a-1)\}/2$, whose change from $a$ to $a+1$ is
$2a+1-n<0$ on the selected range. Thus different splits give different
similarity classes. The assumed exceptional-split bound and integer
counting give the displayed lower bound. Fixing any $\varepsilon_{\mathrm{sp}}>0$
and then taking the lower limit proves a bound $1/2-2\varepsilon_{\mathrm{sp}}$.
Taking the supremum over these fixed choices proves the last assertion;
their thresholds need not be uniform as $\varepsilon_{\mathrm{sp}}\downarrow0$.
\end{proof}

\begin{theorem}[Dimensions in prescribed residue classes]
\label{inf:residue-arithmetic}
For every positive integer $L_{\mathrm{ap}}$ and every integer
$b_{\mathrm{ap}}$, there is an unbounded selection of dimensions
$n\equiv b_{\mathrm{ap}}\pmod{4L_{\mathrm{ap}}}$ with
\[
 |n/2-\widehat p_m|<32L_{\mathrm{ap}}^2/\widehat p_m,
 \qquad 0<\tau<256L_{\mathrm{ap}}^2
\]
eventually. Every sufficiently large selected dimension has the
domains and multiplicity just proved. The threshold may depend on
the fixed modulus, residue, and split/counting parameters.
In particular every residue modulo $L_{\mathrm{ap}}$ occurs
infinitely often. For even $L_{\mathrm{ap}}$ the residue determines
the parity; for odd $L_{\mathrm{ap}}$ both parities occur infinitely
often in each residue class modulo $L_{\mathrm{ap}}$.
\end{theorem}
\begin{proof}
Set $M_{\mathrm{ap}}=4L_{\mathrm{ap}}$. Consecutive convergent vectors
of $\alpha_{\mathrm{ar}}$ form an invertible matrix modulo
$M_{\mathrm{ap}}$. Choose representatives
$A,B\in\{0,\ldots,M_{\mathrm{ap}}-1\}$ satisfying
\[
 P=A P_k^{\mathrm{cf}}+B P_{k-1}^{\mathrm{cf}}
       \equiv b_{\mathrm{ap}}+3\pmod{M_{\mathrm{ap}}},\qquad
 Q=A Q_k^{\mathrm{cf}}+B Q_{k-1}^{\mathrm{cf}}
       \equiv1\pmod{M_{\mathrm{ap}}}.
\]
They cannot both vanish, so $Q\ge Q_{k-1}^{\mathrm{cf}}>0$.
In \eqref{inf:cf-error}, for $T_k^{\mathrm{cf}}\ge A/B$ the expression is at
most $B(A+B)$; when $1<T_k^{\mathrm{cf}}<A/B$ its maximum is attained at
$T_k^{\mathrm{cf}}=1$ and then is at most
$|A-B|\max\{A,(A+B)/2\}$. The cases $A=0$ or $B=0$ follow
directly. Both bounds are at most $2(M_{\mathrm{ap}}-1)^2$.
Set $m=(Q-1)/4$, $n=P-3$, and $p=n/2$. The crossing expansion gives
\[
 \limsup\widehat p_m|p-\widehat p_m|
 \le\frac{\alpha_{\mathrm{ar}}}{2}(M_{\mathrm{ap}}-1)^2
       +d_{\mathrm{ar}}<32L_{\mathrm{ap}}^2.
\]
For the last strict inequality use
$\alpha_{\mathrm{ar}}<363/157$ and $d_{\mathrm{ar}}<1$.
The dimensions tend to infinity and have the prescribed residue.
Equation~\eqref{inf:log-detuning-relation}, with
$C_{\mathrm{ap}}=32L_{\mathrm{ap}}^2$, gives eventually
$\tau<(16\delta_0+o(1))32L_{\mathrm{ap}}^2<256L_{\mathrm{ap}}^2$;
non-coincidence gives $\tau>0$. Apply
Lemma~\ref{inf:fixed-detuning-extension} and
Proposition~\ref{inf:bounded-assembly} with these fixed constants.
Taking $\varepsilon_{\mathrm{sp}}=1/8$, the sublinear exclusions
leave a split in $r\in[1/4,3/8]$ as well.
For a residue modulo an odd $L_{\mathrm{ap}}$, its lifts
$b_{\mathrm{ap}}$ and $b_{\mathrm{ap}}+L_{\mathrm{ap}}$ have opposite
parities, giving the final assertion.
\end{proof}

\begin{corollary}[Almost every fixed split ratio]\label{inf:fixed-ratio}
Use the explicit continued-fraction selections in
Definition~\ref{inf:integer-sequence}, Lemma~\ref{inf:shifted-arithmetic}, and
Theorem~\ref{inf:residue-arithmetic}, retaining their convergent
index $k$. There is a set $\mathcal R_{\mathrm{split}}\subset(0,1/2)$
of full Lebesgue measure such that, for every
$r_0\in\mathcal R_{\mathrm{split}}$, each of these countably many
selected dimension families $(n_k)$ eventually has a constructed
Schiffer domain with split
\[
 a_k=\lfloor n_k r_0\rfloor.
\]
The starting index may depend on $r_0$ and the selected family.
\end{corollary}
\begin{proof}
Fix one selected family and a compact ratio interval inside
$(0,1/2)$. Choose a slightly larger compact interval for
Lemma~\ref{inf:many-splits}; rounding changes the ratio by at most
$1/n_k$, so it remains in that interval eventually.
Take $\eta_{\mathrm{cnt}}=1/6$. At dimension $n_k$ at most
$C n_k^{5/6}$ splits are excluded, with $C$ depending on this
family's fixed detuning bound and the compact interval.
Each integer split corresponds under $a=\lfloor n_k r\rfloor$
to an interval of $r$ of length $1/n_k$. The set of ratios giving
an excluded split therefore has measure at most $C n_k^{-1/6}$.

The selected denominator is at least $Q^{\mathrm{cf}}_{k-3}$
for the first two selections and at least $Q^{\mathrm{cf}}_{k-1}$
for the residue-class selection. Continued-fraction denominators
satisfy $Q^{\mathrm{cf}}_{k+2}\ge Q^{\mathrm{cf}}_{k+1}+Q^{\mathrm{cf}}_k
\ge2Q^{\mathrm{cf}}_k$. Since the selected dimensions are asymptotic
to $\alpha_{\mathrm{ar}}$ times their selected denominators,
$n_k\ge c\,2^{k/2}$ eventually. Hence $\sum_k n_k^{-1/6}<\infty$.
The measure of the union of the excluded-ratio events after index
$K$ is bounded by the tail of this convergent series and tends
to zero. Almost every ratio has only finitely many exclusions.

Take a countable compact exhaustion of $(0,1/2)$ and then the
intersection of these full-measure sets over the two primary
selections and all fixed moduli and residues. Their exceptional
sets have total measure zero. Proposition~\ref{inf:bounded-assembly}
gives the domain at every eventual retained split, proving the
assertion with the stated order of quantifiers.
\end{proof}

\section{Sharper clamped and material estimates}\label{inf:sec-sharp}

The spaces in this section remain $\Cc$ and $\Xc_p=\Gc_p\times\Ec_1$.
The centre and the fixed material field have the membership supplied
by Section~\ref{inf:sec-material}. The estimates below follow from
their exact equations.
For fixed wider parameter ranges use the extensions in
Lemmas~\ref{inf:split-extension} and \ref{inf:fixed-detuning-extension}.

\begin{lemma}[The harmonic gradient coefficient sum]\label{inf:pole-sum}
For endpoint-normalized solid harmonics $S_i,S_j$,
\begin{equation}\label{inf:pole-sum-identity}
 \nabla S_i\cdot\nabla S_j
 =\sum_{k<i+j}2h(h+2k+p-1)c_{ij}^kS_kt^{h-1},\quad h=i+j-k,
 \qquad
 \sum_{k<i+j}2h(h+2k+p-1)c_{ij}^k=4ij.
\end{equation}
Consequently, for $H_u=\HP h_u$,
\begin{equation}\label{inf:sharp-gradient-multiplier}
 \nrm{M_{\nabla H_1\cdot\nabla H_2}}_{\Cc\to\Cc}
 \le4\nrm{h_1}_{\Ec_1}\nrm{h_2}_{\Ec_1}.
\end{equation}
\end{lemma}
\begin{proof}
The product rule for two harmonic polynomials gives
$2\nabla S_i\cdot\nabla S_j=\Delta(S_iS_j)$.
Apply \eqref{inf:profile-laplacian} to each term of the positive solid
product to obtain the first identity. At a unit point of the $x=1$
block axis, $S_i=S_j=1$ and their tangential gradients vanish.
Their full gradients are $2iy$ and $2jy$. Evaluation therefore gives
the second identity; the coefficients are all nonnegative.

Every multiplier $S_kt^{h-1}$ has total degree $2i+2j-2$ and norm
at most $2^k\wt(2i+2j-2)$ by Proposition~\ref{inf:module}.
This is at most $2^{i+j}\wt(2i)\wt(2j)$.
The positive coefficient sum is $4ij\le4(1+i)(1+j)$.
Sum the absolute shape coefficients with their $\Ec_1$ weights
to get \eqref{inf:sharp-gradient-multiplier}. A constant harmonic
has zero gradient, and completion gives the assertion for all shapes.
\end{proof}

\begin{proposition}[Sharp centre and material bounds]\label{inf:sharp-material}
Under the root and centre hypotheses of
Lemma~\ref{inf:detuning-dependencies}, for each fixed centre order
$N\ge4$, uniformly in
$0<\tau\le T\le\log p$ and on the working region
\eqref{inf:nonlinear-working-region},
\begin{equation}\label{inf:sharp-three-bounds}
 \nrm{g_0}_\Cc\le p^3\mathcal B_N(T),\qquad
 \nrm{\Qmat^{-1}}\le p^2\mathcal B_N(T),\qquad
 \sup\nrm{D^2\Fnl}\le p^3\mathcal B_N(T).
\end{equation}
Here $\mathcal B_N(T)=C_N(1+T)^{d_N^{\mathrm{tame}}}e^{C_N\sqrt T}$
is the envelope of Lemma~\ref{inf:detuning-dependencies}, enlarged
if necessary; $C_N,d_N^{\mathrm{tame}}$ depend on the fixed order
and compact split interval. For fixed $T$ it is a constant.
The inverse identities are those of
Theorem~\ref{inf:material-isomorphism}.
\end{proposition}
\begin{proof}
In the proof of Proposition~\ref{inf:nonlinear-map}, replace
\eqref{inf:gradient-multiplier} by
\eqref{inf:sharp-gradient-multiplier}. Every other grouped estimate
there has a constant independent of $p$.
In particular \eqref{inf:second-shape-clamped} and
\eqref{inf:mixed-shape-clamped} are applied to clamped sources, as
their statements require. The inverse multiplier series and their
first two derivatives are bounded on the working ball. Thus
\begin{equation}\label{inf:sharp-nonlinear}
 \nrm{D_f\Fnl(g,f)}\le C\nrm g_\Cc,
 \qquad \nrm{D^2\Fnl(g,f)}\le C(1+\nrm g_\Cc).
\end{equation}
The frequency term is affine in $g$ and independent of the shape.

Put $\mathcal R_{\mathrm{sc}}=(P_{\Phi_N}\Delta-\Delta)\KD$.
Equation~\eqref{inf:relative-scalar} gives
$\nrm{\mathcal R_{\mathrm{sc}}}\le C\tau p^{-2}<1/2$ at the fixed
centre. Since $U_0=1+Kg_0$, its residual equation is exactly
\[
 E_0=(I+\mathcal R_{\mathrm{sc}})g_0+R^2U_0.
\]
Use $\nrm{g_0}\le2\nrm{(I+\mathcal R_{\mathrm{sc}})g_0}$ and
$\nrm{U_0}\le C_Np$ from \eqref{inf:centre-norm-bounds}.
This yields $\nrm{g_0}\le C_Np^3$ and hence the last bound in
\eqref{inf:sharp-three-bounds} by \eqref{inf:sharp-nonlinear}.
This is norm absorption on the already defined clamped source.

The fixed material field is
$\mathfrak f_0=d_0^{-1}\Euler U_0=d_0^{-1}\Euler Kg_0$.
The all-source estimate for $\Euler\KD$ in
\eqref{inf:poisson-operator-bounds} gives
$\nrm{\mathfrak f_0}\le C_Np^2$. Its zero Dirichlet trace and
$\Delta\mathfrak f_0\in\Cc$, already proved in
\eqref{inf:fixed-material-field}, give
$\mathfrak f_0=\KD\Delta\mathfrak f_0$.
Apply the material identity to the constant shape variation one:
\[
 B_{\Phi_N}\mathfrak f_0
 =d_0^{-1}\Euler E_0-D_f\Fnl(z_0)[1].
\]
The residual bound and \eqref{inf:sharp-nonlinear} bound this by
$C_Np^3$. On the other hand it equals
$(I+\mathcal R_{\mathrm{sc}})\Delta\mathfrak f_0+R^2\mathfrak f_0$.
The same norm absorption proves
\begin{equation}\label{inf:sharp-material-field}
 \nrm{\Delta\mathfrak f_0}\le C_Np^4.
\end{equation}
For an unknown shape $\eta$,
\[
 \Delta T_\eta=(\HP\eta)\Delta\mathfrak f_0
 +2\nabla(\HP\eta)\cdot\nabla\KD\Delta\mathfrak f_0.
\]
The module estimate and the all-source gradient bound
\eqref{inf:gradient-C-one} give
$\nrm{\Delta T_\eta}\le C_Np^4\nrm\eta_{\Ec_1}$.
The exact reconstruction \eqref{inf:material-reconstruction} has
$\nrm\eta_{\Ec_1}\le C_Np^{-2}\nrm\psi_\Cc$, by its normal trace
and the normal Hessian bound. It follows that
$\nrm{\delta g}+\nrm\eta\le C_Np^2\nrm\psi$.
Both inverse identities were proved before estimating this expression,
so this proves the middle bound in \eqref{inf:sharp-three-bounds}.
The dependence on $T$ in these same steps is given by row 12 and
the clamped-source and material calculations in the proof of
Lemma~\ref{inf:detuning-dependencies}, yielding the stated envelope.
\end{proof}

\begin{lemma}[Quadratic graph remainder at a polynomial chart]
\label{inf:quadratic-graph}
Let $H$ be a real invariant harmonic polynomial of fixed angular degree,
with size $\delta$ in a fixed multiplier norm including its first
four Euler derivatives. For $p\delta$ sufficiently small put
$q=1+H$, $d=q+\Euler q$, $m=q^{p-1/2}d^{1/2}$, and denote its
unitary Dirichlet operator by $\widehat H_H$. Its first variation is
\[
 V_1(H)=2H\Delta+2\nabla H\cdot\nabla\Euler
                         +2\nabla l_{\mathrm{lin}}\cdot\nabla,
 \qquad l_{\mathrm{lin}}=pH+\tfrac12\Euler H.
\]
In each of the mixed and Hilbert matrix norms at fixed exponential
weight,
\begin{equation}\label{inf:quadratic-graph-bound}
 \nrm{(\widehat H_H-H_0-V_1(H))K_0^{-1}}\le C\delta^2.
\end{equation}
The constant may depend on the fixed degree and weights. This is
an all-source estimate at a fixed polynomial chart.
\end{lemma}
\begin{proof}
Lemma~\ref{inf:pole-sum} also gives a $p$-independent multiplier
bound for fixed-degree gradient products in $\Mmix$: each positive
solid multiplier is bounded by Lemma~\ref{inf:mixed-multipliers}.
In $\Hc$ use its fixed bandwidth and the pointwise gradient bound
as in \eqref{inf:Hilbert-gradient-product}. These bounds apply to
$H$ and each of its fixed Euler derivatives.
With $l=\log m$, harmonicity of $q,d$ gives
\[
 \Delta l=-(p-\tfrac12)q^{-2}|\nabla H|^2
       -\tfrac12d^{-2}|\nabla(1+\Euler)H|^2.
\]
Thus $\Delta l=O(p\delta^2)$,
$|\nabla l|^2=O(p^2\delta^2)$, and
$\Euler l,\Euler^2l=O(p\delta)$ in the multiplier norms.
Subtracting $l_{\mathrm{lin}}$ in its gradient leaves harmonic
gradients times coefficients of size $O(p\delta)$; after
$\nabla\KD$ the result is $O(\delta^2)$.
The exact conjugations are
\begin{align*}
 m\Delta m^{-1}&=\Delta-2\nabla l\cdot\nabla
                            +|\nabla l|^2-\Delta l,\\
 m(\nabla H\cdot\nabla\Euler)m^{-1}
 &=\nabla H\cdot\nabla\Euler-(\nabla H\cdot\nabla l)\Euler
       -(\Euler l)\nabla H\cdot\nabla\\
 &\quad +(\Euler l)(\nabla H\cdot\nabla l)
                         -\nabla H\cdot\nabla\Euler l,\\
 m\Euler(\Euler+1)m^{-1}
 &=\Euler(\Euler+1)-2(\Euler l)\Euler
             +(\Euler l)^2-\Euler^2l-\Euler l,\\
 m\Euler m^{-1}&=\Euler-\Euler l.
\end{align*}
Insert these in the negative scalar pullback
\eqref{inf:scalar-pullback}. Its linear part is exactly $V_1(H)$;
there is no linear potential because $\Delta l_{\mathrm{lin}}=0$.
The following table records every remainder after composition with
$\KD$; all Poisson and gradient bounds are the all-layer bounds in
Lemma~\ref{inf:mixed-Poisson}.
\begin{center}\small\normalfont
\begin{tabular}{@{}>{\raggedright\arraybackslash}p{.60\textwidth}>{\raggedright\arraybackslash}p{.33\textwidth}@{}}
\toprule Remainder & Bound after $\KD$\\\midrule
Quadratic part of $q^{-2}\Delta$ & $C\delta^2$\\
Nonlinear density-gradient terms & $Cp\delta^2\cdot p^{-1}$\\
Density potential & $Cp^2\delta^2\cdot p^{-2}$\\
Nonlinear coefficient of $\nabla H\cdot\nabla\Euler$
 & $C\delta\cdot\delta$\\
$(\nabla H\cdot\nabla l)\Euler$ and
$(\Euler l)\nabla H\cdot\nabla$ & $Cp\delta^2\cdot p^{-1}$\\
Their zero-order terms & $C(p^2\delta^3+p\delta^2)p^{-2}$\\
$|\nabla H|^2\Euler(\Euler+1)$ & $C\delta^2$\\
Its density corrections & $C(\delta^3+\delta^4+\delta^3/p)$\\
The scalar $\Euler$ coefficient involving $C_q$ & $C\delta^2/p$\\
Its density potential & $Cp\delta^3\cdot p^{-2}$\\
\bottomrule
\end{tabular}
\end{center}
Formula~\eqref{inf:Cq-identity} makes $C_q$ quadratic in the fixed
harmonic gradients and Euler derivatives. The inverse multipliers
have convergent series in both matrix algebras. Every row is bounded
by $C\delta^2$, including the possible $p^2$ potentials. The identities
first hold on the polynomial source core and extend by the graph
bounds. Multiplying the $\KD$ version by the angular-diagonal
contraction $H_0K_0^{-1}$ proves \eqref{inf:quadratic-graph-bound}.
\end{proof}

\section{Phase geometry and mixed conditioning}\label{inf:sec-conditioning}

Fix the compact split interval of Section~\ref{inf:sec-widened}
and an integer split $a,b$ in it, with $a+b=2p\in\mathbb N$.
In this section $R\in[2p-1,2p+4]$ is the continued primary Neumann
zero, $0<\tau\le\log p$, and $|p-p_0|=O(\log p/p)$.
All assertions are for sufficiently large $p$, with constants
independent of $\tau$ in that range. Fix
$1<q_{\mathrm{gap}}<2$ and $0<\xi_{\mathrm{inv}}<2-q_{\mathrm{gap}}$.
Assume that the actual quartic domain has invariant gap at least
\[
 \gamma=\kappa_{\mathrm{gap}}\tau p^{-q_{\mathrm{gap}}},\qquad
 \kappa_{\mathrm{gap}}>0\text{ fixed}.
\]
The low-state structure is that of
Lemma~\ref{inf:wide-low-transfer} and the finite recurrence
\eqref{inf:finite-q-recurrence}.

\subsection{One radial pole at each angular index}

\begin{lemma}[Radial projection and complement]\label{inf:radial-poles}
For $0\le j\le\lfloor2p\rfloor$, choose the positive zero $z_j$
of $J_{\nu_j}$ closest in frequency to $R$, choosing the smaller
root in a tie. Let $\Pi_{\mathrm{rad}}$ retain its unit radial vector
and let $Q_{\mathrm{rad}}=I-\Pi_{\mathrm{rad}}$. These are contractions
in $\Mmix$ and $\Hc$. Put $d_j^{\mathrm{pole}}=z_j^2-R^2$. Then
\begin{equation}\label{inf:radial-pole-bounds}
 \nrm{K_0\{Q_{\mathrm{rad}}(H_0-R^2)Q_{\mathrm{rad}}\}^{-1}
                        Q_{\mathrm{rad}}}\le Cp,
 \qquad \nrm{K_0\Pi_{\mathrm{rad}}}\le Cp^2
\end{equation}
in both spaces and their angular-diagonal matrix norms.
On the retained space the mixed norm is exactly $\ell^1(2^j)$;
the Hilbert metric is $\diag(\eta_j^2)$.
\end{lemma}
\begin{proof}
The projections are radial orthogonal projections in each angular
block. Lemma~\ref{inf:mixed-algebra} proves the contraction statements.
Spacing greater than $\pi$ implies that every nonretained zero in a
retained block is at distance at least $\pi/2$ from $R$.
For $j>2p$ one has $\nu_j>5p-1$, so the centrifugal lower bound
puts its entire spectrum above $(5p-1)^2$. The scalar spectral
ratio $(z^2+R^2)/|z^2-R^2|$ is therefore at most $Cp$ on the
complement.

The radial sine trial on $(1/2,1)$ in
\eqref{inf:radial-schrodinger} bounds the first eigenvalue by
$4\nu_j^2+4\pi^2$. Its first zero is thus at most $2\nu_j+7$.
If $R\le2\nu_j$, comparison with that first zero gives
$z_j\le\max\{2R,2\nu_j+7\}=O(p)$.
If $R>2\nu_j$, the differential equation for
$\sqrt xJ_{\nu_j}(x)$ has frequency square
$1-(\nu_j^2-1/4)/x^2\ge3/4$ for $x\ge R$.
Comparison with a sine on an interval of length
$2\pi/\sqrt3$ gives a zero within a bounded distance beyond $R$.
For a direct proof, put $\kappa_{\mathrm{cmp}}=\sqrt3/2$ and
$w(x)=\sin(\kappa_{\mathrm{cmp}}(x-a))$ on
$[a,a+\pi/\kappa_{\mathrm{cmp}}]$. If a solution $v$ of
$v''+Q(x)v=0$, $Q\ge\kappa_{\mathrm{cmp}}^2$, had no zero there,
choose its sign positive. Then
$(v'w-vw')'=-(Q-\kappa_{\mathrm{cmp}}^2)vw\le0$,
whereas the difference of the endpoint values is
$\kappa_{\mathrm{cmp}}(v(a)+v(a+\pi/\kappa_{\mathrm{cmp}}))>0$.
This contradiction proves the comparison.
Again the closest zero has size $Cp$. Thus its $K_0$ eigenvalue is
at most $Cp^2$. The norm identifications follow from the radial
normalization and \eqref{inf:true-H-norm}.
\end{proof}

\begin{lemma}[Linear quartic reduction]\label{inf:centre-form-comparison}
Let $H_{\mathrm q}=z_{\mathrm q}S_2$ be the harmonic extension of
$f_{\mathrm q}$, so $|z_{\mathrm q}|\le C\tau p^{-2}$, and put
$L_{\mathrm q}=H_0+V_1(H_{\mathrm q})-R^2$.
The radial complementary inverse exists and its exact pole Schur
matrix is
\begin{equation}\label{inf:pole-schur-bands}
 \mathcal S_{\mathrm{pole}}
 =\diag(d_j^{\mathrm{pole}})+\mathcal B_{\mathrm{pole},2}
       +p^{-1}\mathcal B_{\mathrm{pole},4}
       +\mathsf R_{\mathrm{pole},6},
\end{equation}
where the first two bands have bandwidths two and four and
\begin{equation}\label{inf:pole-band-bounds}
 \nrm{\mathcal B_{\mathrm{pole},2}}\le C\tau,\quad
 \nrm{\mathcal B_{\mathrm{pole},4}}\le C\tau^2,\quad
 \nrm{\mathsf R_{\mathrm{pole},6}}\le C\tau^3p^{-2}
\end{equation}
in the coefficient column norm and Hilbert row/column norms.
Define the second-order interaction truncation by
\[
 \mathcal S_{\mathrm{pole}}^{(0)}
 =\diag(d_j^{\mathrm{pole}})+\mathcal B_{\mathrm{pole},2}
                         +p^{-1}\mathcal B_{\mathrm{pole},4}.
\]
It deletes the entire remainder, which contains all paths with
at least three interactions and need not have bandwidth six.
The exact matrix and this bandwidth-four truncation have Hilbert
gap at least $c\gamma$.
\end{lemma}
\begin{proof}
For this degree-four harmonic chart the first variation gives
\[
 V_1(H_{\mathrm q})\KD
 =2z_{\mathrm q}M_{S_2}
 +2z_{\mathrm q}\nabla S_2\cdot\nabla\Euler\KD
 +(2p+4)z_{\mathrm q}\nabla S_2\cdot\nabla\KD.
\]
The exact radial profiles in Lemma~\ref{inf:gradient-profiles}
factor each angular block into $c_{2,j}^i$ times a radial operator
of norm at most $C|z_{\mathrm q}|$.
Lemmas~\ref{inf:mixed-multipliers} and \ref{inf:mixed-Poisson}
bound its relative graph norm by $C\tau p^{-2}$ in both matrix
algebras. By \eqref{inf:radial-pole-bounds}, the radial complementary
Neumann ratio is $C\tau/p$. Direct projection has size $C\tau$;
one complementary inverse gives size $C\tau^2/p$ and bandwidth four.
Each further inverse/perturbation factor contributes $C\tau/p$.
Summing the convergent series proves \eqref{inf:pole-band-bounds},
including all complementary radial states.

For the spectral comparison the exact unitary form is
\[
 \int A_q(\nabla u-u\nabla\log m)\cdot
                  \overline{(\nabla u-u\nabla\log m)}.
\]
Its linear density term integrates to
$\int\Delta l_{\mathrm{lin}}|u|^2=0$.
The linear form has principal coefficient
$(1-2H_{\mathrm q})I-y\otimes\nabla H_{\mathrm q}
-\nabla H_{\mathrm q}\otimes y$ and is uniformly comparable
to the ball form. The remaining form error
is bounded by
$C|z_{\mathrm q}|^2(\nrm{\nabla u}_2^2+p^2\nrm u_2^2)$.
The ball Poincar\'e bound absorbs the second term into the first.
Indexwise min--max therefore changes eigenvalues near $R^2$ by at most
$C\tau^2p^{-2}=o(\gamma)$, since
$\tau p^{q_{\mathrm{gap}}-2}\to0$ uniformly here.
Thus $L_{\mathrm q}$ has at least half the actual quartic gap.
For a pole vector reconstruct the radial complementary component;
the full image is exactly its Schur image and its Hilbert norm
is at least that of the pole vector. This gives the Schur gap as
in Proposition~\ref{inf:inverse-Schur}.
Deleting $\mathsf R_{\mathrm{pole},6}$ changes that gap by
$C\tau^3p^{-2}=o(\gamma)$, proving the last assertion.
\end{proof}

\subsection{The two phase estimates}

Use the corrected phase
\begin{equation}\label{inf:corrected-phase}
 \varphi_{\mathrm{Bes}}(\nu,x)=\Xi(\nu,x)-\arctan b(\nu,x)
\end{equation}
with the functions in \eqref{inf:phase-def}. Chebyshev polynomials
of the second kind are denoted by $U_k^{\mathrm{Ch}}$, to distinguish
them from the Debye polynomials $U_k$.

\begin{lemma}[Two-sided phase and energy comparison]
\label{inf:phase-eigenvalue-two-sided}
Fix $0<c_*<1$. For $p-1\le\nu\le c_*R$, let $z$ be the zero of
$J_\nu$ nearest to $R$. If
$d_{\varphi}=\dist(\varphi_{\mathrm{Bes}}(\nu,R),\pi(\Z+1/2))$,
then for all sufficiently large $p$,
\begin{equation}\label{inf:phase-energy-comparison}
 c p(d_{\varphi}-Cp^{-2})
 \le |z^2-R^2|\le C p(d_{\varphi}+Cp^{-2}).
\end{equation}
The constants depend only on $c_*$.
\end{lemma}
\begin{proof}
Put $c_1=\sqrt{1-c_*^2}/2$ and choose a fixed interval
$|x-R|\le X_*$, $X_*=2\pi/c_1$. There $\kappa\ge cp$.
Lemma~\ref{inf:debye-finite} at order one gives
\[
 \frac{J_\nu(x)}{\sqrt{2/\pi}\kappa^{-1/2}}
 =\sqrt{1+b^2}\cos\varphi_{\mathrm{Bes}}(\nu,x)+e(x),
 \qquad |e(x)|\le C p^{-2}.
\]
The explicit phase derivatives give
$c_1\le\partial_x\varphi_{\mathrm{Bes}}\le1$ after increasing
$p$. At a zero the distance of its phase from $\pi(\Z+1/2)$ is
at most $Cp^{-2}$. Conversely, on either side of a lattice phase at
phase distance $2Cp^{-2}$, enlarge $C$ to dominate the error above.
The two values have opposite signs, so the intermediate value
theorem supplies a zero between them. Every lattice phase whose
distance from both endpoints of the phase interval is greater
than $2Cp^{-2}$ has such a zero inside the argument interval.
Since $\partial_x\varphi_{\mathrm{Bes}}\ge c_1$ and
$X_*=2\pi/c_1$, the phase interval contains
$[\varphi_{\mathrm{Bes}}(\nu,R)-2\pi,
\varphi_{\mathrm{Bes}}(\nu,R)+2\pi]$.
The lattice phase nearest to $\varphi_{\mathrm{Bes}}(\nu,R)$ is
at distance at most $\pi/2$, so it has a fixed positive endpoint
margin and satisfies this condition for all sufficiently large $p$.

Choose the lattice phase nearest to the phase at $R$ and a zero
$z'$ associated to it. The derivative lower bound gives
$|z'-R|\le(d_{\varphi}+Cp^{-2})/c_1$. Thus the nearest frequency
zero $z$ also lies in the chosen interval and obeys that upper bound.
For its lower bound, its phase is within $Cp^{-2}$ of some lattice
phase, whence
$d_{\varphi}\le|\varphi_{\mathrm{Bes}}(\nu,z)
-\varphi_{\mathrm{Bes}}(\nu,R)|+Cp^{-2}\le|z-R|+Cp^{-2}$.
Finally $cp\le z+R\le Cp$ on this interval. Multiplication proves
\eqref{inf:phase-energy-comparison}. This proof uses the order-one
finite expansion for all $x/\nu>1$ bounded away from one, including
the ratios below $3/2$ required for the high critical regions.
\end{proof}

\begin{lemma}[Finite-shift resonance through Lommel polynomials]
\label{inf:lommel-resonance}
Fix $0<\alpha_{\mathrm{loc}}<1$. If retained indices $j,j+d$,
$1\le d\le4$, both satisfy
$|d_j^{\mathrm{pole}}|,|d_{j+d}^{\mathrm{pole}}|\le p^{1-\alpha_{\mathrm{loc}}}$,
then
\[
 |U_{2d-1}^{\mathrm{Ch}}(\nu_j/R)|\le Cp^{-\alpha_{\mathrm{loc}}}.
\]
The ratio $\nu_j/R$ is within $Cp^{-\alpha_{\mathrm{loc}}}$ of
\begin{equation}\label{inf:critical-ratios}
 \mathcal T_{\mathrm{crit}}=
 \{1/2,1/\sqrt2,\sqrt3/2,\cos(\pi/8)\}.
\end{equation}
In particular $d=1$ is impossible for large $p$.
\end{lemma}
\begin{proof}
Write $\nu=\nu_j$, $\mu=\nu+2d$, and let $z,z'$ be the two retained
zeros. Their distances from $R$ are at most $Cp^{-\alpha_{\mathrm{loc}}}$.
At $J_\nu(z)=0$, one has $J_{\nu+1}(z)\ne0$ by ODE uniqueness.
Dividing the Bessel recurrence by this nonzero value gives the
finite Lommel recurrence, with initial values zero and one at
shifts zero and one. This identity holds for every real order.
For the finitely many shifts at most nine its recurrence coefficients
are $2\nu/z+O(p^{-1})$. The recurrence
$U_{k+1}^{\mathrm{Ch}}(t)=2tU_k^{\mathrm{Ch}}(t)-U_{k-1}^{\mathrm{Ch}}(t)$
therefore gives, by induction,
\[
 \frac{J_{\nu+m}(z)}{J_{\nu+1}(z)}
 =U_{m-1}^{\mathrm{Ch}}(\nu/z)+O(p^{-1}),\qquad 1\le m\le9.
\]
The ratios and the derivative ratio of $J_\mu$ are bounded by an
absolute constant there, using the derivative recurrence.
On the interval joining $z,z'$, the coefficients of the system for
$(J_\mu,J'_{\mu})$ are uniformly bounded: $x\asymp p$ and
$\mu<z'\le R+Cp^{-\alpha_{\mathrm{loc}}}$.
Its integral equation and Gronwall's inequality give
$\sup|J_\mu'|\le C|J_{\nu+1}(z)|$ on that interval.
Since $J_\mu(z')=0$, integration proves
$|J_\mu(z)|\le Cp^{-\alpha_{\mathrm{loc}}}|J_{\nu+1}(z)|$.
Divide and use the preceding recurrence; replacing $z$ by $R$
changes its polynomial argument by $O(p^{-1-\alpha_{\mathrm{loc}}})$.
This proves the bound, including near the turning region.

The zeros of $U_{2d-1}^{\mathrm{Ch}}$ are the simple numbers
$\cos(h\pi/(2d))$, $1\le h<2d$; this follows from
$U_k^{\mathrm{Ch}}(\cos\beta)=\sin((k+1)\beta)/\sin\beta$.
Here $\nu/R\ge1/2-o(1)$ and is at most $1+o(1)$.
The roots in this range, for $d=1,2,3,4$, are respectively
none, $1/\sqrt2$, $1/2,\sqrt3/2$, and
$1/\sqrt2,\cos(\pi/8)$. Simplicity and
$U_{2d-1}^{\mathrm{Ch}}(1)=2d$ prove the remaining assertions.
\end{proof}

\subsection{Core geometry and Green propagation}

Choose $\alpha_{\mathrm{loc}}>0$ with
$q_{\mathrm{gap}}+\xi_{\mathrm{inv}}<2-2\alpha_{\mathrm{loc}}$.
Set $W_{\mathrm{cond}}=C_{\mathrm{cond}}(1+\tau)$, with a fixed
sufficiently large $C_{\mathrm{cond}}$, and
\[
 \mathcal I_{\mathrm{win}}=\{j:|d_j^{\mathrm{pole}}|\le W_{\mathrm{cond}}\}.
\]
The angular complement of this set in the bandwidth-four matrix
$\mathcal S_{\mathrm{pole}}^{(0)}$ has every principal inverse
bounded by $2/W_{\mathrm{cond}}$. Indeed its band perturbation has
norm $C\tau$ and division by the unperturbed diagonal has norm at most
$1/W_{\mathrm{cond}}$, also after any coordinate restriction.
Let $\mathcal F_{\mathrm{win}}$ be the Schur complement of
$\mathcal S_{\mathrm{pole}}^{(0)}$ onto $\mathcal I_{\mathrm{win}}$.
It is self-adjoint for $\diag(\eta_j^2)$ and has Hilbert gap $c\gamma$
by Lemma~\ref{inf:centre-form-comparison} and angular complementary
reconstruction. The remainder $\mathsf R_{\mathrm{pole},6}$ is
restored after the window inverse in the proof of
Theorem~\ref{inf:mixed-conditioning}.

\begin{lemma}[Critical strips and the true minima]\label{inf:critical-strips}
For $c_{\mathrm{cr}}\in\mathcal T_{\mathrm{crit}}$ let
\[
 \mathcal N_{c_{\mathrm{cr}}}
 =\{j\in\R:|p+2j-1-Rc_{\mathrm{cr}}|
                       \le2p^{1-\alpha_{\mathrm{loc}}/2}\}.
\]
These neighbourhoods are disjoint for large $p$.
Write $c_{\mathrm{cr}}=\cos(\ell_{c_{\mathrm{cr}}}\pi/(2d_{c_{\mathrm{cr}}}))$
with reduced pairs
$(\ell_{c_{\mathrm{cr}}},d_{c_{\mathrm{cr}}})=(2,3),(1,2),(1,3),(1,4)$,
respectively, and put
\[
 f_{c_{\mathrm{cr}}}(j)=\varphi_{\mathrm{Bes}}(p+2j-1,R)
                      +\ell_{c_{\mathrm{cr}}}\pi j/d_{c_{\mathrm{cr}}}.
\]
On $\mathcal N_{c_{\mathrm{cr}}}$ its second derivative is between
$c_{\mathrm{curv}}/p$ and $C_{\mathrm{curv}}/p$, where
$c_{\mathrm{curv}},C_{\mathrm{curv}}>0$ are fixed, and
$|f'_{c_{\mathrm{cr}}}|\le Cp^{-\alpha_{\mathrm{loc}}/2}$.
Its true minimum is attained at
\begin{equation}\label{inf:true-phase-minimum}
 j_{c_{\mathrm{cr}}}^*=\{Rc_{\mathrm{cr}}-(p-1)\}/2+O(p^{-1}).
\end{equation}
Connect successive window vertices in each high neighbourhood
$c_{\mathrm{cr}}>1/2$ if their distance is at most
$L_{\mathrm{core}}(p)=K_{\mathrm{core}}\log p/\log\log p$.
Each resulting component has a common phase level and residue and
diameter at most
\begin{equation}\label{inf:sharp-core-diameter}
 L_{\mathrm{core}}(p)+C\sqrt{W_{\mathrm{cond}}+1}.
\end{equation}
Its indices satisfy $j\ge c_2p$ for a fixed $c_2>0$.
The low core consists of window indices $j\le C\sqrt{W_{\mathrm{cond}}+1}$,
all congruent to two modulo three. Every other low-neighbourhood
window vertex has $j\ge c\sqrt p$ and is isolated at scale
$L_{\mathrm{core}}(p)$.
For the nonseed indices of the low core one also has the uniform formula
\begin{equation}\label{inf:low-core-offsets}
 \begin{gathered}
 d_j^{\mathrm{pole}}=-\frac43(2j-1)(2j-4)
             +O\bigl((1+\log p)^2/p\bigr),\\
 j\equiv2\pmod3,\qquad 5\le j\le C\sqrt{W_{\mathrm{cond}}+1}.
 \end{gathered}
\end{equation}
For $j=5,8,11$, these leading offsets are $-72,-240,-504$, respectively.
\end{lemma}
\begin{proof}
The derivatives in the proof of Lemma~\ref{inf:phase-strip} give
$f_{c_{\mathrm{cr}}}''=4/\kappa+O(p^{-3})$. At
$j_{c_{\mathrm{cr}}}=\{Rc_{\mathrm{cr}}-(p-1)\}/2$ the derivative of the leading phase cancels
the added linear term, while the corrected derivative is
$-2b_\nu/(1+b^2)<0$, of size $O(p^{-2})$.
Its strictly positive second derivative
therefore puts its unique minimum at distance $O(p^{-1})$, proving
\eqref{inf:true-phase-minimum}.

By Lemma~\ref{inf:phase-eigenvalue-two-sided}, window vertices lie
within $C(W_{\mathrm{cond}}+1)/p$ of a level of the lattice
$\pi(\Z+1/2)+\ell_{c_{\mathrm{cr}}}\pi j/d_{c_{\mathrm{cr}}}$.
Levels are spaced by $\pi/d_{c_{\mathrm{cr}}}$.
An edge of length at most $L_{\mathrm{core}}$ changes $f_{c_{\mathrm{cr}}}$ by
$o(1)$, so its endpoints have the same level; reduction of
$(\ell_{c_{\mathrm{cr}}},d_{c_{\mathrm{cr}}})$ makes their residues equal.
Every component thus lies in one strip
$|f_{c_{\mathrm{cr}}}-g|\le C(W_{\mathrm{cond}}+1)/p$.
If $f''\ge m>0$, either monotone side of $|f-g|\le\delta$
has length at most $2\sqrt{\delta/m}$: its variation over length
$h$ is at least $mh^2/2$ and at most $2\delta$.
A connected chain either stays on one side or has an edge of
length at most $L_{\mathrm{core}}$ joining the two sides.
This proves \eqref{inf:sharp-core-diameter}. High neighbourhoods
have $j\asymp p$ by their definitions.

For $c_{\mathrm{cr}}=1/2$ the seed $j=2$ anchors the first level, since
$d_2^{\mathrm{pole}}=-\theta/(2p)+O(\theta^2p^{-3})$.
For $j\le c\sqrt p$ the phase cannot reach another level.
The curvature bound then gives $j\le C\sqrt{W_{\mathrm{cond}}+1}$,
and the first-level residue is $j\equiv2\pmod3$.
Beyond this first group, a window vertex has $j\ge c\sqrt p$.
There $f'_{c_{\mathrm{cr}}}\ge cj/p$; moving by an integer between one and
$L_{\mathrm{core}}$ changes its level by more than
$2C(W_{\mathrm{cond}}+1)/p$ and less than
$\pi/3-2C(W_{\mathrm{cond}}+1)/p$.
Hence such vertices are isolated at that scale.
To obtain \eqref{inf:low-core-offsets}, apply
Lemma~\ref{inf:low-recurrence} at the associated primary crossing
$p_0$. For $j\equiv2\pmod3$, $j\ge5$, its value/normal ratio is
$c_j/p_0+O((j+1)^4/p_0^2)$: the numerator has the stated
first-order recurrence coefficient, and the normal denominator
is its nonzero limiting sign plus $O((j+1)^2/p_0)$.
The collar ODE is uniformly bounded for
$j=O(\sqrt{\log p})$; its nearby zero therefore has squared
offset $-4c_j+O((j+1)^4/p_0)$.
Continue this zero and the Neumann zero to $p$.
Their frequency difference stays $O(\log p/p)$ by the global
order derivative bound, so spacing identifies the continued zero
with the retained pole. Repeating the refined derivative comparison
in the proof of Lemma~\ref{inf:wide-low-transfer} gives
$|z'(t)-R'(t)|\le C(j+1)/p+C(\log p)/p^2$.
The interval has length $O(\log p/p)$; the change of squared
offset is consequently at most
$C(\log p)(j+1)/p+C(\log p)^2/p^2$.
Since $j\le C\sqrt{W_{\mathrm{cond}}+1}=O(\sqrt{\log p})$,
these errors give \eqref{inf:low-core-offsets}.
\end{proof}

\begin{lemma}[Bin-elimination Green bound]\label{inf:bin-green}
Let a matrix on a finite subset of integer sites be
$\diag(d_j)+B$, where $B$ has bandwidth four and norm at
most $b_*\ge0$ in a weighted coordinate $\ell^1$ norm or a diagonal
Hilbert norm. Suppose $W>0$, $|d_j|>W$, $b_*/W<1/4$, and all coordinate
compressions are contractions. Assume the same band bound on
principal restrictions. Divide the sites into consecutive bins of
width four. Every principal inverse has norm at most $2/W$.
If a run of $m$ intervening bins has all $|d_j|\ge L>0$, with
$L\ge4(b_*+2b_*^2/W)$, then its bin-to-bin Green propagation has
bound
\begin{equation}\label{inf:bin-green-bound}
 \frac{C}{W}(2b_*/L)^m.
\end{equation}
Without the stronger diagonals there is exponential decay in the
number of intervening bins. Empty bins disconnect the matrix.
\end{lemma}
\begin{proof}
Diagonal division and a Neumann series bound every principal inverse
by $1/(W-b_*)\le2/W$. The bin matrix is block tridiagonal.
Eliminate bins from the right. The inverse of each right Schur
block is the compression of the corresponding tail inverse, so
has norm at most $2/W$. The diagonal block at a strong bin differs
from the restriction of $\diag(d_j)$ by at most $b_*+2b_*^2/W$.
Another diagonal Neumann series therefore bounds its inverse by
$2/L$. On a homogeneous row, backward elimination gives exactly
$u_b=-G_bT_{b,b-1}u_{b-1}$, with $\nrm{T_{b,b-1}}\le b_*$.
Iteration across the strong run gives $(2b_*/L)^m$.
For a source bin, eliminate on both sides; its resulting solution
operator is a compression of the full inverse and costs $2/W$.
This proves \eqref{inf:bin-green-bound}. Replacing $L$ by the
ordinary Schur bound gives ratio $2b_*/W<1/2$ on other bins.
The argument is valid separately in either specified norm; no
change of coordinates occurs inside the propagation estimate.
\end{proof}

\begin{lemma}[Valuations of low angular paths]\label{inf:path-valuation}
Uniformly for $j\le C\sqrt{1+\tau}+4$, the quartic angular
coefficients satisfy
\begin{equation}\label{inf:low-coefficient-valuations}
 \begin{aligned}
 c_{2,j}^{j+2}&=O(1),&
 c_{2,j}^{j+1},c_{2,j}^j&=O((j+1)/p),\\
 c_{2,j}^{j-1},c_{2,j}^{j-2}&=O((j+1)^2/p^2).
 \end{aligned}
\end{equation}
Every path of quartic interactions and angular-diagonal radial
resolvents carries the product of its angular coefficients.
This factorization applies to the feedback band
$p^{-1}\mathcal B_{\mathrm{pole},4}$ and to every further complementary path.
A two-interaction path starting at a low index $j$ in the stated
range has an extra factor $C(j+1)/p$ in its angular coefficient
product in each of the two cases of net displacement $+3$ and $-3$.
\end{lemma}
\begin{proof}
Put
\[
 \mathfrak u_j=\frac{\mathsf H_{j+1}(1)}{\mathsf H_j(1)},\qquad
 \mathfrak v_j=\mathfrak a_j^2
              \frac{\mathsf H_{j-1}(1)}{\mathsf H_j(1)},\qquad
 a_{\mathrm{quad}}=\frac{2(\ab+1)}{p+2},
\]
with $\mathfrak v_0=0$. The normalized recurrence is
$xG_j=\mathfrak u_jG_{j+1}+\mathfrak b_jG_j+\mathfrak v_jG_{j-1}$.
Two applications give the exact coefficients
\begin{align*}
 \mathsf H_2(1)c_{2,j}^{j+2}&=\mathfrak u_j\mathfrak u_{j+1},\\
 \mathsf H_2(1)c_{2,j}^{j-2}&=\mathfrak v_j\mathfrak v_{j-1},\\
 \mathsf H_2(1)c_{2,j}^{j+1}
 &=\mathfrak u_j(\mathfrak b_j+\mathfrak b_{j+1}-a_{\mathrm{quad}}),\\
 \mathsf H_2(1)c_{2,j}^{j-1}
 &=\mathfrak v_j(\mathfrak b_j+\mathfrak b_{j-1}-a_{\mathrm{quad}}),\\
 \mathsf H_2(1)c_{2,j}^{j}
 &=\mathfrak u_j\mathfrak v_{j+1}+\mathfrak v_j\mathfrak u_{j-1}
                                      +\mathsf H_2(\mathfrak b_j).
\end{align*}
Terms outside the nonnegative index range are zero.
The formulas \eqref{inf:recurrence-a}--\eqref{inf:recurrence-b}
and \eqref{inf:angular-jacobi} give explicitly
\begin{align*}
 \mathfrak u_j&=\frac{(p(1-r)+j)(p+j-1)}{(p+2j)(p+2j-1)},\\
 \mathfrak v_j&=\frac{j(pr+j-1)}{(p+2j-1)(p+2j-2)},\\
 \mathfrak b_j-r&=\frac{2j(p+j-1)(1-2r)}{(p+2j)(p+2j-2)},
 \qquad a_{\mathrm{quad}}-2r=\frac{2(1-2r)}{p+2}.
\end{align*}
Their denominators are comparable to the displayed powers of $p$
uniformly for $j\le C\sqrt{1+\tau}+4$. In particular
$\mathfrak u_j=O(1)$, $\mathfrak v_j=O((j+1)/p)$,
$\mathfrak b_j-r=O((j+1)/p)$, and $a_{\mathrm{quad}}-2r=O(p^{-1})$.
The endpoint ratios are bounded above and below on this range,
and $\mathsf H_2(1)\ge c_{\mathrm{sp}}>0$.
Writing \eqref{inf:H2} in powers of $x-r$ gives
$\mathsf H_2(\mathfrak b_j)=O((j+1)^2/p^2+p^{-1})$.
Since $j=O(\sqrt{\log p})$, every displayed coefficient has exactly
the bound in \eqref{inf:low-coefficient-valuations}.

The exact block formula in the proof of
Lemma~\ref{inf:centre-form-comparison} factors one interaction as
$c_{2,j}^i$ times its radial operator. An angular-diagonal radial
resolvent preserves this scalar factor. Induction on the number
of interactions therefore gives the product on any path, followed
by the product of its bounded radial norms. The Hilbert endpoint
ratio $\eta_i/\eta_j$ telescopes along the path.
Two steps of lengths at most two with net displacement three
must include a one-step edge. Equation~\eqref{inf:low-coefficient-valuations}
gives the asserted extra factor. Positive coefficient sums and
bounded bandwidth convert these bounds to column estimates,
without a factor for the number of low indices.
\end{proof}

\begin{lemma}[The seven off-core block types]\label{inf:off-core-table}
Take the high components of Lemma~\ref{inf:critical-strips}, its
low core, and all remaining window vertices as singleton cores.
For sufficiently large fixed $K_{\mathrm{core}}$, in both the
coefficient operator norm and the Hilbert operator norm,
\begin{equation}\label{inf:global-off-core}
 \nrm{\operatorname{Off}_{\mathrm{cores}}\mathcal F_{\mathrm{win}}}
 \le C\tau(1+\tau)^2p^{-2+2\alpha_{\mathrm{loc}}}.
\end{equation}
The low-core internal off-diagonal part has the stronger bound
$C\tau(1+\tau)^2p^{-2}$ in both norms.
\end{lemma}
\begin{proof}
Call an outside singleton one outside all $\mathcal N_{c_{\mathrm{cr}}}$, and a low
singleton one in $\mathcal N_{1/2}$ beyond its low core.
The following seven types exhaust all pairs of distinct cores;
the last row takes precedence when the neighbourhoods differ.
\begin{center}\small\normalfont
\begin{tabular}{@{}>{\raggedright\arraybackslash}p{.37\textwidth}>{\raggedright\arraybackslash}p{.56\textwidth}@{}}
\toprule Block type & Estimate and mechanism\\\midrule
Outside--outside & $C\tau(1+\tau)^2p^{-2+2\alpha_{\mathrm{loc}}}$;
two attenuated endpoints.\\
Outside--low singleton & Same bound; the low endpoint has the
stronger attenuation $C\tau/p$.\\
Low singleton--low singleton & Same bound; two primary edges
cannot join distinct such vertices.\\
Outside--high component & Faster than any chosen power,
including pairs crossing a critical-neighbourhood boundary.\\
High component--high component in one neighbourhood &
Convex phase barrier; $C\tau p^{-B_*}$ for any prescribed fixed $B_*$.\\
Low core--other core in its neighbourhood or outside &
Generic Green decay across distance at least $c\sqrt p$.\\
Cores in different critical neighbourhoods &
Generic Green decay across distance at least $cp$.\\
\bottomrule
\end{tabular}
\end{center}

Here are the operator estimates behind the table. An outside vertex
has every complementary neighbour within four sites at energy
distance greater than $p^{1-\alpha_{\mathrm{loc}}}$, by
Lemma~\ref{inf:lommel-resonance}; otherwise it would lie in the
smaller critical neighbourhood of radius $Cp^{1-\alpha_{\mathrm{loc}}}$.
Thus either endpoint map, with its adjacent diagonal inverse,
has norm $C\tau p^{-1+\alpha_{\mathrm{loc}}}$.
At a low singleton, for $h=\pm1,\pm2$ the phase change to $j+h$
is $-2h\pi/3+o(1)$, because $j/p\to0$ in the low neighbourhood
and $\partial_\nu\varphi_{\mathrm{Bes}}=-\arccos(\nu/R)+O(p^{-2})$.
Thus primary neighbours have phase separated from the zero lattice
by a fixed amount and denominators at least $cp$.
Feedback edges
already have size $C\tau^2/p$, and their denominator is at least
$W_{\mathrm{cond}}$. They give the endpoint bound $C\tau/p$.
All these maps have fixed bandwidth, so the estimates hold on
whole columns and, in Hilbert coordinates, on rows as well.

No direct band-four entry connects two distinct singleton cores.
Nor can two primary edges connect them: the endpoints would be
two window vertices within four sites and hence in the same
critical component. A length-two path between distinct singletons
therefore includes a feedback edge. With one such edge put the
diagonal inverse at the primary endpoint, giving
$C\tau^3p^{-2+\alpha_{\mathrm{loc}}}$; with two use
$C\tau^4/(p^2W_{\mathrm{cond}})$.
For all paths of length at least three, the resolvent identity
places a diagonal inverse at each endpoint and a bounded
Neumann inverse between them. Their global bound is
$C\tau^3p^{-2+2\alpha_{\mathrm{loc}}}$.
This proves the first three rows as operator norms, with no
sum over the total number of vertices.

For two distinct high components in one neighbourhood, let their
nearest endpoints be $u<v$, with $D=v-u>L_{\mathrm{core}}$, and set
\[
 h=\pi/d_{c_{\mathrm{cr}}},\qquad
 \delta=C(W_{\mathrm{cond}}+1)/p.
\]
Here $h$ is the spacing of the union of the phase levels of all
integer residues. If their phase levels differ,
there is an intervening phase interval a fixed distance from all
levels. Since $|f'_{c_{\mathrm{cr}}}|\le Cp^{-\alpha_{\mathrm{loc}}/2}$,
this supplies at least $cp^{\alpha_{\mathrm{loc}}/2}$ consecutive
sites with denominator $cp$ by
\eqref{inf:phase-energy-comparison}. If both endpoints have level
$g$, define their excursion by
$\mathfrak e=g-\min_{[u,v]}f_{c_{\mathrm{cr}}}$.
The convexity inequality gives
$g-f_{c_{\mathrm{cr}}}(x)\ge cD^2/p-\delta$
on $[u+D/4,v-D/4]$.
If $\mathfrak e\ge h/4$, continuity supplies the phase interval
$[g-h/4,g-h/8]$, a fixed distance from every level;
use the preceding barrier on its preimage.
If $\mathfrak e<h/4$, the middle half has matching-residue
denominators at least $cD^2$, and all other residues have
denominators at least $cp$. This branch forces $D\le C\sqrt p$,
so, after decreasing its fixed constant, take $L=cD^2\le cp$.
Since $D>L_{\mathrm{core}}$ and
$L_{\mathrm{core}}^2/(1+\tau)\to\infty$, one has
$L\ge4(b_*+2b_*^2/W_{\mathrm{cond}})$ with $b_*=C\tau$.
The middle half contains at least $D/8-2$ complete bins of width
four. Their diagonals exceed $W_{\mathrm{cond}}$ for large $p$.
Lemma~\ref{inf:bin-green} bounds propagation
by $\exp\{-cD\log(D^2/[C(1+\tau)])\}$.
Uniformly for $\tau\le\log p$,
\[
 \log\frac{L_{\mathrm{core}}(p)^2}{C(1+\tau)}
 \ge\tfrac12\log\log p
\]
eventually. Choose $K_{\mathrm{core}}$ to give any fixed power
$p^{-B_*-2}$. There are only $O(p)$ retained vertices; summing
the already weighted entry bounds gives all coefficient columns
and Hilbert rows and columns with the desired power.
Each such Schur term has two exterior couplings. Together with
the source factor $1/W_{\mathrm{cond}}$ their size is at most
$C\tau^2/W_{\mathrm{cond}}\le C\tau$, giving the detuning factor
in the high-component rows of the table.

At a boundary of $\mathcal N_{c_{\mathrm{cr}}}$, $|f'_{c_{\mathrm{cr}}}|$ is comparable to
$p^{-\alpha_{\mathrm{loc}}/2}$. Two window points on opposite
sides of that boundary cannot have the same level at integer
distance between one and $cp^{\alpha_{\mathrm{loc}}/2}$:
their phase change exceeds $C(W_{\mathrm{cond}}+1)/p$.
Different levels require at least that distance as well.
Thus generic exponential Green decay proves the outside--high
row, including these boundary pairs. Distinct neighbourhoods are
separated by $cp$. The low core is separated from other vertices
of its neighbourhood by $c\sqrt p$. These facts prove the last
two rows.

It remains to bound the internal low-core off-diagonal part.
Its vertices have residue two modulo three, so there is no direct
primary edge between distinct vertices. The direct feedback of
displacement three has size
$C\tau^2(j+1)p^{-2}$ by Lemma~\ref{inf:path-valuation}.
A primary endpoint lands on a nonresonant index with denominator
$cp$, uniformly for $j\le C\sqrt{1+\tau}$.
Indeed $f_{1/2}(j)-f_{1/2}(2)=O((j+1)^2/p)=O(W_{\mathrm{cond}}/p)$,
while either of the other two residues has its zero lattice shifted
by $\pi/3$ from this level. Lemma~\ref{inf:phase-eigenvalue-two-sided}
therefore supplies that lower bound also for the neighbouring indices.
A feedback endpoint of displacement one, two or four also lands
on a nonresonant index; displacement three retains the extra
coefficient $(j+1)/p$ even if its denominator is only
$W_{\mathrm{cond}}$. Two-primary-edge paths of net displacement
three gain the same extra coefficient and a $cp$ denominator.
A zero-displacement step stays at its window index and is therefore
absent from an endpoint map into the angular complement.
Consequently each endpoint map has coefficient norm at most
$C\tau/p$. A two-edge term containing a feedback edge is bounded by
$C\tau^3p^{-2}+C\tau^4/(p^2W_{\mathrm{cond}})$, placing the inverse
at the primary endpoint when one is present.
The resolvent identity bounds the sum of paths with at least three
edges by $C\tau^3p^{-2}$, using the two endpoint bounds and a
bounded middle Neumann inverse. Bounded-band column sums therefore give
$C\tau(1+\tau)^2p^{-2}$ for the entire low-core off-diagonal matrix.
It is self-adjoint in a positive diagonal metric. Its Hilbert
operator norm equals its spectral radius, which is at most its
coefficient operator norm. This supplies the Hilbert bound without
any factor for the growing low-core cardinality.
\end{proof}

\begin{theorem}[Mixed conditioning and restoration]
\label{inf:mixed-conditioning}
Under the hypotheses of this section the linear quartic operator
has a two-sided mixed graph inverse with
\begin{equation}\label{inf:linear-mixed-inverse}
 \nrm{K_0L_{\mathrm q}^{-1}}_{\Mmix\to\Mmix}
 \le C p^{2+\xi_{\mathrm{inv}}}/\gamma.
\end{equation}
For a fixed order $N$, suppose its true centre has bounded degree,
strong size $C\tau p^{-2}$, and
\begin{equation}\label{inf:conditioning-centre-hypothesis}
 \nrm{f_N-f_{\mathrm q}}_{\mathrm{strong}}
 \le\tau p^{-4}\mathfrak E_N(p),\qquad
 \mathfrak E_N(p)=p^{o(1)}
\end{equation}
uniformly in the parameters. Then its actual scalar graph is a
bounded isomorphism on $\Cc$ and
\begin{equation}\label{inf:conditioning-final-inverse}
 \nrm{J_D^{-1}}_{\Cc\to\Cc}
 \le C_Np^{5/2+\xi_{\mathrm{inv}}}/\gamma.
\end{equation}
The threshold may depend on the compact split interval,
$N,q_{\mathrm{gap}},\xi_{\mathrm{inv}},\kappa_{\mathrm{gap}}$, and
the envelope in \eqref{inf:conditioning-centre-hypothesis}.
\end{theorem}
\begin{proof}
Delete cross-core entries of
$\mathcal F_{\mathrm{win}}$. Their Hilbert norm is $o(\gamma)$ by
\eqref{inf:global-off-core}, so every core retains Hilbert gap
$c\gamma$. In a high core, the exact ratio
\begin{equation}\label{inf:core-conditioning}
 \frac{\eta_{j+1}^2}{\eta_j^2}
 =\frac{2j+p-1}{2j+p+1}
   \frac{(pr+j)(j+1)}{(p(1-r)+j)(p-1+j)}
\end{equation}
is bounded above and below by fixed positive constants since
$j\ge c_2p$ and $j=O(p)$. In coordinates
$\widetilde u_j=2^ju_j$, the coefficient norm is unweighted
$\ell^1$ and the Hilbert weights are $\eta_j/2^j$.
Their consecutive ratios have the same boundedness property.
The diagonal similarity over the core diameter, followed by
$\ell^2$--$\ell^1$ comparison over its cardinality, therefore costs at most
\[
 \exp\{C L_{\mathrm{core}}(p)+C\sqrt{W_{\mathrm{cond}}+1}\}
 \sqrt{L_{\mathrm{core}}(p)+C\sqrt{W_{\mathrm{cond}}+1}+1}
 =p^{o(1)}.
\]
For a singleton the cost is one. On the low core, delete its
internal off-diagonal part in the Hilbert norm. Each diagonal
entry then has modulus at least $c\gamma$. Restoring that part
by a coefficient diagonal Neumann series gives inverse $C/\gamma$,
because $(1+\tau)^2p^{q_{\mathrm{gap}}-2}\to0$.

The direct sum of all core inverses has coefficient norm at most
$Cp^{\xi_{\mathrm{inv}}}/\gamma$, since the displayed subpolynomial
factor is at most $p^{\xi_{\mathrm{inv}}/2}$ for sufficiently large $p$. Its product with the off-core error is bounded by
$C(1+\tau)^2p^{q_{\mathrm{gap}}+\xi_{\mathrm{inv}}-2+2\alpha_{\mathrm{loc}}}=o(1)$.
Thus $\mathcal F_{\mathrm{win}}^{-1}$ has that same bound.
Recovering the angular complement costs bounded factors, since
$\tau/W_{\mathrm{cond}}$ is small. Restoring
$\mathsf R_{\mathrm{pole},6}$ costs
$C\tau^2p^{q_{\mathrm{gap}}+\xi_{\mathrm{inv}}-2}=o(1)$.
Hence the exact pole Schur inverse has the same coefficient bound.

The radial complementary reconstruction has reduced right side
of norm at most $(1+C\tau/p)\nrm f_{\Mmix}$.
The pole component costs $Cp^{\xi_{\mathrm{inv}}}\nrm f/\gamma$;
its $K_0$ norm costs a further $Cp^2$.
The complementary graph norm is bounded by
$Cp(\nrm f+C\tau\nrm{\Pi_{\mathrm{rad}}u})$.
This proves \eqref{inf:linear-mixed-inverse} on the entire graph
domain, with all radial states included. Lemma~\ref{inf:mixed-domain}
applies to the small relative perturbation $V_1(H_{\mathrm q})$;
the bounded scalar shift leaves its graph domain unchanged.
The reconstructions converge in that graph norm. The embeddings and
uniqueness identify this inverse with the Hilbert inverse on
mixed sources, and prove both inverse identities.

For the actual centre put $\Phi=\Phi_N$. Apply
Lemma~\ref{inf:quadratic-graph} with $H=\HP f_N$ and use
\eqref{inf:conditioning-centre-hypothesis}. They give
\[
 \nrm{(\Hu-H_0-V_1(H_{\mathrm q}))K_0^{-1}}
 \le C_N\tau p^{-4}(\mathfrak E_N(p)+\tau)=\tau p^{-4+o(1)}.
\]
Multiply by the mixed graph inverse before any radial coefficient
conversion. Its product is
$p^{q_{\mathrm{gap}}+\xi_{\mathrm{inv}}-2+o(1)}=o(1)$.
A Neumann series restores the true centre on $\Mmix$.
On $\Hc$, its linear graph inverse satisfies
$\nrm{K_0L_{\mathrm q}^{-1}}\le Cp^2/\gamma$.
Indeed, for $L_{\mathrm q}u=f$ write
$K_0u=f+2R^2u-V_1(H_{\mathrm q})u$, use the Hilbert gap
$\nrm u\le C\gamma^{-1}\nrm f$, and absorb
$\nrm{V_1(H_{\mathrm q})K_0^{-1}}<1/2$.
The same centre perturbation therefore has Hilbert product
$p^{q_{\mathrm{gap}}-2+o(1)}=o(1)$.
The restored Hilbert and mixed inverses agree on mixed sources
by the embedding and Hilbert uniqueness.
The density and graph-conjugation proof of
\eqref{inf:density-graph-bounds} uses only $p\delta$ small;
its binomial bound is
$\exp(Cp\delta)=\exp(O(\tau/p))$, hence bounded.
The scalar return therefore preserves the mixed inverse exponent.
The scalar estimates in the proof of Theorem~\ref{inf:relative-graph}
bound the sum of the $\Cc$ and $\Mmix$ relative norms by
$C(\delta+p\delta^2)\le C\tau p^{-2}$.
Thus \eqref{inf:head-tail-hypotheses} holds for large $p$.
The finite-head/infinite-tail proof of Lemma~\ref{inf:head-tail}
now costs exactly
$\sqrt p$. Its tail is solved in $\Cc$ and its head is converted
from the already constructed mixed solution. This proves
coefficient membership, both inverse identities, and
\eqref{inf:conditioning-final-inverse}.
All perturbations retain a factor $\tau$; only the final inverse
has a factor $\tau^{-1}$.
\end{proof}

\section{Uniform dependence on the detuning bound}\label{inf:sec-tame}

Fix the compact split interval, a finite centre accuracy $N$, and
every finite derivative order and complex disc required for that
centre. Constants in this section depend on these choices.
Coefficient and mixed operator estimates are at integer block sizes;
the intervening flow parameters use the weighted realization of
Proposition~\ref{inf:flow-form}.
They are independent of $p$ and of $0<\tau\le T\le\log p$.
We may replace $T$ by $\max(1,T)$ when forming an envelope.

\begin{lemma}[Detuning-dependent estimates]\label{inf:detuning-dependencies}
Suppose $|p-p_0|\le C_{\mathrm{ap}}/p_0$, where
$C_{\mathrm{ap}}\le C_*\log p_0$ with fixed $C_*$, and the continued
zero has $0<\tau\le T\le\log p$. There are finite integers
$k_N,d_N^{\mathrm{tame}}$ and constants $c_N,C_N>0$ such that the
finite-centre estimates hold on
\[
 |\eps|\le c_N(1+T)^{-k_N},\qquad
 \mathfrak P_N(T)=(1+T)^{d_N^{\mathrm{tame}}},\qquad
 \mathcal B_N(T)=C_N\mathfrak P_N(T)e^{C_N\sqrt T}.
\]
Increasing the fixed exponents and constants when necessary, every
bound in the following table holds. Here $E=K_{\mathrm{flow}}(1+\tau)$,
$W=C_{\mathrm{cond}}(1+\tau)$, and $\delta$ is the strong centre size.
The labels in the first column specify the inherited statement
and the quantity whose bound is being extended.

\begingroup\small\normalfont
\renewcommand{\arraystretch}{1.16}
\begin{longtable}{@{}>{\raggedright\arraybackslash}p{.06\textwidth}>{\raggedright\arraybackslash}p{.38\textwidth}>{\raggedright\arraybackslash}p{.48\textwidth}@{}}
\toprule & Quantity and statement & Uniform bound / small parameter\\\midrule
\endfirsthead
\toprule & Quantity and statement & Uniform bound / small parameter\\\midrule
\endhead
1 & Root continuation, \ref{inf:detuning} &
$R-R_{\mathrm{alg}}=-(2\delta_0+O(p^{-1}))(p-p_0)$.\\
2 & Denominator in \eqref{inf:trace-ratio-exact} &
$2R^2-4(p+1)(p+2)=4p^2(1+O(p^{-1}))$.\\
3 & Low-state transfer, \ref{inf:wide-low-transfer} &
$C C_{\mathrm{ap}}(\ell+1)/p+C C_{\mathrm{ap}}(1+C_{\mathrm{ap}})/p^2$.\\
4 & Radius equation, \eqref{inf:rho-equation} &
Analytic in $(\eps,w)$, $w=\theta\eps^2$; $|w|\le T/p^2$.\\
5 & Finite centre jet, \ref{inf:centres} &
Each fixed $\eps$ coefficient is polynomial in $\theta$;
common $\eps$ radius $c_N(1+T)^{-k_N}$.\\
6 & Collar and complex angular disc, \eqref{inf:collar-ode} &
Fixed discs; physical displacement $C\tau/p$.\\
7 & Two physical defects, \eqref{inf:centre-complex-defect} &
$C_N\theta^2p^{-N}\mathfrak P_N(T)$, also in the weighted
coefficient norms of \ref{inf:centre-coefficient-defect}.\\
8 & Quartic comparison, \eqref{inf:centre-profile} &
Physical error $C_N\tau p^{-3}\mathfrak P_N(T)$;
normalized error $C_N\tau p^{-4}\mathfrak P_N(T)$.\\
9 & Density, \eqref{inf:density-graph-bounds} &
$\nrm{m^{\pm1}}\le e^{Cp\delta}$,
$\delta\le C\tau/p^2$, so $p\delta\le CT/p$.\\
10 & Radial modes, \ref{inf:radial-centre-norm} &
$C_{N,k}p^{k+1}$; normalization units remain separated from zero.\\
11 & Known-mode composition, \eqref{inf:composition-bound} &
Head factor $e^{C_N\sqrt T}$; tail step factor
$e^{C_N\sqrt T/p}$ against Bessel ratio $1/4$.\\
12 & Clamped centre and material,
\ref{inf:centre-source-norm}, \ref{inf:sharp-material} &
$\nrm{U_0}\le p\mathcal B_N(T)$;
$\nrm{g_0}\le p^3\mathcal B_N(T)$;
$\nrm{\Qmat^{-1}}\le p^2\mathcal B_N(T)$;
$\sup\nrm{D^2\Fnl}\le p^3\mathcal B_N(T)$.\\
13 & Relative flow form, \ref{inf:flow-relative} &
$C\tau p^{-2}$ and $C\tau^2p^{-3}$;
series ratio $C\tau/p^2$.\\
14 & Flow complement, \ref{inf:flow-complement} &
Choose $E=K_{\mathrm{flow}}(1+\tau)$; ratio $C\tau/E$.\\
15 & Schur derivative, \eqref{inf:flow-schur-remainder} &
$C\tau^2(E^{-1}+p^{-1})$;
$\tau^3/E^2\le\tau^2/K_{\mathrm{flow}}E$.\\
16 & Seed diagonal, \ref{inf:flow-low-block} &
$\theta/(2p)+O(\theta^2p^{-3})$;
weighted low remainder $C\tau^2(E^{-1}+p^{-1})p^{-2}$.\\
17 & Index-five diagonal, \eqref{inf:low-offsets} &
$-72+O((1+T)/p)$; primary perturbation $O(\tau/p)$.\\
18 & Low feedback, \eqref{inf:weighted-low-inverse} &
For $a_*=C\tau^2(E^{-1}+p^{-1})$, bound $Ca_*^2/(p\tau)$;
ratio to $\tau$ tends to zero.\\
19 & Transport, \eqref{inf:transport-bounds} &
$C\tau^2/p$ in the weighted derivative form;
ratio to the main margin $O(\tau/p)$.\\
20 & Norm recovery, \ref{inf:flow-theorem} &
$\nrm u\le C(\sqrt p+\tau)\nrm{\mathsf Wv}$;
slope at least $c\tau/p$.\\
21 & Phase/count, \ref{inf:window-isolation}--\ref{inf:global-count} &
Phase thickness $C(1+W)/p$;
count $C\sqrt{1+W}\,p^{2/3}$ for $W\le C\log p$.\\
22 & Shell and split count, \ref{inf:quartic-count-window},
\ref{inf:many-splits} &
Shell $C\tau/p^2$; count window $C(1+\tau)$;
bad splits $C\sqrt{1+\tau}(p^{2/3}+p^{8/3-q_{\mathrm{gap}}})$.\\
23 & Radial conditioning, \ref{inf:radial-poles} &
Complementary graph $Cp$, Neumann ratio $C\tau/p$.\\
24 & Linear spectral comparison, \ref{inf:centre-form-comparison} &
Error $C\tau^2p^{-2}$; divided by $\gamma$ it is
$C\tau p^{q_{\mathrm{gap}}-2}$.\\
25 & Angular conditioning window &
$W=C_{\mathrm{cond}}(1+\tau)$; every principal inverse $2/W$.\\
26 & High core, \ref{inf:critical-strips} &
Diameter $L_{\mathrm{core}}+C\sqrt{W+1}$;
condition factor $e^{C(L_{\mathrm{core}}+\sqrt{W+1})}
\sqrt{L_{\mathrm{core}}+C\sqrt{W+1}+1}$.\\
27 & High barrier, \ref{inf:off-core-table} &
$\log(L_{\mathrm{core}}^2/[C(1+\tau)])\ge
\tfrac12\log\log p$; choose $K_{\mathrm{core}}$ before $p$.\\
28 & Low core, \ref{inf:path-valuation} &
$j\le C\sqrt{1+\tau}$; one-step coefficient
$C(j+1)/p$, downward coefficients $C(j+1)^2/p^2$.\\
29 & Low-core off-diagonal, \ref{inf:off-core-table} &
$C\tau(1+\tau)^2p^{-2}$ in coefficient and Hilbert norms.\\
30 & Singleton off-core, \eqref{inf:global-off-core} &
$C\tau(1+\tau)^2p^{-2+2\alpha_{\mathrm{loc}}}$.\\
31 & True-centre graph, \ref{inf:quadratic-graph} &
$C_N\tau p^{-4}(\mathfrak P_N(T)+\tau)$;
restore in $\Mmix$ before coefficient conversion.\\
32 & Radial conversion, \ref{inf:head-tail} &
One factor $\sqrt p$; tail inverse uniformly bounded.\\
33 & Convexity, \ref{inf:geometry-convexity} &
Normalized $C^2$ shape $C\tau p^{-2}$ plus Newton correction.\\
34 & Non-ball witness, \ref{inf:geometry-nonball} &
Leading physical coefficient $-\theta/(16r(1-r)p)$;
centre error $C_N\tau p^{-2}\mathfrak P_N(T)$.\\
\bottomrule
\end{longtable}
\endgroup
The double lift, residual, and material-defect bounds are
\begin{equation}\label{inf:tame-residuals}
 \nrm{\ell_N}\le C_N\theta^2p^{-N}\mathfrak P_N(T),\quad
 \nrm{\Fnl(z_0)}\le C_N\theta^2p^{2-N}\mathfrak P_N(T),\quad
 \nrm{\mathcal V_{\mathrm{res}}}\le
                  C_N\theta^2p^{3-N}\mathfrak P_N(T).
\end{equation}
The normal Hessian satisfies
$\nrm{D_p^{-1}}_{\Ec_1}\le C_Np^{-2}\mathcal B_N(T)$.
The scalar fixed-centre graph difference is $C\tau p^{-2}$ and
the unitary relative bound is $C(\delta+p\delta^2)$.
Every $o(1)$ obtained from these envelopes is uniform in
$0<\tau\le\log p$. The local cluster/cofactor bound is used only
with fixed $T$; the growing-window inverse is
Theorem~\ref{inf:mixed-conditioning}.
\end{lemma}
\begin{proof}
We give the dependence calculations in the order of their inputs.

\emph{Roots and low states.}
Equation~\eqref{inf:order-derivative} has error $O(p^{-2})$ and its
comparison function is smooth near $1/2$. Throughout the stated
interval, $p/R=1/2+O(p^{-1})$ and $R\to\infty$.
Thus $R'=2\pi/(3\sqrt3)+O(p^{-1})$ and
$R'_{\mathrm{alg}}=2+O(p^{-2})$.
Integrating their difference retains exactly the factor $p-p_0$.
Substitution in \eqref{inf:trace-ratio-exact} proves rows 1--2 and
\begin{equation}\label{inf:log-detuning-relation}
 \theta=(16\delta_0+O(p^{-1}))p_0(p-p_0),\qquad
 R=2p+3+O((1+C_{\mathrm{ap}})/p).
\end{equation}
There is no additive detuning error. Row 3 is
\eqref{inf:wide-low-error}.

\emph{Finite jets and their remainders.}
The radius equation is analytic in $(\eps,w)$ with
$w=\theta\eps^2$. Coefficient extraction in the finite recurrence,
the collar ODE, the monic products, and the update
\eqref{inf:formal-update} shows inductively that every fixed
$\eps$ coefficient is a polynomial in $\theta$.
All leading divisors are independent of $\theta$: they are the
normalization units, $-1/8$, $c_j$, and nonzero functions of $r$.
For this fixed finite jet choose $k_N$ larger than the finitely
many ratios of positive $\theta$ degrees to positive $\eps$ orders.
On $|\eps|\le c_N(1+T)^{-k_N}$, $|\theta|\le2(T+1)$, its
amplitudes and displacement are uniformly small by choosing $c_N$.
The variable $w$, the collar denominators $\rho+\eps y$, and every
finite normalization remain in fixed compact analytic sets.
The same choice works on the prescribed fixed complex angular disc.

The exact defects of the evaluated finite jet vanish through the
required order in $\eps$. They have an analytic factor $\theta^2$
by the identical field/boundary linear seed, as proved in
Theorem~\ref{inf:centres}. Divide by this factor; the maximum
principle on the larger $\theta$ circle bounds the quotient
uniformly on the smaller disc. The Cauchy remainder estimate in
$\eps$ costs at most a fixed power $(1+T)^{d_N^{\mathrm{tame}}}$.
The same argument applied to $(h_N-h_{\mathrm q})/\theta$ uses its
vanishing first three $\eps$ coefficients and gives row 8.
For $T\le\log p$ the condition
$p^{-1}(1+T)^{k_N}\to0$ puts the real input inside the chosen disc.
Lemma~\ref{inf:analytic-coefficients} transfers the bounds to every
fixed coefficient moment and radius. This proves rows 4--8.

The first term of the normalized shape has size $C\tau p^{-2}$,
with a constant independent of the fixed order. Its remainder is
$C_N\tau p^{-4}\mathfrak P_N(T)$; its ratio to the leading size
tends to zero. Thus $\delta\le C\tau p^{-2}$ in the strong norm.
The scalar estimates in the proof of Theorem~\ref{inf:relative-graph}
are $C(\delta+p\delta^2)$, with no fixed-$T$ compactness constant.
Since $p\delta=O(T/p)\to0$, the density series give row 9 and the
stated graph bounds.

The mode-normalization ratios in
Lemma~\ref{inf:radial-centre-norm} are finite recurrence functions
of $(\eps,w)$ and tend to the same nonzero constants. Its Parseval
head bound and the real-order tail ratio therefore hold uniformly.
For composition, the argument of the square root in
\eqref{inf:composition-bound} is $C_N\tau p^{-2}s(2s+\beta+1)$.
It is $O_N(T)$ on the $s=O_N(p)$ head, and on the tail gives an
extra geometric step $e^{C_N\sqrt T/p}$, absorbed by the ratio
$1/4$. The finite amplitudes stay bounded on the real input.
This proves rows 10--11 and the old membership/moment estimates
with envelope $\mathcal B_N(T)$.
The exact lift columns and their Laplacians give
\eqref{inf:tame-residuals} by precisely the two-layer calculation
of Proposition~\ref{inf:residual-estimates}.
The inverse multiplier series preserve every fixed moment since
$C\tau/p^2\to0$.
The normal Hessian identity \eqref{inf:normal-hessian} now gives
its stated inverse bound. Repeating the two norm absorptions in
Proposition~\ref{inf:sharp-material} yields row 12; finite products
of its envelopes are included by increasing $d_N^{\mathrm{tame}},C_N$.
These arguments use composition only for the known finite modes,
clamped gains only on $\Gc_p$, and all-layer Poisson estimates
for $\mathfrak f_0$.

\emph{Uniform flow.}
The series in Lemma~\ref{inf:flow-relative} is dominated by
$\sum h^6(C\tau/p^2)^h$, proving row 13 with an absolute constant
for the fixed split interval. In the complementary proof choose
$E=K_{\mathrm{flow}}(1+\tau)$ independently of $\varsigma$.
Then $C\tau/E\le C/K_{\mathrm{flow}}$ and $E/p^2\to0$.
The differentiated complementary inverse adds
$C\tau^3/E^2$, bounded as in row 15. Thus rows 14--15 hold
without differentiating a moving window projection.
The finite recurrence and the collar root argument give rows
16--17 with their displayed factors of $\theta$.
In the low feedback calculation
\[
 \frac{a_*^2}{p\tau}
 \le C\left(\frac{\tau}{pK_{\mathrm{flow}}^2}
       +\frac{\tau^2}{p^2K_{\mathrm{flow}}}
       +\frac{\tau^3}{p^3}\right)=o(\tau).
\]
The transport error is $C\tau^2/p=o(\tau)$.
Choose $K_{\mathrm{flow}}$ to absorb $C\tau/K_{\mathrm{flow}}$
into the positive angular margin and then take $p$ large.
Low-state transfer from row 3 gives the same $32$ exclusion on
indices below $\lfloor\sqrt p/16\rfloor$.
The norm recovery in the last paragraph of
Theorem~\ref{inf:flow-theorem} is
$\nrm u\le C(\sqrt p+\tau)\nrm{\mathsf Wv}$;
$\tau\le\log p=o(\sqrt p)$ proves rows 18--20.

\emph{Counting.}
The proof of Lemma~\ref{inf:window-isolation} needs only
$2W/R<1/10$ and a frequency interval shorter than $\pi$.
These hold for $W\le C\log p$.
In Lemma~\ref{inf:phase-strip}, the Debye error is independent of
$W$ and replacing the zero by $R$ costs at most $C W/p$.
Its thickness is therefore $C(1+W)/p$.
On each dyadic band the determinant error in
Lemma~\ref{inf:lattice-strip} has its same fixed curvature term
and additional error $O(Wp^{-2/3})=o(1)$.
Its number of collinear points is at most $C\sqrt{1+W}$.
Summing the dyadic intervals as in Theorem~\ref{inf:global-count}
gives row 21; the turning strip has only $O(p^{2/3})$ indices
uniformly in this range of $W$.
The shell width $C\tau/p^2$ gives count window $C(1+\tau)$ by
indexwise min--max. The grid proof of Lemma~\ref{inf:many-splits}
then gives row 22 with its explicit square-root factor.

\emph{Conditioning and geometry.}
Rows 23--32 are the estimates proved in
Section~\ref{inf:sec-conditioning}, which kept $\tau\le\log p$
throughout. Row 8 supplies its true-centre hypothesis with
$\mathfrak E_N(p)=C_N\mathfrak P_N(\log p)$.
The high-core diameter in row 26 is linear in $\sqrt W$;
its exponential is therefore subpolynomial. Its barrier gain is
row 27. The three low-core facts used in rows 28--29 are: its
indices have residue two modulo three; every primary edge lands
on a nonresonant index with denominator $cp$; a feedback edge
that lands on another resonant index has displacement three and
the extra coefficient $(j+1)/p$. These are the statements proved
in Lemmas~\ref{inf:critical-strips} and \ref{inf:path-valuation}.
The full off-diagonal estimates are operator norms as in
Lemma~\ref{inf:off-core-table}. The factor in row 31 is multiplied
by the mixed graph inverse, yielding
$p^{q_{\mathrm{gap}}+\xi_{\mathrm{inv}}-2+o(1)}\to0$.
The final radial conversion costs just row 32.
Rows 33--34 now follow from the fixed-order shape comparisons and
the dimension-independent $C^2$ embedding.

Finally
\[
 \mathcal B_N(\log p)
 =\exp\{O_N(\log\log p+\sqrt{\log p})\}=p^{o(1)}.
\]
The other envelopes in the table are fixed powers of $\log p$ or
have the form
\[
 \exp\{O(\log p/\log\log p+\sqrt{\log p})\}.
\]
The strict power margins absorb all of them. The two factors in
\eqref{inf:tame-residuals} are $\theta^2$, the inverse has only
$\tau^{-1}$, and the radius and witness carry $\tau$.
No positive lower bound for $\tau$ has been used.
\end{proof}

\begin{corollary}[Uniform logarithmic flow and grid count]
\label{inf:log-flow-count}
For the inputs above, the signed slope is at least $c\tau/p$ in
$|\lambda|<\tau/(16p)$. For any fixed $5/3<q_{\mathrm{gap}}<2$,
all but
\begin{equation}\label{inf:log-bad-splits}
 C\sqrt{1+\tau}\bigl(p^{2/3}+p^{8/3-q_{\mathrm{gap}}}\bigr)
\end{equation}
interior integer splits have actual quartic gap at least
$\kappa_{\mathrm{gap}}\tau p^{-q_{\mathrm{gap}}}$.
\end{corollary}
\begin{proof}
Rows 13--22 of Lemma~\ref{inf:detuning-dependencies} establish the
flow, count, and spacing hypotheses. Apply the clipped-curve
argument of Lemma~\ref{inf:many-splits} with its explicitly stated
number of curves. The expression in \eqref{inf:log-bad-splits}
is $o(p)$ because $q_{\mathrm{gap}}>5/3$.
\end{proof}

\section{The logarithmic family and fixed-order closure}
\label{inf:sec-logarithmic}

\begin{definition}[The logarithmic input family]\label{inf:log-family}
With $m_0$ as in Definition~\ref{inf:input-families}, put
\[
 \mathcal P_{\log}=\left\{p\in\tfrac12\Z:p>1,\quad
 \exists m\ge m_0:\ |p-\widehat p_m|\le\frac{\log p}{8p}\right\}.
\]
The root $R$ and detuning are obtained by continuation from this
crossing. For sufficiently large members the crossing is unique,
and \eqref{inf:log-detuning-relation} gives $0<\tau\le\log p$.
Indeed $16\delta_0<8$, and non-coincidence excludes $p=\widehat p_m$.
Every fixed-width family $\mathcal P_{\mathrm{fix}}(C_{\mathrm{ap}})$ is contained
in $\mathcal P_{\log}$ from some point on.
\end{definition}

\begin{lemma}[Multipliers giving unbounded detuning]\label{inf:log-multiplier}
In each parity, $\mathcal P_{\log}$ has an unbounded selection with
$\widehat p_m|p-\widehat p_m|\to\infty$ and $\tau\to\infty$.
\end{lemma}
\begin{proof}
The congruence and shifted congruence lemmas give infinitely many
pairs $(P,Q)$ with $Q\equiv1\pmod4$, either desired parity of $P$,
and $0<e=Q|Q\alpha_{\mathrm{ar}}-P|<9/2$.
For each sufficiently large $Q$ choose the largest positive
$l\equiv1\pmod4$ with
\[
 l^2e\le\tfrac1{10}\log(lQ).
\]
There is such a multiplier eventually, and the set is finite since
$l^2/\log(lQ)\to\infty$ at fixed $Q$. Moreover $l\to\infty$:
every fixed admissible multiplier eventually satisfies the inequality
because $e<9/2$. Failure at $l+4$ gives the explicit two-sided bound
\begin{equation}\label{inf:multiplier-two-sided}
 \left(\frac l{l+4}\right)^2\frac{\log(lQ)}{10}
 <l^2e\le\frac{\log(lQ)}{10}.
\end{equation}
Set $P'=lP$, $Q'=lQ$, $p'=(P'-3)/2$ and let $p_0'$ be the
crossing with index $(Q'-1)/4$. Then $Q'\equiv1\pmod4$ and the
parity of $2p'$ is preserved. The crossing expansion gives
\[
 p_0'(p'-p_0')
 =-\frac{\alpha_{\mathrm{ar}}}{4}l^2Q(Q\alpha_{\mathrm{ar}}-P)
          +d_{\mathrm{ar}}+o(1).
\]
Since $l/(l+4)\ge1/5$, it follows that
\[
 \frac{\alpha_{\mathrm{ar}}}{1000}\log Q'-C
 \le p_0'|p'-p_0'|
 \le\frac{\alpha_{\mathrm{ar}}}{40}\log Q'+C.
\]
Here $p'/Q'\to\alpha_{\mathrm{ar}}/2$ and
$p_0'/p'\to1$. As $\alpha_{\mathrm{ar}}/40<1/8$, the upper bound
puts these inputs in $\mathcal P_{\log}$ eventually; the lower
bound tends to infinity. Equation~\eqref{inf:log-detuning-relation}
then gives $\tau\to\infty$. Successive primary crossings are
separated by $2\alpha_{\mathrm{ar}}+o(1)$, so another crossing
cannot place these same inputs in a fixed-width family.
\end{proof}

\begin{theorem}[Newton closure at order twenty]\label{inf:fixed-order-newton}
Fix $5/3<q_{\mathrm{gap}}<2$ and
$0<\xi_{\mathrm{inv}}<2-q_{\mathrm{gap}}$, a compact split interval,
and a positive gap constant. At every sufficiently large logarithmic
input whose quartic domain has gap at least
$\gamma=\kappa_{\mathrm{gap}}\tau p^{-q_{\mathrm{gap}}}$,
the order-$20$ centre has an exact real zero within
\[
 \nrm{z_*-z_0}_{\Xc_p}\le\tau p^{-11}.
\]
The zero gives a bounded strictly convex non-ball analytic domain
with the frequency-one Schiffer data.
\end{theorem}
\begin{proof}
Fix the parameters in the statement. Choose
$\alpha_{\mathrm{loc}}>0$ and an internal loss smaller than
$\xi_{\mathrm{inv}}$ as in Theorem~\ref{inf:mixed-conditioning}.
Fix the angular window constant and then $K_{\mathrm{core}}$ to
control the global Green sums. Fix the radial head cutoff.
Take $N=20$ and fix the complex angular disc and coefficient moments
needed by the centre, clamping and material estimates.
These choices all precede the threshold in $p$.

Lemma~\ref{inf:detuning-dependencies} supplies the true-centre
hypothesis of Theorem~\ref{inf:mixed-conditioning} and the
same-space residual and material factorization. Its row 12 supplies
the sharp material and nonlinear bounds uniformly in this range.
They give
\[
 \nrm{\AN}\le\tau^{-1}p^{B_{\log}+o(1)},\qquad
 B_{\log}=\tfrac92+q_{\mathrm{gap}}+\xi_{\mathrm{inv}}<\tfrac{13}2,
 \qquad \sup\nrm{D^2\Fnl}\le p^{3+o(1)},
\]
where $\AN=\Qmat^{-1}J_D^{-1}$ is a bounded isomorphism.
On the closed ball of radius $\mathfrak r_p=\tau p^{-11}$,
the three contraction quantities are bounded by
\begin{align}
 \nrm{\AN\Fnl(z_0)}/\mathfrak r_p
 &\le p^{B_{\log}-7+o(1)}\le p^{-1/2+o(1)},
 \label{inf:log-gate-one}\\
 \nrm{I-\AN D\Fnl(z_0)}
 &\le\tau p^{B_{\log}-17+o(1)}
       \le(\log p)p^{-21/2+o(1)},\label{inf:log-gate-two}\\
 \sup\nrm{\AN D^2\Fnl}\,\mathfrak r_p
 &\le p^{B_{\log}-8+o(1)}\le p^{-3/2+o(1)}.
 \label{inf:log-gate-three}
\end{align}
All tend to zero uniformly in $0<\tau\le\log p$.
The source displacement is below one and the total shape norm
below $1/4$, so the entire closed ball is in the working region.
For large $p$ the map $z\mapsto z-\AN\Fnl(z)$ maps that ball
to itself and has Lipschitz constant at most $1/2$, by the same
Taylor integral inequalities used in Theorem~\ref{inf:newton-threshold}.
Completeness and injectivity of $\AN$ give an exact real zero.

Its normalized $C^2$ shape is $O(\tau p^{-2})+O(\tau p^{-11})$,
which tends to zero. The physical seed coefficient is
$-\theta/(16r(1-r)p)+O(\tau p^{-2}\mathfrak P_{20}(\log p))$;
the correction changes it by at most $C\tau p^{-10}$.
Division by $\tau/p$ makes both errors tend to zero.
Thus the curvature and non-ball arguments in
Section~\ref{inf:sec-geometry} apply. The completed-space
identification in Proposition~\ref{inf:classical-solution} supplies
the classical field and its two boundary traces in the integer
dimension $2p$.
\end{proof}

\begin{proposition}[Assembly of the logarithmic family]\label{inf:log-assembly}
For every fixed $0<\varepsilon_{\mathrm{sp}}<1/4$ and $0<\eta_{\mathrm{cnt}}<1/3$,
every sufficiently large $p\in\mathcal P_{\log}$ has the domains of
Theorem~\ref{inf:fixed-order-newton} at all but
$O_{\mathrm{sp},\eta_{\mathrm{cnt}}}(p^{2/3+\eta_{\mathrm{cnt}}})$ interior integer splits.
Distinct retained splits give pairwise non-similar domains.
The similarity-class lower limit is at least $1/2$ after division
by $n=2p$, along this family's dimensions.
\end{proposition}
\begin{proof}
Choose
\[
 q_{\mathrm{gap}}=2-\eta_{\mathrm{cnt}}/2,\qquad
 \xi_{\mathrm{inv}}=\eta_{\mathrm{cnt}}/4,\qquad
 \alpha_{\mathrm{loc}}=\eta_{\mathrm{cnt}}/16.
\]
These obey all strict inequalities in the conditioning argument.
Definition~\ref{inf:log-family} and
Lemma~\ref{inf:detuning-dependencies} supply its arithmetic,
low-state and uniform-constant hypotheses.
Corollary~\ref{inf:log-flow-count} gives the selected quartic gaps
and at most $C\sqrt{\log p}\,p^{2/3+\eta_{\mathrm{cnt}}/2}
=O(p^{2/3+\eta_{\mathrm{cnt}}})$ exceptions.
For every retained split apply Theorem~\ref{inf:fixed-order-newton}
at the single fixed accuracy twenty. All its thresholds are
uniform over the compact split interval. Lemma~\ref{inf:similarity-classes}
proves non-similarity and the lower-limit conclusion by the same
integer count and then by letting the fixed margin tend to zero.
\end{proof}

\section{Shape and scope}\label{inf:sec-shape}

\begin{proposition}[Quartic shape and Hermite limit]\label{inf:shape}
Fix a construction with a compact interior split interval and a fixed
centre order. For its domains
write the physical boundary as $\varrho=R+h(x)$ and set
$v_r=r(1-r)$. In either family, uniformly in every fixed $C_x^k$ norm,
\begin{equation}\label{inf:shape-leading}
 h(x)=-\frac{\theta}{16v_rp}\mathsf H_2(x)+O(\tau p^{-2}),
 \qquad
 \frac{p^2}{\theta}\left(\frac{\varrho(x)}R-1\right)
 =-\frac{(x-r)^2}{32v_r}+O(p^{-1}).
\end{equation}
For $y$ in a fixed compact interval the concentration-scale limit is
\begin{equation}\label{inf:Hermite-limit}
 \frac{p^2}{\theta}h(r+\sqrt{v_r/p}\,y)
 =-\frac{\operatorname{He}_2(y)}{16}+O(p^{-1/2}),
 \qquad \operatorname{He}_2(y)=y^2-1.
\end{equation}
The constants may depend on the fixed centre order in the bounded
family; the logarithmic family uses order twenty.
\end{proposition}
\begin{proof}
For fixed bounded detuning, \eqref{inf:centre-profile} and the
cofactor Newton radius give
\[
 h=-(2\chi/g_2)\mathsf H_2+O(\tau p^{-3}).
\]
For logarithmic detuning the same formula holds with remainder
$C\tau p^{-3}\mathfrak P_{20}(\log p)$ by row 8 of
Lemma~\ref{inf:detuning-dependencies}; the physical Newton correction
is $O(\tau p^{-10})$.
The ellipse estimate in the proof of Lemma~\ref{inf:space-embedding}
is $|G_j(z)|\le2(j+1)\rho_{\mathrm{ell}}^j$ for fixed
$1<\rho_{\mathrm{ell}}<2$. Cauchy's formula on a smaller ellipse
therefore bounds each fixed $x$ derivative by
$C_k(j+1)\rho_{\mathrm{ell}}^j$, absorbed by the coefficient
weight $(j+1)2^j$. Thus these comparisons hold in every fixed
$C_x^k$ norm.
Now $\chi=\theta/(4Rp)$, $R=2p+3+O(\log p/p)$ and
$g_2=4v_r/p+O(p^{-2})$.
Fixed powers of $\log p$ divided by $p$ tend to zero.
These estimates prove \eqref{inf:shape-leading}.

The exact identity following from \eqref{inf:H2} is
\[
 \mathsf H_2(x)=(x-r)^2-
 \frac{2(1-2r)}{p+2}(x-r)-\frac{v_r}{p+1}.
\]
At $x=r+\sqrt{v_r/p}\,y$ it equals
$v_rp^{-1}(y^2-1)+O(p^{-3/2})$ on compact $y$ intervals.
The sharper quartic comparison above, multiplied by $p^2/\theta$,
has error at most $C\mathfrak P_{20}(\log p)/p=O(p^{-1/2})$ in
the logarithmic family, and $O(p^{-1})$ in the bounded family.
This proves \eqref{inf:Hermite-limit}. Every error before division
contains $\tau$, so the assertions are uniform as $\tau\downarrow0$.
\end{proof}

\begin{remark}[Trace differentiation on two zero curves]
\label{inf:trace-curve-derivative}
Fix an odd integer $m\ge3$. Suppose two actual simple Bessel zero
curves $R_N(p)=j_{p,k_N}$ and $R_D(p)=j_{p+m,k_D}$ meet at
$(p_0,R_0)$, where $p_0\to\infty$ and
$R_0/p_0\to\sec\beta_{\mathrm{tr}}$,
$0<\beta_{\mathrm{tr}}<\pi/2$. The indices $k_N,k_D$ are fixed on each local
continuation and may vary with the chosen crossing.
Define along the Neumann curve
\[
 \chi_m(p)=\frac{J_{p+m}(R_N(p))}
 {J'_{p+m}(R_N(p))-(p-1)J_{p+m}(R_N(p))/R_N(p)}.
\]
Then
\begin{equation}\label{inf:true-zero-derivative}
 \begin{aligned}
 \chi_m'(p_0)&=R_N'(p_0)-R_D'(p_0),\\
 \left.\frac d{dp}(4R_N(p)p\chi_m(p))\right|_{p=p_0}
 &=\left(\frac{4m(\tan\beta_{\mathrm{tr}}-\beta_{\mathrm{tr}})
                    \cos\beta_{\mathrm{tr}}}{\sin^3\beta_{\mathrm{tr}}}
        +o(1)\right)p_0.
 \end{aligned}
\end{equation}
For the primary quartic crossing this coefficient is $16\delta_0$.
The detuning comparison in Lemma~\ref{inf:detuning} follows
from the exact quotient \eqref{inf:trace-ratio-exact}.
\end{remark}
\begin{proof}
At a Dirichlet zero the denominator of the trace quotient equals
$J'_{p+m}(R_D(p))\ne0$. Its partial derivative with respect to
its radius argument is therefore one. Differentiating its zero
value along the actual Dirichlet curve gives partial order
derivative $-R_D'(p_0)$. The chain rule on $R_N$ proves the first
identity.
By \eqref{inf:order-derivative},
\[
 R_N'-R_D'
 =f_{\mathrm{ord}}(p_0/R_0)
       -f_{\mathrm{ord}}((p_0+m)/R_0)+O(p_0^{-2}).
\]
For $x=\cos\beta_{\mathrm{tr}}$ one has
$f'_{\mathrm{ord}}(x)=-\cos\beta_{\mathrm{tr}}
(\tan\beta_{\mathrm{tr}}-\beta_{\mathrm{tr}})/\sin^3\beta_{\mathrm{tr}}$.
Taylor expansion therefore makes the difference
$\{m(\tan\beta_{\mathrm{tr}}-\beta_{\mathrm{tr}})
\cos^2\beta_{\mathrm{tr}}/\sin^3\beta_{\mathrm{tr}}+o(1)\}/p_0$.
Since $\chi_m(p_0)=0$, differentiating its prefactor $4R_Np$
adds zero. Multiplication by $4R_0p_0$ proves the second identity.
For $m=3$, $\beta_{\mathrm{tr}}=\pi/3$, it reduces to $16\delta_0$.
\end{proof}

\begin{proposition}[Density of the arithmetic families]
\label{inf:arithmetic-density}
Each fixed-width candidate family $\mathcal P_{\mathrm{fix}}(C_{\mathrm{ap}})$ and
the candidate family $\mathcal P_{\log}$ has natural density zero
in either parity lattice. Their dimension sets have density zero
in $\mathbb N$.
\end{proposition}
\begin{proof}
The crossing expansion shows that membership implies
\[
 \left\|\frac{c_{\mathrm{ar}}}{\pi}p+b_{\mathrm{ph}}\right\|_{\R/\Z}
 \le C\frac{1+\log p}{p}
\]
with $b_{\mathrm{ph}}=3\sqrt3/(2\pi)-3/4$; a fixed-width family
has the stronger numerator $C$.
The frequency $c_{\mathrm{ar}}/\pi=\sqrt3/\pi-1/3$ is irrational
by transcendence of $\pi$. On either lattice $p=k$ or $k+1/2$,
each nonconstant integer Fourier mode of this rotation has a bounded
geometric partial sum, so its average tends to zero.
Continuous trigonometric majorants and minorants of a fixed arc
\cite{inf:zygmund} then give its frequency equal to its length. For any fixed
$0<\zeta<1/2$, all sufficiently late candidate points lie in
the fixed arc of radius $\zeta$ because $(1+\log p)/p\to0$.
Their upper density is at most $2\zeta$. Letting $\zeta$ tend to
zero proves the assertion on both lattices; combining them and
rescaling by two proves the dimension statement.
\end{proof}

The existence assertions use the quartic angular seed. Extending this
construction to other angular seeds requires a centre construction,
a signed-flow estimate for its wider angular interaction, and a
closure in the resulting amplitude scale. The density statement
above describes the listed arithmetic families; it leaves the set
of all dimensions supporting Schiffer domains undetermined.

\begin{lemma}[The spectral cap is inactive]\label{inf:cap-inactive}
For every sufficiently large input in either family and every
candidate split in a fixed compact interval, the actual quartic
invariant Dirichlet gap is less than $R^2$. Thus
\[
 \min\{R^2,\dist(R^2,\spec H_a)\}
                  =\dist(R^2,\spec H_a),
\]
where $H_a$ is the quartic Dirichlet operator in unit-ball normalization.
\end{lemma}
\begin{proof}
In the radial angular-zero sector of the ball, the sine trial on
$(1/2,1)$ gives
$\lambda_1(B)\le4(p-1)^2+4\pi^2$.
The quartic shell has relative width $\delta\le C_{\mathrm{sp}}\tau/p^2$.
Min--max and dilation imply
\[
 0<\lambda_1(H_a)\le4(p-1)^2+4\pi^2+C_{\mathrm{sp}}\tau
       <(2p-1)^2\le R^2
\]
for fixed $T$ or $\tau\le\log p$ and sufficiently large $p$.
The gap is at most $R^2-\lambda_1(H_a)<R^2$, proving the claim.
\end{proof}

\begin{remark}[Choosing a split from strict gap enclosures]
\label{inf:split-selection}
For the finite set of candidate splits, first enclose the strict
conditions $\eta_{\mathrm{alg}}<1/2$ of
Lemma~\ref{inf:quartic-computable-relative}. They hold for all
candidates at sufficiently large inputs. Lemma~\ref{inf:gap-enclosures}
then supplies rational enclosures $[L_a,U_a]$ of their capped gaps
with widths tending to zero. At each precision set $U=\max_aU_a$
and retain the least $a$
with $L_a>0$ and $2L_a>U$, when one exists.
For sufficiently large family inputs, let $G$ be the maximum actual
gap. Lemma~\ref{inf:cap-inactive}
makes the caps inactive. Whenever $G>0$, this selection terminates:
an index attaining $G$ has $L_a\to G$, while $U\to G$.
Its output has actual gap at least $L_a>G/2$.
For fixed-detuning inputs Theorem~\ref{inf:integer-gap}, with the
fixed-parameter extension, gives $G\ge c_T\tau p^{-5/3}$.
For logarithmic inputs Corollary~\ref{inf:log-flow-count} gives,
for each fixed $5/3<q_{\mathrm{gap}}<2$,
$G\ge c_q\tau p^{-q_{\mathrm{gap}}}$. Either bound supplies the
positive gap and the appropriate closure. This strict-comparison
rule needs no decision of equality between two spectral gaps.
\end{remark}

\appendix
\section{Notation index}\label{inf:sec-notation}

The split parameter is $\varsigma$; the radial Jacobi index is $s$.
Monic angular polynomials use the font $\mathsf H$, while $H_0$
is an operator. Physical radii use $\varrho$; the squared reference
radius is $t$ and the nonsquared one is $\zeta$.
Constants written as $C$ may increase within a proof. A subscript
records permitted dependence. Unless that dependence is stated,
constants in the bounded base chain are independent of $p$, the split,
and $0<|\theta|\le20$. Section~\ref{inf:sec-widened} allows dependence
on its fixed compact split interval and detuning bound;
Lemma~\ref{inf:detuning-dependencies} gives the explicit envelopes
for the logarithmic chain.

\begingroup\small
\renewcommand{\arraystretch}{1.16}
\begin{longtable}{@{}>{\raggedright\arraybackslash}p{0.29\textwidth}>{\raggedright\arraybackslash}p{0.65\textwidth}@{}}
\toprule Symbol & Meaning and first definition\\\midrule
\endfirsthead
\toprule Symbol & Meaning and first definition\\\midrule
\endhead
$n=2p$ & Ambient dimension; $p\in\tfrac12\Z$ and $n\in\mathbb N$ in the final construction.\\
$\mathcal P,\ (p_m)$ & Selected integer set and its increasing enumeration; Definition~\ref{inf:integer-sequence}.\\
$\widehat p_m,\ p_0$ & Real primary crossing orders for $m\ge m_0$; $p_0$ is associated to the current half-lattice input.\\
$P_j^{\mathrm{cf}},Q_j^{\mathrm{cf}}$ & Continued-fraction numerator and denominator; selected pairs without an index are $(P,Q)$.\\
$T_j^{\mathrm{cf}},\rho_j$ & Continued-fraction tail $[a_{j+1};a_{j+2},\ldots]$ and denominator ratio $Q_{j-1}^{\mathrm{cf}}/Q_j^{\mathrm{cf}}$ in \eqref{inf:cf-error}.\\
$c_{\mathrm{ar}},\alpha_{\mathrm{ar}},d_{\mathrm{ar}},\delta_0$ & Arithmetic and detuning constants in \eqref{inf:arithmetic-constants}.\\
$R_{\mathrm{alg}}(v),A_m$ & Algebraic crossing radius and phase centre in \eqref{inf:algebraic-radius}.\\
$a,b$ & Integer block sizes, $a+b=2p$; the base range is $p/2\le a\le3p/4$, extended in Section~\ref{inf:sec-widened}.\\
$\ab,\bb$ & Beta parameters $a/2,b/2$.\\
$r,\varsigma$ & Split ratio $a/(2p)$ and flow parameter $(r-1/2)^2$.\\
$X,Y$ & Cartesian block vectors in $\R^a,\R^b$.\\
$x,\ t,\ \zeta,\ \varrho$ & Angular coordinate $|X|^2/|y|^2$, reference squared radius, reference radius, and physical radius.\\
$q_X,q_Y$ & The invariant polynomials $|X|^2,|Y|^2$.\\
$\mu_r$ & Normalized beta measure, \eqref{inf:beta-measure}.\\
$\mathsf H_j,G_j,P_j$ & Monic, endpoint-normalized and orthonormal angular Jacobi polynomials; \eqref{inf:angular-normalization}.\\
$\eta_j$ & Angular Hilbert norm of $G_j$; \eqref{inf:angular-norm}.\\
$\mathfrak a_j,\mathfrak b_j$ & Jacobi multiplication recurrence coefficients; \eqref{inf:recurrence-a}--\eqref{inf:recurrence-b}.\\
$\mathsf a_k,\mathsf b_k,\mathsf e_k$ & Leading sequence, first correction and remainder of Lemma~\ref{inf:low-recurrence}.\\
$Q_{\mathrm e},Q_{\mathrm o},E(A,B)$ & Even and odd denominators, and the error product, in the proof of Lemma~\ref{inf:congruence}.\\
$S_j$ & Solid harmonic $t^jG_j(x)$ of degree $2j$.\\
$\Lambda,\lambda_j$ & Positive angular operator and eigenvalues $4j(j+p-1)$.\\
$\nu_j,\beta_j$ & The common number $p-1+2j$, used respectively as Bessel order and radial Jacobi parameter.\\
$J_\nu,Y_\nu,H^{(1)}_\nu$ & Ordinary Bessel functions of the first and second kind, and their outgoing Hankel combination.\\
$\mathrm K_\nu$ & Modified Bessel function of the second kind; its roman font distinguishes it from the graph operator $K_0$.\\
$j_{\nu,k},z_{j,k}$ & Positive Bessel zeros and the separated ball radii $j_{p-1+2j,k}$.\\
$J_s^\beta,\Psi_{j,s}$ & Radial Jacobi polynomial and full invariant polynomial basis, \eqref{inf:radial-basis}.\\
$A_{\mathrm{Deb}},\mathcal A_\nu$ & Scalar Debye amplitude and its normalization; \eqref{inf:debye-definitions}, Lemma~\ref{inf:debye-first}.\\
$\mathsf R_\nu^{\mathrm{Deb}}$ & Finite Debye remainder bound; \eqref{inf:debye-remainder}.\\
$\Xi(\nu,z),b(\nu,z)$ & Corrected-phase ingredients in \eqref{inf:phase-def}.\\
$\varphi_p(j)$ & Lattice phase graph in \eqref{inf:phase-graph}.\\
$R$ & Continued radial Neumann zero $J_p(R)=0$; error $O(p^{-1})$ at fixed detuning and $O(\log p/p)$ in the logarithmic family.\\
$\chi,\theta,\tau$ & Seed trace ratio, signed detuning $4Rp\chi$, and magnitude $|\theta|$.\\
$\mathfrak c_{ij}^k,c_{ij}^k$ & Monic and endpoint-normalized angular product coefficients, respectively.\\
$g_2$ & Monic quartic self-coupling $[\mathsf H_2]\mathsf H_2^2$.\\
$h_{\mathrm q},f_{\mathrm q}$ & Physical and normalized quartic boundary displacement; \eqref{inf:quartic-domain}.\\
$h_N,f_N$ & Physical finite-centre displacement and its normalized version $h_N/R$.\\
$N,\eps$ & Fixed accuracy order and the formal small parameter $1/p$.\\
$R_{\mathrm{an}},\rho',\rho_{\mathrm{ell}}$ & Prescribed analytic defect disc radius, stronger coefficient radius $\rho'>2$, and boundary analyticity ellipse parameter $1<\rho_{\mathrm{ell}}<2$.\\
$\rho,Z,S_0$ & Scaled centre radius $R/p$, its square, and the polynomial in \eqref{inf:rho-equation}.\\
$\SeedA,\SeedA_v,a_j$ & Formal scaled seed amplitude, its coefficients, and finite-mode field amplitudes.\\
$\mathcal I_{\mathrm{res}},c_j$ & Fixed resonant index set $\{2,5,8,\ldots\}$ and positive nonseed divisors in \eqref{inf:resonant-trace}.\\
$T_\varsigma,\epsilon_{\mathrm q},\rho_\varsigma$ & Quartic angular multiplier, $-\theta/(2R^2)$, and normalized trial radius $1+\epsilon_{\mathrm q}T_\varsigma$.\\
$\mathsf W,w_j$ & Flow degree-weight matrix and $(j+1)/(p+j+1)$.\\
$\mathsf D,\Euler$ & Mass-unitary radial generator $\zeta\partial_\zeta+1/2$, and ordinary Euler derivative $y\cdot\nabla$.\\
$H_\varsigma$ & Exact unitary quartic Dirichlet operator in the flow representation.\\
$\mathfrak q_\varsigma$ & Quartic Dirichlet form; \eqref{inf:flow-formula}.\\
$\mathfrak g_\gamma,\mathcal G(\gamma)$ & Quadratic correction form and its weak form operator; Proposition~\ref{inf:flow-form}.\\
$c_j^{\mathrm{diag}},d_j^{\mathrm{diag}}$ & Rational diagonal coefficients in \eqref{inf:T-diagonal}.\\
$\mathcal N_j,\mathcal E_j$ & Angular variation numerators; \eqref{inf:N-angular}, \eqref{inf:E-angular}.\\
$j_{\mathrm{pos}},v,D_{\mathrm{pos}}$ & The shifts $j-7$ in Lemma~\ref{inf:angular-positive}, $j-8$ in Appendix~\ref{inf:sec-positivity}, and its positive denominator.\\
$\PP,\PP_H,L$ & Flow-window spectral projection, its high block, and the two-state block $j=2,5$.\\
$T_{\PP}^{\mathrm{ang}}$ & Angular principal submatrix, without radial overlaps.\\
$\mathcal S_\varsigma$ & Actual compressed transport matrix, including radial overlaps.\\
$D_\PP,Z_\PP,D_H,Z_H$ & Window energy offsets and normalized boundary frequencies, and their high-block restrictions; Lemmas~\ref{inf:transport}, \ref{inf:flow-low-block}.\\
$\mathfrak C_\varsigma$ & Normalized flow perturbation $K_0^{-1/2}(H_\varsigma-H_0)K_0^{-1/2}$; proof of Lemma~\ref{inf:flow-complement}.\\
$\mathfrak U_0,\mathfrak U$ & Unperturbed and exact normalized complementary inverses in that proof.\\
$\mathcal R_{2,\varsigma},\mathcal R_\varsigma$ & Quadratic form remainder and exact flow Schur remainder; Lemmas~\ref{inf:flow-relative}, \ref{inf:flow-complement}.\\
$F_{\PP},F_H$ & Exact flow Schur matrices, at spectral parameter $\lambda$.\\
$\omega_p$ & Local-flow half-bandwidth $\tau/(16p)$.\\
$\epsilon_{\mathrm{sgn}}$ & Sign of the eigenvalue flow, fixed to $\operatorname{sign}\theta$ in the grid argument.\\
$\gamma_p,c_{\mathrm{gap}}$ & Centre Dirichlet gap $c_{\mathrm{gap}}\tau p^{-5/3}$, with $c_{\mathrm{gap}}=c_{\mathrm q}/2$.\\
$\Cc,\Cc^{[k]},\wt$ & Source coefficient space, its degree moments, and grade weight $(1+\ell/p^2)^p$.\\
$\rho_w$ & Fixed angular coefficient radius two in Section~\ref{inf:sec-spaces}.\\
$\Gc_p,\Ec_1,\Xc_p$ & Clamped source subspace, one-moment shape space, and their product.\\
$\HP$ & Harmonic extension from the sphere by $G_j\mapsto S_j$.\\
$\Mmix,\Hc$ & Mixed radial Hilbert/angular coefficient space and true invariant $L^2$ space.\\
$\Hc_{\mathrm{flow}}$ & Mass-unitary representation of $\Hc$ with radial measure $d\zeta$.\\
$\sigma$ & Exponential angular matrix-norm parameter; fixed to one in the inverse proof.\\
$\KD,K$ & Dirichlet inverse of $\Delta$, and its restriction to clamped sources.\\
$d_s,b_0,k_s$ & Poisson denominators $\beta+2s,\beta+1$ and radial column profile; Lemma~\ref{inf:poisson-columns}.\\
$\ell_{\mathrm{out}},Z_{\mathrm{grad}}$ & Output harmonic degree and factor $h+\ell_{\mathrm{out}}+p-1$; Lemma~\ref{inf:gradient-profiles}.\\
$L_j^D,L_j^N$ & Unit Dirichlet and normal double-trace lift columns; \eqref{inf:lift-columns}.\\
$\kf,\mathfrak h$ & General radial Poisson profile in Lemma~\ref{inf:mixed-Poisson}, and $q-1$ in Proposition~\ref{inf:Hilbert-relative}.\\
$\delta$ & Strong centre size; \eqref{inf:strong-centre-norm}.\\
$H_0,K_0$ & Positive ball Dirichlet Laplacian and shift $H_0+R^2$.\\
$\Phi_f,q,d$ & Scalar harmonic chart and its factors in \eqref{inf:scalar-chart}.\\
$P_\Phi\Delta,B_\Phi,B_{\Phi,D}$ & Scalar pullback Laplacian, its Helmholtz shift, and its restriction to the Dirichlet graph domain; \eqref{inf:scalar-pullback-definition} and following paragraph.\\
$J_D$ & Scalar Dirichlet graph operator $B_{\Phi_N,D}\KD$ at the fixed centre.\\
$\varepsilon_0,\varepsilon_{\mathrm{work}}$ & Analytic chart radius $1/2$ and closed working shape radius $1/4$.\\
$A_q,M_q,C_q$ & Pullback metric, gradient square, and contracted Hessian; \eqref{inf:pullback-metric}.\\
$m,H_\Phi,\Hu,\mathcal V$ & Jacobian square root, scalar Dirichlet pullback, its unitary realization, and $\Hu-H_0$.\\
$D_N,N_N^{\mathrm{phys}}$ & Physical boundary defects of a finite centre.\\
$N_N^{\mathrm{ref}},\ell_N$ & Reference normal defect and exact double-trace lift.\\
$U_*,U_0,g_0,z_0$ & Pulled-back finite field, clamped field, its source and centre pair; \eqref{inf:clamped-centre}.\\
$d_j^{\mathrm{mode}}$ & Value or normal mode normalization; Lemma~\ref{inf:radial-centre-norm}.\\
$\epsilon_{\mathrm{cmp},p}$ & Known-centre composition multiplier size in \eqref{inf:composition-bound}.\\
$\epsilon_{\mathrm{rec}}$ & Fixed-centre inverse-series ratio in the proof of Proposition~\ref{inf:residual-estimates}.\\
$\Fnl,E_0,\mathcal V_{\mathrm{res}}$ & Nonlinear coefficient map, centre residual and material residual map $\Xc_p\to\Cc$; $\mathcal V=\Hu-H_0$ instead acts on the Dirichlet graph domain.\\
$D_p,\mathfrak f_0,T_\eta$ & Normal Hessian, fixed material field and material shape variation; \eqref{inf:material-definitions}.\\
$\Qmat$ & Exact material isomorphism, \eqref{inf:Qmat-definition}.\\
$\mathcal A,\mathcal D_0$ & Centred unitary and ball operators $\Hu-R^2,H_0-R^2$.\\
$\Pi,Q,\mathcal A_Q$ & Inverse-window projection, its complement, and exact complementary operator.\\
$\Feff$ & Exact Schur matrix for the coefficient inverse, at spectral parameter zero.\\
$M,D_{\mathrm{sep}}$ & Cluster cardinality bound and logarithmic separation constant.\\
$\Pi_T,Q_T,K_{\mathrm{head}}$ & Coefficient head/tail projections and their fixed degree cutoff.\\
$b_{\mathrm{inv}},B_{\mathrm{Newt}},L_{\mathrm{Newt}}$ & Inverse and Newton exponents; \eqref{inf:constant-order}, \eqref{inf:Newton-correction}.\\
$m_{\mathrm{mat}},d_{\mathrm{nl}}$ & Material-inverse and nonlinear second-derivative exponents in Theorem~\ref{inf:newton-threshold}; both equal six in the application.\\
$\AN,\mathfrak r_p$ & Frozen Newton inverse and contraction radius $\tau p^{-L_{\mathrm{Newt}}}$.\\
$C_0,C_{\mathrm{rel}}$ & Absolute centre-size and relative graph constants; \eqref{inf:uniform-delta}, \eqref{inf:relative-two-norms}.\\
$E_{\mathrm{inv}},c_t$ & Fixed inverse-window factor, and fixed turning-spacing margin $0<c_t<\pi/\sqrt2$ in \eqref{inf:constant-turning}.\\
$\mathfrak W_p$ & Physical monic non-ball witness $R[\mathsf H_2]f$.\\
$\mathcal P_{\mathrm{fix}}(C_{\mathrm{ap}}),\mathcal P_{\mathrm{bd}}$ & Full half-lattice fixed-width input sets, with $\mathcal P_{\mathrm{bd}}=\mathcal P_{\mathrm{fix}}(31/10)$; Definition~\ref{inf:input-families}.\\
$\mathcal P_{\log}$ & Logarithmic input family, Definition~\ref{inf:log-family}.\\
$\varepsilon_{\mathrm{sp}},I_{\mathrm{sp}}$ & Fixed interior split margin and its $\varsigma$ interval. The formal centre parameter remains $\eps=1/p$.\\
$K_{\mathrm{sp}},W_{\mathrm{sp}}$ & Quartic multiplier and enlarged counting-window constants, Lemma~\ref{inf:split-extension}.\\
$C_{\mathrm{ap}},T$ & Fixed scaled-distance and detuning bounds; $T$ is allowed to grow only under Lemma~\ref{inf:detuning-dependencies}.\\
$L_{\mathrm{ap}},b_{\mathrm{ap}},M_{\mathrm{ap}}$ & Arithmetic modulus, residue and $4L_{\mathrm{ap}}$, Theorem~\ref{inf:residue-arithmetic}.\\
$\mathcal R_{\mathrm{split}},r_0,n_k$ & Full-measure set of fixed ratios, a chosen ratio, and a continued-fraction-indexed selected dimension family; Corollary~\ref{inf:fixed-ratio}.\\
$q_{\mathrm{gap}},\xi_{\mathrm{inv}},\alpha_{\mathrm{loc}}$ & Gap exponent, conditioning loss and critical-neighbourhood exponent.\\
$\eta_{\mathrm{cnt}}$ & Positive loss in the bad-split counting exponent; the material shape variation remains $\eta$.\\
$\gamma,\kappa_{\mathrm{gap}}$ & Gap lower bound $\kappa_{\mathrm{gap}}\tau p^{-q_{\mathrm{gap}}}$ and its fixed constant.\\
$B_{\mathrm{bd}},L_{\mathrm{bd}},B_{\log}$ & Newton inverse/radius exponents for the extended bounded chain and inverse exponent for the logarithmic chain.\\
$\mathcal N_{\mathrm{sim}}(n)$ & Number of integer splits $0<a<n/2$ admitting a bounded strictly convex non-ball analytic domain with the block symmetry and data \eqref{inf:main-equations}; one representative per split gives that many similarity classes by Lemma~\ref{inf:similarity-classes}.\\
$\mathsf F_l,\mathsf G_l$ & Rational sine--cosine coefficients for half-integer Bessel functions.\\
$\mathsf L^{\mathrm{Lom}}_{k,\mu},U_k^{\mathrm{Ch}}$ & Lommel polynomials and second-kind Chebyshev polynomials; $U_k$ without the superscript remains the Debye polynomial.\\
$\varphi_{\mathrm{Bes}}$ & Corrected Bessel phase $\Xi-\arctan b$, \eqref{inf:corrected-phase}.\\
$\Pi_{\mathrm{rad}},Q_{\mathrm{rad}},z_j,d_j^{\mathrm{pole}}$ & One-pole-per-angular-block projection, its complement, retained zero, and retained squared-energy offset.\\
$z_{\mathrm q},H_{\mathrm q},L_{\mathrm q}$ & Quartic harmonic coefficient, harmonic lift $z_{\mathrm q}S_2$, and linear quartic unitary operator minus $R^2$.\\
$\widehat H_H,V_1(H),l,l_{\mathrm{lin}}$ & Fixed polynomial-chart unitary operator, its first variation, log density and its linear part; Lemma~\ref{inf:quadratic-graph}.\\
$\mathcal R_{\mathrm{sc}}$ & Fixed-centre unshifted scalar graph perturbation in Proposition~\ref{inf:sharp-material}.\\
$\mathcal S_{\mathrm{pole}},\mathcal S_{\mathrm{pole}}^{(0)}$ & Exact radial-pole Schur matrix and its second-order interaction truncation of bandwidth four.\\
$\mathcal B_{\mathrm{pole},2},\mathcal B_{\mathrm{pole},4},\mathsf R_{\mathrm{pole},6}$ & Pole bands and all paths of at least three interactions in \eqref{inf:pole-schur-bands}; the remainder need not have bandwidth six.\\
$W_{\mathrm{cond}},\mathcal I_{\mathrm{win}},\mathcal F_{\mathrm{win}}$ & Conditioning angular window width, retained index set, and Schur complement of $\mathcal S_{\mathrm{pole}}^{(0)}$ onto that set.\\
$\mathcal T_{\mathrm{crit}},c_{\mathrm{cr}}$ & Four critical phase ratios and a member of this set; $c$ without a subscript remains a generic positive constant.\\
$\mathcal N_{c_{\mathrm{cr}}},f_{c_{\mathrm{cr}}}$ & Critical angular neighbourhood and phase plus its rational slope.\\
$\ell_{c_{\mathrm{cr}}},d_{c_{\mathrm{cr}}},j_{c_{\mathrm{cr}}},j_{c_{\mathrm{cr}}}^*$ & Reduced critical angle pair, uncorrected stationary index, and true corrected-phase minimum.\\
$h,\mathfrak e$ & In the same-level high-core argument, the spacing $\pi/d_{c_{\mathrm{cr}}}$ of the union of residue levels and excursion $g-\min_{[u,v]}f_{c_{\mathrm{cr}}}$.\\
$c_{\mathrm{curv}},C_{\mathrm{curv}}$ & Fixed lower and upper phase-curvature constants in Lemma~\ref{inf:critical-strips}.\\
$K_{\mathrm{core}},L_{\mathrm{core}}(p)$ & Fixed separation constant and $K_{\mathrm{core}}\log p/\log\log p$; unrelated to the normal Hessian $D_p$.\\
$\mathfrak u_j,\mathfrak v_j,a_{\mathrm{quad}}$ & Endpoint-normalized Jacobi recurrence coefficients and $a_{\mathrm{quad}}=2(\ab+1)/(p+2)$, Lemma~\ref{inf:path-valuation}.\\
$\mathfrak P_N(T),\mathcal B_N(T),d_N^{\mathrm{tame}},k_N$ & Polynomial and exponential-square-root envelopes, their finite exponent, and the finite-jet disc exponent.\\
$\mathfrak E_N(p)$ & Subpolynomial strong-centre comparison envelope in \eqref{inf:conditioning-centre-hypothesis}.\\
$v_r,\operatorname{He}_2$ & Variance factor $r(1-r)$ and monic Hermite polynomial $y^2-1$.\\
$f_{\mathrm{ord}},R_N,R_D,\chi_m$ & Watson comparison function, actual Neumann/Dirichlet zero curves, and their trace ratio.\\
$\beta_{\mathrm{tr}}$ & Limiting trace-curve angle with $R_0/p_0\to\sec\beta_{\mathrm{tr}}$ in Remark~\ref{inf:trace-curve-derivative}.\\
$n_{\uparrow},n_{\downarrow}$ & Numbers of radial parameter-raising and lowering steps in Lemma~\ref{inf:clamped-smoothing}.\\
$\gamma_F$ & Hilbert gap in Lemma~\ref{inf:cluster-inverse}.\\
$F_{\nu,j},C_{j,\mathrm{rad}}$ & Entire fixed-mode radial factor and its radial norm constant in Section~\ref{inf:sec-material}.\\
$K_{0,\PP},K_{0,Q}$ & Restrictions of $K_0$ to the two flow spectral blocks.\\
$\mathsf W_L,F_L,u_L,\mathcal R_{HL},\mathcal R_{LH}$ & Low flow weight, Schur block, eigenvector block, and off-diagonal remainder blocks.\\
$\mathsf O,x_{\mathrm{tr}}$ & Orthogonal rigid-motion factor and translation vector in Proposition~\ref{inf:Pompeiu}.\\
$H_a,\mathcal K_a,\mathcal R_{\mathrm{alg},a}$ & Quartic unitary Dirichlet operator at split $a$, its compact positive inverse, and $(H_a-H_0)H_0^{-1}$; Appendix~\ref{inf:sec-gap-enclosures}.\\
$\mu_{\mathrm{alg}},\eta_{\mathrm{alg}}$ & Absolute quartic coefficient $|z_{\mathrm q}|$ and computable relative bound $4000p\mu_{\mathrm{alg}}$.\\
$\Pi^{\mathrm{alg}}_{J,K_{\mathrm{rad}}},\Lambda_{\mathrm{miss}}$ & Rectangular Bessel--Jacobi projection and lower bound for omitted ball eigenvalues.\\
$L_a,U_a,f_R$ & Rational capped-gap enclosures and the Lipschitz function transferring the compact inverse spectrum to those gaps.\\
$\epsilon_{\mathrm{cut}}$ & Endpoint cutoff for the positive-weight integrations in Lemma~\ref{inf:gap-enclosures}.\\
\bottomrule
\end{longtable}
\endgroup

\noindent\begin{minipage}{\textwidth}
\subsection*{The spectral windows}
\begin{center}
\begin{tabular}{@{}lll@{}}
\toprule Name & Width about $R^2$ & Purpose\\\midrule
$W_{\mathrm{flow}}$ & fixed & flow complement and derivative estimates\\
$W_{\mathrm{count}}$ & fixed & count all curves that may enter the flow band\\
$W_{\mathrm{loc}}$ & arbitrary fixed & local lattice count and clusters\\
$E_{\mathrm{inv}}\tau$ & proportional to $\tau$ & coefficient inverse Schur matrix\\
$E$ & $K_{\mathrm{flow}}(1+\tau)$ & uniform logarithmic flow\\
$W_{\mathrm{cond}}$ & $C_{\mathrm{cond}}(1+\tau)$ & mixed-conditioning cores\\
\bottomrule
\end{tabular}
\end{center}
\end{minipage}\par

In the bounded base inverse $W_{\mathrm{loc}}=20E_{\mathrm{inv}}$;
Section~\ref{inf:sec-widened} uses $W_{\mathrm{loc}}=TE_{\mathrm{inv}}$
for its fixed detuning bound $T$.
The flow and curve-count constants are fixed independently of
$E_{\mathrm{inv}}$. The small flow band $\omega_p=\tau/(16p)$
and the resulting gap $\gamma_p$ shrink with the detuning.

\subsection*{Dependence of constants}

Each row below records the uniformity convention used in its
statements. A listed threshold may increase to accommodate a
finite number of earlier thresholds. Unlisted free parameters
are quantified uniformly over the ranges in the corresponding
statement. These base-chain rows are uniform for $0<\tau\le20$.

\begingroup\small
\renewcommand{\arraystretch}{1.2}
\begin{longtable}{@{}>{\raggedright\arraybackslash}p{0.31\textwidth}>{\raggedright\arraybackslash}p{0.30\textwidth}>{\raggedright\arraybackslash}p{0.32\textwidth}@{}}
\toprule Estimate & Constants may depend on & Threshold may depend on\\\midrule
\endfirsthead
\toprule Estimate & Constants may depend on & Threshold may depend on\\\midrule
\endhead
Uniform Debye bounds & explicit $\nu,z,L$ as displayed & none beyond the stated inequalities\\
Phase strip and global count & fixed window $W$ & $W$\\
Logarithmic local count & $W_{\mathrm{loc}}$; the integer $M_W$ is independent of $D$ & $W_{\mathrm{loc}},D$\\
Cluster cardinality & $M_W$ only & $W_{\mathrm{loc}},M_WD$\\
Finite centre defects & fixed $N$, complex neighbourhood $U$, requested fixed derivative order & the same fixed choices\\
Coefficient defect transfer & $N,k,\rho'$ & $N,k,\rho'$; uniform in $\theta$ and the split\\
Quartic angular matrices & derivative order $k\le2$ and the fixed interval $I$ & $p\ge64$\\
Band form estimate & explicit bandwidth factor $(m+1)^3$ & none depending on the radial index\\
Flow complement remainder & absolute constants; independent of $W_{\mathrm{flow}}$ & fixed $W_{\mathrm{flow}}$, with $p^2\ge W_{\mathrm{flow}}$\\
Local flow margin & absolute constants fixed before $p$ & $W_{\mathrm{flow}}$ and low-degree exclusion\\
Curve count & $C_{40}$ for $W_{\mathrm{count}}=40$ & the threshold of the fixed-window count\\
Integer-split gap & $c_{\mathrm{flow}},C_{40}$, both absolute & their fixed values\\
Centre gap transfer & $N,c_{\mathrm q}$ & $N,c_{\mathrm q}$\\
Source module and clamped estimates & fixed angular radius two; displayed shape degrees only & $p\ge64$\\
Higher smooth chart norms & fixed derivative order and fixed dimension & fixed-domain estimates\\
Centre source moments & $N,k$ & $N,k$\\
Material inverse and $D^2\Fnl$ & $N$; exponents six independent of $N$ & $N$ and the fixed working radius $1/4$\\
Strong centre-size bound & absolute $C_0$, independent of $N$ & $N$ absorbs its finite-degree remainder\\
Relative graph bounds & $C_0$, fixed moment ten and fixed exponential weights & $N$, plus absolute smallness conditions\\
Complementary inverse & fixed $E_{\mathrm{inv}}$ and relative constant & those constants and $N$\\
Cofactor inverse & $M,c_{\mathrm{gap}},E_{\mathrm{inv}},D_{\mathrm{sep}}$ & the same fixed constants\\
Head/tail conversion & fixed $K_{\mathrm{head}}$ & the relative bound threshold\\
Full coefficient inverse & preceding absolute constants & $N$ and the fixed preceding constants\\
Newton ball & fixed $N=4M+40$ and preceding constants & their finite maximum\\
Convexity and non-ball witness & absolute shape embeddings and fixed-order centre constants & $N$ and the Newton thresholds\\
Classical regularity after the zero & fixed resulting domain and dimension & fixed-domain estimates\\
\bottomrule
\end{longtable}
\endgroup

The auxiliary real-parameter extensions used to construct centres
and vary the split are separate from the Euclidean conclusion.
Every final domain is formed at $p\in\tfrac12\Z$ and integer
block sizes $a,b=2p-a$. The accuracy $N$ is fixed before those
sufficiently large inputs are considered and remains fixed as $p$ grows.

\section{A polynomial positivity identity}\label{inf:sec-positivity}
The rational functions $\mathcal N_j,\mathcal E_j$ in the angular variation lemma obey
\[\frac{\mathcal N_j}{2}-\frac23\sum_{h=-1}^{2}\mathcal E_{j+h}
=\frac{4p^2}{D_{\mathrm{pos}}(p,j)}\sum_{k=0}^{12}v^k Q_k(u),
\qquad u=p-64,\quad v=j-8.\]
Here the positive denominator is
\begin{align*}D_{\mathrm{pos}}(p,j)={}&3(2j+p)^2(p+1)(p+2)^2(p+4)(p+5)\\
&\cdot(2j+p-5)(2j+p-4)^2(2j+p-3)(2j+p-2)^2\\
&\cdot(2j+p+1)(2j+p+2)^2(2j+p+3).\end{align*}
The following lists give the coefficients of $Q_k(u)$ in increasing
powers of $u$, starting with its constant coefficient. Their lengths
for $k=0,\ldots,12$ are $13,13,13,12,11,10,9,8,7,6,5,4,3$, a
total of $114$ coefficients. Higher powers have coefficient zero.
Multiplying the closed forms \eqref{inf:N-angular} and
\eqref{inf:E-angular} by $D_{\mathrm{pos}}(p,j)$ and substituting
$p=64+u$, $j=8+v$ gives the identity. Every listed coefficient is
positive, so the numerator is strictly positive for $u,v\ge0$.
\par\smallskip\noindent$k=0$: \begingroup\small\raggedright
$7753073057355915455846400$,\allowbreak{} 
$1215360332373242887229952$,\allowbreak{} 
$87272032201419163433472$,\allowbreak{} 
$3795893130561640498368$,\allowbreak{} 
$111379277105283072768$,\allowbreak{} 
$2322600629405384192$,\allowbreak{} 
$35294561281256064$,\allowbreak{} 
$393801439328684$,\allowbreak{} 
$3201816890896$,\allowbreak{} 
$18499863104$,\allowbreak{} 
$72102712$,\allowbreak{} 
$170196$,\allowbreak{} 
$184$.\par\endgroup
\par\smallskip\noindent$k=1$: \begingroup\small\raggedright
$3180378375750299550050304$,\allowbreak{} 
$473937198500360715148800$,\allowbreak{} 
$32259215635963195792896$,\allowbreak{} 
$1325755805929596036096$,\allowbreak{} 
$36623628139750000400$,\allowbreak{} 
$716093978925466680$,\allowbreak{} 
$10156021734940612$,\allowbreak{} 
$105193508825770$,\allowbreak{} 
$789036501246$,\allowbreak{} 
$4175089100$,\allowbreak{} 
$14771472$,\allowbreak{} 
$31314$,\allowbreak{} 
$30$.\par\endgroup
\par\smallskip\noindent$k=2$: \begingroup\small\raggedright
$562664300050453696883712$,\allowbreak{} 
$78738823661925490820352$,\allowbreak{} 
$5007990979423659644544$,\allowbreak{} 
$191197382688847752160$,\allowbreak{} 
$4872666916298666520$,\allowbreak{} 
$87154787492568924$,\allowbreak{} 
$1118925565579542$,\allowbreak{} 
$10351501782387$,\allowbreak{} 
$68133328557$,\allowbreak{} 
$308708398$,\allowbreak{} 
$902196$,\allowbreak{} 
$1491$,\allowbreak{} 
$1$.\par\endgroup
\par\smallskip\noindent$k=3$: \begingroup\small\raggedright
$57679238717080616484864$,\allowbreak{} 
$7495244567904683655936$,\allowbreak{} 
$439582532126405531840$,\allowbreak{} 
$15344545497976190928$,\allowbreak{} 
$353819688334375200$,\allowbreak{} 
$5650365815615444$,\allowbreak{} 
$63648016435252$,\allowbreak{} 
$504433412122$,\allowbreak{} 
$2746860336$,\allowbreak{} 
$9738624$,\allowbreak{} 
$20076$,\allowbreak{} 
$18$.\par\endgroup
\par\smallskip\noindent$k=4$: \begingroup\small\raggedright
$3852565146000679166976$,\allowbreak{} 
$459625726775754912896$,\allowbreak{} 
$24519727131647080608$,\allowbreak{} 
$769653952044921056$,\allowbreak{} 
$15727048908327064$,\allowbreak{} 
$218338672283856$,\allowbreak{} 
$2082487959758$,\allowbreak{} 
$13447668765$,\allowbreak{} 
$56116245$,\allowbreak{} 
$136135$,\allowbreak{} 
$145$.\par\endgroup
\par\smallskip\noindent$k=5$: \begingroup\small\raggedright
$177747049322645907456$,\allowbreak{} 
$19228553890224334656$,\allowbreak{} 
$918983144616089664$,\allowbreak{} 
$25449989920330928$,\allowbreak{} 
$449695578923952$,\allowbreak{} 
$5252148923508$,\allowbreak{} 
$40491107304$,\allowbreak{} 
$198350292$,\allowbreak{} 
$558912$,\allowbreak{} 
$688$.\par\endgroup
\par\smallskip\noindent$k=6$: \begingroup\small\raggedright
$5833725019983771648$,\allowbreak{} 
$564050096670576064$,\allowbreak{} 
$23716561720515776$,\allowbreak{} 
$566030886527952$,\allowbreak{} 
$8379892988992$,\allowbreak{} 
$78720883076$,\allowbreak{} 
$457623948$,\allowbreak{} 
$1502480$,\allowbreak{} 
$2128$.\par\endgroup
\par\smallskip\noindent$k=7$: \begingroup\small\raggedright
$137658610846236672$,\allowbreak{} 
$11688344788927104$,\allowbreak{} 
$422564419955328$,\allowbreak{} 
$8425693173440$,\allowbreak{} 
$99978826496$,\allowbreak{} 
$705144000$,\allowbreak{} 
$2732800$,\allowbreak{} 
$4480$.\par\endgroup
\par\smallskip\noindent$k=8$: \begingroup\small\raggedright
$2323153433438208$,\allowbreak{} 
$169347834087168$,\allowbreak{} 
$5104046345952$,\allowbreak{} 
$81336054512$,\allowbreak{} 
$721895088$,\allowbreak{} 
$3377952$,\allowbreak{} 
$6496$.\par\endgroup
\par\smallskip\noindent$k=9$: \begingroup\small\raggedright
$27391092756480$,\allowbreak{} 
$1662337753920$,\allowbreak{} 
$39942732160$,\allowbreak{} 
$474290240$,\allowbreak{} 
$2777600$,\allowbreak{} 
$6400$.\par\endgroup
\par\smallskip\noindent$k=10$: \begingroup\small\raggedright
$214440898560$,\allowbreak{} 
$10360755904$,\allowbreak{} 
$184765248$,\allowbreak{} 
$1437184$,\allowbreak{} 
$4096$.\par\endgroup
\par\smallskip\noindent$k=11$: \begingroup\small\raggedright
$1001816064$,\allowbreak{} 
$35857920$,\allowbreak{} 
$414720$,\allowbreak{} 
$1536$.\par\endgroup
\par\smallskip\noindent$k=12$: \begingroup\small\raggedright
$2113536$,\allowbreak{} 
$48896$,\allowbreak{} 
$256$.\par\endgroup

\section{Computable spectral gap enclosures}\label{inf:sec-gap-enclosures}

This appendix supplies the optional selection rule of
Remark~\ref{inf:split-selection} and is not part of the proof
chain of Theorems~\ref{inf:theorem-A} and \ref{inf:theorem-B}.

Fix an input $p,R,\theta$ and the finite set of integer splits in a
compact interval of Section~\ref{inf:sec-widened}. Represent each
quartic Dirichlet operator $H_a$ on $\Hc$ by the harmonic chart
\[
 q=1+z_{\mathrm q}S_2,\qquad d=1+5z_{\mathrm q}S_2,\qquad
 z_{\mathrm q}=-\frac{\theta\mathsf H_2(1)}{2R^2pg_2}.
\]
All parameters of this chart are computable from an isolated
Neumann zero $R$ and the rational block parameters. The following
bound supplies a strict starting condition for the spectral enclosure.

\begin{lemma}[A computable quartic relative bound]
\label{inf:quartic-computable-relative}
Assume integer block sizes with $a+b=2p$, $\bb\ge\ab\ge16$,
$p\ge64$, and
$\mu_{\mathrm{alg}}=|z_{\mathrm q}|\le1/(100p)$.
Then the positive unitary quartic operator satisfies
\begin{equation}\label{inf:algorithm-relative-bound}
 \nrm{(H_a-H_0)H_0^{-1}}_{\Hc\to\Hc}
 \le\eta_{\mathrm{alg}},\qquad
 \eta_{\mathrm{alg}}=4000p\mu_{\mathrm{alg}}.
\end{equation}
For each fixed compact split interval, the condition
$\eta_{\mathrm{alg}}<1/2$ holds at every sufficiently large
input of either theorem, uniformly over its candidate splits.
\end{lemma}
\begin{proof}
Write $\mu=\mu_{\mathrm{alg}}$ within this proof.
The endpoint proof in Lemma~\ref{inf:angular-sup} uses
$\bb-1\ge\ab-1\ge0$ and therefore applies to these parameters.
Homogeneity and its derivative bounds, at degree four, give
\[
 \nrm{S_2}_\infty\le1,\qquad
 \nrm{\nabla S_2}_\infty\le4\sqrt2<6,
 \qquad \nrm{D^2S_2}_{\mathrm{op},\infty}\le32.
\]
For the last bound, at a unit point the radial-radial entry is
$12G_2$, the radial-tangential entries are $3\nabla_SG_2$, and
the tangential block is $\nabla_S^2G_2+4G_2g_S$.
The diagonal blocks have norm at most twenty and the off-diagonal
block at most twelve. Homogeneity gives the same bound inside the ball.

Thus $q,d>0$, $q^{-1},d^{-1}\le2$, and
\[
 \nrm{D\Phi^{-1}}_{\mathrm{op},\infty}\le3,\qquad
 \nrm{A_q}_{\mathrm{op},\infty}\le9,
 \qquad \nrm{A_q-I}_{\mathrm{op},\infty}\le200\mu.
\]
Indeed the inverse is $q^{-1}(I-y\otimes\nabla q/d)$.
In \eqref{inf:pullback-metric}, the last bound follows from
\[
 |q^{-2}-1|\le8\mu,\qquad
 q^{-2}\left(2|\nabla q|/d+|\nabla q|^2/d^2\right)
 \le96\mu+576\mu^2.
\]
The drift vector in the scalar coordinate formula is
$-D\Phi^{-1}(2A_q\nabla q+y(A_q:D^2q))$.
Its norm is at most
$3(108+288n)\mu\le2000p\mu$, using $n=2p$.
For $u=H_0^{-1}f$, the ball estimates give
$\nrm{D^2u}_2\le\nrm f_2$ and
$\nrm{\nabla u}_2\le(p-1)^{-1}\nrm f_2$.
The Frobenius norm of a symmetric $n$-matrix is at most $\sqrt n$
times its operator norm. Consequently, for $H_\Phi=-P_\Phi\Delta$,
\begin{equation}\label{inf:algorithm-scalar-bound}
 \nrm{(H_\Phi-H_0)H_0^{-1}}
 \le200\sqrt{2p}\,\mu+\frac{2000p\mu}{p-1}
 \le500p\mu.
\end{equation}

Put $m=q^{p-1/2}d^{1/2}$ and $l=\log m$.
Since $p\mu\le1/100$, the inequalities
$\log(1+t)\le t$ and $-\log(1-t)\le2t$ on $0\le t\le1/2$
give $\nrm m_\infty,\nrm{m^{-1}}_\infty\le2$.
Harmonicity of $q,d$ gives
\[
 \nrm{\nabla l}_\infty\le13p\mu,
 \qquad \nrm{\Delta l}_\infty\le200p\mu^2.
\]
For example the gradient bound follows from
$2(p-1/2)6\mu+5\cdot6\mu\le13p\mu$;
the Laplacian bound follows from
$4p(6\mu)^2+2(30\mu)^2\le200p\mu^2$.
The identity
$\Delta m^{-1}=m^{-1}(|\nabla l|^2-\Delta l)$ therefore yields
\[
 \nrm{[H_0,m^{-1}]H_0^{-1}}
 \le2\left\{\frac{26p\mu}{p-1}
      +\frac{169p^2\mu^2+200p\mu^2}{(p-1)^2}\right\}
 \le60\mu.
\]
Here $p\ge64$ and $\mu\le1/(100p)$ bound the braces by $30\mu$.
It follows that $\nrm{H_0m^{-1}H_0^{-1}}\le3$.
Finally use $H_a=mH_\Phi m^{-1}$ and the exact factorization
\begin{align*}
 (H_a-H_0)H_0^{-1}
 ={}&m(H_\Phi-H_0)H_0^{-1}(H_0m^{-1}H_0^{-1})\\
    &+m[H_0,m^{-1}]H_0^{-1}.
\end{align*}
Its norm is at most $2(1500p\mu+60\mu)\le4000p\mu$.
All estimates hold on the Dirichlet graph core and extend by its
density. In the families under consideration,
$\mu_{\mathrm{alg}}\le C_{\mathrm{sp}}\tau p^{-2}$, hence
$\eta_{\mathrm{alg}}\le C_{\mathrm{sp}}\tau/p\to0$ uniformly for
fixed detuning bounds and for $\tau\le\log p$.
\end{proof}

\begin{lemma}[Convergent gap enclosures]\label{inf:gap-enclosures}
For a computable input satisfying the hypotheses of
Lemma~\ref{inf:quartic-computable-relative} and
$\eta_{\mathrm{alg}}<1/2$, the capped gap
\[
 g_a=\min\{R^2,\dist(R^2,\spec H_a)\}
\]
has rational lower and upper enclosures of arbitrarily small
prescribed width.
\end{lemma}
\begin{proof}
Set $\mathcal R_{\mathrm{alg},a}=(H_a-H_0)H_0^{-1}$ and
$\mathcal K_a=H_a^{-1}$. The strict relative bound gives
\begin{equation}\label{inf:compact-inverse-series}
 \mathcal K_a=H_0^{-1}\sum_{v=0}^{\infty}(-\mathcal R_{\mathrm{alg},a})^v,
 \qquad
 \left\|H_0^{-1}\sum_{v>L}(-\mathcal R_{\mathrm{alg},a})^v\right\|
 \le\frac{\eta_{\mathrm{alg}}^{L+1}}
 {(p-1)^2(1-\eta_{\mathrm{alg}})}.
\end{equation}
The operator $\mathcal K_a$ is compact, positive and self-adjoint.
Use the orthonormal Bessel--Jacobi basis of $H_0$ and the rectangular
projection $\Pi^{\mathrm{alg}}_{J,K_{\mathrm{rad}}}$ onto
$0\le j\le J$, $1\le k\le K_{\mathrm{rad}}$.
The radial min--max principle gives
\[
 z_{j,k}^2\ge\nu_j^2-1/4+\pi^2 k^2.
\]
Thus one may take the omitted-eigenvalue lower bound
\[
 \Lambda_{\mathrm{miss}}=\min\{\nu_{J+1}^2-1/4+\pi^2,
                   \ \nu_0^2-1/4+\pi^2(K_{\mathrm{rad}}+1)^2\}.
\]
It tends to infinity as $J,K_{\mathrm{rad}}\to\infty$.
For this projection, abbreviated $\Pi^{\mathrm{alg}}$, commutation
with $H_0^{-1}$ and self-adjointness give
\begin{equation}\label{inf:compact-inverse-tail}
 \nrm{\mathcal K_a-\Pi^{\mathrm{alg}}\mathcal K_a\Pi^{\mathrm{alg}}}
 \le\frac{2}{(1-\eta_{\mathrm{alg}})\Lambda_{\mathrm{miss}}}.
\end{equation}
Indeed $(I-\Pi^{\mathrm{alg}})\mathcal K_a$ is bounded by the
right side without its factor two, using
$\mathcal K_a=H_0^{-1}(I+\mathcal R_{\mathrm{alg},a})^{-1}$.
Its adjoint has the same bound, and their sum gives
\eqref{inf:compact-inverse-tail}.

We describe the finite matrix computations in these two truncations.
The basis values and derivatives are computable from the convergent
Bessel series, finite Jacobi polynomials and their normalizations.
The positive roots are simple and separated by more than $\pi$.
The sine comparison in Lemma~\ref{inf:radial-poles}, applied on
disjoint length-four intervals beyond $2\nu$, gives the coarse
bound $j_{\nu,k}\le2\nu+8k$. Since the first root exceeds
$\nu>1$, a rational mesh starting at one, with gaps below $\pi$,
finds every root up to a prescribed index by sign changes.
Mesh points can be chosen with certified nonzero values: of two
candidate points less than $\pi$ apart, at least one is not a zero,
and successive series enclosures detect that value. Choosing them
in small disjoint neighbourhoods of an initial mesh preserves
its order and gap bound. The same procedure refines each isolated
root interval. Simplicity then gives a nonzero derivative bound
after finitely many refinements, supplying its normalization.

For a finite basis sum $v$, the function
$\mathcal R_{\mathrm{alg},a}v$ is an explicit smooth invariant
function. Its coefficients and its squared norm are integrals in
the radial and angular variables with positive weights.
Every chart denominator is bounded away from zero by
$q\ge1-\mu_{\mathrm{alg}}$, $d\ge1-5\mu_{\mathrm{alg}}$.
The regular Bessel series times solid harmonics gives computable
uniform bounds for the function and its required derivatives on
the closed ball: each derivative series has factorial tails and
only finitely many input modes occur.
After the coordinate substitution, use an endpoint cutoff
$0<\epsilon_{\mathrm{cut}}<1/2$. The normalized radial and angular
densities are $2p\zeta^{2p-1}$ and
$x^{\ab-1}(1-x)^{\bb-1}/B(\ab,\bb)$. Their omitted endpoint
measures are bounded, respectively, by
\[
 \frac{2p\epsilon_{\mathrm{cut}}^{2p}}{2p},\qquad
 \frac{\epsilon_{\mathrm{cut}}^{\ab}}{\ab B(\ab,\bb)},\qquad
 \frac{\epsilon_{\mathrm{cut}}^{\bb}}{\bb B(\ab,\bb)}
\]
at the origin and the two angular endpoints. Multiplying by the
computed supremum bounds controls those tails.
On the remaining compact rectangle, derivative bounds give a
computable modulus of continuity and rational interval Riemann
sums converge from above and below. Here
$B(\ab,\bb)=\Gamma(\ab)\Gamma(\bb)/\Gamma(p)$.
For these integer or half-integer parameters the Gamma formulas
give convergent rational enclosures; a positive lower enclosure
of $B(\ab,\bb)$ gives the required upper bound for its reciprocal.
Equivalently one can obtain this lower enclosure from its positive
integral. Thus the tail constants $2p$ and $1/B(\ab,\bb)$ have
computable rational upper bounds. This constructs convergent
rational enclosures of every required integral.

To approximate the image of a finite vector, use Parseval in this
complete basis. For a larger finite projection $\Pi'$, its squared
residual is
\[
 \nrm{\mathcal R_{\mathrm{alg},a}v}^2
 -\sum_{i\in\Pi'}
       |\langle\mathcal R_{\mathrm{alg},a}v,e_i\rangle|^2.
\]
Its upper enclosure eventually falls below any positive target by
completeness, increasing the projection and integration precision.
This supplies a terminating precision search for each application
of $\mathcal R_{\mathrm{alg},a}$. An input error is multiplied by at
most $\eta_{\mathrm{alg}}$ and added to the next approximation error.
There are only finitely many applications in the truncated series
\eqref{inf:compact-inverse-series}. Taking errors sufficiently small
for its finitely many input columns gives a rational approximation
to the finite compression in operator norm. A column error at most
the requested tolerance divided by the number of columns suffices,
by Cauchy--Schwarz. Symmetrization preserves
that error bound.

The spectrum of a finite rational symmetric matrix can be enclosed
by rational arithmetic. One construction enumerates rational
orthogonal matrices formed from coordinate reflections and plane
rotations with entries
$(1-t^2)/(1+t^2)$ and $2t/(1+t^2)$, $t\in\mathbb Q$.
This family is dense in the orthogonal group. Continue until the
off-diagonal Frobenius norm of the conjugated matrix is below a
prescribed tolerance. The spectral theorem guarantees termination,
and the diagonal entries then enclose its spectrum with that error.
For self-adjoint operators an operator-norm error bounds their
spectral Hausdorff distance: outside that neighbourhood the
resolvent identity and a Neumann series give invertibility, in both
directions. Combining this fact with
\eqref{inf:compact-inverse-series}--\eqref{inf:compact-inverse-tail}
gives arbitrary norm-accurate finite spectral enclosures of
$\mathcal K_a$, with zero adjoined for the infinite complement.

Finally put
\[
 f_R(t)=\min\{R^2,|t^{-1}-R^2|\}\quad(t>0),\qquad
 f_R(t)=R^2\quad(t\le0).
\]
It is continuous and Lipschitz with constant $4R^4$: its nonconstant
part has $t\ge1/(2R^2)$ and derivative of absolute value at most
$1/t^2$. Also it is Lipschitz with constant one as a function of
$R^2$, so the input radius can be enclosed at the same time.
Spectral inversion gives
$g_a=\min f_R(\spec\mathcal K_a)$.
The preceding finite spectral enclosures therefore give rational
bounds on $g_a$ with error tending to zero. Clipping the lower bound
at zero preserves the enclosure. This proves the assertion.
\end{proof}

\section*{Use of AI}
AI systems, including GPT-6 Astra, GPT-6.1 Sol and Claude Opus~5.5, were
used in developing the proofs and preparing the manuscript. The author
takes responsibility for the content.

\begin{thebibliography}{99}
\bibitem{inf:dlmf}
F. W. J. Olver, A. B. Olde Daalhuis, D. W. Lozier, B. I. Schneider,
R. F. Boisvert, C. W. Clark, B. R. Miller, B. V. Saunders, H. S. Cohl,
and M. A. McClain (eds.),
\emph{NIST Digital Library of Mathematical Functions},
\url{https://dlmf.nist.gov/}, Chapter 10,
\S\S10.2, 10.6, 10.16, 10.17, and 10.21, in particular
Eqs.~10.2.2, 10.16.1, and 10.21.17.
\bibitem{inf:watson}
G. N. Watson, \emph{A Treatise on the Theory of Bessel Functions},
second edition, Cambridge University Press, 1944.
\bibitem{inf:BakerWustholz}
A. Baker and G. W\"ustholz,
\emph{Logarithmic Forms and Diophantine Geometry},
New Mathematical Monographs 9, Cambridge University Press, 2008,
\S1.2.
\url{https://doi.org/10.1017/CBO9780511542862}.
\bibitem{inf:gasper}
G. Gasper, Linearization of the product of Jacobi polynomials. II,
\emph{Canadian Journal of Mathematics} \textbf{22} (1970), 582--593,
Theorem 1. \url{https://doi.org/10.4153/CJM-1970-065-4}.
\bibitem{inf:siegel}
C. L. Siegel, \emph{\"Uber einige Anwendungen diophantischer Approximationen}
(1929), reprinted in U. Zannier (ed.), \emph{On Some Applications of
Diophantine Approximations}, Edizioni della Normale, 2014, pp.~81--138.
\url{https://doi.org/10.1007/978-88-7642-520-2_2}.
\bibitem{inf:lindemann}
F. Lindemann, Ueber die Zahl $\pi$,
\emph{Mathematische Annalen} \textbf{20} (1882), 213--225.
\url{https://doi.org/10.1007/BF01446522}.
\bibitem{inf:kato}
T. Kato, \emph{Perturbation Theory for Linear Operators},
second edition, Classics in Mathematics, Springer, 1995,
Chapters II and IV--VII.
\url{https://doi.org/10.1007/978-3-642-66282-9}.
\bibitem{inf:GT}
D. Gilbarg and N. S. Trudinger,
\emph{Elliptic Partial Differential Equations of Second Order},
Classics in Mathematics, Springer, 2001, Chapters 6--9.
\url{https://doi.org/10.1007/978-3-642-61798-0}.
\bibitem{inf:MorreyNirenberg}
C. B. Morrey, Jr., and L. Nirenberg,
On the analyticity of the solutions of linear elliptic systems of
partial differential equations,
\emph{Communications on Pure and Applied Mathematics}
\textbf{10} (1957), 271--290.
\url{https://doi.org/10.1002/cpa.3160100204}.
\bibitem{inf:MorreyBoundary}
C. B. Morrey, Jr., On the analyticity of the solutions of analytic
non-linear elliptic systems of partial differential equations.
Part II. Analyticity at the boundary,
\emph{American Journal of Mathematics} \textbf{80} (1958), 219--237.
\url{https://doi.org/10.2307/2372831}.
\bibitem{inf:sacksteder}
R. Sacksteder, On hypersurfaces with no negative sectional curvatures,
\emph{American Journal of Mathematics} \textbf{82} (1960), 609--630.
\url{https://doi.org/10.2307/2372973}.
\bibitem{inf:zygmund}
A. Zygmund, \emph{Trigonometric Series}, second edition, Volumes I and II,
Cambridge University Press, 1959, Vol.~II, Ch.~X, Theorem~3.13.
\bibitem{inf:banach}
S. Banach, Sur les op\'erations dans les ensembles abstraits et leur
application aux \'equations int\'egrales,
\emph{Fundamenta Mathematicae} \textbf{3} (1922), 133--181.
\url{https://doi.org/10.4064/fm-3-1-133-181}.
\bibitem{inf:BST73}
L. Brown, B. M. Schreiber, and B. A. Taylor,
Spectral synthesis and the Pompeiu problem,
\emph{Annales de l'Institut Fourier (Grenoble)} \textbf{23} (1973),
no.~3, 125--154. \url{https://doi.org/10.5802/aif.474}.
\bibitem{inf:Wil76}
S. A. Williams, A partial solution of the Pompeiu problem,
\emph{Mathematische Annalen} \textbf{223} (1976), no.~2, 183--190.
\url{https://doi.org/10.1007/BF01360881}.
\bibitem{inf:Zal92}
L. Zalcman, A bibliographic survey of the Pompeiu problem,
in \emph{Approximation by Solutions of Partial Differential Equations}
(B. Fuglede et al., eds.), NATO ASI Series C, vol.~365,
Kluwer, Dordrecht, 1992, 185--194.
\url{https://doi.org/10.1007/978-94-011-2436-2_17}.
\bibitem{inf:CLdDP26}
G. Cao-Labora and J. de Dios Pont,
Counterexamples to Schiffer's conjecture, arXiv:2608.05114 (2026).
\url{https://arxiv.org/abs/2608.05114}.
\bibitem{inf:CS26}
M. J. Colbrook and G. Stepaniants,
A computer-assisted counterexample to the planar Pompeiu and Schiffer
conjectures, arXiv:2608.01579 (2026).
\url{https://arxiv.org/abs/2608.01579}.
\bibitem{inf:Guo26}
J. Guo,
Convex counterexamples to the Schiffer and Pompeiu conjectures in
dimensions two to eighteen, arXiv:2609.35419 (2026).
\url{https://arxiv.org/abs/2609.35419}.
\bibitem{inf:doCarmoLima}
M. P. do Carmo and E. Lima,
Isometric immersions with semi-definite second quadratic forms,
\emph{Archiv der Mathematik} \textbf{20} (1969), 173--175.
\url{https://doi.org/10.1007/BF01899010}.
Reprinted in K. Tenenblat (ed.), \emph{Manfredo P. do Carmo---Selected Papers},
Springer, Berlin, Heidelberg, 2012, 43--45.
\url{https://doi.org/10.1007/978-3-642-25588-5_4}.
\end{thebibliography}

\end{document}
